% 本文件由 LaTeX 排版 Agent 根据有效 translation.json 自动生成。
% author=Hilary of Poitiers; work_id=3302; title=论天主圣三

\chapter{论天主圣三}

% structure: p0001 source=h1 target=omitted reason=page-title-already-rendered-by-parent

\SourceRecord{Translated by E.W. Watson and L. Pullan.；Nicene and Post-Nicene Fathers, Second Series；Vol. 9.；Edited by Philip Schaff and Henry Wace.；Buffalo, NY: Christian Literature Publishing Co.,；1899.；<http://www.newadvent.org/fathers/3302.htm>.}

% structure: p0001 source=h1 target=section reason=page-title

\section{论天主圣三（卷一）}

卷一。这部论著以圣依拉利自己的灵修历史开始，其中事件在叙述中的呈现无疑比作者经历中实际发生得更合乎逻辑且对称。他讲述了一个纯洁而高尚的灵魂——据我们所知，既未受卑劣欲望的阻碍，也未因冷漠而停滞——如何寻求人生的适当目的和目标。他首先上升到古代哲学家的概念，随后随着他在圣经中逐渐学习更多天主启示，他通过认识大公信仰中所启示的天主而达成了他寻求的目标。但这种幸福不仅仅是纯粹理智知识的结果，也是信仰的结果。在第一至十四节中，我们看到了这种从无知和恐惧到知识与平安的进步。在此他本可安歇，若非他被授予司铎（即当时所说的主教）职务，这使他承担了关心他人得救的责任。而这种关怀是必要的，因为（第十五、十六节）异端四起，主要有两种：撒贝利乌斯派认为圣父和圣子只是一个天主位格的名称或面貌，因此圣子并无真正的诞生；而亚流派（然而依拉利很少以倡导者之名称呼它，更喜欢称之为\Quote{新异端}）或多或寡公开断言圣子是被造的，而非受生，因此在本性上与圣父不同，不是真正意义上的天主。依拉利声明（第十七节）他的目的是驳斥这些异端，并通过圣经的证据来证明真信仰。他要求听众对他将引用的圣经忠心相信；没有这种信仰，他的论证对他们将无益（第十八节）；在第十九节中，他警告他们类比论证的局限性，他必须使用这种论证，尽管它因他必须用来表达无限者的有限比喻而不充分。然后在第二十节中，他以适度的自豪谈到他精心安排论证，以引导读者得出正确结论；在第二十一至三十六节中，他总结了作品的内容。他以一个祈祷结束了第一卷（第三十七、三十八节），这个祈祷表达了他对自己所持真理的确信，并恳求圣父与圣子赐予他必要的语言雄辩和推理说服力，以便值得地呈现关于他们的真理。\par

1. 当我寻求一种与人类生命能力相称且本身公义的事业——无论是出于本性的推动，还是受智者研究的启发——以便我能达到配得上所赐予我们的神圣理解恩赐的某些结果时，许多在一般观念中被认为使生活既有用又可取的事物浮现在我的脑海中。尤其是那现在以至今仍被视为最值得向往的、闲暇与财富的结合，进入了我的思绪。二者缺一似乎更像是邪恶的根源而非行善的机会，因为贫穷中的闲暇几乎被感觉为生命本身的放逐，而焦虑中的财富本身就是一种折磨，且因失去最被渴望和追求之物后所遭受的更深羞辱而更为恶化。然而，尽管这两者包含了生命中最高最好的安逸，它们似乎与野兽的正常享乐相去不远：野兽在绿草丰茂的阴凉处游荡，同时享受着免于劳役的安全和食物的丰裕。因为如果这被视为人类生命的最佳、最完美的行为方式，那么结果就是我们与整个无理性的动物世界有一个共同的目标，尽管感受的范围不同：因为所有动物，在自然界赋予的丰富供应和绝对闲暇中，都有充分的机会享受而不必为占有忧虑。\par

2. 我相信，绝大多数人之所以摒弃并谴责他人满足于这种无思想、动物般的生活，唯一的原因是自然本身教导他们：认为人活着只是为满足贪婪和怠惰，出生没有高尚的功业或美好成就的目标，并且这生命本身被赐予却没有进步到永生的能力——这样的生命，我们甚至应断言它不配被视为天主的恩赐，因为它充满了痛苦和麻烦，从空白的婴儿期到老年的迷惘，都浪费在自身之上——这是与人性不相称的。我相信，人们受自然本身的驱使，通过教导和实践提升到了我们所称为忍耐、节制和宽容的德行，因为他们深信正确的活着意味着正确的行动和正确的思想，并且不朽的天主并没有赐予生命只以死亡告终；因为没有人能相信良善的赐予者赐予了生命令人愉悦的感觉，却要让它被死亡的阴郁恐惧所遮蔽。\par

3. 然而，尽管我不能指责他们这使灵魂免于过错并预知、巧避或耐心忍受生活烦恼的劝诫为愚蠢无用，我仍不能视这些人为足以引导我达到良善幸福生命的向导。他们的格言不过是人类冲动的老生常谈；动物本能本就能理解它们，而理解却不服从的人，将陷入比动物无理性更低级的疯狂。此外，我的灵魂不仅渴望做那些忽略会带来羞耻和痛苦的事，更渴望认识那赐予这伟大恩赐的天主父，它感到将自己完全奉献给祂是当行的，事奉祂是它真正的荣耀，它所寄望于祂，在祂的慈爱中，如同在安全的家中和港湾里，它能在焦虑生命的一切烦恼中得安息。它被一种炽热的渴望燃烧着，要领悟祂或认识祂。\par

4. 这些教师中，有些人搬出了可疑神祇的大家族，并因相信神灵有性活动而叙述了从一位神到另一位神的出生与谱系。另一些人则断言有大小不一的神明，其区别与他们的能力相称。有些人否认任何神祇的存在，将崇拜局限于一种他们认为是偶然的震动与碰撞所产生的自然界。另一方面，许多人遵循共同信仰，断言有一位天主存在，却宣称祂对人事漠不关心、毫不在意。还有一些人崇拜大地和空气中的元素，即受造物的实际有形体与可见形式；最后，有些人将他们的神祇安置在人或野兽、驯养或野生的、鸟类或蛇类的图像中，将宇宙的主宰和无限的圣父囚禁在这些金属、石头或木头的狭窄牢笼中。我确信这些人不能是真理的解释者，因为尽管他们在其仪式的荒谬、污秽和不虔敬上一致，但在他们愚蠢信仰的基本信条上却存在分歧。我的灵魂在这些主张中感到困惑，却仍然沿着那条必然导向对天主真正认识的益处良多的道路前行。它不能认为对祂自己所创造的世界漠不关心值得归于天主，或者认为赋有性别及生育者与被生者的谱系与神性的纯全、强大本性相容。不，反而它确信：那神圣而永恒者必须是没有性别之分的，因为那自存者不可能在自身之外留下任何高于自己的东西。因此，全能和水恒只属于一位，因为全能不能有强弱之分，永恒不能有先后之分。在天主内，我们必须崇拜绝对的永恒和绝对的权能。\par

5. 当我的心灵思忖这些以及许多类似的思想时，我偶然看到了那些按照希伯来信仰的传统由梅瑟和先知们所写的书卷，并在其中发现了天主创造者亲自为自己作证的话语：\BibleQuote{我是自有者}，又说：\BibleQuote{那自有的打发我到你们这里来}\BibleRef{出 3:14}。我承认我很惊讶，发现其中对天主的指示如此精确，以致用最适合人类理解的术语表达了关于神圣本性奥秘的不可企及的洞见。因为心灵所能把握的天主属性中，没有比存在更能表征祂的了，因为存在，在绝对意义上，不能用来指称那将要终结的，或那有开始的，而那位现在把存在的延续与拥有圆满福乐结合在一起的，在过去不可能、在未来也不可能不存在；因为凡是神圣的，既不能生起，也不能消亡。因此，既然天主的永恒与祂自身不可分离，那么祂揭示这一件事，即祂是，作为祂绝对永恒的保证，是配得上祂的。\par

6. 为了指示天主的无限，\BibleQuote{我是自有者}这句话显然足够了；但是，我们还需要理解祂的威严与权能的运作。因为虽然绝对存在是唯独属于那永恒存在、在遥远的过去也没有开始的祂，但我们又听到从永恒而神圣的天主口中发出的一句与祂相称的话，祂说：\BibleQuote{谁用手心量过诸天，用手掌测过大地？}\EditorialNote{依撒意亚 40:12}，又说：\BibleQuote{天是我的宝座，地是我的脚凳。你们要为我建造怎样的殿宇？或是我安息的地方在哪里？}整个天被握在天主的手心中，整个地被握在祂的手掌中。现在，天主的话对虔诚心灵粗略的思考是有益的，但对耐心的学生比对片刻的听者揭示出更深的意义。因为这个被握在天主手中的天也是祂的宝座，而被抓在祂手中的地也是祂脚下的脚凳。这并非记载下来要我们从宝座和脚凳——坐姿的比喻——得出祂有像身体一样的空间延伸，因为那作为祂宝座和脚凳的，也被那无限的全能者握在手掌和手心中。它被记载，为使天主在所有被生和受造物中被认知为在内也在外，荫庇也内住，环绕一切并渗透一切，因为手掌和手，它们握着，显示祂外在控制的权能，而宝座和脚凳，通过它们对坐着者的支撑，显示外在事物对内在一位的服侍，这位内在者祂自己在外，却将一切握在掌握中，却住在属于祂自己的外在世界中。如此，天主从内到外，控制和对应着宇宙；由于祂是无限的，祂临在于万物之中，在无限的祂内，万事万物都包含在内。在这些虔诚的思想中，我的灵魂，专注于追求真理，以此为乐。因为似乎天主的伟大远远超越祂手工作品的心智能力，以至于无论有限的人心在冒险定义祂时如何尽力伸展，那挣扎的有限本性与遥不可及的无限无限之间的差距并未缩小。我通过自身的虔敬反思已经明白了这一点，但发现它被先知的话所证实：\BibleQuote{我往何处去躲避你的神？或往何处逃避你的面？我若升到天上，你就在那里；我若下到阴间，你也在那里；我若乘清晨的翅膀，住在海的极处（你就在那里），因为在那里你的手必引导我，你的右手必扶持我。}没有一处空间天主不在；空间并非独立于祂存在。祂在天上，在阴间，在海外；住在万物中，又包围一切。因此，祂拥抱宇宙，也被宇宙拥抱，不局限于任何部分，却渗透万有。\par

7. 因此，虽然我的灵魂因领会到这庄严而不可测度的智慧而欢喜，因为它能敬拜这无限的无穷性作为自己的父和创造者，但更热切地渴望认识其无限而永恒之主的真容，以便能相信那不可测量的神性穿着与他智慧之美相称的光辉。然后，当虔诚的灵魂因自身的软弱而困惑迷惘时，它从先知的声音中听到了对天主的这种绝妙的比较尺度：\BibleQuote{由伟大及受造物的美丽，造物主自可看出。}\BibleRef{智 13:5}。伟大事物的创造者至高至大，美丽事物的创造者至美至善。既然那工程超越我们的思想，则那创造者必然超越一切思想。因此，天空、空气、大地和海洋都是美丽的：整个宇宙也是美丽的，正如希腊人所承认的，他们因其美丽的秩序而称之为\OriginalTerm{宇宙}{\GreekText{κόσμος}}，即\Emph{秩序}。但如果我们的思想能凭自然本能估计这宇宙之美——正如我们在某些鸟类和兽类身上所见的那种本能，它们的叫声虽然低于我们理解的水平，但对它们自己却有清晰的意义，虽然它们无法表达；而在其中，既然一切言语都是某种思想的表达，就存在着对它们自己显明的含义——那么，这宇宙之美的主岂不应当被承认为在环绕他的所有美丽中他自己是最美的吗？因为虽然他永恒荣耀的光辉超出了我们思想的最佳能力，但却不能看不见他是美的。我们必须诚然承认天主是最美的，并且是一种美，虽然他超越我们的理解，却迫使自己成为我们感知的对象。\par

8. 因此，我的心灵充满了这些通过自身反思和圣经教导而获得的结果，便安心地如同在某座平安的瞭望塔上，停留在那荣耀的结论上，承认其本性使其只能对造物主怀有一种崇敬，而不能有其他或大或小的崇敬；即相信他的伟大超越我们的理解，但不超越我们的信德。因为理性的信德与理性相近，并接受其帮助，即使这同一理性无法应对永恒全能者的浩瀚。\par

9. 所有这些思想之下，隐藏着一种本能的希望，它加强了我对信德的肯定，即在将来，通过虔诚的思想和正直的生活，可以赢得某种完美的幸福；这仿佛是等待着胜利战士的奖赏。因为对天主的真信德若未得赏报，倘若灵魂被死亡摧毁，在肉身的生命熄灭中熄灭，那就不合理了。即使单凭理性也恳求说，天主将人引入一种分享他思想和智慧的存在，却只是为了等待生命被收回、永恒死亡的判决，这是不配于天主的；他从虚无中造出人，使他在世界上占据一席之地，却只为了让他占据后便消亡。因为，根据唯一合理的受造理论，其目的是使非存在的东西成为存在，而不是使存在的东西停止存在。\par

10. 然而，我的灵魂为自己也为身体感到沉重担忧。它坚持对天主的真信德有坚定的信念和虔诚的忠诚，但开始对自身以及那必须共享其毁灭的身体居所深感焦虑。在这种状态中，除了对律法和先知教导的认识之外，它又学会了宗徒在福音中所教导的真理——\BibleQuote{在起初已有圣言，圣言与天主同在，圣言就是天主。这圣言在起初就与天主同在。万物是藉着他而造的；凡受造的，没有一样不是藉着他而造的。在他内有生命，这生命是人的光。光在黑暗中照耀，黑暗决不能胜过他。曾有一人，由天主派遣，名叫若翰。这人来，是为作证，为给光作证。那普照每个人的真光，正在进入这世界。他在世界，世界是藉着他造的，世界却不认识他。他来到自己的领域，自己的人却没有接受他。但是，凡接受他的，他给他们权利成为天主的子女，就是那些信他名字的人；他们不是由血气，也不是由肉情，也不是由男欲，而是由天主生的。于是，圣言成了血肉，居住在我们中间；我们看见了他的光荣，正如父独生者的光荣，充满恩宠和真理。}\BibleRef{若 1:1} 在这里，灵魂超越其自然能力的成就，对天主有了比它所梦想的更多的教导。因为它得知它的创造者是出自天主的天主；它听见圣言是天主，并在起初与天主同在。它开始明白世界之光常在世界，世界却不认识他；他来到自己的领域，自己的人却没有接受他；但那些因信德接受他的人得以成为天主的子女，不是由肉体的拥抱而生，也不是由血气的孕育而生，也不是由肉身的欲望而生，而是由天主而生；最后，它得知圣言成了血肉，居住在我们中间，并且看见了他的光荣，这光荣，如同父独生者的光荣，藉着恩宠和真理而圆满。\par

11. 在此，我那战栗而忧伤的灵魂找到了比它所想象的更宽广的希望。首先，它被引入对\Person{圣父}的认识。然后，它得知那凭自然理性之光归之于造物主的永恒、无限与美，也同样属于\Person{独生子天主}。它没有使自己的信德分散于多位神祇之中，因它听说祂是\Quote{出自天主的天主}；也没有陷入将本质的差异归之于这\Quote{出自天主的天主}的错误，因它得知祂是\Quote{满溢恩宠和真理}。我的灵魂也未在\Quote{出自天主的天主}中察觉到任何违背理性之处，因为祂被启示为\Quote{在起初就与天主同在的天主}。它看到，虽然这救恩信德的赏报是伟大的，但能获得这知识的人极少，因为\Quote{自己的人却没有接受祂}，尽管那些接受祂的人，藉着并非肉体的出生，而是信德的出生，被提升为天主的子女。它也得知，成为天主的子女并非强迫，而是一种可能，因为天主的恩赐虽赐予众人，却并非不可避免的遗传，而是赐给自愿选择的奖赏。唯恐\Quote{凡愿意的都能成为天主的子女}这一事实会使我们软弱的信德动摇（因为我们最渴望、却最不期待那因其伟大而难以盼望的事），\Person{圣言}成了血肉，好藉祂的降生成人，我们的血肉得以与\Person{圣言}结合。唯恐我们以为这成了血肉的\Person{圣言}是另一位\Person{圣言}，或祂的血肉来自与我们不同的身体，祂\Quote{居住在我们中间}，好藉祂的居住显明祂是内住的天主，并藉祂在我们中间的居住，显明祂是降生成人的天主，且取了与我们相同的血肉，更甚者，尽管祂屈尊接受了我们的血肉，却仍不缺乏祂自己的属性；因为祂，\Person{圣父}的\Person{独生子}，\Quote{满溢恩宠和真理}，完全拥有祂自己的属性，并真实地享有我们的属性。\par

12. 我的灵魂欣然接受了这天主奥秘的教诲，它如今藉着血肉接近天主，藉信德蒙召获得新生，被赋予了自由和能力去赢取天上的重生，意识到其\Person{圣父}和造物主的爱，确信祂不会消灭一个祂从虚无中召唤到生命的受造物。它能估量这些真理是多么超越人的心智视野；因为处理普通思想对象的理性，无法设想任何超越它自身所感知或能自创的东西是存在的。我的灵魂衡量天主的伟大化工——这些化工是按照祂永恒的全能尺度进行——不是凭它自己的感知能力，而是凭一种无限的信德；因此它拒绝因不理解而怀疑\Quote{圣言在起初就与天主同在}以及\Quote{圣言成了血肉，居住在我们中间}，而是铭记这一真理：怀着愿相信的心，就会获得理解的能力。\par

13. 为了不让灵魂误入歧途、沉迷于外教哲学的某种幻象，它从宗徒受启发的言语中接受了这进一步的全备信德教训：—— \BibleQuote{你们要谨慎，免得有人用哲学和虚妄的诡计，按照人的传授，按照世俗的言论，而不是按照基督，把你们掳去；因为在祂内住着整个天主性圆满地有形有体地，你们在祂内也得到了圆满，祂是一切率领者和掌权者的元首；在祂内你们也受了不是人手所行的割损，而是基督的割损，脱去肉体的身躯；你们在圣洗圣事中与祂一同埋葬了，也因着那使祂从死者中复活的天主的运作，与祂一同复活了。你们从前因罪恶和未受割损的肉体是死的，祂却使你们与祂一同生活，赦免了你们的一切过犯，抹去了那藉规条反对我们、与我们为敌的债务，把它从中间撤去，钉在十字架上；祂脱去了肉体，公然展示权势，靠着自信战胜了它们。} \BibleRef{哥 2:8} 坚定的信德拒绝哲学探究的虚幻巧辩；真理拒绝被这些人类愚蠢的诡诈手段所击败，并被谎言奴役。它不将天主局限于我们共同理性的界限之内，也不按照世俗的言论来论基督，在祂内住着整个天主性圆满地有形有体地，以至于地上心智尽力去理解祂，却被那无穷的永恒和全能所挫败。我的灵魂判断祂是那一位，吸引我们向上分享祂自己的天主性，从此解除了身体礼仪的束缚；与象征性的法律不同，祂没有引入任何割损肉体的仪式，而是意图使我们的精神，从恶习中受割损，通过戒绝罪恶来净化身体的一切自然官能，使我们在圣洗圣事中与祂同死，回归永恒的生命（因为重生就是旧生命的死亡），并死于罪恶，得以重生为不朽，正如祂放弃自己的不朽为我们而死，我们也要与祂一同从死亡中醒来而得不朽。因为祂取了我们有罪的肉体，好藉着披上我们的肉体而赦免罪恶；祂与我们共享这肉体，是藉着披上它，而不是在它内犯罪。祂藉着死亡涂抹了死亡的判决，好藉着在祂内对我们族类的新创造，扫除旧法律所定的惩罚。祂让人把自己钉在十字架上，好把世界所应受的一切诅咒钉在十字架的诅咒上并予以消灭。祂作为人受尽了苦，好使权势蒙羞。因为圣经曾预言，那位是天主的应当死去；那信靠祂之人的胜利和凯旋在于：那位不朽的、不能被死亡战胜者，要死去，好使必死的人获得永生。这些天主的作为，以超越我们理解的方式完成，我再说一次，不能为我们的自然官能所理解，因为无限永恒者的工作只能被无限的智力所把握。因此，正如天主降生成人、不朽者死去、永恒者被埋葬这些真理不属于理性秩序，而是独特的权能工作，同样，那既是人又是天主、那死去的是不朽的、那被埋葬的是永恒的，这不是理智的效果，而是全能的效果。我们，因此，是由天主在基督内藉着祂的死而一同复活。但是，既然在基督内有整个天主性的圆满，我们在此就有天主圣父的启示，要与那死去者一同复活我们；我们必须承认基督耶稣不是别的，正是具有整个天主性圆满的天主。\par

14. 在这平安稳妥的保障中，我的灵魂欣然且满怀希望地安息，如此不畏惧死亡的打断，以致死亡似乎只是永生的别名。现在这肉身的生命远非重担或苦楚，而是被视为如同儿童看待字母、病人看待药水、遇难水手看待游泳、年轻人看待职业训练、未来统帅看待初次出征一样；也就是说，作为一种对当前需要的可忍受的顺从，却带着幸福永生的应许。此外，我开始宣讲那些我灵魂亲身相信的真理，作为托付给我的主教职务的责任，运用我的职位促进众人的得救。\par

15. 当我这样进行时，一些鲁莽邪恶之人的谬误暴露了出来——他们对自己绝望，对他人毫不怜悯，以自己软弱的本性衡量天主本性的威能。他们声称，不是他们已上升到对无限事物的无限认识，而是他们已将以前未定义的一切知识归结在普通理性的范围内，并限定了信德的界限。然而，宗教的真正工作是服从的服务；而这些人忽视自己的软弱，藐视神圣的实体，竟着手改进天主的教导。\par

16. 姑且不论其他异端者无谓的探究——然而，当我的论证过程有机会时，我不会保持沉默——有些人篡改福音的信德，他们借口忠于独一天主，否认天主独生子的诞生。他们声称天主延伸到人里，而非下降；祂曾暂时取了我们的血肉而成为人子，但先前并非、现在也不是天主子；在祂身上并无神圣的诞生，而生者与受生者同一；并且（为了维护他们视为对天主合一完全忠信的主张）降生成人中存在无中断的连续，圣父延伸到童贞女中，祂自己作为自己的儿子而诞生。相反，另一些人（因为是异端者，因为除基督外无救恩，基督在起初就是与天主同在的天主圣言）否认祂是生而有的，宣称祂只是受造的。他们认为，承认诞生就是承认祂为真天主，而受造证明祂的天主性不真实；虽然这种解释是对合一信德的欺骗，但若视作准确的定义，他们却认为可以当作比喻说法蒙混过去。他们在名称和信仰上，将祂真实的诞生降格为受造，为将祂从神圣的合一中割裂，如此，作为一个被召叫而存在的受造物，祂便不能要求天主性的圆满——这圆满并非借着真实的诞生而属祂。\par

17. 我的灵魂渴望回应这些疯狂的攻击。我记起，得救信德的核心不仅在于信天主，更在于信天主为圣父；不仅在于信基督，更在于信基督为天主子；信祂，不是作为受造物，而是作为生于天主的天主造物主。我的首要目标，是借着先知与福音作者明确的宣告，驳斥那些人的疯狂与无知，他们利用天主的合一（这本身是一个虔敬而有益的宣认）作为掩饰，否认或是在基督内天主诞生了，或是祂是真天主。他们的目的是在信德中心孤立一位孤独的天主，使基督——尽管大能——仅成为一个受造物；因为他们声称，天主的诞生将信徒的信德扩展为对多个天主的信赖。但我们蒙神圣教导，既不自认两个天主，也不自认一个孤独的天主；我们将引证福音与先知来支持我们宣认天主圣父与天主圣子，在信德中联合而非混淆。我们不承认他们的同一性，也不允许作为妥协，认为基督在某种不完善的意义上是天主；因为天主生于天主，不能与祂的父相同，因为祂是圣子，也不能在本性上相异。\par

18. 而你，因信德的热忱和对真理的热忱——这真理不为世界及其哲学家所知——促使你阅读我的人，必须记住要避免世俗头脑中虚弱无根据的臆测，并以虔敬的受教之心打破偏见和半知半解的障碍。需要重生理智的新官能；每个人都必须借着赐予灵魂的天上恩赐，使他的理智得到光照。首先，他必须站在天主那确定的基础\OriginalTerm{实体}{substantia = \GreekText{ὑποστάσει}}上，如圣\Person{耶肋米亚}所说，既然他要聆听关于那本性（实体），他应扩展其思想直至与主题相称，不与自身设定任意标准，而是以无限性来判断。再者，尽管他知道自己有份于神圣本性，如圣宗徒\Person{伯多禄}在\WorkTitle{第二封书信}中所说，但他不能以自身局限性来衡量神圣本性，而应以天主关于祂自身荣耀启示的尺度来衡量祂所说的。因为最好的学生并不将自己的思想读入书中，而是让书揭示自身；他从书中汲取其意义，而非将自己的意义导入其中，也不在打开书页之前就确定了一个意思而强加于其言语。既然我们将要谈论天主的事，让我们假定天主对自己有全备的知识，并谦恭地俯伏于祂的言语。因为只有借着祂自己的话我们才能认识的那一位，是有关祂自己的合适见证。\par

19. 如果在讨论天主的本性与诞生时，我们提出某些类比，但愿无人认为这些比较是完美而完全的。天主与世间事物之间不可能有比较，然而我们理智的软弱迫使我们去寻找来自较低领域的例证，以解释我们关于更高主题的意思。日常生活表明，我们在普通事物上的经验如何使我们能够对不熟悉的主题形成结论。因此，我们必须将任何比较视为对人有益，而非对天主的描述，因为它提示而非穷尽我们所寻求的意义。这样的比较，当它将血肉的本性与属神的本性、不可见与可触并置时，也不可被视为过于大胆，因为它承认自己是人类心智软弱所必需之帮助，并恳求对不完善类比可能的谴责。本着这一原则，我着手进行我的任务，打算使用天主提供的术语，但用人生的实例来为我的论证着色。\par

20. 首先，我已统筹整个著作的计划，借着各卷册的逻辑顺序来考虑读者的利益。我决意不发表未完成或考虑不周的论文，免得其杂乱无章的排列类似于一群农民混乱的喧嚷。既然无人能攀上峭壁，除非有突出的岩架帮助其登顶，我在此按次序设下我们攀登的初级轮廓，将我们困难的论证过程引向最易行的路径；不是在崖面上凿出台阶，而是将其削成平缓的斜坡，如此，行者几乎不感到努力便可抵达高处。\par

21. 因此，在第一卷之后，第二卷阐明天主圣子出生之奥迹，为使那些因圣父、圣子及圣神之名而受洗的人，得以认识真实的名号，不被其含义困扰，而是准确了解事实与意义，从而充分确信：在所使用的言辞中，他们拥有真实的名号，而这些名号蕴含着真理。\par

22. 在此关于天主圣三的简短浅显论述之后，第三卷虽缓却稳地前进。它引用祂权能的最大实例，将那本超乎领悟的言语——\BibleQuote{我在父内，父在我内} \BibleRef{若 10:38}——祂论及自身所说的话，带入信德理解的范围内。如此，超越世人迟钝心智的真理，便是配备理性与知识的信德之奖赏；因为我们既不能怀疑天主关于自身的圣言，也不能以为虔诚的理性无力领悟祂的大能。\par

23. 第四卷从异端者的教义开始，否认与他们在诋毁教会信德时所使用的谬误有任何牵连。它公布他们中有些人最近颁布的悖信信条，揭露他们从法律为所谓天主的合一而提出的论证之虚伪，因而也是邪恶。它摆出法律和先知的所有证据，以证明主张天主的合一而排除基督的天主性是不敬，以及声称如果基督是独生天主，那么天主就不是一体，这是叛逆。\par

24. 第五卷接着按照异端主张的顺序予以回应。他们曾错误地宣称，他们遵循法律，并赋予法律关于天主合一的意义，且从中证明真天主是独一位格；这是为了借他们关于独一真天主的结论，剥夺主基督的出生，因为出生是来源的证据。作为回应，我逐步肯定他们所否认的；因为从法律和先知我证明：并非有两个神，也非一个孤立无援的真天主，既没有歪曲对神圣合一的信德，也没有否认基督的出生。由于他们说主耶稣基督是被造的而非出生的，藉恩赐而拥有神圣名号，并非藉权利，我已从先知证明祂真实的天主性，以致祂被承认为真天主，祂固有天主性的确证将使我们坚持确信天主是独一的。\par

25. 第六卷揭示此异端教义的彻底欺骗性。为使其主张获得信任，他们谴责异端者的不敬教义：即\Person{瓦伦提诺}、\Person{撒伯里乌}、\Person{摩尼}、\Person{希拉加斯}。他们将教会虔诚的言辞盗用为自己的亵渎作掩护。他们挑剔并篡改这些异端者的言辞，将其修正得模糊近似正统，以便在表面上谴责异端的同时压制圣教信德。但我们清楚说明这些人的言辞和教义各自为何，并宣告教会与判罪的异端者毫无牵连或共融。应受谴责的，我们谴责；应谦卑接受的，我们接受。因此，耶稣基督的天主性圣子身份——他们最极力否认的对象——我们藉圣父的见证、基督自身的宣称、宗徒的宣讲、信徒的信德、魔鬼的哀号、犹太人的反对（其本身即是一种承认）、以及不认识天主的外邦人的认可来证明；这一切都是为了将基督未曾留给我们无知借口的真理，从争议中解救出来。\par

26. 接着，第七卷从已获致真信德的基础上，在大辩论中宣布裁决。首先，它以其对坚固信德的稳妥无可辩驳之证明，投入\Person{撒伯里乌}与\Person{厄比雍}以及这些真天主性反对者之间的激烈冲突。它质问\Person{撒伯里乌}关于否认基督先存性的立场，也质问他的攻击者们关于声称基督是受造物的断言。\Person{撒伯里乌}忽略了圣子的永恒性，但相信真天主在人体内工作。我们当前的对手否认祂是生的，断言祂是被造的，且未在祂的作为中看到真天主的工程。双方争议的，我们相信。\Person{撒伯里乌}否认那是圣子在作工，他错了；但他声称所作的工程是真天主的，则证明了他的论点有力。教会分享他对那些否认基督内真有天主之人的胜利。但当\Person{撒伯里乌}否认基督在万世之前已存在时，他的对手们令人信服地证明基督的活动是永恒的，我们站在他们一边，共同驳斥\Person{撒伯里乌}——他虽在此活动中认出真天主，却未认出是天主子。而我们先前两位对手联合起来反驳\Person{厄比雍}：第二方证明基督的永恒存在，第一方则证明祂的作为是真天主的。因此，异端者互相推翻，而教会——反对\Person{撒伯里乌}、反对那些称基督是受造物者、反对\Person{厄比雍}——见证主耶稣基督是\OriginalTerm{真天主出自真天主}{very God of very God}，在万世之前出生，又在后期作为人出生。\par

27. 没有人可以怀疑，我们遵循了真正恭敬和纯正道理的途径。我们先从法律和先知证明基督是天主子，然后证明祂是真天主，且无损于奥秘的一体性；继而我们用福音的见证来支持法律和先知，并从福音中也证明祂是天主子，且祂自己就是真实的天主。在证明祂有\Quote{子}的名号之后，最轻而易举的事就是表明这名号真实地描述了祂与圣父的关系，尽管普遍的用法认为赐予儿子的名号就是儿子身份的确凿证据。但是，为了不给那些诬蔑天主独生子真正出生的人留下任何诡诈和欺骗的漏洞，我们使用了祂真实的天主性作为祂真实儿子身份的证据；为了表明祂（众人所公认）拥有天主子名号的就是天主，我们援引了祂的名号、祂的出生、祂的本性、祂的能力、祂的声明。我们证明了祂的名号是对祂自身的准确描述，\Quote{子}的称号是出生的证据，在祂的出生中祂保持了祂的神性，并随同祂的本性保持了祂的能力，而且那能力在自觉而有意的自我启示中彰显出来。我逐一陈述了每一点的福音证据，说明祂的自我启示如何展示祂的能力，祂的能力如何揭示祂的本性，祂的本性如何因出生而属于祂，并从祂的出生而来的是祂称为\Quote{子}的称号。如此，一切亵渎的耳语都被止息了，因为主耶稣基督亲自以祂自己的口吻教导我们：正如祂的名号、祂的出生、祂的本性、祂的能力所宣告的，祂就是真实意义上的神性，从真天主而来的真天主。\par

28. 前两卷致力于确认对基督作为天主子及真天主的信仰，第八卷则专事证明天主的唯一性，表明这唯一性是与圣子的诞生并行不悖的，并且圣子的诞生并不导致神格中的二元性。首先，它揭露了这些异端者所运用的诡辩，他们虽不能否认天主圣父和圣子的真实存在，却企图回避对这种存在的宣认；它击碎了他们那无力和荒谬的辩解，即宣称在诸如\Quote{众信徒都一心一意}以及\BibleQuote{栽种的和浇灌的，都是一事}\BibleRef{格前 3:8}，以及\BibleQuote{我不但为他们祈求，而且也为那些因他们的话而信从我的人祈求，愿众人都合而为一！父啊！愿他们在我们内合而为一，就如你在我内，我在你内，使他们也得以在我们内}\BibleRef{若 17:20}这样的经文中，表达的是意志和心意的合一，而非神性的合一。通过对这些经文真实意义的思考，我们证明它们涉及神圣诞生的真实性；然后，展示我们主的一系列自我启示，我们用宗徒的语言和圣神的话语，呈现了天主作为圣父和独生子的全部而完美的荣耀奥秘。因为有圣父，我们知道有圣子；在那圣子中，圣父向我们彰显，因此我们确信祂是作为独生子而生的，并且祂是真天主。\par

29. 在有关救恩的事上，仅仅提出信德所提供并认为充分的证据是不够的。未经我们检验的论证可能会误导我们误解我们自己话语的含义，除非我们主动出击，揭露敌人证据的空洞，从而在证明其荒谬的基础上建立我们自己的信仰。因此，第九卷致力于驳斥异端者试图否定天主独生子诞生的论证；这些异端者无视自创世之初就隐藏的启示奥秘，忘记了福音信仰宣布了天主与人的结合。因为他们否认我们的主耶稣基督是天主，是与天主相似、与天主平等的圣子与圣父，由天主所生，并因其出生之权而作为真实的神灵存在，他们惯常引用我们主的话，如：\BibleQuote{你为什么称我善？除了惟一天主外，没有谁是善的}\BibleRef{路 18:19}。他们辩称，主藉着对那称祂为善的人的责备，以及祂只承认天主是善的断言，将祂自己排除在那唯一良善的天主的良善之外，也排除在那只属于唯一者的真神性之外。结合这段经文，他们亵渎性的推理又联系到另一处：\BibleQuote{永生就是：认识你，唯一的真天主，和你所派遣来的耶稣基督}\BibleRef{若 17:3}。他们说，在这里，祂承认圣父是唯一的真天主，而祂自己既不是真的也不是天主，因为对唯一真天主的认识仅限于那拥有所指定属性者。他们还自称非常清楚祂在这段经文中的意思，因为祂也说：\Quote{子不能由自己作什么，除非看见父作什么}。他们声称，祂只能仿效这一事实证明了祂本性的局限。全能者和其行动依赖于另一位先前活动者之间不能有任何比较；理性本身在能力与无能之间划下绝对的界限。这界限如此清晰，以至于祂自己关于天主圣父曾宣告说：\Quote{父比我大}。如此坦白的承认止息了所有异议；将天主的尊荣和本性归于一个否认它们的存在，乃是亵渎和疯狂。祂如此彻底地缺乏真天主的特质，以至于祂自己竟为祂自己作证说：\BibleQuote{至于那日子或那时刻，除了父以外，谁也不知道，连天上的天使和子都不知道}\BibleRef{谷 13:32}。一个不知道父亲秘密的儿子，因其无知，必然与知道秘密的父亲疏远；一个知识有限的本性不能分享那唯独免于无知统治的威严和能力。\par

30. 因此，我们揭露他们因歪曲和曲解基督话语的意义而陷入的亵渎性误解。我们通过说明祂当时在回答何种问题、在何种时刻说话、祂愿意向我们传授何种有限知识，来解释那些话语；我们让语境说明话语，而不是强制让前者与后者一致。因此，每一处差异，例如在\BibleQuote{父比我大}与\BibleQuote{我与父原是一体}\BibleRef{若 10:30}之间，或在\BibleQuote{除了天主一个外，没有谁是良善的}\BibleRef{路 18:19}与\BibleQuote{看见了我，就是看见了父}\BibleRef{若 14:9}之间，或如\BibleQuote{父啊！一切都是我的，也都是你的}与\BibleQuote{他们认识你，唯一的真天主}之间这种极大的不同，或在\BibleQuote{我在父内，父在我内}与\BibleQuote{至于那日子和那时刻，除了父以外，谁也不知道，连天上的天使和子都不知道，只有父知道}\BibleRef{谷 13:32}之间，都可以通过区分渐进的启示与祂本性及权能的完全表达来说明。两者都是同一位发言者所说的话，而对每组话语真实力量的解释将表明，基督的真天主性丝毫未受损害，因为为了形成福音信德的奥迹，基督的诞生和名号是逐渐启示的，并且是在祂选择的时间和条件下进行的。\par

31. 第十卷的目的与信德一致。因为，异端者以他们自以为有的智慧之愚昧，歪曲了受难的一些情节和言语，对主耶稣基督的天主性及权能做出无礼的矛盾论断，我被迫使证明这是亵渎性的误解，并且这些事是由主自己记录下来，作为祂真实而绝对威严的证据。在他们对信德的戏仿中，他们用这样的话自欺欺人：\BibleQuote{我的心灵忧闷得要死。}\BibleRef{玛 26:38} 他们认为，祂必定远非天主无苦无乐的生命，因为祂的心灵笼罩着对即将来临灾祸的沉重恐惧，祂在苦压下甚至谦卑自己祈祷：\BibleQuote{父啊！如果可能，请把这杯从我面前撤去}，并且确实表现出害怕承受祂祈祷求免的试炼的样子；祂的全部本性被痛苦压倒，以至于在十字架上的那一刻，祂呼喊：\BibleQuote{我的天主，我的天主，你为什么舍弃了我？} 被痛苦的苦涩所迫，抱怨自己被遗弃：祂缺乏圣父的帮助，在断气时说：\BibleQuote{父啊！我把我的灵魂交在你手里。}\BibleRef{路 23:46} 他们说，临终时刻缠绕着祂的恐惧，使祂将自己的灵魂托付给天主圣父的照顾：祂自己处境的绝对无望迫使祂将自己的灵魂交托给另一位保管。\par

32. 他们的愚昧与亵渎同样巨大，他们未能注意到，在相似情况下所说的基督的话总是一致的；他们拘泥于字句，而忽略了祂言语的目的。在\BibleQuote{我的心灵忧闷得要死}\BibleRef{玛 26:38}与\BibleQuote{从此你们要看见人子坐在权能的右边}\BibleRef{玛 26:64}之间存在着极大的差异；同样，在\BibleQuote{父啊！如果可能，请把这杯从我面前撤去}与\BibleQuote{父所给我的杯，我岂能不喝呢？}\BibleRef{若 18:11}之间；此外，在\BibleQuote{我的天主，我的天主，你为什么舍弃了我？}\BibleRef{玛 27:46}与\BibleQuote{我实在告诉你，今天你就要与我一同在乐园里}\BibleRef{路 23:43}之间；以及在\BibleQuote{父啊！我把我的灵魂交在你手里}与\BibleQuote{父啊！宽赦他们吧！因为他们不知道他们做的是什么}之间；他们狭隘的心智，无法把握神圣的意义，在试图解释时陷入亵渎。在焦虑与平静之间，在匆忙与请求拖延之间，在痛苦之言与鼓励之言之间，在绝望于己与为他人充满信心的恳求之间，存在着广泛区别；异端者显露出他们的不虔，因为他们忽视了关于天主性和基督天主性的宣称——这些解释了祂言语中的一类——而将注意力集中于那些只关乎祂在世职务的言行。因此，我已经列举了主耶稣基督灵魂与肉身奥迹中所包含的所有元素；一切都已找出，无一遗漏。接下来，以理性的宁静之光照亮这个问题，我将祂的每一句话按其意义归属到相应的类别，从而表明祂也拥有永不摇动的信心、永不犹豫的意志、永不怨言的保证；当祂将自己的灵魂托付给父时，其中也包含为他人祈求宽恕的祈祷。因此，福音教导的完整呈现解释并证实了基督的所有（而不仅是部分）话语。\par

33. 为此，连复活的荣耀也不能开启这些丧亡者的眼睛，使他们保持在信德的明确界限内，他们竟以伪装的敬畏为亵渎锻造了武器，甚至将奥迹的启示歪曲为对天主的侮辱。根据这句话：\BibleQuote{我升到我的父和你们的父那里去，升到我的天主和你们的天主那里去}\BibleRef{若 20:17}，他们辩称：既然那位父同样是我们和祂的父，那位天主也是我们的和祂的天主，祂自己承认与我们共享那与父及与天主的关系，便将祂排斥于真天主性之外，并使祂隶属于天主造物主，是祂的受造物和下属，正如我们一样，尽管祂已获得儿子的义子名分。不仅如此，我们也不可认为祂具有天主性本性的任何特质，因为宗徒说：\BibleQuote{但当祂说\InnerQuote{万物都屈服于祂}时，显然那使万物屈服于祂的，是不算在内的。万物既屈服于祂，那时，子自己也要屈服于那使万物屈服于祂的，好叫天主成为万物之中的万有}\BibleRef{格前 15:27}。他们说，因为屈服证明了被屈服者的无能以及屈服于主权者。第十一卷恭敬地讨论了这个论点；它从宗徒的这些话本身证明：屈服不但不是基督无能的证据，反而是祂作为天主子的真天主性的标记；祂的父和天主也是我们的父和天主这一事实，对我们来说是无限的优势，对祂来说绝非贬低，因为那作为人而生、并承受了我们肉身所有痛苦的那一位，已升到我们的天主和父那里，作为我们的代表去接受祂的荣耀。\par

34. 在这篇论著中，我们遵循每一门教育的进程。首先是简单的课程，通过练习慢慢熟悉学科的基础；然后学生可以在实际生活中检验他所接受的训练。同样，士兵在操练纯熟后可以上战场；辩护律师在熟悉修辞学校的辩论后，敢于进入法庭的冲突；水手在学会了在故乡的有屏障港湾中驾驶船舶后，才能被信任在公海和远方的风暴中航行。这就是我们在这一最严肃、最困难的科学——即全部信德的教导——中所采取的进程。首先是对未受教的信徒进行关于基督的诞生、名字、天主性、真天主性的简单教导；此后我们安静而坚定地前进，直到读者能够摧毁异端者的每一个论点；最后，我们终于在这伟大而光荣的冲突中让他们与敌人对阵。单靠普通资源的人类理性无力把握永恒生育的概念，但通过研究神圣事物，他们得以领悟超越平常思维范围的奥迹。他们能够驳斥关于主耶稣的那个悖论——该悖论的全部力量和貌似合理来自一种盲目的异教哲学：即断言\Quote{祂曾有一个时期不存在}、\Quote{祂在出生之前并不存在}以及\Quote{祂是从无中造成的}；仿佛祂的出生证明祂先前不存在，并在某一时刻进入存在，而天主独生子因此可以受制于时间的概念，好像信德本身（通过授予\InnerQuote{子}的称号）和出生的本性本身就证明了祂曾有一个时期不存在。因此，他们主张祂是从无中出生的，理由是出生意味着给予先前不存在的东西存在。我们回应，根据宗徒和福音作者的见证，宣认父是永恒的，子也是永恒的，并证明子是万物的天主，具有绝对而非有限的前在性；他们亵渎逻辑的这些大胆攻击——\Quote{祂是从无中出生的}和\Quote{祂在出生之前并不存在}——对祂无效；祂的永恒与子身份一致，子身份与永恒一致；在祂身上没有独特的免除出生，而是一种从永远就有的出生，因为出生意味着父，而天主性是与永恒不可分的。\par

35. 对先知语言的无知和对解释圣经的不熟练，使他们歪曲了这段经文的要点和含义：\BibleQuote{主在祂道路的起点、在祂诸工之先，就创造了为元始的我}\BibleRef{箴 8:22}，他们竭力从中确立基督作为天主是被创造的，而非被生的，因此祂分享受造物的本性，尽管祂在受造的方式上超越它们，却没有神圣出生的荣耀，只有一种超然受造物的能力。我们回应，不引入任何新的想法或先入为主的意见，而是让这段智慧本身的经文展现其真实含义和目的。我们将指出，祂为天主道路的起点和为祂的诸工而被创造这一事实，不能被曲解为关于神圣永恒出生的证据，因为为这些目的受造与从永远出生是两种完全不同的事情。提到出生，只谈出生；提到创造，首先要提及创造的原因。有一种在万有之先出生的智慧，又有为特定目的而创造的智慧；那从永远就有的智慧是一回事，那在时间流逝中出现的智慧是另一回事。\par

36. 既然我们已经得出结论，必须从我们对天主独生子的信仰宣认中排除\Quote{受造物}一词，我们现在继续提出关于圣神的理性与虔诚的教导，以便那些通过耐心而认真的研读前几卷而确立信念的读者，能够获得完整的信仰阐述。当异端教导关于这主题的亵渎也被扫除，且那使我们重生的、纯洁无玷的天主圣三奥迹，以宗徒和福音传道者的权威，用救恩的精确措辞固定下来时，这一目标就达到了。人将不再敢仅凭人的理性，将那神圣之神——我们领受祂作为不朽的保证和与天主无罪的本质共融的泉源——列于受造物之中。\par

37. 全能的上主天主，我知道作为我生命的首要职责，我应将我所有言语和思想都奉献于祢。祢所赐予的言谈之恩，除却服务祢、向盲目的世界和悖逆的异端者宣讲祢并显示祢是圣父、是独生子天主的圣父之外，不能再给我更高的赏报。但这仅是我愿望的表达；我还必须祈求祢的帮助和慈悲之恩，愿祢圣神的气息充满我所展开的信仰与宣认之帆，并赐顺风助我在教导的航程中前进。我们可以信赖那说\Emph{\BibleQuote{你们求，必要给你们；你们找，必要找着；你们敲门，必要给你们开门。}}\BibleRef{路 11:9}者所许的诺言；我们将在匮乏中祈求所需之物。我们将不知疲倦地研读祢的先知和宗徒的著作，并叩响每一扇隐秘知识之门，但回答祈祷、赐予所求、打开我们所叩之门的，是祢。我们的心灵天生迟钝而视线模糊，我们软弱的理智被禁锢在对神圣事物不可逾越的无知之墙内；但研读祢的启示将我们的灵魂提升至理解神圣真理，而顺从信仰是通往超越无助理性之确定性的道路。\par

38. 因此，我们指望祢支持这项事业的初始颤抖步伐，指望祢的援助，使它得以增强并繁荣。我们指望祢赐给我们引领先知和宗徒的那位圣神的共融，使我们能按他们说话的原意理解他们的言语，并为每一句话赋予正确的含义。因为我们将要讲述他们以奥迹所宣讲的事：关于祢，永恒的天主，永恒独生子天主的圣父，唯独无生者；以及关于唯一主耶稣基督，由祢从永远所生。我们不可因本性差异的任何借口，将祂与祢分开，或使祂成为多位神祇之一。我们不可说祂并非由祢所生，因为祢是独一。我们不可不宣认祂为真实的天主，因为祂是由祢——真实的天主，祂的圣父——所生。因此，求祢赐给我们精确的言语、稳健的论证、优雅的风格、对真理的忠诚。使我们能说出我们所相信的，以便我们能像先知和宗徒所教导的那样，宣认祢——独一天主我们的圣父，以及唯一主耶稣基督——并且封住异端者反对的口，宣告祢是天主，却非孤独的，而祂是天主，绝非虚妄的意义。\par

\SourceRecord{Translated by E.W. Watson and L. Pullan.；Nicene and Post-Nicene Fathers, Second Series；Vol. 9.；Edited by Philip Schaff and Henry Wace.；Buffalo, NY: Christian Literature Publishing Co.,；1899.；<http://www.newadvent.org/fathers/330201.htm>.}

% structure: p0001 source=h1 target=section reason=page-title

\section{论天主圣三（第二卷）}

他以命令给万民授洗 \BibleRef{玛 28:19} 作为信德的摘要开始；这本身已经足够，若不是异端者对它的含义的歪曲迫使需要解释的话。因为（第3、4节）异端是误解圣经的结果；这里我们必须注意，圣经被认为是双方共同的根据。所有人都认为它字面真实，并按照对自己最有利的方式组合经文。依拉略认为所有异端是对真理的联合反对，提出两个反对意见：他们的论证互相破坏，并且他们是近代的。然后在第5节中，他表达了自己接近这一主题时的敬畏。他必须使用的语言完全不足，但他仍被迫使用它。在第6、7节中，他开始谈论天主作为父的概念；在第8至11节中，他继而讨论天主圣子。他陈述了信德当如何相信；仅仅接受基督奇迹的真实性是不够的（第12、13节）。那奥秘，正如在 \BibleBook{若望}{福音} \BibleRef{若 1:1} 所启示的（ \BibleBook{若望}{福音} 第一章1至4节），必须是信德的对象。在第14至21节中，他针对当前的反对意见解释了这段话，然后胜利地宣称异端的一切努力都是徒劳的（第22节）。他在第23节中提出圣经证明经文反对每一个反对者，然后在第24、25节中指出我们因所启示的无限天主屈尊就卑而欠负的债。因为，在基督所降卑的一切羞辱中，天主的威严仍是不可分割地属于祂，并在祂诞生的环境及在世的生活中彰显出来（第26至28节）。本书在第29至35节中，以对圣神道理的陈述结束，这个陈述在当时该道理尚未发展的状态下尽可能完美。\par

1. \Emph{信徒们}始终在那神圣言语中寻得满足，当我们领受作为见证的大能时，我们的耳朵从福音中听到了这言语：—— \BibleQuote{你们要去使万民成为门徒，因父、及子、及圣神之名给他们授洗，教训他们遵守我所吩咐你们的一切；看！我同你们天天在一起，直到今世的终结。} \BibleRef{玛 28:19} 这段话包含了人类救恩奥秘的哪一部分呢？什么被遗忘，什么被留下在黑暗中？一切都丰盈，如同出自天主的丰盈；完美，如同出自天主的完美。这段话包含了当用的确切言词、本质的行动、过程的顺序、以及对天主本性的洞察。他命令他们因 \Emph{圣父、圣子、圣神}之名授洗，即宣认造物主、独生者及 \OriginalTerm{恩赐}{Gift}。因为天主圣父是独一的，万有都出于祂；我们的主耶稣基督，独生者，万有都是藉着祂而有的，也是独一的；圣神，天主赐予我们的恩赐，充满万有，也是独一的。这样，一切都按照所拥有的能力和所赐予的恩惠被排列——唯一大能从祂万有，唯一后裔藉着祂万有，唯一恩赐赐予我们圆满的望德。在那至高的合一中，在圣父、圣子、圣神中，包含了永恒中的无限、祂在清楚肖像中的相似、我们在恩赐中享用的祂，毫无所缺。\par

2. 但是，异端者和亵渎者的错误迫使我们处理不当之事，攀登危险的高峰，说出不可言说的话语，踏入禁地。信德本应在静默中履行诫命，崇拜圣父，与祂一起崇敬圣子，在圣神内丰盈，但我们不得不竭尽语言之贫乏来表达思想，这些思想远超言语。别人的错误迫使我们犯错，竟敢将本应静默地藏在心中敬拜的真理用人类言语表达出来。\par

3. 因为有许多人兴起，对圣经的清楚言语随意作出自己的解释，而非其真正和唯一的含义，这公然违背了言语的明确意义。异端在于所赋予的含义，而非所写的字句；罪责在于解释者，而非文本。真理难道不是不可摧毁的吗？当我们听到 \Emph{父} 这名字时，这名号不是已经隐含了子性吗？圣神被提名，祂怎能不存在？我们不能将父性从父分离，或将子性从子分离，如同我们不能否认我们所领受的恩赐在圣神中的存在一样。然而，思想扭曲的人将整个事情投入疑问和困难中，愚蠢地颠倒词语的明确含义，剥夺父的父性，因为他们想剥夺子的子性。他们拿走父性，断言子不是本性上的子；因为若生者和被生者没有相同的属性，子就不是出于父的本性；若子的存有与父的不同且不像，他就不是子。然而，天主在什么意义上是父（正如祂是），如果祂没有在祂的子中生了祂自己相同的实体和本性？\par

4. 因此，既然他们不能改变所记载的事实，便将人所发明的新原则和理论施加于其上。例如，\Person{撒伯流}使圣子成为圣父的延伸，而这方面的信德只是名辞而非现实，因为他使同一个位格既是儿子，同时也是父亲。\Person{赫比翁}不承认天主圣子有任何起源，除了来自玛利亚；他表明祂不是先为天主后为人，而是先为人后为天主；宣称童贞玛利亚并未接纳一位先存者到自己内，那在起初就是与天主同在的圣言天主，而是藉着圣言的行动生了肉身；他心目中的\Quote{圣言}不是指先存的\OriginalTerm{独生子天主}{Only-begotten God}的性体，而只是提起来的声音。同样，今日某些教师断言，天主的\OriginalTerm{肖像、智慧与德能}{Image and Wisdom and Power of God}被造成于无中，且有时间性。他们这样做，为了拯救被视为圣子之父的天主，免于被降低到圣子的水准。他们害怕圣子从祂所生的这个出生会使祂失去光荣，因此他们来援救天主，把祂的圣子说成是从无中受造的受造物，为使天主能独自完美地生活，没有一个由祂所生并分享祂性体的圣子。那么他们的圣神道理与我们不同又何足为奇呢？当他们竟敢把那圣神的赐予者置于受造、变化与不存在之下时。如此，他们破坏了信德奥迹的一致性与完整性。他们打碎了天主的绝对合一，在一切显然为每位所共有的地方指派性体上的差异；他们否认圣父，因他们剥夺了圣子的真正子性；他们否认圣神，因他们盲目于我们拥有祂、基督赐予祂的事实。他们藉夸耀其教理的逻辑完美而将未经训练的灵践出卖于毁灭；他们藉掏空术语的含义来欺骗听众，虽然名号仍为真理作证。我略过其他异端的陷阱：\Person{瓦伦廷}派、\Person{马西昂}派、\Person{摩尼}派等等。它们不时引起一些愚蠢灵魂的注意，并因接触的感染本身而致命；一旦他们教义的毒素进入听众的思想，一种瘟疫就与另一种同样具有破坏性。\par

5. 他们的背叛使我们陷入困难而危险的处境，必须在此深奥难解之事上作出超越圣经陈述的明确宣示。主曾说过，万民要\Emph{\BibleQuote{因圣父、圣子及圣神之名}}受洗。信德的言辞是清楚的；异端者却竭力使含义陷入疑惑。我们不得因此增删既定的格式，但必须限制他们解释的自由。由于他们受魔鬼狡猾的煽动，以恶毒掏空了教义的含义，同时却保留传达真理的名号，我们必须强调这些名号所传达的真理。我们必须按圣经言辞中实际找到的那样，宣告圣父、圣子及圣神的尊威与功能，从而阻止异端者剥夺这些名号所含的天主性体含义，并借这些名号本身迫使他们将术语的使用限定于其恰当含义。我无法想象我们对手的心智是怎样的，他们歪曲真理，幽暗光明，分裂不可分者，撕扯无伤者，消解完美的合一。在他们眼中，撕碎完美、为全能制定律法、限制无限也许是小事；至于我，回答他们的任务使我充满焦虑；我的头脑眩晕，理智震惊，我自己的言辞必是一种告白，并非我言词软弱，而是我哑口无言。然而，承担此任务的渴望强加于我；这意味着抵挡骄傲者、引导迷途者、警告无知者。但这主题无穷尽；我冒险以比天主自己所使用的更精确的言辞谈论天主，看不到尽头。祂已指定了名号——圣父、圣子及圣神——这些是有关天主性体的知识。言词不能表达，情感不能怀抱，理智不能理解进一步探究的结果；一切都是不可言说、不可企及、不可理解的。语言因主题的宏大而耗尽，其光辉的炫丽使凝视的眼睛失明，理智无法包容其无边的范围。尽管如此，在我们所负有的必要之下，向祂祈求赦免，这些属性都属于祂，我们将冒险、探究和谈论；而且——这是在如此严肃的事情上我们敢于作出的唯一承诺——我们将接受祂所指示的任何结论。\par

6. 一切存有都源于圣父。在基督内并借着基督，祂是万有的根源。与一切其他存在相反，祂是自存的。祂的存在并非来自外在，而是从自身并在自身内拥有。祂是无限的，因为没有任何事物能包含祂，而祂包含一切；祂永恒地不受空间限制，因为祂是无限的；永恒地先于时间，因为时间是祂的创造。任凭想象驰骋到你所认为的天主极限，你会发现祂就在那里；无论你如何竭力延伸，总有一个更远的地平线让你竭力去够。无限是祂的属性，正如你拥有做出这种努力的能力。言语会枯竭，但祂的存在不会受限于边界。或者，翻阅历史书页，你会发现祂始终在场；即使数字无法表达你所追溯的远古，天主的永恒也不会减少。振作你的理智去领悟祂的整体，祂却避开你。天主作为整体，在你掌握中留下了某物，但这某物不可分割地与祂的整体相连。因此你错过了整体，因为只有一部分留在你手中；不，甚至不是一部分，因为你是在与一个你未能分割的整体打交道。因为部分意味着分割，整体是不可分割的，而天主无处不在，且无论在哪里，祂都是完全临在。因此，理性无法应对祂，因为在祂自身之外找不到任何静观点，而永恒永远是祂的。这就是那不可测度之本质的奥迹的真实陈述，由\Quote{圣父}之名表达：天主是不可见的、不可言说的、无限的。让我们以沉默承认言语无法描述祂；让感官承认试图领悟的失败，理性承认界定的失败。然而，如我们所说，祂在\Quote{圣父}中有一个名字指示祂的本质；祂是无条件的圣父。祂不像人那样从外在源泉获得父权。祂是无始的、永存的、本质永恒的。只有圣子认识祂，因为除了圣子，并除了圣子愿意启示的人，没有人认识圣父；也不认识圣子，除非圣父。\BibleRef{玛 11:27} 二者彼此完全知悉。因此，既然\BibleQuote{除圣子外，没有人认识圣父}，我们关于圣父的思想应与圣子（唯一忠信的见证，祂向我们启示圣父）的思想一致。\par

7. 我感觉关于圣父的这些事，比说出来更容易。我很清楚没有言语足以描述祂的属性。我们必须感到祂是不可见的、不可领悟的、永恒的。但说祂是自存的、自源的、自持的，说祂是不可见的、不可领悟的、不朽的；这一切是对祂荣耀的承认，我们含义的暗示，我们思想的素描，但言语无力告诉我们天主是什么，词语无法表达实在。你听说祂是自存的；人的理性无法解释这样的独立。我们能找到支持的对象和被支持的对象，但这样存在的东西显然有别于其存在的原因。再者，如果你听说祂是自源的，找不到任何实例中生命恩赐的赐予者与所赐的生命是同一的。如果你听说祂是不朽的，那么就有某种不是出自祂的东西，并且祂凭其本性与之毫无接触；事实上，死亡并不是\Quote{不朽}这个词声称独立于天主的唯一事物。如果你听说祂是不可领悟的，这就等于说祂不存在，因为与祂接触是不可能的。如果你说祂是不可见的，一个不能有形存在的东西不能确定自己的存在。因此，我们对天主的宣认因语言的缺陷而失败；我们所能组合的最佳词语也不能指明天主的实在和伟大。对天主的完美认识，是这样认识祂：我们确信我们不可对祂无知，却又无法描述祂。我们必须相信，必须领悟，必须敬拜；这样的虔诚行为必须代替定义。\par

8. 我们现在从一个无港海岸的危险换成了开阔海洋的风暴。我们既不能安全前进，也不能安全后退，而面前的道路比身后的更加艰难。圣父是怎样的，我们便应如何相信。心思在谈论圣子时恐惧畏缩；每句话我都战栗，唯恐陷入不忠。因为祂是那无始者的苗裔，由一而生一，由真而生真，由生者而生，由完美而生完美；能力的能力、智慧的智慧、荣耀的荣耀、不可见天主的肖像、无始圣父的肖像。然而，我们如何能设想独生子是无始者的苗裔呢？圣父多次从天上呼喊：\BibleQuote{这是我的爱子，我所喜悦的}。这不是撕裂或分离，因为生者没有情欲，所生者是不可见天主的肖像，并作证：\BibleQuote{父在我内，我在父内}。\BibleRef{若 10:38} 这不是单纯的义子，因为祂是天主真正的儿子，并呼喊：\BibleQuote{看见了我，就是看见了父}。祂并非如受造物一样因命令而存在，因为祂是独一天主的独生子；祂在自己内拥有生命，如同生祂的那位拥有生命，因为祂说：\BibleQuote{正如父在自己内拥有生命，同样祂赐予子在自己内拥有生命}。圣父的部分并不在圣子内，因为圣子作证：\BibleQuote{凡父所有的，都是我的}，又说：\BibleQuote{凡是我的，都是你的；你的，也是我的}，宗徒也证明：\BibleQuote{因为在天主性的一切圆满，有形有体地居住在祂内}。\BibleRef{哥 2:9} 而按事物本性，部分不能拥有整体。祂是完美圣父的完美圣子，因为拥有一切的那位已将一切赐予了祂。然而我们不可想象圣父没有给予，因为祂仍拥有；或祂失去了，因为祂给予了圣子。\par

9. 这种受生的方式因此是仅为二者所保守的秘密。如果有人因自身能力不足而未能解开这奥秘，尽管确知圣父与圣子彼此处于那些关系中，他仍会因我所承认的无知而更加痛苦。我也茫然无知，但我并不追问。我因总领天使与我同样无知、天使未曾听闻解释、世界不能容纳它、没有先知窥见、没有宗徒追寻、圣子亲自没有启示这事实而得安慰。愿此类可怜的抱怨止息。无论你是谁，探究这些奥秘的人，我并非命你重新探索高处、宽处和深处；我只要求你耐心地安于对天主性受生方式的不知，因你尚且不知道祂的受造物如何生成。请回答我这个问题：你的感官给你任何你自己是如何被生的证据吗？你能解释你成为父亲的经过吗？我不是问你从哪里获得知觉，如何获得生命，你的理性从哪里来，你的嗅觉、触觉、视觉、听觉的本性是什么；我们拥有这一切的功用这一事实就是它们存在的证据。我所问的是：你如何将这些给予你的子女？你如何接植感官、照亮眼睛、植入心智？你若能，就告诉我。那么，你拥有你不了解的能力，你赋予你不理解的天赋。你对自己的存在的奥秘泰然处之，却对天主奥秘的无知亵渎地不耐烦。\par

10. 那么，请听无始之父，请听独生子。听祂的话：\Quote{父比我大}\BibleRef{若 14:28}，\Quote{我与父原是一体}，\Quote{看见我的，就看见了父}，\Quote{父在我内，我在父内}，\Quote{我从父出来}，\Quote{那在父怀里的}，\Quote{凡父所有的，都交付给了子}，\Quote{子在自己内有生命，如同父在自己内有生命一样}\BibleRef{若 5:26}。在这些话中，请听子、肖像、智慧、能力、天主的荣耀。其次请注意圣神宣告：\Quote{谁能述说祂的世代？}\EditorialNote{依撒意亚 53:8}。注意主的保证：\Quote{除了父，没有人知道子；除了子和子所愿意启示的人，也没有人知道父}\BibleRef{玛 11:27}。进入这奥秘，投入笼罩那受生的黑暗，在那里你将单独与无始的天主和独生的天主在一起。开始，继续，坚持。我知道你不会达到终点，但我将为你的进步而喜乐。因为虔诚地行走无尽之路的人，虽无结论，却会因自己的努力而获益。理性因缺乏言语而衰竭，但当它停下时，因所付出的努力而变得更好。\par

11. 子从那真正有生命的父领受生命；独生者从无始者，后代从父母，生命从生命。\Quote{父怎样在自己内有生命，也照样赐给子在自己内有生命}\BibleRef{若 5:26}。子从那完美者而完美，因祂从那完整者而完整。这不是分割或分离，因为彼此在对方内，而天主性的圆满在子内。不可理解者由不可理解者而生，因为除了彼此，无人认识他们；不可见者由不可见者而生，因为子是无形天主的肖像，看见子的就看见了父。有区别，因为他们是父与子；并非他们的天主性在种类上不同，因为二者是一体，天主出自天主，独一天主出自独一无始的天主。他们不是两位天主，而是出自一位的一位；不是两个无始的，因为子是从无生者所生。没有差异，因为活天主的生活在活着的基督内。关于他们天主性的本性，我决定说这些；并非认为我已成功地对信仰作了总结，而是认识到这主题是取之不尽的。你反驳说，信仰因此没有用处，因为它什么也不能领悟。并非如此；信仰的恰当服事是领悟并宣认真理，即它无力理解它的对象。\par

12. 还剩下一些关于子奥秘受生的事要说；更确切地说，这更多的事就是一切。我颤抖，我迟疑，我能力衰竭，我不知道从何开始。我无法说出子受生的时间；不确定这事实便是不敬。我该恳求谁？我该呼求谁？我该从什么书籍借取术语来陈述如此困难的问题？我该搜遍希腊哲学吗？不！我曾读过：\Quote{智者在哪里？今世的搜寻者在哪里？}\BibleRef{格前 1:20}那么，在这事上，今世的哲学家，外邦人的智者，是哑口无言的：因为他们拒绝了天主的智慧。我该转向律法的经师吗？他在黑暗中，因为基督的十字架为他成了绊脚石。我或许该命你闭眼不看异端，沉默地放过它，理由是：我们若相信癞病人得洁净、聋子听见、瘸子奔跑、瘫痪者站立、瞎子（一般地）看见、生来瞎眼的人被赐予眼睛、魔鬼被驱逐、病人痊愈、死人复活，便对我们所宣讲的那位表示了足够的尊敬。异端者承认这一切，却丧亡。\par

13. 现在请看一件不比瘸子行走、瞎子看见、魔鬼逃遁、死者复活更不奇妙的事。站在我身边，引导我穿越我所阐述的种种难题的，是一位贫穷的渔夫。他目不识丁，未受教育，手持钓线，衣衫湿透，两脚泥泞，十足的水手气。请思考并判断：是使死人复活更伟大，还是使一个未经训练的心灵领悟如此深奥的奥秘——正如他所说：\BibleQuote{在起初已有圣言}。\BibleRef{若 1:1}这\Emph{在起初已有}是什么意思？他回溯时间的空间，世纪被抛在身后，年代被取消。无论你为这个\Emph{起初}指定什么日期，你都错过了目标，因为即使在那时，我们所说的祂就已存在。环顾宇宙，注意它所记载的：\BibleQuote{在起初天主创造了天地}。\BibleRef{创 1:1}这个\Emph{起初}一词确定了创造的瞬间；你可以为一件明确说成\Emph{在起初}发生的事件指定日期。但我的这位渔夫，没有学问，未读经书，却不受时间束缚，不被其广大吓倒；他穿透了\Emph{起初}之前。因为他的\Emph{有}既无时间限制，也无开始；那\OriginalTerm{非受造的}{uncreated}\OriginalTerm{圣言}{Word}\Emph{在起初已有}。\par

14. 但或许我们会发现，我们的渔夫已经偏离了所要解决问题的条件。他已使圣言摆脱时间的限制；那自由之物过着它自己的生活，不受任何服从的约束。因此，让我们最仔细地注意接下来的话：\Emph{\BibleQuote{圣言与天主同在}}。我们发现，那在起初之前就\Emph{有}的圣言，是\Emph{与天主同在}，不受时间条件限制。那\Emph{有}的圣言，\Emph{与天主同在}。当我们寻求祂在时间中的起源时，祂是缺席的，却始终与时间的\OriginalTerm{创造者}{Creator}同在。这一次，我们的渔夫逃脱了；也许他将屈服于等待他的难题。\par

15. 因为你将辩称，言语是声音的响动；它是事物的命名，思想的表达。但这圣言与天主同在，并在起初已有；永恒思想者的思想表达必然是永恒的。现在我暂且以渔夫的名义给你一个简短的答复，直到我们看他如何为自己的朴实辩护。那么，言语的本质是：它首先是一种潜能，之后是过去的事件；只有在被听到的瞬间才存在。我们怎能说\Emph{\BibleQuote{在起初已有圣言}}，而一个言语既不存在于某一确定时间点之前，也不存活于其后？我们甚至能说存在一个时间点，其中有一个言语存在吗？不仅说话者口中的言语在说出来之前不存在，在说出的瞬间就消逝了，甚至在说出的时刻，从开始的声音到结束的声音也有变化。这就是我作为旁观者所想到的答复。但你的对手渔夫有自己的回答。他将首先责备你的疏忽。即使你未经训练的耳朵没有听清开头的话\Emph{\BibleQuote{在起初已有圣言}}，为什么抱怨下一句\Emph{\BibleQuote{圣言与天主同在}}？你所听到的是\Emph{\BibleQuote{圣言在天主内}}吗——那是某种深奥哲学的主张？或者你的方言不分\Emph{在……内}和\Emph{与……同}？断言是：那在起初已有的，是与另一位同在，而不是在另一位内。但我不从句子开头辩论；后续部分自会说明。现在请听圣言的品位和名字：\Emph{\BibleQuote{圣言就是天主}}。你的辩解——圣言是声音的响动、思想的表达——便不攻自破。圣言是实在，不是声音；是存有，不是言语；是天主，不是虚无。\par

16. 但我战栗而不敢言；这胆量使我震惊。我听见\Emph{\BibleQuote{圣言就是天主}}，而先知教导我天主是唯一的。为使我免于进一步的恐惧，渔夫朋友，请更充分地传授这一伟大奥秘。表明这些断言与天主唯一性一致；其中没有亵渎，没有曲解，没有否认永恒。他继续说：\Emph{\BibleQuote{祂在起初与天主同在}}。这\Emph{祂在起初}消除了时间的限制；\Emph{天主}一词表明祂不仅是一个声音；祂\Emph{与天主同在}证明祂既不侵犯也不被侵犯，因为祂的身份没有被另一位吞没，并且清楚地说明祂与唯一的\OriginalTerm{无始}{Unbegotten}天主同在，作为天主、祂的\OriginalTerm{独生子}{Only-begotten Son}。\par

17. 我们仍在等待，渔夫，你对圣言的完整描述。可以这样说，祂在起初已有，但也许祂不在起初之前。对此我也将以渔夫的名义作出答复。圣言只能是祂所\Emph{是}的那个；那\Emph{是}是无条件、无限的。但渔夫自己怎么说？\BibleQuote{万物是藉着祂造成的}。因此，既然万物——宇宙藉祂而成——没有祂就不存在，祂作为万有的创造者，必然有无可估量的存在。因为时间是可知的、可分的延展度量，不是空间上的，而是持续上的。万物都是来自祂，毫无例外；那么时间本身也就是祂的受造物。\par

18. 但是，我的渔夫，有人会反对说，你的言语鲁莽而过分；\Emph{\BibleQuote{万物是藉着祂造成的}}需要限定。有那\OriginalTerm{无始}{Unbegotten}者，非受造；也有那\OriginalTerm{圣子}{Son}，由\OriginalTerm{无生父}{Unborn Father}所生。这个\Emph{万物}是一个不加防备的陈述，不承认任何例外。当我们沉默不语，不敢回答或试图思考某种答复时，你却插话，\Emph{\BibleQuote{没有祂，什么也没有造成}}。你使神性的创造者回归其位，同时宣告祂有一位同伴。从你所说\Emph{没有祂}就没有造成任何事物，我得知祂并非独处。藉着祂完成工程的那一位是一；没有祂就没有完成的那一位是另一位：创造者与同伴之间划出了区别。\par

19. 对无始的创造者的敬畏曾令我困扰，惟恐你在那一切皆由圣言造成的广阔宣称中将他也包括了。你用\BibleQuote{\Emph{没有他，什么也没有造成}}消除了我的恐惧。然而正是这同一句\BibleQuote{\Emph{没有他，什么也没有造成}}又带来了烦恼和困惑。那么，确有某物是由那另一位造成的；诚然，它不是\BibleQuote{\Emph{没有他}}而造成的。如果那另一位的确造成了什么，即使圣言在造物时在场，那么\BibleQuote{\Emph{万物是借着他造成的}}便不真实。作创造者的同伴是一回事，作创造者本身是另一回事。对先前的异议我能找到自己的回答；但在此事上，渔夫，我只能立刻转向你的话\BibleQuote{\Emph{万物是借着他造成的}}。如今我明白了，因为宗徒已光照了我：\BibleQuote{\Emph{无论看得见的还是看不见的，是宝座，是宰制者，是率领者，是掌权者，一切都是借着他，并在他内所造的}}\BibleRef{哥 1:16}。\par

20. 既然万物都是借着他造成的，请来帮助我们，告诉我们什么是不借着他造成的。\BibleQuote{\Emph{那在他内造成的就是生命}}。那在他内造成的，当然不是不借着他造成的；因为那在他内造成的也是借着他造成的。万物是在他内并借着他造成的。\BibleRef{哥 1:16}他们是在他内造成的，因为他作为创造者天主而诞生。再者，凡在他内造成的，没有一样是不借着他造成的，因为受生的天主就是生命，并且作为生命而诞生，并非在出生后才成为生命；在他内并没有两个要素，一个天生，一个后来赋予。在他身上，出生与成熟之间没有间隔。凡在他内受造的，没有一样是不借着他造成的，因为他就是使他们的创造成为可能的生命。此外，天主，天主子，因其出生而成为天主，并非在出生之后。他是从永生者生为永生者，从真实者生为真实者，从完美者生为完美者，他出生时即拥有圆满的能力。他无需在日后学习他的出生是什么，而是因着他是从天主生的天主而意识到他的天主性。\BibleQuote{\Emph{我与父原是一体}}\BibleRef{若 10:30}——这是无始者的独生圣子的话。这是独一天主宣告自己为父与子的声音；父在子内发言，子在父内发言。因此也有\BibleQuote{\Emph{看见我的，就是看见了父}}；因此也有\BibleQuote{\Emph{凡父所有的，都赐给了子}}；因此也有\BibleQuote{\Emph{父怎样在自己内有生命，也赐给子在自己内有生命}}；因此也有\BibleQuote{\Emph{除了子，没有人认识父，除了父，没有人认识子}}\BibleRef{玛 11:27}；因此也有\BibleQuote{\Emph{在他内住有整个圆满的天主性，具有身体}}。\par

21. 这生命是人的光，那光照亮黑暗。为了安慰我们，因为我们无力描述先知所说的他的受生（\EditorialNote{依撒意亚 53:8}），渔夫接着说：\BibleQuote{\Emph{黑暗没有胜过他}}。\BibleRef{若 1:4}单纯理性的言语被挫败而缄默；那依偎在主怀里的渔夫得以教导表达这奥迹。他的言语不是世上的言语，因为他所处理的是不属于这世界的事物。让我们知道这是什么，如果你能从他的话中提取出任何超出字面含义的教导；如果你能将我们推导出的真理用其他词语表达，请公布出来。如果没有——确实是没有——让我们虔诚地接受渔夫的这个教导，并在他的话中认出天主的圣言。让我们满怀崇敬地紧守对圣父和圣子的真实宣认，父是无始的，子是独生的，不可言喻，其尊威超越一切表达和一切感知。让我们如同若望一样，依偎在主耶稣的怀里，好使我们也能理解并宣示这奥迹。\par

22. 这一信德及其每一部分，都借着福音的证据、宗徒的教导以及异端者四面八方的诡诈攻击的徒劳，印刻在我们身上。这基础面对风雨和洪流，坚定不移；风暴不能推翻它，滴水不能穿空它，洪水不能冲走它。其卓越性由无数次未能损害它的攻击所证明。有些药物的配方不仅对单一疾病有效，而且对一切都有用；它们具有普遍效力。公教信仰亦然。它不是一种专门针对某种疾病的药物，而是针对一切病症；毒性不能制胜它，数量不能击败它，复杂性不能困扰它。它始终如一，面对并战胜所有敌人。奇妙的是，一种言词形式竟能包含对一切疾病的治疗，一种真理陈述竟能对抗一切谎言的诡计。让异端招兵买马，每个教派都来战斗吧。我们对他们挑战的回答是：有一位无始的天主圣父，和一位独生的天主子，完美父母所生的完美后裔；圣子受生并非由于圣父的减少或其本体的减损，而是因为拥有一切的那一位生了拥有一切的圣子；圣子不是从圣父流溢或发出，而是包含且内在于整个天主性中，属于那位无论在哪里临在都永远临在的那一位；他不受时间限制，没有时间限制，因为万物是由他造成的，而且他确实不能被他自己所创造的界限所限制。这就是公教与宗徒信仰，福音教导我们并我们所宣认的。\par

23. 倘若\Person{Sabellius}胆敢将圣父与圣子混淆为两个名称同一含义，使之并非合一，而是一位；他将从福音中立刻得到答复——并非一次两次，而是再三重复：\Quote{这是我的爱子，我所喜悦的。}他将听到：\BibleQuote{父比我大}\BibleRef{若 14:28}，以及\Quote{我往父那里去}，\Quote{父啊！我感谢你}，\Quote{父啊！光荣我}，\BibleQuote{你是生活的天主之子}\BibleRef{玛 16:17}。让\Person{厄维翁}尝试侵蚀信仰吧，他否认天主之子在童贞母腹之前有任何生命，仅在肉身生命开始之后才承认圣言。我们将请他再读：\BibleQuote{父啊！在创造世界之前，我与您同享的光荣，求您现在光荣我}\BibleRef{若 17:5}，以及\Quote{在起初已有圣言，圣言与天主同在，圣言就是天主}，\Quote{万物是藉着祂而造成的}，\Quote{祂已在世界上，世界原是藉着祂造成的，但世界却不认识祂}。让那些自称宗徒的新潮传道者——假基督的宗徒——上前来，对天主之子嘲弄侮辱吧。他们必须听：\Quote{我从父出来}，\Quote{子在父怀里}，\Quote{我与父原是一体}，\Quote{我在父内，父在我内}。最后，倘若他们像犹太人一样愤怒，因基督自称天主为父，使自己与天主平等，他们必须接受祂对犹太人的回答：\Quote{你们要凭我的作为相信：父在我内，我在父内。}因此，我们唯一不动摇的根基，我们唯一幸福的信仰磐石，是来自\Person{伯多禄}口中的宣认：\BibleQuote{你是生活的天主之子}\BibleRef{玛 16:16}。在此基础上，我们可以回答每一个因乖僻的聪明或苦涩的背叛而攻击真理的异议。\par

24. 其余部分涉及父旨意的安排。童贞玛利亚、诞生、身体、然后十字架、死亡、下降阴府——这一切是我们的救恩。为了人类，天主之子由童贞玛利亚和圣神而生。在这一过程中，祂亲自服务自己；凭祂自己的大能——天主的能力——荫庇她，播下祂身体的开端，开始了祂在肉身生命的第一阶段。祂这样做，是为了藉着降生成人，从童贞玛利亚那里为自己取得血肉之性，并通过这种结合，产生一个全人类的圣化身体；以便通过祂乐意取用的那个身体，全人类都隐藏在祂内，而祂反过来，通过祂那不可见的存在，在所有人内被重现。因此，无形的天主肖像并没有鄙视人类生命开端所标志的耻辱。祂经历了每一阶段：通过成胎、诞生、啼哭、摇篮以及随后的每一次屈尊。\par

25. 对于如此伟大的屈尊，我们能作出何等相称的回报？独一无二的独生天主，无可言喻地由天主所生，进入了童贞女的子宫，成长并取了卑微人性的框架。祂维系宇宙，万有在祂内并藉着祂而存在，却由寻常分娩而生；祂的声音使总领天使和天使战栗，天地和这世界的一切元素都融化，竟被听到发出婴儿的啼哭。那不可见、不可理解，目视、感觉和触摸都无法测度的，竟被包裹在摇篮中。若有人以为这一切不配于天主，那么，恩惠越是不配于天主的尊威，他就越应承认自己因所受恩惠而欠的债。人被造所藉的那位，成为人并无所得；我们的收益是天主降生成人并住在我们中间，藉着取了那一位的肉身，使一切血肉成为祂的住所。我们因祂被贬抑而被高举；耻辱对祂是光荣对我们。祂是神，以肉身为居所，而我们反过来，从肉身被重新提升到天主那里。\par

26. 但或许挑剔的心灵会因摇篮、啼哭、诞生和成胎而困扰，我们必须将每一件事所含的光荣归于天主，好让我们以充份被祂统治要求所充满的灵魂，走近祂的自我贬抑，在祂的屈尊中不忘记祂的尊威。因此，让我们注意谁侍奉在祂的成胎周围。一位天使对\Person{匝加利亚}说话；不育者得到生育能力；司祭从焚香处出来成为哑巴；\Person{若翰}还在母腹中就爆发出言语；一位天使祝福\Person{玛利亚}，并许给她，一位童贞女，将成为天主之子的母亲。她意识到自己的童贞，为这艰难的事而忧虑；天使向她解释天主的大能作为，说：\BibleQuote{圣神要从上面降到你内，至高者的能力要荫庇你}\BibleRef{路 1:35}。圣神从上面降下，圣化了童贞玛利亚的子宫，并在其中嘘气（因为\BibleQuote{风随意向那里吹}\BibleRef{若 3:8}），祂将自己与人的血肉之性混合，用强力夺取那陌生的领地。而且，为避免因人类结构的软弱而发生失败，至高者的能力荫庇了童贞女，以一层云似的东西包围她，加强她的软弱，使那阴影（即天主的能力）坚固她的身体框架，以接受圣神的生育能力。这就是成胎的光荣。\par

27. 现在让我们来思量伴随祂诞生、啼哭和摇篮的荣耀。天使告诉若瑟，童贞女要生一个儿子，祂的名要叫厄玛奴耳，意思是\Quote{天主与我们同在}。圣神藉先知预言，天使作证；那诞生的就是\Quote{天主与我们同在}。一颗新星的光芒为贤士们显现；属天的记号护送着天上的主。天使向牧羊人报告说，主基督诞生了，世界的救主。天军众天使聚集赞美这诞生；天主亲族的喜乐宣告了这伟大工程的完成。于是，他们宣告：\Quote{天主在天上受光荣，主爱的人在世享平安。}如今，贤士们前来朝拜那裹在襁褓中的祂；在毕生从事虚妄哲学的奥秘仪式之后，他们向躺在摇篮里的婴孩屈膝。这样，贤士们弯腰尊敬幼年的软弱；祂的啼哭被天使们属天的喜乐所迎接；感动先知的那位圣神、报信的天使、新星的光芒，都围绕祂服务。就这样，圣神的降临和至高者的庇荫使祂诞生。内在的实体与外表截然不同；眼睛看到一样，灵魂看到另一样。童贞女生子；她的孩子是天主。一个婴孩啼哭；天使们被听见在赞美。有粗糙的襁褓；天主被朝拜。祂尊威的荣耀并未因取了血肉的卑微而丧失。\par

28. 在祂此后的人世生活中也是如此。祂以人形度过的全部时间都用在天主的事工上。我无暇详述；只需说，在祂所行的各样异能和医治中，一个事实显而易见：祂因所取的血肉而为人，因所行的事迹证据确凿而为天主。\par

29. 关于圣神，我不应缄默，却也无需多言；但为了那些无知的人，我不能不说。无需多言，因为我们必当宣认祂是由圣父与圣子所发。就我而言，我认为讨论祂是否存在是不恰当的。祂确实存在，因为祂被赐予、被领受、被保有；祂在我们的信德宣认中与圣父和圣子结合，并且不能被排除在圣父和圣子的真实宣认之外；如果除去一部分，整个信德就受损了。若有人问我们对此结论持何种意义，他与我们都已读到宗徒的话：\Quote{因为你们是天主的子女，天主就派遣了自己儿子的圣神，到我们心内喊说：\InnerQuote{阿爸，父啊！}}\BibleRef{迦 4:6}，以及\Quote{你们不要叫天主的圣神忧郁，因为你们是在祂内受了印证}\BibleRef{弗 4:30}，又说：\Quote{我们领受的不是这世界的精神，而是由天主来的圣神，为使我们能明了天主所赐与我们的一切}\BibleRef{格前 2:12}，还有：\Quote{你们已不属于血肉，而是属于圣神，只要天主的圣神住在你们内。谁若没有基督的圣神，谁就不属于基督}\BibleRef{罗 8:9}，以及\Quote{如果那使耶稣从死者中复活的圣神住在你们内，那么那使基督从死者中复活的，也必要藉那住在你们内的圣神，使你们有死的身体复活}。因此，既然祂存在，且被赐予、被保有，并出自天主，愿那些诽谤祂的人保持沉默。当他们问：祂是通过谁？祂为何存在？祂的本性如何？我们回答：祂是万物藉以存在、万物由之而来的那一位，并且祂是天主之神，是天主赐予信友的恩赐。若我们的回答令他们不悦，他们必定也不悦于宗徒和先知，因为他们所说的正如我们所说。此外，圣父和圣子也必受同样的不悦。\par

30. 我相信，某些人之所以继续无知或怀疑，是因为他们看到这第三个名称——圣神，常被用来指圣父或圣子。对此不必质疑；无论是指圣父还是圣子，祂都是神，且是圣的。\par

31. 但是福音的话，\BibleQuote{天主是神}\BibleRef{若 4:24}，需要谨慎审查其含义和目的。因为每一句话都有前因和目的，必须通过研究意义来确定。我们必须记住这一点，免得因了\Quote{天主是神}这句话，不仅否认了圣神的名号，也否认了祂的工作和恩赐。主曾与一位撒玛黎雅妇人说话，因为祂来是为作全人类的救赎者。祂详细谈论了活水、她的五个丈夫，以及她当时所有的并不是她丈夫的那人之后，妇人回答说：\Quote{主，我看出祢是位先知。我们的祖先在这座山上朝拜；而你们却说，人应当朝拜的地方是在耶路撒冷。}主回答说：\Quote{女人，你信我吧！时候将到，你们在山上或是在耶路撒冷，将不再朝拜父。你们朝拜你们所不认识的；我们朝拜我们所认识的，因为救恩出自犹太人。但时候将到，且现在就是，那些真正朝拜的人将以心神和真理朝拜父，因为父正寻找这样的人朝拜祂。天主是神，朝拜祂的人必须以心神和真理朝拜祂，因为天主是神。}我们看到，那妇人满脑子遗传传统，认为天主必须在山上（如撒玛黎雅）或在圣殿（如耶路撒冷）受朝拜；因为撒玛黎雅违背法律，在山丘上选定了一个朝拜地点，而犹太人则认为所罗门所建的圣殿是他们宗教的所在，两者的偏见都将无所不包、无限的天主局限在山顶或建筑的拱顶之下。天主是看不见的、不可理解的、不可测量的；主说时候已到，天主不应再在山上或圣殿中受朝拜。因为神不能被局限或束缚；祂在空间和时间中无所不在，在任何条件下都充满地临在。因此，祂说，那些以心神和真理朝拜的人才是真正的朝拜者。而这些将要在心神中朝拜天主之神的人，将以一位作为方法，另一位作为朝拜的对象：因为两者与朝拜者的关系各不相同。\Quote{天主是神}这句话并不改变圣神有祂自己的名字，以及祂是赐予我们的恩赐这一事实。那位将天主局限在山丘或圣殿的妇人被告知，天主包罗万有，自给自足；祂，那看不见、不可理解的，必须以看不见、不可理解的方式来朝拜。当基督教导说，天主既是神，就必须在圣神中朝拜祂，并揭示了在这对天主，那神，在圣神中的朝拜里，有多么大的自由和知识、多么无限的崇拜境界时，那所赐予的恩赐和朝拜的对象便清楚地显明了。\par

32. 宗徒的话也有同样的含义：\BibleQuote{主就是神，主的神在哪里，那里就有自由}\BibleRef{格后 3:17}。为了清楚说明他的意思，他区分了存在的神，以及祂的神所属的那一位：\Emph{祂}和\Emph{祂的}在意义上不同。因此，当他说\Quote{主就是神}时，他揭示了天主的无限；当他补充说\Quote{主的神在哪里，那里就有自由}时，他指出了属于天主的那一位；因为祂是主的神，并且\Quote{主的神在哪里，那里就有自由}。宗徒作此陈述，并非出于他自己论证的需要，而是为了清晰。因为圣神处处是一，光照所有圣祖、先知和法律的全班，甚至在母腹中感动若翰，在适当的时候赐予宗徒和其他信徒，使他们能认识所赐予的真理。\par

33. 让我们从主亲口的话来听圣神在我们内的工程。祂说：\BibleQuote{我还有许多事要告诉你们，但你们现在不能承载}\BibleRef{若 16:12}。因为\Quote{为你们有益的是我离去；我若离去，我必派遣护慰者到你们这里来}。又说：\Quote{我要求父，祂必赐给你们另一位护慰者，使祂永远与你们同在，就是真理之神}。\Quote{祂要引导你们进入一切真理，因为祂不凭自己说话，而只把所听到的说出来，并把将来的事宣告给你们。祂要光荣我，因为祂要从我那里领受。}这些话是为显示群众应如何进入天国；它们包含了对施予者善意的保证，以及对恩赐的方式和条件的说明。它们告诉我们，因为我们软弱的心智无法理解父或子，我们那觉得天主的降生成人难以相信的信德，将由圣神——联合之纽带与光明之源——的恩赐所光照。\par

34. 下一步自然是听宗徒讲述这恩赐的能力和功用。他说：\BibleQuote{凡被天主圣神引导的，都是天主的子女。因为你们没有领受恐惧的奴仆之神，而是领受子女之神的，藉着我们喊说：\Quote{阿爸，父呀！}}\BibleRef{罗 8:14}；又说：\BibleQuote{没有人藉着天主圣神说\InnerQuote{耶稣当受诅咒}，也没有一个能说\InnerQuote{耶稣是主}，除非在圣神内}\BibleRef{格前 12:3}；他还补充说：\Quote{恩赐各有不同，却是同一圣神；职务各有不同，却是同一主；功效各有不同，却是同一天主，在一切人身上行一切事。圣神赐给每个人显明的恩宠，是为人的好处。这人藉圣神获得智慧的言语，那人按照同一圣神获得知识的言语，另一人因同一圣神获得信德，另一人获得治病的奇恩，另一人行奇迹，另一人说预言，另一人辨别神恩，另一人说各种语言，另一人解释语言。可是这一切都是同一圣神所行的。}这里我们有了对恩赐的目的和结果的陈述；在如此清楚地定义了祂的起源、行动和能力之后，我实在不能想象还会有什么疑惑留存。\par

因此，让我们善用这伟大的恩惠，寻求对这最需要的恩赐的亲身体验。因为宗徒说，在我已引用过的话中，\Emph{\BibleQuote{但我们并没有领受此世的精神，而是领受了那出于天主的圣神，为使我们知道天主所白白赐给我们的事物。}}我们领受祂，正是为知道。人体官能若被剥夺运用，便会陷于休眠。眼睛若无自然或人工之光，便不能尽其职责；耳朵若不闻声音或响动，便不知其功用；鼻孔若不闻气味，便不知其目的。并非官能不存在，只因从未被使用，而是对其存在毫无经验。同样，人的灵魂，除非因着信德领受了圣神的恩赐，便具有认识天主的先天官能，却缺乏知识之光。那在基督内的恩赐是同一的，且被白白赐予，完全赐予众人；不被拒绝给任何人，按每人愿意接受的度量而给予；恩赐的储藏越丰富，获取它的渴望就越热切。这恩赐与我们同在直到世界终穷，是我们等待的安慰，因祂所赐的恩惠，是我们将来希望的保证，是我们心灵的光明，我们灵魂的太阳。这圣神我们必须寻求并赢得，然后以信德和服从天主的诫命牢牢持守。\par

\SourceRecord{Translated by E.W. Watson and L. Pullan.；Nicene and Post-Nicene Fathers, Second Series；Vol. 9.；Edited by Philip Schaff and Henry Wace.；Buffalo, NY: Christian Literature Publishing Co.,；1899.；<http://www.newadvent.org/fathers/330202.htm>.}

% structure: p0001 source=h1 target=section reason=page-title

\section{\WorkTitle{论天主圣三（第三卷）}}

第三卷。在§§1-4中，主的话\Quote{我在父内，父在我内}被视为典型。人不能理解，只能领悟。希拉里在可解释的范围内解释了它们。但天主的自我启示总是神秘的。基督的奇迹是不可解释的（§§5-8）；这是天主的方式，旨在抑制人的僭越。人的智慧是有限的，当它越界侵入信德领域时，就变成了愚妄。接着，在§§9-17中，经文\BibleRef{若 17:1}被解释为证明在独一的天主内有圣父和圣子两位，并揭示天主作为父的形像。然后，在§§18-21中，基督的奇妙事迹被提出为祂奇妙诞生的证据。我们不可问祂如何能与父同永，因为我们问祂如何能穿过紧闭的门也是徒劳。这两种问题都是僭越。基督所给予的启示（§§22-23）是：天主是祂的父；\Emph{\OriginalTerm{他们是\InnerQuote{一}，而非\InnerQuote{一个}}{Unum sunt, non Unus}}。最后，在§§25-26中，他回到推理的徒劳。真正的智慧是在我们不能理解时相信；我们必须信赖信德，而非证明。\par

1. 主的话：\Quote{我在父内，父在我内}\BibleRef{若 14:11}使许多人困惑，这并不奇怪，因为人类理性的力量无法提供任何可理解的意义。似乎不可能有某物既在另一物之内又在它之外，或者（既然我们讨论的存在尽管不相分离，但仍保持各自的存在和状态）这些存在能互相包含，以致一个永久地包裹另一个，同时又永久地被另一个包裹，而这个另一个却是它包裹的对象。这是人的智慧永远无法解决的问题，人类的探究也永远无法找到这种神圣存在状态的类比。但人不能理解的，天主却能是。我并不是说，既然这是天主所做的断言，它就立刻为我们所理解。我们必须自己思考，并逐渐认识这句话：\Quote{我在父内，父在我内}的含义；但这取决于我们能否把握这一真理：基于神圣真理的推理能够建立其结论，即使它们似乎与宇宙的规律相矛盾。\par

2. 为了尽可能容易地解决这个极其困难的问题，我们首先必须掌握圣经关于圣父和圣子的知识，以便我们能够更精确地谈论这些熟悉和惯常的事情。如我们在上一卷书中详细讨论后所断定的，圣父的永恒超越空间、时间、现象以及人类思想的一切形式。祂在万有之外又在万有之内，祂包含万有却不被任何所包含，不能因增加或减少而改变，不可见、不可理解、圆满、完美、永恒，不从另一处获得祂所拥有的一切，但若有从祂而来的，祂仍然完整且自足。\par

3. 因此，那\OriginalTerm{未生者}{Unbegotten}在时间之前从祂自身生了一个圣子；不是从任何预先存在的质料，因为万有都是藉着圣子；不是从无，因为圣子来自圣父自身；不是藉着生育，因为在天主内既无变化也无虚空；不是作为祂自己的一块被切割、撕下或拉长，因为天主是无情无体的，只有可能和有身体的存在才能这样被对待，而且正如宗徒所说，在基督内\Quote{蕴藏着天主圆满的\Emph{天主性}}\BibleRef{哥 2:9}。以不可理解、不可言说的方式，在时间和世界之前，祂从自己的未生实体生了独生子，藉着爱和能力将祂全部的天主性赋予那诞生。因此，祂是未生、完美、永恒圣父的完美、永恒、\OriginalTerm{独生子}{Only-begotten}。但祂因所取的身体而具有的那些属性，是祂为我们的救恩所怀的善意之果。因为祂作为天主子，本是无形无体、不可理解的，却取了那样程度的有形体和卑微，以便能进入我们的理解、感知和默观范围。这是对我们软弱的俯就，而非放弃祂自己的固有属性。\par

4. 因此，祂是完美圣父的完美圣子，未生天主的独生子，从拥有一切的那一位领受一切，是出自天主的天主、出自圣神的圣神、出自光的光，大胆地说：\Quote{父在我内，我在父内}\BibleRef{若 10:38}。因为正如父是圣神，子也是圣神；正如父是天主，子也是天主；正如父是光，子也是光。因此，那些在父内的属性是子所赋有的属性的源头；也就是说，祂是完全由完全之父而来的完全之子；不是从外部引入的，因为在子之前一无所有；不是从无造成的，因为子来自天主；不是局部之子，因为天主性的圆满在子内；不是在某些方面是子，而是在一切方面；是按照那有能力者的旨意，以祂唯一知道的方式成为子。凡在父内的，也在子内；凡在未生者内的，也在独生子内。一位出自另一位，两者是一体；不是两个成为一体，而是一位在另一位内，因为两者内的是同一的。父在子内，因为子来自祂；子在父内，因为父是祂的唯一根源；独生子在未生者内，因为祂是出自未生者的独生子。因此，彼此内互相包含，因为正如在未生父内一切圆满，在独生子内也一切圆满。这就是子在父内、父在子内的合一，这就是能力，这就是爱；我们的希望、信德、真理、道路和生命，不是争论父的能力或贬低子，而是敬畏祂诞生的奥秘和威严；使未生父超乎一切较量，承认独生子在永恒和能力上与祂相等，宣认关于天主子——祂来自天主。\par

5. 天主具有这样的权能，我们的理性方法无法理解，但我们的信德基于祂行动的确切证据而深信不疑。我们将在肉身领域和灵性领域中找到这种行动的实例，其彰显并非以能说明圣子诞生之比喻的形式，而是以一件可理解却奇妙的事迹。在加里肋亚的婚宴上，水变成了酒。我们是否有言语可述、感官可查明：究竟是何种方法产生了这一变化，使水的无味消失，代之以酒的浓郁醇香？这不是混合，而是一种创造，这创造不是开端，而是转化。并非一种较弱的液体通过混入较强的元素而获得；而是现有之物消失，新物诞生。新郎焦虑，家人慌乱，婚宴的和睦岌岌可危。耶稣被请求帮助。祂没有起身或忙乱，祂毫不费力地完成了这工作。水被倒入缸中，酒从杯中被取出。倒水之人的感官与取酒之人的感官相矛盾。倒水的人期望倒出的是水，取酒的人认为定是倒入了酒。中间的时间无法解释液体性质的变化。行动的方式避开了视觉和感官，但天主的权能在所成就的结果中彰显。\par

6. 在五个饼的事例中，同类型的奇迹令我们惊叹。通过饼的增加，五千男人以及无数的妇女儿童免于饥饿；其方法避开了我们观察的能力。五个饼被献上并被擘开；当宗徒们分饼时，一连串新创造的碎块不知不觉地经过他们的手。他们正在分的饼并未变小，但他们的手却始终满握碎块。过程的迅疾避开了视线；你眼看着一只手满握碎块，同时看到另一只手中的饼并未减少，而碎块的堆一直在增长。切饼的人忙于工作，吃的人努力吃；饥饿的人得到饱足，剩下的碎块装满了十二筐。视觉或感官无法发现如此引人注目的奇迹之方式。那不存在的被创造出来；我们所见的超越了我们的理解。我们唯一的资源就是信靠天主的全能。\par

7. 在这些天主的奇迹中，没有欺骗，没有为取悦或欺哄而作的微妙伪装。天主之子的这些行为并非出于自我炫耀的欲望；那有无数的天使事奉的祂，从未欺哄过人。我们所拥有的一切都是借祂而创造的，祂能需要我们的什么？祂需要我们这些或沉睡、或纵欲、或充满暴行和流血的罪债、或因狂欢而沉醉之人的赞美吗？祂是那总领天使、宰制者、率领者、掌权者在天上以永恒不倦的声音无眠、无休、无罪地赞美的。他们赞美祂，因为祂是不可见的天主的肖像，在祂自己内创造了他们全体，创造了世界，立定了诸天，安置了星辰，稳固了大地，奠定了深渊的根基；因为后来祂降生，战胜了死亡，打破了地狱之门，为自己赢得了同为继承者的子民，将肉身从腐朽中提升到永生的荣耀。因此，祂不可能从我们这里获取什么，能促使祂以这些神秘而无法解释的作为（好像祂需要我们的赞美）来披上辉煌。但天主预见人类的罪恶和愚妄将如何被误导，并知道不信将敢于对天主的事作出判断，因此祂用祂权能的记号（这必使我们最胆大妄为的人踌躇）战胜了僭越。\par

8. 因为世上许多智者（他们的智慧在天主前是愚妄）否认我们所宣告的天主出自天主、真出自真、完美出自完美、一出自一，好像我们教导不可能的事。他们坚信通过逻辑推理所达到的某些结论：\Quote{无不能生于一，因为每一次生育都需要两位父母}；\Quote{如果这子生于一，祂已领受了生者的一部分；若祂是一部分，那么两者都不完美，因为那从之生出子的祂缺少了某物，并且那由另一部分构成的祂不可能有圆满。因此，两者都不完美，因为生者失去了祂的圆满，而受生者并未获得它。}这就是那世上的智慧，天主甚至在先知的时代就已预见，并通过先知谴责说：\BibleQuote{我要消灭智者的智慧，废弃明达者的明达。}\BibleRef{依 29:14}宗徒也说：\BibleQuote{智者在哪里？经师在哪里？这世代的辩士在哪里？天主岂不是使这世上的智慧成了愚拙吗？因为世人凭自己的智慧，既然不认识天主，天主就乐意借着愚拙的宣讲，拯救那些相信的人。犹太人要求神迹，希腊人寻求智慧；我们却宣讲被钉十字架的基督：这为犹太人是绊脚石，为外邦人是愚拙；但那蒙召的，无论是犹太人或是希腊人，基督是天主的能力、天主的智慧。因为天主的愚拙总比人聪明，天主的软弱总比人强壮。}\BibleRef{格前 1:20}\par

9. 因此，天主子负有管理人类的职责，首先成了人，好使人能信祂；祂从我们中间兴起，为我们作属神之事的见证，并藉着肉身的软弱，向我们这些软弱而属血肉的人宣讲天主父，如此承行父的旨意，正如祂所说：\BibleQuote{我来不是为执行我的旨意，而是为执行派遣我来的那位的旨意}。\BibleRef{若 6:38} 这并不是祂自己不愿意，而是为彰显祂因父的旨意而有的服从，因为祂自己的旨意就是承行父的旨意。这就是那承行父旨意的意志，祂以这些话作证：\BibleQuote{父啊，时辰来到了，请光荣你的子，使子也光荣你；正如你赐给祂权柄掌管一切血肉，使祂将你所赐给祂的人，赐给他们永生。永生就是：认识你，唯一的真天主，和你所派遣的耶稣基督。我在地上已光荣了你，完成了你所委托我作的工作。现在，父啊，请在你面前光荣我，赐给我在创世之前，我与你所享有的光荣。我已将你的名显示给你赐给我的人们}。祂用简短而少的话语揭示了祂被委派和指定的全部任务。然而，这些话语虽简短而少，却是真信仰抵御魔鬼一切诡计建议的保障。让我们简要思考每一短语的力度。\par

10. 祂说：\BibleQuote{父啊，时辰来到了；光荣你的子，使子也光荣你}。祂说时辰来到了，而非日子或时间。时辰是一日的分划。这必定是哪个时辰？当然是祂在受难时，为坚强门徒所说的那个时辰：\BibleQuote{人子应受光荣的时辰来到了}。\BibleRef{若 12:23} 因此，这就是祂祈求受父光荣的时辰，为使祂自己光荣父。但这是什么意思？将要给予光荣的祂反要求接受光荣？将要赐予尊荣的祂反为自己恳求？祂会缺少即将偿还的那事物吗？在此，让世上的哲学家、希腊的智者来拦阻我们的道路，撒下他们的三段论之网来缠住真理。让他们问：如何？来自何处？为何？当他们找不到答案时，让我们告诉他们，这是因为\BibleQuote{\Emph{天主拣选了世上愚妄的，为羞辱那有智慧的}}。\BibleRef{格前 1:27} 这就是为什么我们这些愚妄的人能理解世上哲学家所不能理解的事。主曾说：\BibleQuote{\Emph{父啊，时辰来到了；}}祂揭示了祂受难的时辰，因为这些话正是在那一刻说的；然后祂又加上：\BibleQuote{光荣你的子}。但是子要如何受光荣呢？祂由童贞女诞生，从摇篮和童年长大成人，经过睡眠、饥饿、口渴、疲乏和眼泪，活出了人的生活；现在祂还要被吐唾沫、被鞭打、被钉十字架。为什么呢？这些事是为保证我们在基督内看到纯全的人性。但十字架的羞辱不属于我们；我们没有被判受鞭打，也没有被吐唾沫所玷污。父光荣子，如何光荣？接下来祂被钉在十字架上。然后发生了什么？太阳没有落山，反而逃走了。怎么回事？它不是隐在云后，而是离开了既定的轨道，世界的一切元素都感受到了基督之死的同一冲击。星辰在它们的轨道上，为避免同谋，自行熄灭以逃避目睹那场景。大地做了什么？它在悬挂于树上的主的重负下颤抖，抗议自己无力束缚那垂死者。然而，岩石和石头岂能不拒绝给祂安息之所？是的，它们裂开、破碎，力量耗尽。它们必须承认，凿于岩石中的坟墓不能囚禁那等待埋葬的身体。\par

11. 接下来呢？大队的百夫长、十字架的看守者喊道：\BibleQuote{\Emph{这人真是天主子}}。\BibleRef{玛 27:54} 因这\OriginalTerm{赎罪祭}{Sin-offering}的中介，受造物被释放；连岩石也失去坚固和力量。那些把祂钉在十字架上的人承认，这人真是天主子。结果证实了这断言。主曾说：\BibleQuote{光荣你的子}。祂藉着\Emph{你的}这个词断言，祂是天主子，不仅在名号上，而且在本质上。我们众人都是天主的子女；祂却是另一种意义上的子。因为祂是天主\OriginalTerm{真实而独一的圣子}{true and own Son}，出于本源而非因义子名分，不仅在名号上，而且在真理上，是生而非受造的。因此，在祂受光荣之后，这宣认触及了真理；百夫长宣认祂是真正的天主子，为使任何信徒不要怀疑一个甚至连迫害者的仆役也不能否认的事实。\par

12. 但或许有人会以为，祂为了祈求的那荣耀是祂所缺乏的，而祂期望被一位更伟大者所荣耀，这正显明了祂能力的不足。确实，谁会否认圣父是更伟大的呢？无始的比受生的更伟大，圣父比圣子更伟大，派遣者比被派遣者更伟大，愿意者比服从者更伟大？祂自己要作自己的见证：——\BibleQuote{圣父比我大}。这是我们必须承认的事实，但我们必须谨慎，免得在未经深思的思想中，圣父的尊威遮掩了圣子的荣耀。这种遮掩是圣子所祈求的同一荣耀所禁止的；因为祈祷\Quote{圣父，光荣祢的圣子}，是由\Quote{使圣子也光荣祢}来完成的。因此，圣子并不缺乏能力，祂在领受了这荣耀之后，将以荣耀回报。然而，如果祂并不缺乏，为何祂还要祈祷呢？没有人会祈求他所不需要的东西。难道圣父也有所缺乏吗？或者圣父如此轻率地将祂的荣耀给予，以致需要圣子归还吗？不；一位从未缺乏，另一位也无需祈求，然而两者都要彼此给予。因此，祈求荣耀的给予和回报，既不是对圣父的掠夺，也不是对圣子的贬低，而是证实了同一天主性住在两者之中。圣子祈求自己被圣父所荣耀；圣父不认为被圣子所荣耀是屈辱。荣耀的互换——给予和接受——宣告了圣父与圣子权能的合一。\par

13. 接下来我们必须确定这荣耀是什么，从何而来。我确信，天主是不会改变的；祂的永恒不容有缺陷或增补，也不容有增益或损失。惟独祂有这样的特性：祂从永远就是祂所是的。祂从永远就是的，出于祂的本性，祂不可能不再是的。那么，祂怎么能接受荣耀呢？祂已完全拥有荣耀，祂的宝藏并不减少；没有新的荣耀可以获取，也没有失去的荣耀可以恢复。我们陷入了困境。但福音作者并没有让我们失望，尽管我们的理性已显明其无助。为告诉我们圣子应归还给圣父的是何种荣耀，祂说出这些话：\BibleQuote{正如祢赐给祂权柄掌管凡有血肉的人，使祂将所赐给祂的人，赐给他们永生。永生就是：认识祢，惟一的真天主，和祢所派遣的耶稣基督}。因此，圣父因圣子而受光荣，是通过祂被我们认识。这荣耀就在于：圣子降生成人，从圣父领受了掌管凡有血肉之人的权柄，并负责将永生赐给我们这些短暂、为身体所累的人。永生对于我们，不是工作的结果，而是内在权能的结果；不是通过新的创造，而只是通过对天主的认识，就取得了那永恒的荣耀。天主的荣耀没有增加；它不曾减少，所以也无法补充。但是，祂因圣子而在我们这些无知、流放、被玷污、活在绝望的死亡和无法无天的黑暗中的人眼中，受到光荣；光荣在于圣子因圣父所赐的掌管凡有血肉之人的权柄，将永生赐予我们。正是通过圣子的这工作，圣父受到光荣。因此，当圣子从圣父领受万有时，圣父光荣了祂；反之，当万有通过圣子而受造时，祂光荣了圣父。所给予的荣耀的回报在于：圣子所有的荣耀都是圣父的荣耀，因为祂所拥有的一切都是圣父的恩赐。因为执行使命者的荣耀归于授予使命者的荣耀，受生者的荣耀归于生者的荣耀。\par

14. 但是，永生在于什么呢？祂自己的话告诉我们：——\BibleQuote{使他们认识祢，惟一的真天主，和祢所派遣的耶稣基督}。这里有什么疑问或困难，或矛盾之处吗？认识真天主就是永生；但仅仅认识祂并不赐予永生。那么，祂又加了什么？\Quote{和祢所派遣的耶稣基督}。在\Quote{祢，惟一的真天主}中，圣子对圣父表达了应有的敬意；而通过加上\Quote{和祢所派遣的耶稣基督}，祂将自己与真天主性联系在一起。信徒在宣认时，并不在两者之间画线，因为他的永生希望在于两者，事实上，真天主与祂在信经中紧随其名的那一位是不可分离的。因此，当我们读到\BibleQuote{使他们认识祢，惟一的真天主，和祢所派遣的耶稣基督}时，这些派遣者和被派遣者的措辞，并非意在藉着某种区别或分辨的外表，传达圣父与圣子真天主性之间的差异，而是引导人虔诚地宣认他们为生者和受生者。\par

15. 因此，圣子完全而终末地光荣了圣父，在随后的话语中：\Quote{我在世上已光荣了你，完成了你所委托给我的工作。}所有对圣父的赞美都来自圣子，因为对圣子的每一赞美都是对圣父的赞美，因为圣子所成就的一切都是圣父所愿意的。天主之子生于人形；但天主的能力在于童贞生育。天主之子显现为人；但天主临在于祂的人性行动中。天主之子被钉在十字架上；但在十字架上，天主战胜了人类的死亡。基督，天主之子，死了；但所有血肉在基督内都得以复生。天主之子在地狱；但人被带回天堂。我们对基督这些工作的赞美有多大，我们带给那基督的天主性所从出之者的赞美也就有多大。这些是圣父在地上光荣圣子的方式；作为回报，圣子以大能的工作向外邦人的无知和世界的愚昧启示祂所从出者。这种给予和接受的荣光交换，并不意味天主性的增加，而是对那些曾活在无知于天主的人因获得认识而献上的赞美。事实上，那从祂而有万物的圣父，有什么不是丰盛地拥有呢？圣子缺少什么呢？在祂内岂不是天主性的圆满乐意居住吗？圣父在地上受光荣，因为所吩咐的工作完成了。\par

16. 接下来让我们看看圣子期望从圣父那里得到的这种光荣是什么；这样我们的阐述就完整了。随后的经文是：\Quote{我在世上已光荣了你，完成了你所委托给我的工作。现在，父啊，请以你自身光荣我，就是我在世界以先与你同有的光荣。我已将你的名显扬给人。}因此，圣子通过祂的工作使圣父受光荣，因为祂被承认为天主，是独生子的父，祂为了我们的救恩，愿意祂的圣子由童贞女所生而成为人；这位圣子的一生，在受难中完成，与童贞生育的卑谦一致。因此，因为天主之子，全美善，从永远生于天主性的圆满中，如今藉着降生成人成为人，并准备好受死，祂祈求能同天主一起受光荣，正如祂在地上光荣圣父一样；因为在那一刻，天主的能力在血肉中，在不认识祂的世界眼前正受到光荣。但祂所期望的与圣父同有的光荣是什么呢？当然就是祂在世界以先与圣父同有的光荣。祂曾拥有天主性的圆满；祂现在仍然拥有，因为祂是天主之子。但那位原是圣子的，也成为了人子，因为\BibleQuote{圣言成了血肉。}祂并没有失去祂原有的存在，而是成为了祂以前不是的；祂没有放弃自己的地位，却接纳了我们的地位；祂祈求祂所取的本性能被提升到祂从未放弃的光荣中。因此，既然圣子是圣言，圣言成了血肉，圣言是天主，在起初就与天主同在，而且在世界奠基之前圣子就是圣言；如今降生成人的这位圣子祈祷，使血肉能对圣父成为圣子之所是。祂祈祷生于时间的血肉能获得永恒光荣的光辉，使血肉的腐朽被吞没，转化为天主的能力和圣神的纯洁。这是祂向天主的祈祷，是圣子对圣父的宣认，是那将在审判之日所有人都看见被刺穿并带着十字架印记的血肉之身躯的恳求；是那在山上显其光荣的血肉，是那升天并坐在天主右边、保禄看见过、斯德望曾朝拜过的血肉。\par

17. \Emph{圣父}这个名号已显扬给人；问题是，圣父自己的名号是什么？然而，天主的名号从未不为人知。梅瑟曾在荆棘丛中听到过，\WorkTitle{创世记}在创造史的起头就宣告了，法律已宣扬，先知已颂扬，世界历史已使人类熟悉它；甚至外邦人也在虚假的面纱下朝拜它。人类从未对天主的名号一无所知。然而他们确实一无所知。因为除非人宣认祂为圣父，即独生子的父，并宣认圣子不是藉着\OriginalTerm{分割}{partition}、\OriginalTerm{延伸}{extension}或\OriginalTerm{由来}{procession}，而是由祂所生，作为圣父的儿子，不可言喻、不可理解地，并且保留了那天主性的圆满，祂就是从这圆满中并在这圆满中作为真实、无限、完美的天主而出生，没有人能认识天主。这就是\Emph{天主性的圆满}的意义。如果这些事中缺少任何一项，就不会有那乐意居住在祂内的圆满。这就是圣子的讯息，祂向无知的人类的启示。当人们认出祂是这么神圣的一位圣子的父时，圣父就藉着圣子受光荣。\par

18. 圣子为要使我们确信祂这神圣诞生之真实，便以祂的工程作为比方，好叫我们借着不可言喻之工程中所显示的不可言喻之能力，学习那不可言喻之诞生的教训。例如，当水变成酒，五个饼使五千人（妇女儿童除外）得以饱足，且碎块装满了十二筐时，我们看见了事实，却无法理解；事成了，却难倒我们的理性；过程无法追随，但结果是显而易见的。以吹毛求疵的精神闯入，而我们所探究之事又是我们无法深究到底的，这是愚昧。因为如同圣父因无始而不可言喻，同样圣子因是独生者而不可言喻，因为受生者乃是无始者的肖像。现在，我们必须借感官和言语来形成对肖像的概念；也必须是借同样的方式来形成对肖像所代表之物的观念。但在我们这里，我们这些只能处理可见可触之物的官能，却在追求不可见之物，努力把握不可捉摸之物。然而我们并不羞愧，不因自己的不敬而自责，反而怀疑并批评天主的奥迹和能力。祂如何是子？祂从何而来？圣父因祂的诞生失去了什么？祂是从圣父的哪一部分诞生的？我们这样问；然而一直有工程的证据摆在我们面前，向我们保证天主的行动并不受我们理解其方式能力的限制。\par

19. 你问，正如圣神所教导的，圣子是以何种方式诞生的？我要问你一个关于肉身的事。我不问祂如何由童贞女诞生；我只问：在她孕育祂的肉身使之预备诞生的过程中，她的肉身是否遭受了任何损失？确实，她不是以寻常方式怀了祂，也没有经受人类交合的羞耻而生祂：然而她生了祂，祂的整个人体完美无缺，而她自己完全无损。虔诚必定要求我们相信，我们在一个人身上借祂的能力成为可能的事，对天主来说也是可能的。\par

20. 但是，你——无论你是谁——想要探究那不可探究之事，并认真地对天主的奥迹和能力作出判断——我转向你求教，请你启发我，一个无知的、单纯相信天主所说一切的人，关于我即将提及的一件情况。我听从主的话，并且由于我相信所记载的，我确信在祂复活之后，祂曾多次以肉身向许多不信者的群众显现。至少，祂向多默这样做了，多默曾抗议说，除非他摸到祂的伤口，他决不相信。祂的话是：\BibleQuote{除非我看见祂手上的钉痕，把指头探入钉痕，并把手探入祂的肋旁，我决不相信。} \BibleRef{若 20:25} 主甚至屈就于我们软弱的理解力；为了满足不信之心的疑惑，祂行了一件祂无形能力的神迹。你，我这位天上之事的批评者，请若可能，解释祂的行动。门徒们在一间关着门的屋子里；自从主受难以后，他们秘密地聚集开会。主显现出来，通过应战多默的挑战来坚固他的信德；祂把自己的身体给他摸，把自己的伤口给他触。的确，那要被人认出受过伤的人，必须出示那受伤的身体。我问，主是以肉身穿过那间关着门的屋子的哪一部分墙壁进入的？宗徒仔细而精确地记述了当时的情形：\Quote{耶稣来了，门关着，站在中间。} 因为祂站在那里，是肉身的显现；没有任何欺骗的嫌疑。让你心灵的眼睛追随祂进入的路径；让你理智的目光伴随祂进入那紧闭的居所。墙壁没有破口，门没有打开；然而，看哪，祂站在中间，祂的权能无可阻挡。你是不可见之事的批评者；我要求你解释一件可见的事。一切都保持原样；没有身体能穿过木石的空隙。主的身体并不消散，然后消散后重新聚合；然而，那站在中间的是从哪里来的？你的感官和言语无法解释它；事实是确定的，但它超越了人类解释的领域。如果如你所说，我们关于神圣诞生的叙述是谎言，那么请证明这关于主进入的叙述是虚构。如果我们因为无法发现一件事是如何发生的就假定它没有发生，那我们就是把自己的理解力的限度当作现实的限度。但证据的确定性证明了我们反驳的虚假。主站在关着门的屋子里，在门徒中间；圣子是由圣父所生。不要因为你的小才智无法确定祂是如何来到那里的，就否认祂站在那里；不要基于感官和言语无法理解那诞生之超越奇迹，而拒绝相信天主——独生者、完满的圣子，出于无始者、完满的圣父。\par

21. 不仅如此，整个自然的结构都会支持我们，驳斥怀疑天主之工程和能力的渎神之罪。然而我们的不信甚至与明显的真理作对；我们在狂怒中甚至试图把天主从祂的宝座上拉下来。如果我们能，我们会用身体的力量攀上天空，会打乱日月星辰的秩序，会扰乱潮汐的涨落，会阻止江河的源头或使它们倒流，会动摇世界的根基，在我们反对天主父工程的极端不敬中。好在我们身体的局限将我们限制在更适度的范围内。当然，我们若能做到，会造成的危害是毫不掩饰的。在一个方面我们是自由的；因此，我们用亵渎的傲慢歪曲真理，把我们的武器转向天主的话语。\par

22. 圣子说：\Quote{父啊，我已将祢的名彰显于人。}这里有什么可谴责或暴怒的理由呢？难道你否认圣父吗？圣子首要的目的正是使我们认识圣父。但事实上，当你认为圣子并非由祂所生时，你就是在否认祂。然而，如果祂像其他受造物一样是天主任意创造的，祂又为何有子的名号？我能敬畏天主作为基督的创造者，如同祂作为宇宙的建立者；使祂成为那创造总领天使、天使、有形无形万物、天地及我们周围一切受造者之物的创造者，这原是祂权能的合宜运用。但主来所要成就的工作，不是要你认识天主作为万物的创造者的全能，而是要你认识祂是那对你说话的圣子的父。在天上，有祂以外的掌权者，强大而永恒；但只有一位独生子，祂与它们的区别不仅仅是能力的程度，而是它们都藉着祂受造。既然祂是真而唯一的圣子，我们就不该宣称祂是从无中造成的，而把祂当作私生子。你听见\Quote{子}的名字；要相信祂是子。你听见\Quote{父}的名字；要铭记祂是父。为什么要在这些名字周围布满怀疑、敌意和仇恨？天主的事有名字，真实地指示了实际；为什么要强行给它们明显的含义加上任意的解释？经上提到父与子；不要怀疑这些话的意思。圣子启示的目的和目标，就是要你认识圣父。为什么要挫伤先知们的辛劳、圣言的降生成人、童贞女的产痛、奇迹的效果、基督的十字架？这一切都为你而施行，这一切都赐予你，好使你能通过这一切认识父与子。你却用任意行动、创造或义子的理论取代了真理。请思考基督所进行的争战。祂这样描述：\Quote{父啊，我已将祢的名彰显于人。}祂并没有说：\Quote{祢创造了诸天的创造者}，或\Quote{祢造了全地的制造者。}祂说：\Quote{父啊，我已将祢的名彰显于人。}接受你救主的智慧之恩赐。要确信有一位生了子的父，一位被生的子；祂在祂本性的真理中由那实有的父所生。要记住，这启示不是将父显明为天主，而是将天主显明为父。\par

23. 你听见这些话：\BibleQuote{我与父是一体}。\BibleRef{若 10:30} 你为什么把圣子从圣父那里撕裂分开？他们是合一：一个绝对的存在，一切事物都完美地与那绝对的存在共融，祂由那绝对的存在而来。当你听见圣子说\BibleQuote{我与父是一体}时，要使你对事实的看法与位格相符；接受生者与受生者关于祂们自己的声明。要相信祂们是一体，正如祂们也是生者与受生者。为什么要否认共同的本性？为什么要攻击真正的天主性？你又听见：\BibleQuote{父在我内，我在父内}。\BibleRef{若 10:38} 这事关乎父与子，由圣子的工程所证明。我们的科学不能将身体包裹于身体内，也不能像水倒进酒里那样把一物注入另一物；但我们宣认在两者之内，有同等的权能和天主性的圆满。因为圣子从圣父领受了一切；祂是天主的肖像，祂本体的肖像。\BibleQuote{祂本体的肖像} \BibleRef{希 1:3} 这句话将基督与祂所来自的那位区分开来，但这只是为了确立祂们各自的实存，而不是教导本性的差异；而\Quote{父在子内，子在父内}的意思是：在两者之内有天主性的圆满。圣父不因圣子的实存而受损，圣子也不是圣父残缺的部分。肖像意味着原型；肖像是相对术语。任何事物除非源自天主，否则不能肖似天主；完美的肖像只能从其代表的事物中反射出来；精确的相似排除了任何差异因素的假设。不要扰乱这肖像；不要在真理显明无差异的地方制造分离，因为那说\BibleQuote{让我们按我们的肖像，照我们的模样造人} \BibleRef{创 1:26} 的那一位，藉着\Quote{我们的模样}这话启示了彼此相像的存在。不要触动，不要触摸，不要歪曲。要持守教导真理的名号，持守圣子关于祂自己的声明。我不愿你用自己发明的赞美来谄媚圣子；你若满足于圣经，就已足够。\par

24. 再者，我们不可如此盲目地信赖人的理智，以致以为我们对所思想的事物有完全的知识，或以为一旦形成一套匀称而一致的理论，终极问题便已解决。有限的心思不能领悟无限；一个依赖他者而存在的存有，不能达到对其造物主或自身的圆满知识，因为它对自我的意识受到其环境的渲染，并且设有其知觉不能逾越的界限。它的活动并非自因，而是出于造物主，一个依赖造物主的存有对其任何官能都没有完全的拥有，因为它的根源在于自身之外。因此，凭着一条不可抗拒的法则，这存有若说它对任何事有完全的知识，便是愚妄；它的能力有其无法改变的限度，只有当它受幻觉蒙蔽，以为自己的渺小界限与无限同界时，它才能空夸拥有智慧。因为它对于智慧是无能为力的，它的知识仅限于其知觉的范围，并分享其依赖存在的无能。因此，这有限本性自夸拥有唯独源于无限知识的智慧，这种伪装招致了宗徒的轻蔑和嘲笑，他称其智慧为愚妄。他说：\BibleQuote{因为基督派遣我，不是为施洗，而是为宣传福音，且不用智慧的言论，免得基督的十字架失去效力。因为十字架的道理，为丧亡的人是愚妄，为我们得救的人，却是天主的德能。因为经上记载：\InnerQuote{我要摧毁智者的智慧，废弃贤者的聪明。}智者在哪裡？经师在哪裡？今世的辩士在哪裡？天主岂不是使今世的智慧变成愚妄吗？因为世人没有凭自己的智慧认识天主，天主就乐意用愚妄的道理来拯救那些相信的人。原来犹太人要求神蹟，希腊人寻求智慧，而我们却宣讲被钉在十字架上的基督：这为犹太人固然是绊脚石，为外邦人是愚妄，但为那些蒙召的，不拘是犹太人或希腊人，基督却是天主的德能和天主的智慧。因为天主的愚妄总比人明智，天主的软弱也总比人刚强。}\BibleRef{格前 1:17} 因此，一切不信都是愚妄，因为它只取用自己有限知觉所能获得的智慧，并以此渺小的尺度衡量无限，从而断定它不能理解的事必然不可能。不信是无能参与辩论的结果。人们确信某事从未发生，因为他们已认定它不可能发生。\par

25. 因此，宗徒熟悉人的思想那种狭窄的假设，即它不知道的就不是真理，便说他不是用智慧的言論宣讲，免得他的宣讲落空。为避免被视为宣讲愚妄的人，他补充说：十字架的道理，为丧亡的人是愚妄。他知道不信的人认为唯一真正的知识就是构成他们自己智慧的知识，而由于他们的智慧只认识他们狭窄视野内的事，那另一种唯独神圣而圆满的智慧在他们看来便是愚妄。因此，他们的愚妄实际上在于他们误以为是智慧的软弱想象。由此，那些为丧亡之人是愚妄的事，对得救之人却是天主的德能；因为后者从不以自己的不足官能为尺度，而是将上天的全能归于天主的行动。天主按此意义废弃智者的智慧和贤者的聪明，正因为他们承认自己的愚妄，救恩才赐给相信的人。不信的人对一切超出他们认识范围的事都判以愚妄的裁决，而信的人将救恩奥秘的选择留给天主的德能和威严。天主的事中没有愚妄；愚妄在于那种要求天主以神蹟和智慧作为相信条件的人间智慧。犹太人要求神蹟是愚妄；他们因长期熟悉法律而对天主的圣名有一定认识，但十字架的绊脚石使他们退避。希腊人寻求智慧是愚妄；他们带着外邦人的愚妄和人的哲学，追问天主为何被高举在十字架上。并且，由于顾念我们心思能力的软弱，这些事被隐藏在奥秘中，犹太人和希腊人的这种愚妄便变为不信；因为他们指责那些凭其本性无法理解的真理不值得合理相信。但是，因为世上的智慧如此愚妄——先前它藉着天主的智慧未能认识天主，即宇宙的光辉和祂为祂手工作品所计划的奇妙秩序并未教导它敬畏造物主——天主乐意藉着宣讲愚妄来拯救那些相信的人，即藉着十字架的信德使永生成为有死之人的分份；这样，人智慧的自信便蒙羞，救恩在人以为愚妄所在之处被找到。因为基督，对外邦人是愚妄，对犹太人是绊脚石，却是天主的德能和天主的智慧；因为在神性事物中人看来软弱愚妄的，在真智慧和能力上超越地上的思想和力量。\par

26. 因此，天主的行为不可凭人的官能来考察；创造主不可受祂亲手所造之物的判断。我们应披上愚拙，好得智慧；不是那种轻率结论的愚拙，而是谦逊地意识到自己软弱的愚拙，为叫天主大能的证据教导我们那些世俗哲学论证达不到的真理。因为当我们完全意识到自己的愚拙，并感受到自己理智的无助和贫乏时，借着上智的谋略，我们将被引入天主的智慧；不为无限威严和威力设限，也不把自然的主束缚在自然律之下；确信对我们而言，有关天主的唯一真信德，就是祂同时是创始者和见证者。\par

\SourceRecord{Translated by E.W. Watson and L. Pullan.；Nicene and Post-Nicene Fathers, Second Series；Vol. 9.；Edited by Philip Schaff and Henry Wace.；Buffalo, NY: Christian Literature Publishing Co.,；1899.；<http://www.newadvent.org/fathers/330203.htm>.}

% structure: p0001 source=h1 target=section reason=page-title

\section{论天主圣三（卷四）}

卷四．从某种意义上说，这一卷是论着的开始，后来有时也被引为第一卷。他在§1中说，前三卷是先前写成的。它们陈述了关于基督天主性的真理，并简要驳斥了各种异端。现在他开始对亚流主义发动主攻。首先（§2）他重申他的困难所在：人的语言和思想无法应对无限者。然后（§3）他讲述亚流人如何曲解基督的永恒子性。为抵御这种篡改真理的行为，教会采用了\OriginalTerm{同质}{Homoousion}一词（§§4-7）；希拉里解释并辩护了它的使用。在§8中，他通过收集亚流人歪曲为己用的圣经经文，表明这种定义是必要的，并在§9-10指出他们对这些经文的利用是不诚实的。在§11中，他确切告诉我们亚流教导是什么，并用他们自己的一个信条，即\OriginalTerm{亚流致亚历山大书信}{Epistola Arii ad Alexandrum}（§§12-13）将其摆出。在§14中，这一教义被谴责；它并没有解释，而是曲解。通过梅瑟所宣告的\BibleQuote{以色列，你要听，上主你的天主是唯一的上主}\BibleRef{申 6:4}，亚流人以此为依据，但这只揭示了真理的一个方面（§15）。它并未穷尽真理；因为天主在创造的历史（§§16-22）、在亚巴郎的生命（§§23-31）以及在梅瑟的生命（§§32-34）中都被描绘为并非一位孤独的位格。而这也是先知们的教导，如选自依撒意亚、欧瑟亚和耶肋米亚的经文所示（§§35-42）。所有这些收集的证据都表明，在天主性中有圣父和圣子，并且圣子是天主。\par

1.\ 我相信，先前写成的这部论着的前几卷，已经提供了无可辩驳的证明：我们持守并宣认由福音作者和宗徒所教导的、对圣父、圣子和圣神的信仰，并且我们与异端者之间不可能有任何交往，因为他们无条件地、无理地、鲁莽地否认我们的主耶稣基督的天主性。然而，仍有一些要点，我觉得有义务纳于本卷和随后各卷中，以便通过揭露他们的每一个谎言和亵渎，使我们信仰的确信更为确定。因此，我们将首先探究他们教导的危险何在，这种不敬所带来的风险；其次，他们持守什么原则，他们提出什么论据来反对我们所坚持的宗徒信仰，以及他们用什么言语诡计来欺骗听众的坦诚；最后，他们用什么注释方法剥夺圣经话语的力量和意义。\par

2.\ 我们很清楚，无论是人的言语还是人性的类比，都不能使我们完全洞察天主的事。不可言说者不能屈从于定义的界限和限制；属灵的事物与每一类或每一个有形之物的实例都是截然不同的。然而，由于我们的主题是关于天上本性的事，我们必须运用普通的本性和普通的言语，作为表达我们心智所领会之物的手段；这手段无疑不配天主的尊威，但因我们理智的软弱而被迫使用，它只能用我们自身的情况和我们自己的话语，向他人传达我们的感知和结论。这一真理在本书中已经强调过，但在此重复，是为了使我们在引用来自人事的任何类比时，不致使我们认为天主类似有形的本性，或将属灵的存有与我们有情的自身相比，而是被视为将可见事物的外在形象作为不可见事物内在含义的线索。\par

3.\ 因为异端者说，基督不是出于天主，即圣子不是由圣父所生，并且祂不是出于本性而是出于任命而为天主；换句话说，祂接受了由赐名构成的义子名分，是如同许多人成为天主的儿子那样的天主子；再者，基督的尊威是天主广施恩惠的证据，祂是诸神中的神，如同有许多神一样；虽然他们承认，在祂被立为子和被命名为天主时，所显示的爱比在其他情况下更为丰盛，祂的义子在时间顺序上是首位的，祂比其他被立的儿子更大，并因祂受造时所伴随的更大荣光而在受造物中居于首位。有些人为了宣认天主的全能，又补充说，祂是按天主的肖像受造的，并且祂像其他受造物一样，是从无中被提升起来，成为永远造物主的肖像，因天主的一声命令，藉天主的权能从无变为有，而天主能从无中塑造祂自己的肖像。\par

4. 此外，他们利用自身对历史事实的了解，即早期的主教们曾教导圣父与圣子同体，便以狡猾的借口说这是一个异端观念，从而颠覆真理。他们声称，这个术语\Quote{同体}——希腊文为\OriginalTerm{同体}{homoousion}——用以意指并表达圣父就是圣子；也就是说，圣父从无限中延伸出来，进入童贞玛利亚，从她取了身体，并在所取的身体中给自己取了\Quote{圣子}的名号。这是他们对\OriginalTerm{同体}{homoousion}的第一个谎言。他们的第二个谎言是，这个词\OriginalTerm{同体}{homoousion}暗示圣父与圣子共同分有某种先于两者、且不同于两者之物，并且某种想象的实体或\OriginalTerm{实体}{ousia}，先于一切物质存在，自从太古以来就已存在，并被分割并完全分给两者；他们声称，这证明两者各自具有先于自身的本性，且两者在物质上同一。因此，他们公开谴责对\OriginalTerm{同体}{homoousion}的宣认，理由是这个词没有区分圣父与圣子，并使得圣父在时间上后于那与圣子共有的物质。他们还针对\OriginalTerm{同体}{homoousion}一词捏造了第三种反对意见，即按其解释，该词的含义是圣子源自圣父实体的一种分割，好像一个物体被切成两半，圣子就是被割下的一部分。他们说，\Quote{同体}的意思是，从整体切下的部分继续分享被切离之物的本性；但天主是不可受难的，不能被分割，因为若祂必须遭受因分割而缩小，祂就会发生变化，并且如果祂完美的实体离开祂而寓于切下的部分中，祂就会变得不完美。\par

5. 他们还认为，他们有一个对先知、福音作者和宗徒都具有概括性驳斥的说法，即断言圣子是在时间之内受生的。他们指责我们不合逻辑，因为我们说圣子从永远就已存在；既然他们否定祂永恒的可能性，他们就不得不相信祂是在某个时间点受生的。因为如果祂并非总是存在，那么就有祂不在的时候；如果存在一个祂不在的时候，那么时间就先于祂。那位并非永远存在者，是在时间之内开始存在的；而那位不受时间限制者，必然是永恒的。他们否定圣子永恒的理由是，祂的永远存在与对祂受生的信仰相矛盾；仿佛我们宣认祂永恒存在，就使得祂的受生不可能似的。\par

6. 何等愚蠢而无敬畏之心的恐惧！何等不敬的为天主担忧！他们自称在\OriginalTerm{同体}{homoousion}一词和圣子永恒的断言中发现的含义，是被教会所憎恶、拒绝和谴责的。教会宣认一位天主，万物由祂而来；宣认一位主耶稣基督，万物藉祂而有；一位从祂而来，一位藉祂而有；一位是万有的根源，一位是万物藉以受造的媒介。在万物由祂而来的那位身上，她认出无始的尊威；在万物藉祂而成的另一位身上，她认出与祂根源同等的权能；因为两者在创世工程和统治受造物方面同为至高。在圣神内，她认出天主是神，不可受难、不可分割，因为她从主那里学到，神没有血肉和骨头\BibleRef{路 24:39}；这警告她免于以为天主既然是神，就可能承受肉身的痛苦和损失。她宣认一位永远不生的天主；也宣认一位天主的独生子。她宣认圣父是永恒的、无始的；也宣认圣子的开端是从永远而来的。不是说祂没有开端，而是说祂是那位无始之父的圣子；不是说祂自生，而是说祂来自那位从永远非受生者；从永远而生，即是从圣父的永恒中领受祂的诞生。因此，我们的信仰摆脱了异端邪说的猜测；它以固定而公开的词语表达，尽管至今尚未提出对我们宣认的理性辩护。然而，为免对教父们使用\OriginalTerm{同体}{homoousion}一词的含义以及我们对圣子永恒的宣认有任何疑虑，我已列出证据，使我们确信圣子永远存留于祂由圣父所生的那实体中，而且圣父生圣子并未减损圣父所存留的实体丝毫；那些圣人，受天主教导的启示，当他们说圣子与圣父是\OriginalTerm{同体}{homoousios}时，并未指向我提到的那些缺陷或缺失。我的目的是要抵消一种印象，即这个\OriginalTerm{实体}{ousia}，这个声称祂与圣父同体的断言，是对独生子受生的否定。\par

7. 为确信这两个短语的必要性——它们被采纳和使用作为对抗当时异端暴民的最佳保障——我认为最好是回应我们现今异端者的固执不信，并藉着福音作者和宗徒们的见证驳斥他们虚妄而有害的教导。他们自诩可以为自己的每一个命题提供证据；事实上，他们已为每个命题附上圣经中的某些章节；这些章节被如此粗暴地曲解，以至于唯有不识字的人才会被曲解的才智在其解释上所掩盖的似是而非的真理所诱骗。\par

8. 因为他们企图通过只赞美圣父的天主性，来剥夺圣子的天主性，他们辩称经上记载：\Quote{以色列，请听！上主我们的天主是唯一的上主}\BibleRef{申 6:4}，并且主在回答询问法律最大诫命的法学士时也重复了这句话：\Quote{以色列，请听！上主我们的天主是唯一的上主}\BibleRef{谷 12:29}。他们又说保禄宣告说：\Quote{因为天主只有一个，在天主与人之间的中保也只有一个}\BibleRef{弟前 2:5}。此外，他们坚持只有天主是智慧的，为的是不给圣子留下任何智慧，他们依据宗徒的话说：\Quote{那能坚固你们按我的福音和耶稣基督的宣讲，并按照永恒的天主——独一智慧的天主——的命令，藉着先知的经卷所启示的奥秘，如今彰显于万国，为使人信服；愿荣耀藉着耶稣基督归于祂，直到永远。}\BibleRef{罗 16:25} 他们还辩称只有祂是真实的，因为依撒意亚说：\Quote{他们必祝福你，真实的天主}\BibleRef{依 65:16}，并且主自己在福音中作证说：\Quote{永生就是：认识祢，唯一的真实天主，和祢所派遣的耶稣基督}\BibleRef{若 17:3}。他们又推理说只有祂是善的，为的是不给圣子留下任何善，因为经上藉着祂说：\Quote{除了天主一个外，没有善的}\BibleRef{谷 10:18}；而且只有祂有大能，因为保禄说：\Quote{在预定的时期，祂要显示给我们——那蒙福的、唯一的全能者，万王之王，万主之主}\BibleRef{弟前 6:15}。并且，他们确信在父那里没有改变也没有转变，因为祂藉着先知说：\Quote{我是上主你们的天主，我决不改变}\BibleRef{拉 3:6}，以及雅各伯宗徒说：\Quote{在祂没有改变}；他们也确信祂是正义的审判者，因为经上记载：\Quote{天主是正义的审判者，坚强而忍耐}；祂眷顾一切，因为主论到飞鸟时说：\Quote{你们的圣父尚且养活它们}\BibleRef{玛 6:26}，以及：\Quote{两只麻雀不是卖一个铜钱吗？但没有一只掉在地上，不经过你们圣父的允许；就是你们的头发也都被数过了}。他们说父预知一切，正如有福的苏撒纳所说：\Quote{永恒的天主，祢知道隐藏的事，在一切事发生以前，祢就预知}；祂是不可理解的，正如经上记载：\Quote{天是我的宝座，地是我的脚凳。你们要为我建造什么样的殿宇？或者哪里是我安息的地方？这一切都是我的手所造，这一切都是我的}；祂包含万有，正如保禄作证：\Quote{因为我们在祂内生活、行动、存在}\BibleRef{宗 17:28}，以及圣咏作者说：\Quote{我往何处去，才能躲避祢的神？我往何处逃，才能躲避祢的面？我若升到天上，祢在那里；我若下到阴府，祢也在那里。我若乘晨曦之翼，远居海极，在那里，祢的手仍引导我，祢的右手仍扶持我}；祂是无形的，因为经上记载：\Quote{天主是神，朝拜祂的人当以心神和真理朝拜祂}\BibleRef{若 4:24}；祂是不死不灭、不可见的，正如保禄说：\Quote{那唯一具有不死不灭者，住于不可接近的光中，没有人看见过，也不能看见}\BibleRef{弟前 6:16}，以及福音作者说：\Quote{从来没有人见过天主，只有那在圣父怀里的独生子，将祂彰显出来}\BibleRef{若 1:18}；只有祂永远保持无生而存在，因为经上记载：\Quote{我是自有者}，以及\Quote{你要这样对以色列子民说：\InnerQuote{那自有者}派遣我到你们这里来}\BibleRef{出 3:14}，又藉着耶肋米亚说：\Quote{上主！祢是主。}\par

9. 谁能看不出这些陈述充满欺诈和谬误？尽管论题被巧妙地混淆、经文被组合，但恶意与愚蠢却深深地铭刻在这番劳心费力的狡猾而笨拙的努力上。例如，在他们的信仰要理中，他们承认只有圣父是无生的；好像我们这边有人会认为，那位生了万有借以受生的祂的天主，其存有来自任何外在源头似的。事实上，祂被称为\Quote{圣父}这一事实本身就揭示祂是圣子存在的原因。\Quote{圣父}这个名称并不暗示祂自己来自另一位，却清楚地告诉我们圣子是由谁而生。因此，让我们将圣父自己的特殊而不可分享的特性归给祂，承认在祂内住有那无始的全能之永恒能力。我确信，无人会怀疑，他们在承认天主圣父时特别强调某些特性作为祂所专有且不可让与的理由，是为了让祂独自保有这些特性。因为当他们说只有祂是真实的、只有祂是正义的、只有祂是智慧的、只有祂是不可见的、只有祂是善的、只有祂是大能的、只有祂是不死不灭的，他们是在用这个\Quote{只有}一词筑起一道屏障，将圣子排除在这些特性的分享之外。他们说，那唯一的，在祂的特性中没有伙伴。但如果我们认为这些特性只存在于圣父内，而不也存在于圣子内，那么我们就必须相信天主圣子既无真理也无智慧；祂是一个由可见的、物质的元素构成的有形之体，性情不善、软弱无力、缺少不死不灭；因为我们把圣父当作这些特性的唯一拥有者，从而将圣子排除在所有这些特性之外。\par

10. 我们如今有意论述那属于天主独生子之至极尊威与圆满天主性，并不担心读者以为，以丰沛言辞论及圣子，会减损天主圣父的光荣，仿佛每一归于圣子的赞美都已先从他那里被剥夺。因为相反，圣子的尊威正是圣父的光荣；那生出如此可敬之子的源头，必是光荣的。圣子除藉出生外一无所有；圣父分享因那出生之权利所接受的一切崇敬。如此，我们损减圣父尊荣的暗示便被平息，因为我们将教导的一切属于圣子的光荣，都将反射回去，归光荣于那生了如此伟大圣子的圣父。\par

11. 如今我们已揭露了他们以尊崇圣父为掩饰而贬低圣子的计谋，下一步是聆听他们关于圣子信仰的确切措辞。因为，我们必须逐一回击他们的每一项主张，并藉圣经的证据显明他们教义的不虔敬，所以必须将他们论圣父的话与他们对圣子所记录的决定附上，以便通过比较他们对圣父的宣认与对圣子的宣认，在我们解决问题时遵循统一的秩序。他们陈述其结论说：圣子并非从预先存在的质料而来，因为万物都是藉他而受造；也不是从天主而生，因为没有任何事物能从天主那里被抽出；而是从无中造成，也就是说，他是天主的完美受造物，虽然与其他受造物不同。他们论证他是受造物，因为经上记载：\BibleQuote{上主在祂道路的起点就创造了我}\BibleRef{箴 8:22}；他虽是完美的天主工程，却与其他工程不同，他们首先引用\BibleBook{保禄}{致希伯来人书}的话：\BibleQuote{他是比众天使更为优越的，如同他承受的名字比他们更尊贵}\BibleRef{希 1:4}，以及：\BibleQuote{因此，诸位圣洁的弟兄，同享天召的人，应当注目我们宣认的宗徒和司祭长耶稣基督，他对那造他的保持忠信}。为了贬低圣子的权能、尊威和天主性，他们主要依靠他自己的话：\BibleQuote{父比我大}\BibleRef{若 14:28}。但他们承认他不属于普通受造物之列，根据\BibleQuote{万物是藉着他而造成的}\BibleRef{若 1:3}。于是他们将全部亵渎的教义总结为以下这些话：——\par

12. 我们宣认独一天主，唯一非受造，唯一永恒，唯一无始，唯一真实，唯一拥有不死，唯一美善，唯一全能，创造者、安排者和管理者，不变不易，公义良善，属法律、先知和新约。我们相信这位天主在万世之前生了独生子，藉着他创造了世界和万物；他生他不是虚像，而是真实，按他自己的意愿，使他不变不易，是天主的完美受造物，却不像他的其他受造物；是他的工程，却不是其他的工程；并非如华伦提努所主张的，圣子是圣父的一种发展；也不是如摩尼对圣子所宣称的，是圣父的一个\OriginalTerm{同质}{consubstantial}部分；也并非如撒伯流所说，从一造出两个，同时是子和父；也非如希拉卡所认为，是光中之光，或双焰的灯；也非如他先前存在，后来才新生或被再造为子，这种观念常被你，蒙福的教宗，在教会和弟兄的聚会中公开谴责。但正如我们所肯定的，我们相信他是在时间与世界之前，由天主的意愿受造，并从圣父获得生命与存在，圣父使他分享自己光荣的完美。因为，当圣父将万有的继承赐给他时，并未使自己失去那无始的属性，他本是万有的源头。\par

13. \Quote{所以有三个位格：圣父、圣子和圣神。天主，就其自身而言，是万有的原因，完全无始，与一切分离；而圣子，由圣父在时间之外发出，在万世之前被创造和建立，生之前并不存在，但于时间之外、万世之前诞生，作为唯一圣父的唯一圣子而存在。因为他既非永恒，也非与圣父\OriginalTerm{同永恒}{co-eternal}，也非与圣父\OriginalTerm{同非受造}{co-uncreate}，也不与圣父并存，如有些人所说，他们假设两个无始的原则。而天主在万有之前，作为不可分割者和万有的起始。因此他也在圣子之前，正如我们从你公开的宣讲中所学到的。那么，既然他的存在、光荣的完美、生命，以及万有的托管都来自天主，故此天主是他的源头，并统治他，作为他的天主，因他在他之前。至于诸如\Emph{\Quote{从他}}、\Emph{\Quote{从母腹}}、\Emph{\Quote{我从父出来，来到世上}}这样的说法，若被理解为表示圣父延伸一部分，也就是那独一实体的某种发展，那么圣父将具有复合性、可分性、可变性和形体性，按照他们的说法；因此，就他们的话而言，无形体的天主将被归于物质的特性。}\par

这就是他们的错误，这就是他们那毒害人的教导；为支持它，他们借用圣经的话，歪曲其含义，并利用人的无知作为机会使他们的谎言获得信任。然而，我们确实必须通过天主自己的这些话来理解天主的事。因为人的软弱无法凭自身的力量获得天上知识的认识；处理身体事物的官能无法形成对不可见世界的概念。既非我们受造的身体实体，也非天主为日常生活的目的所赐予的理性，能够确定并断言天主的本性和工作。我们的才智不能上升到天上知识的层次，我们的感知能力缺乏领会那无限大能的力量。我们必须相信天主关于祂自己的话，并谦卑地接受祂恩赐的洞见。我们必须在拒绝祂的见证（如外邦人所为）与相信祂的本来面目之间做出选择，而后者唯一可能的方式，就是按照祂向我们显示自己的面向来思想祂。因此，让私人的判断终止，让人的理性止步于神圣设立的界限。怀着这种精神，我们避免一切关于天主的亵渎和鲁莽的断言，并紧紧依附于启示的字句。我们探究的每一点都应按我们主题——祂的教导——来考量；绝不可将孤立短语拼凑起来，压制上下文，以欺骗和误导缺乏经验的听众。词语的含义应通过考虑它们被说出的情境来确定；词语必须由情境来解释，而不是强迫情境符合词语。无论如何，我们将完全处理我们的主题；我们将陈述词语被说出时的情境以及词语的真实含义。每一点都将按顺序考虑。\par

他们的出发点是：他们说，我们只承认一位天主，因为梅瑟说：\Quote{\BibleQuote{以色列，请听！上主你的天主是唯一的上主}}。\BibleRef{申 6:4} 但这是一条任何人曾敢于怀疑的真理吗？或者曾有哪位信徒承认过别的，即那一位天主，万有出自祂，那一位没有出生的大能，祂就是那\OriginalTerm{无源之能}{unoriginated Power}？但是，天主的唯一性这一事实并不提供否认祂圣子天主性的机会。因为梅瑟，或者不如说天主通过梅瑟，将这条作为首要诫命赐给那人民——他们在埃及和旷野都专注于偶像和假神的崇拜——他们必须相信一位天主。这条诫命有其真理和理由，因为万有出自的天主是唯一的。但是，让我们看看这位梅瑟是否没有承认那万有藉以存在的那一位也是天主。如果祂的圣子分享\OriginalTerm{天主性}{Godhead}，天主并未被剥夺，祂仍然是天主。因为情形是天主来自天主，一位来自一位，天主是唯一的一，因为天主来自祂。并且相反地，圣子并不因为天主圣父是唯一而成为较少的天主，因为祂是天主的\OriginalTerm{独生子}{Only-begotten Son}；并非永恒地无生，以致剥夺圣父的唯一性，也不与天主不同，因为祂是从祂所生。我们不可怀疑：祂是天主，是由于那天主所生的事实，这向我们相信的人证明天主是唯一；然而，让我们看看那向以色列宣布\Quote{\BibleQuote{上主你的天主是唯一}}的梅瑟，是否也宣告了圣子的天主性。为了证实我们对主耶稣基督天主性的宣认，我们必须使用那同一见证者的证据——异端者正是依赖他来宣认一位天主，他们想象这必然涉及否认圣子的天主性。\par

16. 因此，既然宗徒的话：\BibleQuote{一位天主圣父，万物都出自祂；一位主耶稣基督，万物都藉着祂} \BibleRef{格前 8:6}，是关于天主的准确而完整的宣认，让我们看看梅瑟关于世界之初说了什么。他的话是：\BibleQuote{天主说：在水中间要有穹苍，好将水与水分开。事就这样成了，天主就造了穹苍，并将水从中间分开。} \BibleRef{创 1:6} 那么，你在此便有了\Quote{万物出自祂}的天主与\Quote{万物藉着祂}的天主。你若否认这一点，就必须告诉我们天主的创造之工是藉着谁完成的，否则就得指出一种尚未受造之物中的服从，即当天主说\InnerQuote{要有穹苍}时，这种服从促使穹苍自行建立。这样的推测与圣经的明确含义不符。因为万有，如先知所说\BibleRef{加下 7:28}，是从无中被造的；不是现有之物的变形，而是对非存在者创造为完美的形态。是藉着谁呢？请听圣史的话：\BibleQuote{万有是藉着祂造成的}。你若问这是谁，同一位圣史会告诉你：\BibleQuote{在起初已有圣言，圣言与天主同在，圣言就是天主。祂在起初就与天主同在。万有是藉着祂造成的} \BibleRef{若 1:1}。你若想要反驳那是圣父说\InnerQuote{要有穹苍}的观点，先知会回答你：\BibleQuote{祂一发言，万有形成；祂一下令，万物受造}。所记载的话\InnerQuote{要有穹苍}向我们启示圣父说了话。但在随后的\InnerQuote{事就这样成了}中，在记载天主行了这事时，我们必须认出行事者的位格。\BibleQuote{祂一发言，万有形成}；圣经没有说祂愿意了，就做了。\BibleQuote{祂一下令，万物受造}；你注意到它没有说它们因祂的喜悦而存在。那样的话，在天主与等待受造的世界之间就没有中保的职分了。万物所出自的天主发出创造的命令，而万物所藉着成为的天主执行命令。在同一个名号下，我们宣认发令者与执行命令者。你若是胆敢否认\InnerQuote{天主造了}是指圣子说的，那么你如何解释\BibleQuote{万有是藉着祂造成的}？或宗徒的话\BibleQuote{一位主耶稣基督，万物都藉着祂}？或\BibleQuote{祂一发言，万有形成}？如果这些受灵感的话能说服你固执的心，你就不会再认为\BibleQuote{以色列，请听！主你的天主是唯一的主}这句话拒绝了圣子的天主性，因为在世界奠基之时，说话的那位就宣告祂的圣子也是天主。但让我们看看从这发令的天主与执行的天主之间的区分中，我们能获得什么更大的益处。因为即使在我们的自然理性中，也难以认同在\InnerQuote{祂下令，它们形成}这些话里，所指的是一位单一而孤立的一位，但为了避免一切疑惑，我们必须阐述世界创造之后所发生的事。\par

17. 当世界完成，要造它的居民时，关于他所说的话是：\BibleQuote{让我们按我们的肖像，照我们的模样造人} \BibleRef{创 1:26}。我问你，你认为天主是对自己说这些话吗？难道不明显祂不是在对自己说话，而是对另一位说话？你若是回答说祂是独自的，那么祂亲口反驳你，因为祂说：\BibleQuote{让我们按我们的肖像，照我们的模样造人}。天主藉着立法者以我们能理解的方式向我们说话；也就是说，祂藉着语言——祂乐意赐予我们的官能——使我们认识祂的行动。的确，在\InnerQuote{天主说：要有穹苍}以及接着的\InnerQuote{天主造了穹苍}这些话中，已有了圣子——万有藉着祂而造——的迹象；但恐怕我们以为天主这些话是多余且无意义的，以为祂向自己发出创造的命令，又以自己服从——因为对一个孤独的天主来说，还有什么比向自己发出口头命令更荒谬的，当时只需运用祂的意志即可？——祂决定给我们更完美的保证，即这些话是指祂以外的另一位。当祂说\InnerQuote{让我们按我们的肖像，照我们的模样造人}时，祂提及一位伙伴，推翻了祂孤独的理论。因为一个孤立的存在不能做自己的伙伴；再者，\InnerQuote{让我们造}与孤独不相容，而\InnerQuote{我们的}只能用于同伴。\Emph{我们}和\Emph{我们的}这两个词都与孤独的天主对自己说话的概念不相容，也同样与对一位与说话者毫无共同之处的陌生者说话的概念不相容。你若解释这段经文是说祂是孤独的，我问你，你是否认为祂是在对自己说话？你若不明白祂是在对自己说话，你怎能假定祂是孤独的？如果祂是孤独的，我们应会看到祂被描述为孤独的；如果祂有同伴，则不是孤独的。\Emph{我}和\Emph{我的}会描述前一种状态；后一种则由\Emph{我们}和\Emph{我们的}表明。\par

18. 因此，当我们读到\Emph{\Quote{让我们照我们的肖像，按我们的模样造人}}时，这两个词\Emph{\Quote{我们}}和\Emph{\Quote{我们的}}揭示了既不是一位孤立的天主，也不是一位天主在两个不同的位格中；我们的宣信必须与这两个真理相符。因为\Emph{\Quote{我们的肖像}}——不是\Emph{\Quote{我们的肖像们}}——证明两者拥有同一性体。但仅凭词语论证是不够的，除非其结果被事实的证据所证实；因此经上记载：\BibleQuote{天主造了人；按天主的肖像造了他}\BibleRef{创 1:27}。如果祂所说的话，我问，是一个孤立天主的自言自语，那么我们对这最后一段陈述该作何解释？因为我在其中看到三重指涉：造物者、被造之物和肖像。被造之物是人；天主造了他，并照天主的肖像造了他。如果\BibleBook{}{创世记}是在谈论一个孤立的天主，它肯定会说\Emph{\Quote{照祂自己的肖像造了他}}。但由于该书预示了福音的奥迹，它没有谈论两个天主，而是谈论天主和天主，因为它说人是通过天主而照天主的肖像造的。因此我们发现天主照一个与祂自己和天主共有的肖像和模样造了人；提及行动者阻止我们假设祂是孤立的；而作品，按照两者的肖像和模样完成，证明一位和另一位的神性之间没有本质差异。\par

19. 提出进一步的论点似乎是在浪费时间，因为关于天主的真理并不因重复而增强力量；单一陈述足以确立它们。然而，了解所有已经启示的内容对我们有益，因为虽然我们不为圣经的话语负责，但我们却要为赋予它们的意义交账。天主给诺厄的众多诫命之一是\Quote{\Emph{凡流人血的，他的血也必被人流，因为天主照自己的肖像造了人}}。这里再次有肖像、受造物和造物主之间的区别。天主见证祂照天主的肖像造了人。当祂将要造人时，因为祂是谈及自己，却不对自己说，天主说\Emph{\Quote{照我们的肖像}}；而人造成后，又说\Emph{\Quote{天主照天主的肖像造了人}}。如果祂对自己说\Emph{\Quote{我照我的肖像造了人}}，在语言上并无不当，因为祂已经通过\Emph{\Quote{让我们照我们的肖像造人}}表明两位在性体上是合一的。但为了更彻底地消除一切关于天主是否是孤立存在的疑虑，当祂造人时，我们被告知，祂造了他\Emph{\Quote{照天主的肖像}}。\par

20. 若你仍想主张天主圣父在独处中对祂自己说了这些话，我可以与你一同承认一种可能性，即祂可能在独处中自言自语，仿佛在与同伴交谈，并且可以相信祂希望\Emph{\Quote{我照天主的肖像造了人}}等同于\Emph{\Quote{我照我自己的肖像造了人}}。但你自己信德的宣认将会反驳你。因为你已宣认：万物来自圣父，但都是通过圣子；而\Emph{\Quote{让我们造人}}这句话表明，那说话的就是万物之源；而\Emph{\Quote{天主照天主的肖像造了他}}则清楚指向那通过祂完成工作的那一位。\par

21. 再者，为使一切自欺成为不可能，那你自己已宣认为基督的智慧将以这些话面对你：\BibleQuote{当祂确立诸天以下的泉水，当祂坚固大地的根基时，我已在祂身边，布置一切。我是祂所喜悦的。此外，我日日在他面前欢跃，在祂因祂所造的世界和世人而欢欣的时候}\BibleRef{箴 8:28}。一切困难都被解除；错误本身也必须承认真理。有一个智慧与天主同在，在万世之前受生；不仅与祂同在，而且布置，因为\Emph{\Quote{她与他同在，布置一切}}。请注意这布置或安排的工作。圣父通过祂的命令是原因；圣子通过执行所受命之事而布置。两位之间的区别由各自分派的工作标明。当经上说\Emph{\Quote{让我们造}}时，受造与命令的话语认同；但当记载\Emph{\Quote{我与他同在，布置一切}}时，天主揭示祂没有孤立地做这工作。因为祂在祂面前欢跃，祂告诉我们祂也欢喜回报；\BibleQuote{此外，我日日在他面前欢跃，在祂因祂所造的世界和世人而欢欣的时候}\BibleRef{箴 8:28}。智慧已教导我们她喜乐的原因。她因圣父的喜乐而喜乐，圣父为世界的完成和世人而喜乐。因为经上记载\Emph{\Quote{天主看了认为好}}。她喜乐，因为天主对祂通过她、按祂命令所造的工作感到满意。她承认她的喜乐来自圣父因世界的完成和世人的欢欣；世人是由于在一个人亚当里整个人类开始了其进程。因此，在创世中，并非孤立圣父的独白；祂的智慧是祂工作中的伙伴，并在它们联合劳动结束时与祂一同喜乐。\par

22. 我知道对这些话语的全面解释涉及许多重大问题的讨论。我并不回避，而是暂时推迟，将它们留待以后阶段的探讨。目前我专心对付亵渎者的信条，或者更确切地说，不信的条款，即梅瑟宣告天主的孤独。我们不会忘记，这个断言在某种意义上是真实的，即有一个天主，万物来自祂；但我们也没有忘记，这个真理不能成为否认圣子神性的借口，因为梅瑟在其著作的整个过程中清楚地表明天主与天主的存在。我们必须考察关于天主的拣选和颁赐法律的历史如何宣告天主与天主是并列的。\par

23. 天主曾多次与亚巴郎说话，撒辣对哈加尔动了怒，因她主母不生育，而婢女却怀了儿子，便起了嫉妒之心。当哈加尔离开她面前时，圣神就论她这样说：\BibleQuote{上主的天使对哈加尔说：回到你主母那里，屈服在她手下。上主的天使又对她说：我要使你的后裔极其繁衍，多到不可胜数。}又说：\BibleQuote{她称那与她说话的上主之名为：你是看顾我的天主。}是上主的天使在说话，所说的内容却远超一位使者（因为使者一词的本义即如此）所能做到的。祂说：\BibleQuote{我要使你的后裔极其繁衍，多到不可胜数。}使万民繁衍的能力不属于天使的职务。然而，对于这位被称为上主的天使、却说出唯独属于天主的话语者，圣经如何说呢？\BibleQuote{她称那与她说话的上主之名为：你是看顾我的天主。}起初祂是上主的天使；然后祂是上主，因为\BibleQuote{她称上主之名；}第三，祂是天主，因为\BibleQuote{你是看顾我的天主。}那被称为上主的天使，也是上主和天主。天主子依照先知的话，也是\BibleQuote{伟大的谋士天使}。为清楚区分位格，祂被称为上主的天使；那出自天主的天主也是上主的天使，但为让祂享有应有的尊荣，祂也被称作上主和天主。\par

24. 在此段经文中，同一天主首先是上主的天使，然后相继是上主和天主。但对亚巴郎，祂只是天主。因为当位格的区分首先被标示出来，作为防止人错觉天主是独一存有的保障之后，祂真实而无保留的名号便可安全地宣告了。经上这样记载说：\BibleQuote{天主对亚巴郎说：看，你的妻子撒辣必给你生一个儿子，你要给他起名叫依撒格；我要与他立定永远的盟约，与他以后的苗裔。至于依市玛耳，看，我听了你的祈求，我已祝福了他，要使他极其繁衍；他要生十二个族长，我要使他成为一个大民族。} \BibleRef{创 17:19} 难道还能怀疑那先前被称为上主的天使，后来在这里被说是天主吗？在这两处，祂都在谈论依市玛耳；在两者中，都是同一位要使他繁衍。为避免我们以为这是另一位对哈加尔说话者不同的发言者，天主的话明确地证实了同一性，说：\BibleQuote{我已祝福了他，要使他繁衍。}这祝福是重复先前的事件，因为哈加尔已经被告知；这繁衍是应许将来要成就的，因为这是天主对亚巴郎首次提及依市玛耳。现在是天主对亚巴郎说话；对哈加尔说话的是上主的天使。因此，天主和上主的天使是同一的；那为天主的天使也是天主子、天主。祂被称为天使，因为祂是\BibleQuote{伟大的谋士天使}；但后来祂被说成是天主，以免我们以为那为天主者仅仅是一位天使。现在让我们按顺序重述这些事实。上主的天使对哈加尔说话；祂也以天主的身份对亚巴郎说话。同一位发言者向两者说话。祝福赐给了依市玛耳，应许他要成为一个大民族。\par

25. 在另一个事例中，圣经通过亚巴郎启示出说话的是天主。他得到进一步的应许，要生一个儿子依撒格。随后有三人向他显现。亚巴郎虽看见三人，却只朝拜一位，并承认祂为主。圣经说三人站在他面前，但他清楚地知道哪一位是他必须朝拜和宣认的。外表上没有任何区别，但凭着信德之眼、灵魂的视线，他认出了他的主。然后圣经接着说：\BibleQuote{他对他说：明年这时候，我必要回到你这里，你的妻子撒辣要生一个儿子。} \BibleRef{创 18:10} 随后上主对他说：\BibleQuote{我岂能隐瞒亚巴郎我所要做的事呢？}又说：\BibleQuote{上主又说：索多玛和哈摩辣的喊声甚大，他们的罪恶极重。}然后经过一段长谈（为求简洁略过），亚巴郎因无辜者与有罪者一同遭毁灭而忧伤，说：\BibleQuote{审判全地的主，岂不行公义吗？}上主说：\BibleQuote{我若在索多玛城里见到五十个义人，我就为他们的缘故饶恕那地方。}后来，当对亚巴郎的兄弟罗特的警告结束时，圣经说：\BibleQuote{上主从天上从主那里，将硫磺与火，降与索多玛和哈摩辣。}不久之后，\BibleQuote{上主照祂所说的眷顾了撒辣，上主对撒辣行了祂所说的话，撒辣就怀孕，在亚巴郎年老的时候，照天主对他所说的时候，给亚巴郎生了一个儿子。}后来，当婢女和她的儿子被赶出亚巴郎的家，因害怕孩子会在旷野缺水而死时，同一圣经说：\BibleQuote{天主听见童子的声音，上主的天使从天上呼唤哈加尔，对她说：哈加尔，有什么事？不要害怕，因为天主已经听见童子的声音，在他所在的地方。起来，把童子抱起，拉住他的手，因为我要使他成为大民族。}\par

这是什么盲目的无信，何等不信之心的迟钝，何等任性的不虔，竟对这一切浑然不知，或知晓却置若罔闻！确实，这些事被记录下来，正是为了使谬误或遗忘不致阻碍对真理的认知。如果我们如将要证明的那样，不可能逃避对事实的认识，那么否认它们就无异于亵渎。这段记载始于天使对哈加尔说的话，应许她依市玛耳将繁衍成为大国，并赐给他不可计数的后裔。她听了，并藉着她的宣认揭示了祂是主和天主。故事以祂以天主的天使形象显现开始；结束时，祂被明认为天主自己。因此，那位在宣布伟大计画时行使天使职务的，在名称和本性上就是天主。名称与本性相符；本性并未被扭曲以迎合名称。再者，天主对亚巴郎也提及同样的事：他被告知依市玛耳已蒙祝福，并将繁衍成为一个民族；天主说：\BibleQuote{我已祝福了他}。这与先前所指示的位格并无变化；祂表明就是祂自己已赐予了祝福。圣经从奥秘到清晰的启示，显然始终一贯；它始于天主的天使，进而揭示正是天主亲自在这件事上说了话。\par

神圣叙事的进程伴随着教义逐步的发展。在我们已讨论的段落中，天主对亚巴郎说话，应许撒辣将要生一个儿子。之后有三个人站在他面前；他敬拜其中一位，并宣认祂是主。在亚巴郎这样敬拜和宣认之后，那一位应许祂将在同一季节回来，那时撒辣将得到她的儿子。那一位又再次以人的形像被亚巴郎看见，并以同样的应许向他问候。变化的只是名称；亚巴郎每次的宣认都是一样的。他看见的是一位人，但亚巴郎却敬拜祂为主；无疑，他以奥秘的方式看见了将要来临的降生成人。如此坚强的信德没有失却它的认出；主在福音中说：\BibleQuote{你们的父亲亚巴郎欢欣踊跃，渴望看见我的日子；他看见了，极其喜乐。}\BibleRef{若 8:56} 继续历史；他看见的那位人应许祂将在同一季节回来。注意应许的实现，同时记住，许下应许的是一位人。圣经怎么说？\BibleQuote{主眷顾了撒辣。} 所以这位人就是主，履行祂自己的应许。接下来是什么？\BibleQuote{天主照祂所说的，对待了撒辣。} 叙述将祂的话称为人的话，记载撒辣被主眷顾，宣称结果是天主的工作。你确信说话的是一个人，因为亚巴郎不仅听见，而且看见了他。当同一部圣经，曾称祂为人，却又承认祂是天主时，你还能不那么确定祂是天主吗？因为它的原话是：\BibleQuote{撒辣怀孕，在老年给亚巴郎生了一个儿子，正是天主对亚巴郎所说的那个时期。} 但正是那一位人曾应许祂会来。如果你相信祂仅仅是人，那么事实上，除非那来者是主、是天主，否则无法解释。将这些事件联系起来。众所周知，是那一位人应许祂会来，使撒辣怀孕生子。现在接受教训，宣认信德：是主天主亲自降临，使她怀孕生子。那一位人以天主的能力许下应许；同一天主以同样的能力实现了应许。因此，天主借言语和行动启示自己。接着，亚巴郎看见的三个人中有两位离开；留下来的那一位是主、是天主。并且不仅是主和天主，还是审判者，因为亚巴郎站在主面前说：\BibleQuote{你决不会这样做，将义人同恶人一并消灭；那样的话，义人将如恶人一样。决不会这样，审判全地的主，你不行这样的审判。}\BibleRef{创 18:25} 因此，亚巴郎藉着他所有的话语，教导我们那使他成义的信德；他在三人中认出主，唯独敬拜祂，并宣认祂是主和审判者。\par

免得你陷入错误，以为亚巴郎对那一位的宣认是对他所看见的三个人共同的尊敬，请注意罗特看见那两位离开时所讲的话：\BibleQuote{罗特看见他们，就起身迎接，俯伏在地，说：\InnerQuote{我的主人们，请到仆人家里住下。}} 这里的复数形式\Quote{我的主人们}表明这不过是一次天使显现；而在另一情况下，这位信德的祖宗只将应得的尊敬归于那一位。因此，神圣叙述清楚表明，三人中的两人仅是天使；它先前已通过以下话语宣示那一位是主、是天主：\BibleQuote{主对亚巴郎说：\InnerQuote{撒辣为什么笑，说：我年老了，还能生育吗？难道天主有办不到的事吗？下一年这个季节，我要回到你这里，那时撒辣将要生一个儿子。}} 圣经是准确而一致的；我们察觉不到像将复数用于独一天主和主那样的混乱，也没有将神圣尊荣归于两位天使。毫无疑问，罗特称他们为\Quote{我的主人们}，而圣经称他们为天使。前者是人的尊敬，后者是真实的事实。\par

现在，公义审判的惩罚降在索多玛和哥摩辣身上。为了我们探究的目的，我们能从中得到什么教训？\BibleQuote{主从主那里降下硫磺与火。} 那就是\Quote{从主来的主}；圣经并未因名称的不同而对他们的本性加以区分，而是区分了他们本身。因为我们在福音中读到：\BibleQuote{圣父不审判任何人，却将审判权全部交给了子。}\BibleRef{若 5:22} 因此，主所赐予的，主已从主领受了。\par

30. 你们现在已经看到了天主作为审判者既是主，又是主；接下来要学习在天主与天主之间也有同样的名字共有。雅各伯因害怕他的兄弟而逃遁，在梦中看见一个梯子立在地上，梯顶直达天上，天主的使者在梯子上上去下来，而主在梯子之上安息，将祂赐给亚巴郎和依撒格的一切祝福都赐给了他。后来天主这样对他说：\BibleQuote{天主对雅各伯说：起来，上贝特耳去，住在那里，并在那里给那在你逃避你兄弟的面目时显现给你的天主，筑一座祭坛。}（\BibleRef{创 35:1}）天主要求尊敬天主，并清楚表明这要求是为另一位而非祂自己。\Emph{那在你逃跑时显现给你的祂}是祂的话；祂小心翼翼地防止任何位格的混淆。说话的是天主，祂所谈论的是天主。两者的尊威通过它们在真正的名号\OriginalTerm{天主}{God}下的结合而被彰显，而话语则清楚地宣告它们各自的存在。\par

31. 这里又有些考虑涌上我的心头，必须在全面论述这个主题时加以注意。但辩护的顺序必须适应攻击的顺序，我保留这些悬而未决的问题，留待下一卷讨论。当前，关于那要求尊敬天主的天主，我只指出一点就够了：那与哈加尔说话的天主的使者，在与亚巴郎谈论同一件事时，就是天主和主；那与亚巴郎说话的人也是天主和主；而与主一同被看见、并被主派往罗特那里的两位天使，先知只称他们为天使，没有别的。天主不仅以人形显现给亚巴郎，祂也以人形显现给雅各伯。不仅如此，我们还被告知，祂还摔跤；不仅摔跤，而且被对手制服。我没有时间，也不打算讨论这摔跤的预表意义。那摔跤的确实是天主，因为雅各伯战胜了天主，以色列看见了天主。\par

32. 现在让我们查考一下，除了哈加尔的事例，天主的使者是否也被发现就是天主自己。祂确实被发现，不仅是一位天主，而且是亚巴郎、依撒格和雅各伯的天主。因为主的使者从荆棘丛中显现给梅瑟；你认为我们该推断听到的是谁的声音？是那被看见者的声音，还是另一位的声音？没有欺骗的余地；圣经的话是清楚的：\BibleQuote{主的使者从荆棘里火焰中向他显现}，又说：\BibleQuote{主从荆棘里呼叫他，说：梅瑟，梅瑟！他回答：什么事？主说：不要走近这里，将你脚上的鞋脱下，因为你站的地方是圣地。主又说：我是亚巴郎的天主，依撒格的天主，雅各伯的天主。}那在荆棘丛中显现的，从荆棘丛中说话；显现与声音的场所是同一处；说话的不是别人，正是那被看见的。当眼睛看见祂时，祂是天主的使者；当耳朵听见祂时，祂是主；那被听见的主被认出来是亚巴郎、依撒格和雅各伯的天主。当祂被称为天主的使者时，事实显明祂并非自足而孤立的存有，因为祂是天主的使者。当祂被称为主和天主时，祂接受了因祂的本性和名字所应得的全部称号。因此，从荆棘丛中显现的使者，就是主和天主。\par

33. 请继续研究梅瑟所作的见证；留意他是何等勤勉地抓住每一个机会宣扬主和天主。你注意到这段经文：\BibleQuote{以色列，你要听，上主你的天主是唯一的上主。}\BibleRef{申 6:4} 也注意他这首神圣诗歌中的话：\Quote{看，看，我是上主，除我以外没有别的天主。}虽然整首诗中天主一直是发言者，但他以这样的话结束：\Quote{诸天哪，你们要同祂一同欢乐，愿天主的众子都赞美祂。外邦人哪，你们要与祂的子民一同欢乐，愿天主的众天使都尊崇祂。}天主应受天主众天使的荣耀，祂说：\Quote{因为我是上主，除我以外没有别的天主。}因为祂是天主独生子，\Quote{独生}这一称号排除了在此特性上的任何分享，正如\Quote{无源}这一称号否认了在这方面有任何分享无源圣父特性者。圣子是独一者出自独一者。除了无源的天主外，没有无源者；同样，除了天主的独生子外，没有独生者。他们各自单独而独特，分别是那一位无源者和那一位独生者。所以他们两位是独一的天主，因为在永恒的天主性中，独一者与祂所出的独一者之间没有本性差异的鸿沟。因此，祂应受天主众子的敬拜，应得天主众天使的荣耀。天主从祂的众子和众天使那里要求尊荣和敬畏。请注意是谁要接受这尊荣，以及由谁来呈献。祂是天主，而他们是天主的众子和众天使。免得你以为这尊荣不是要求赐予那位分享我们本性的天主，而是梅瑟在此想到应给予天主圣父的敬畏——尽管确实，圣父必须在圣子内受尊荣——请查考在本书末尾天主祝福若瑟的话。它们是：\BibleQuote{愿那在荆棘中显现者所喜悦的福祉，临到若瑟的头上和头顶。}\BibleRef{申 33:16} 因此，天主要由天主的众子来敬拜；但这位天主本身就是天主子。天主要由天主的众天使来尊崇；但这位天主本身就是天主的使者。因为天主曾作为天主的使者从荆棘中显现，而祈祷若瑟能领受祂所喜悦的祝福。祂并不因为是天使就不是天主；祂也不因为是天主就不是天使。这是天主位格的清晰指示；在无源者和受生者之间划出了明确的界限。天国的奥秘得以启示，我们被教导不要幻想天主独居，当天使和天主的众子要敬拜那位是天主使者和祂圣子的天主时。\par

34. 让我们将此视为从梅瑟诸书中获得的答案，或者不如说是梅瑟本人的答案。异端者幻想他们可以利用梅瑟对天主独一性的断言来反驳天主圣子的天主性；这是一种亵渎，公然违背了他们自己见证的明确警告，因为每当梅瑟承认天主是独一的时候，他从未停止教导圣子的天主性。我们的下一步必须引用先知们关于同一圣子的多方面的言论。\par

35. 你知道这些话：\Quote{以色列，你要听，上主你的天主是唯一的上主；}但愿你能正确理解它们！按你的解释，我徒然寻找它们的意义。在圣咏中说：\Quote{天主，你的天主，用喜油傅抹了你。}请向读者的心中灌输傅油者和受傅者之间的区别；区分\Quote{你}和\Quote{你的}：清楚指明这些话是向谁和关于谁说的。因为这一明确的宣认是前面段落的结论，那段话是这样：\Quote{天主，你的宝座永远常存；你王国的权杖是正直的权杖。你喜爱正义，憎恨不义。}然后他继续说：\Quote{因此，天主，你的天主，用喜油傅抹了你。}因此，永恒国度的天主，因祂喜爱正义和憎恨不义，被祂的天主所傅抹。这些名号之间难道没有某种巨大的差异、某种超越我们心智跨越的鸿沟吗？不；位格的区别由\Quote{你}和\Quote{你的}表明，但没有暗示本性的差异。\Quote{你的}指向那根源，\Quote{你}指向那根源的后裔。因为祂是由天主而来的天主，正如先知同样的话所宣告的：\Quote{天主，你的天主，用喜油傅抹了你。}祂自己的话作证，在无源的天主之前没有别的天主：\BibleQuote{你们是我的见证，我也作见证，——上主天主说——并我所拣选的仆人，为使你们知道、相信并明白我是那一位，在我以前没有别的天主，在我以后也没有。}\BibleRef{依 43:10} 这样，那无始者的尊威被宣告，那从无源者而来者的荣耀被维护；因为\Quote{天主，你的天主，用喜油傅抹了你。}\Quote{你的}这个词宣告了祂的诞生，却并不与祂的本性矛盾；\Quote{你的天主}意味着圣子是从祂而生，分享天主性。但圣父是天主这一事实并不妨碍圣子也是天主，因为\Quote{天主，你的天主，用喜油傅抹了你。}圣父和圣子都被提及；天主这一共同称号保证了他们在本性和尊荣上是独一的。\par

36. 但为免这些话——\Quote{\Emph{我是自有者，在我以前没有别的天主，在我以后也不会有}}——成为亵渎傲慢的藉口，藉以证明圣子不是天主，因为在没有天主在他以前的那位天主之后，没有别的天主，就必须考虑这段经文的意图。天主是他自己最好的解释者，但他所拣选的仆人与他一同向我们保证：在他以前没有天主，在他以后也不会有。他自己对自己的见证确实足够，但他又加上了他所拣选的仆人的见证。如此，我们有了两者一致的见证：在他以前没有天主；我们接受这真理，因为万有都出自他。我们也有他们二者的见证：在他以后也不会有天主；但他们并不否认在以往曾从他有天主诞生。已经有过仆人如此说话，并为圣父作证；仆人是从那个民族中诞生的，天主的选民要从那民族中兴起。他在福音中也陈述同样的真理：\BibleQuote{\Emph{请看，我所拣选的仆人，我所钟爱的，我的心灵所喜悦的}}。\BibleRef{玛 12:18}因此，这就是天主说\Quote{\Emph{在我以前没有别的天主，在我以后也不会有}}的含义。他藉着这断言——他自己以前或以后都没有天主——启示了他永恒不变威严的无限。但他使他的仆人分享作证和拥有天主之名的权利。\par

37. 这从他自己的话中显而易见。因为他向先知欧瑟亚说：\BibleQuote{\Emph{我不再怜爱以色列家，反而完全与他们为敌。但我要怜爱犹大家，并藉他们的天主上主拯救他们}}。\BibleRef{欧 1:6}在这里，天父将天主的名字毫无歧义地给予了圣子，在圣子内他在无数世代前拣选了我们。他说\Quote{\Emph{他们的天主}}，因为圣父既然是无始的，独立于万物，就将我们赐予他的圣子作为产业。同样，我们读到：\Quote{\Emph{你向我请求，我必将外邦人赐给你作产业}}。对于万有都从他而来的那一位，没有人能成为他的天主，因为他是永恒的，没有开始；但圣子有天主作为他的父，他从父而生。然而对我们来说，圣父是天主，圣子也是天主；圣父启示我们：圣子是我们天主，而圣子教导我们：圣父是我们之上的天主。我们要记住的关键是，在这段经文中，圣父将天主的名字——他自己无始威严的称号——给予了圣子。但我对欧瑟亚的这些话已做了充分的评论。\par

38. 再者，天父藉依撒意亚论及我等上主所作的宣告是何等清晰！他说：\BibleQuote{\Emph{因为上主，以色列的圣者天主，那创造未来之事的这样说了：关于你们的儿女你问我，关于我手的作为你们可吩咐我。我造了大地和其上的人，我命令了所有星辰，我按正义兴起了一位君王，他的所有道路都是正直的。他要建造我的城邑，使我的子民从被掳中归回，不用代价也不需报酬——万军的上主说。埃及要劳苦，以及雇士人和舍巴人的货物。身材魁梧的人要来归顺你，作你的仆人，带着锁链跟随你，并要敬拜你，向你祈求，因为天主在你内，除你以外没有别的天主。你真是天主，我们却不认识，以色列的天主，救主。凡反抗他的都要蒙羞受辱，在混乱中行走}}。\BibleRef{依 45:11}还留有反驳的余地或无知藉口吗？若亵渎继续，它岂非在无耻的对抗中生存？万有都来自他、他藉命令创造万物的天主，宣称自己是宇宙的创造者，因为除非他说话，没有任何事物被造。他宣称他兴起了一位正义的君王，这位君王为他——即为天主——建造一座城，并使他的人民从被掳中返回，不用礼物也不需报酬，因为我们都是白白得救的。接着，他述说在埃及的劳苦之后，在雇士人和舍巴人的货物之后，身材魁梧的人要归顺他。我们该如何理解这埃及的劳苦、这雇士人和舍巴人的货物？让我们回想东方的贤士如何敬拜并向上主进献贡物；让我们估量那前往犹大白冷的漫长朝圣之旅的疲惫。在贤士王侯们的艰辛旅程中，我们看到了先知所暗示的埃及的劳苦。因为当贤士以他们虚假、物质的方式执行天主的能力为他们所定的职责时，整个外邦世界藉他们奉献了其敬拜所能达到的最深敬意。这些贤士也从雇士人和舍巴人的货物中献上了黄金、乳香和没药的礼物；这是另一位先知所预言的，他曾说：\Quote{\Emph{雇士人要在他的面前俯伏，他的仇敌要舔尘土。塔尔史士的众王要献礼，阿剌伯和舍巴的众王要进贡，必有阿剌伯的黄金献给他}}。贤士和他们的供物代表了埃及的劳苦和雇士人与舍巴人的货物；朝拜的贤士代表了外邦世界，并向上主献上外邦人中最精选的礼物。\par

39. 至于那些高大的人，他们将归向祂，带着锁链跟随祂，毫无疑问他们是谁。翻开福音：伯多禄要跟随他的主时，被束上腰带。读宗徒：基督的仆人保禄夸耀他的锁链。让我们看看这位\Quote{耶稣基督的囚犯}是否在自己的教导中符合天主关于天主祂圣子所说的预言。天主曾说过：\BibleQuote{他们要恳求，因为天主在祢内。}现在请注意并消化宗徒的这些话：\BibleQuote{天主在基督内，使世界与自己和好。}\BibleRef{格后 5:19} 然后预言继续说：\BibleQuote{除了祢以外，没有天主。}宗徒立即用这话配合：\BibleQuote{因为只有一个主耶稣基督，万物都藉着祂而有。}\BibleRef{格前 8:6} 显然，除祂以外没有别的，因为祂是唯一的。第三个预言是：\BibleQuote{祢是天主，我们却不认识祢。}但保禄，曾是教会的迫害者，说：\BibleQuote{圣祖们也是他们的，并且基督按肉身说也是从他们来的，祂是在万有之上，永远受赞美的天主。}\BibleRef{罗 9:5} 这就是那些带着锁链的人所要传报的信息；他们确实是高大的人，要坐在十二个宝座上审判以色列十二支派，并跟随他们的主，在教导和殉道上为祂作证。\par

40. 因此，天主在天主内，而天主正是天主所居住的那位。但既然天主在祂内，\BibleQuote{除了祢以外，没有天主}这句话如何为真呢？异端者！为了支持你承认一位独一的父，你使用了\BibleQuote{除我以外，没有天主}这句话；如果你对\BibleQuote{除我以外，没有天主}的解释是否认圣子的天主性，那么你能赋予天主圣父庄严宣告的\BibleQuote{除了祢以外，没有天主}什么意义呢？在那情况下，天主能向谁说\BibleQuote{除了祢以外，没有天主}呢？你无法暗示这位独一者是对自己说的。主藉着那些敬拜并恳求的高大之人的口，对祂所召来的君王说：\BibleQuote{因为天主在祢内。}事实与单独存在不一致。\BibleQuote{在祢内}暗示有一位在说话者声音可及的范围内（如果可以这样说）。完整的句子\BibleQuote{天主在祢内}不仅揭示了天主临在，也揭示了天主居住在那位临在者内。这些话区分了居住者与祂所居住的那位，但这只是位格上的区分，而非本性上的。天主在祂内，而祂——天主在其中者——就是天主。天主的居所不可能是在一个与祂本性相异或陌生的本性内。祂居住在一位属于祂自己的、从祂所生者内。天主在天主内，因为天主来自天主。\BibleQuote{因为祢是天主，我们却不认识祢，以色列的天主，拯救者。}\par

41. 我下一本书致力于反驳你否认天主在天主内的说法；因为先知继续说：\BibleQuote{凡抗拒祂的都必羞愧、惊惶，并在混乱中行走。}这是天主对你不信的判决。你与基督对立，正是为了祂的缘故，圣父的声音发出庄严的斥责；因为你否认其天主性的那位就是天主。而你以敬畏天主为借口否认它，因为祂说：\BibleQuote{除我以外，没有别的神。}你当羞愧、混乱；那无始的天主不需要你所提供的尊严；祂从未要求你归于祂的这种孤立威严。祂拒绝你多管闲事的解释，即曲解祂的话\BibleQuote{除我以外，没有别的神}为否认祂从自己所生的圣子的天主性。为了挫败你通过将天主性以某种特殊意义归于自己而破坏圣子天主性的目的，祂以绝对天主性的归赋来完善独生子的荣耀：\BibleQuote{除了祢以外，没有天主。}为什么要在完全相等者之间作区分？为什么要分开完全匹配的？天主圣子的独特特征是没有别的神在祂之外；天主圣父的独特特征是没有别的神离开祂。使用祂关于自己的话；用祂自己的术语承认祂，并呼求祂为君王；\BibleQuote{因为天主在祢内，除了祢以外，没有天主。因为祢是天主，我们却不认识祢，以色列的天主，拯救者。}用如此恭敬的话所作的承认，没有僭越之嫌；其措辞不会引起反感。最重要的是，我们必须记住，拒绝它意味着羞愧和耻辱。在心里沉思这些话，并用它们来承认祂，这样就能逃避所威胁的羞愧。因为如果你否认天主圣子的天主性，你就不是通过以孤独威严崇拜祂来增加天主的荣耀，而是因拒绝尊敬圣子而轻视了圣父。要凭着信德和敬畏承认那无始的天主：除了祂以外，没有天主；为天主独生子主张：除了祂以外，没有天主。\par

42. 如同你们已听了\Person{梅瑟}和\Person{依撒意亚}的话，现在也听\Person{耶肋米亚}所教导的同一真理：\BibleQuote{这是我们的天主，没有别的能与他相比；他寻出了全部知识的途径，将它赐给了他的仆人\Person{雅各伯}和他心爱的\Place{以色列}。此后他在地上显现，并与人同居。} 因为此前他曾说：\BibleQuote{他是人，谁能认识他？} 这样，你看见天主在地上显现，并在人中间居住。现在我问你，当\Person{耶肋米亚}宣告天主在地上显现并与人同居时，你将如何解释\BibleQuote{从来没有人见过天主，只有那在父怀里的独生子，将他表明出来}？\BibleRef{若 1:18} 显然，除了圣子，没有人能看见圣父；那么，这位被看见并与人同居的是谁？他必定是我们的天主，因为他是以人的形态可见的天主，人能够触摸他。并且要牢记先知的话：\BibleQuote{没有别的能与他相比。} 如果你问这怎么可能，请听这句话的其余部分，以免你被诱惑而否认圣父在信仰宣认中的份位。\BibleQuote{\Place{以色列}，你要听！主你的天主是唯一的。} 整段经文是：\BibleQuote{没有别的能与他相比，他寻出了全部知识的途径，将它赐给了他的仆人\Person{雅各伯}和他心爱的\Place{以色列}。此后他在地上显现，并与人同居。} 因为在天主与人之间只有一位中保，他既是天主也是人；他既在赐予法律上，又在取了我们的身体上作中保。因此没有别的能与他相比，因为他是唯一的，从天主而生于天主，万物是借着他而造的，无论是天上的还是地上的，时间与世界也是借着他造成的。总之，一切存在之物都因他的行动而存在。就是他教训\Person{亚巴郎}，与\Person{梅瑟}说话，为\Place{以色列}作证，居住在先知内，因圣神由童贞女诞生，以他的苦难将我们的仇敌权势钉在十字架上，在地狱中杀死死亡，以他的复活坚固我们希望的確证，以他身体的荣耀摧毁人类肉身的腐败。因此没有别的能与他相比。因为这些都是天主独生子的独特权能；唯有他是从天主而生，是这些伟大特权的幸福拥有者。没有第二位神能与他相比，因为他是天主从天主，不是从任何外来的存在所生。在他里面没有新造、奇异或近代的事物。当\Place{以色列}听见它的天主是唯一的，并且没有第二位神能相比，而人们可能把他视为神——那身为天主子之神——这启示的意思是：圣父天主和圣子天主完全是一体的，不是位格的混淆，而是本质的同一。因为先知禁止我们，既然天主圣子是神，就不应将他比作某第二位神祇。\par

\SourceRecord{Translated by E.W. Watson and L. Pullan.；Nicene and Post-Nicene Fathers, Second Series；Vol. 9.；Edited by Philip Schaff and Henry Wace.；Buffalo, NY: Christian Literature Publishing Co.,；1899.；<http://www.newadvent.org/fathers/330204.htm>.}

% structure: p0001 source=h1 target=section reason=page-title

\section{论天主圣三（第五卷）}

第五卷。\Person{希拉利}现在指出（§ 1）\OriginalTerm{亚略派}{Arian}论点的争议力量。若他面对他们的断言保持沉默，他们便会声称他同意他们，即圣子只以某种较低层次的意义是天主。另一方面，若他反对他们，他似乎就会与\Person{梅瑟}关于天主唯一性的启示相矛盾。在§ 2中，他概述了第四卷的论点，即\WorkTitle{圣经}的见证证明天主并非孤独的一位；正如他所说的，确有\Emph{天主和天主}。然而亚略派还有另一个漏洞；他们的信经（§ 3）宣称只有一位真天主。他们可能会争辩说，基督确实是天主，但其本性不同于圣父。为驳斥这一点，\Person{希拉利}再次回顾创造的历史（§§ 4-10），证明叙述不仅揭示了圣子在这项工作中有份，而且揭示了祂与圣父的平等和本性的同一；换句话说，祂不仅是天主，而是真天主。同一真理从\Person{亚巴郎}的生平中得到证明（§§ 11-16）。此外，圣子（作为\OriginalTerm{天使}{Angel}，多次）的这些自我启示是降生成人的预兆。祂首先在肉身中被看见，随后由肉身而生。亚略派专注于基督人性生活的卑微条件，因此，由于缺乏全面的视野，无法辨别祂真正的天主性。但\Person{希拉利}不会预先引用\WorkTitle{福音}的证据（§§ 17, 18）。他回到\WorkTitle{旧约}，从\Person{雅各伯}的异象（§§ 19, 20）以及向\Person{梅瑟}的启示（§§ 21-23）证明了他的观点。在总结和巩固前面的论点（§§ 24, 25）之后，他继续从\Person{依撒意亚}的某些段落证明先知承认圣子是真正的天主（§§ 26-31），并且\Person{圣保禄}也是这样理解（§§ 32, 33）。然后，在§§ 34, 35中，着重论述了已经达到的结果。\Person{希拉利}表明，正是亚略派未能承认唯一的真天主；因为基督是真天主，却并非第二位天主。最后，在§§ 36-39中，引用了\Person{梅瑟}、\Person{依撒意亚}和\Person{耶肋米亚}为证，证明基督是出自天主的天主，并在天主内的天主。\par

1. 我们在前几卷中对异端者疯狂而亵渎的教义所作的答辩，使我们清醒地陷入了一种困境：无论我们驳斥对手，还是保持沉默，我们的读者都面临同等的危险。因为当不信者以喧嚣的无礼将天主的唯一性强加给我们——这种唯一性是虔诚而合理的信仰所不能否认的——时，谨慎的心灵陷入了两难：无论它肯定还是否定这一命题，都同样会有亵渎的危险。对于人类的逻辑来说，说肯定同一教义可能是不虔敬的，否定它也可能是不虔敬的，这似乎荒谬而不合理，因为凡是虔敬地坚持的，必定是不虔敬地争论的；若推翻一个主张有好的目的，那么梦想支持它会有好处，就显得愚蠢了。但人类的逻辑在天主的计划面前是谬误，在应对天上的智慧时是愚拙；它的思想受其局限束缚，它的哲学被自然理性的软弱所限制。它必须在自己眼中成为愚拙，才能在天主面前成为智慧；也就是说，它必须认识到自身能力的贫乏，并寻求天主的智慧。它必须变得智慧，不是按照人类哲学的标准，而是按照那上升到天主的智慧，才能进入祂的智慧，其眼睛才能看清世界的愚拙。异端者巧妙地策划，让这种被视为智慧的愚拙成为他们的引擎。他们利用对唯一天主的宣认，为此他们援引\WorkTitle{法律}和\WorkTitle{福音}的见证，说：\BibleQuote{以色列，你要听！主你的天主是唯一的。}他们很清楚涉及的风险，无论他们的断言是遭到反对还是被默然放过；无论发生哪种情况，他们都看到促进其异端的机会。若怀着亵渎意图所提出的神圣真理遭受沉默，那沉默就被解释为同意；就好像在承认，因为天主是唯一的，所以祂的圣子不是天主，天主永远独处。另一方面，若他们大胆论据中所包含的异端遭到反对，这种反对就被打上背离真正福音信仰的烙印，因为福音信仰精确地陈述了天主的唯一性；或者他们反过来指责对手陷入了相反的异端，即只承认圣父和圣子是一个位格。这就是那致命的诡计，披着诱人无害的外衣，是世上的智慧，在天主面前是愚拙的，锻造出来欺骗我们的，在他们信仰的第一条上，我们既不能宣认也不能否认，而毫无亵渎的危险。我们行走在两方面的危险之间；天主的唯一性可能迫使我们否认圣子的天主性，或者，若我们宣认圣父是天主、圣子是天主，我们可能被驱入以\OriginalTerm{撒伯流主义}{Sabellian}的方式来解释圣父和圣子之合一的异端。因此，他们坚持\Emph{唯一天主}的诡计，要么将第二位格排除在天主性之外，要么通过承认祂为第二位天主而摧毁合一，要么使唯一性仅仅成为名义上的。因为唯一性，他们会辩称，排斥第二位；第二位存在会破坏合一；二不能成为一。\par

2. 但我们这些已获得天主的上智——这在世界看来是愚妄——并立意藉着对主之真信德的健全与拯救性宣认，揭露他们教导中蛇一般的诡诈；我们如此安排撰写计划，以便获得展示真理的有利据点，而不致陷入异端论断的危险。我们谨慎避免两种极端：既不否认天主是唯一，又依据那宣告天主唯一性的立法者的明证，清晰地陈述有\Quote{天主}与\Quote{天主}的真理。我们教导说，天主之所以为唯一，绝非由于二者的混淆；我们并不因宣讲多位天主而将祂撕裂，也不只承认名称上的区别。而是将祂呈现为\Quote{天主}与\Quote{天主}，将天主的唯一性问题暂时搁置，留待以后更充分的讨论。因为福音告诉我们，梅瑟在宣告天主唯一时教导了真理；而梅瑟藉着宣告唯一天主，也证实了福音关于\Quote{天主}与\Quote{天主}的教诲。因此，我们并不与权威矛盾，而是将教导建立于其上，证明向以色列启示的天主唯一性，并不容许否认天主圣子的天主性；因为那作为我们断言存在唯一天主的权威，也是我们承认祂圣子天主性的权威。\par

3. 因此，我们论文的安排紧密遵循所提出异议的顺序。由于他们亵渎与不诚实的宣认中下一项是\Quote{我们宣认唯一真天主}，这整卷第二卷便致力于探讨天主圣子是否为真天主的问题。因为显然，异端者巧妙设计了这一顺序：先提及\Quote{唯一天主}，再提及\Quote{唯一真天主}，目的是将圣子从天主的名号与本性中分离；因为必然产生这样的想法：既然真理内在于唯一天主，就必须严格限于祂。既然梅瑟在宣告天主唯一性时毫不动摇地包含了圣子的天主性，那么让我们回到他教导所传达的主要段落，探究他是否希望我们相信那一位——他已教导我们为天主——的子也是真天主。显然，一事物的真确性或真实性关乎其本性与能力。例如，真小麦是那些长穗、麦芒环绕、经簸扬去糠、磨成粉、做成饼、作为食物，并表现面包的本性与用途的。因此，自然能力是真理的证据；让我们以此检验梅瑟称为天主的祂是否为真天主。关于这位唯一也是真天主，我将在当下暂缓讨论，以免我若未能立即应战并拥护那在圣父与圣子两位格中的唯一真天主，急切焦虑的灵魂就会被危险的疑虑所压抑。\par

4. 既然我们都承认天主视祂的圣子为天主这一共同基础，我问你：世界的创造何以证明我们断言圣子是真天主为错误？毫无疑问，万有都是藉着圣子，因为如宗徒所说：\BibleQuote{\OriginalTerm{万有都是藉着祂，并在祂内}{All things are through Him, and in Him}}\BibleRef{哥 1:16}。如果万有都是藉着祂，且万有都是出自无有，且只有藉着祂才能产生，那么拥有天主的本性与能力的祂，在真天主性的哪个元素上有所欠缺？祂拥有天主本性的能力，能将不存在之物带入存在，并随心创造。因为\Quote{天主看了，认为好}。\par

5. 法律说：\BibleQuote{天主说：要有穹苍}，随后又补充：\BibleQuote{于是天主造了穹苍}，它没有引入其他区别，只有位格的区别。它没有显示能力或本性的差异，也没有改变名称。在同一天主的名号下，它首先启示了发言者的思想，然后揭示了创造者的行动。叙述者的语言丝毫未剥夺祂的天主本性和能力；恰恰相反，它多么精确地强调了祂的真天主性。使创造之言生效的能力只属于那样的本性：说话即成就。那么，如果发令者是真天主，那么创造者如何不是真天主？若所说的言辞是真正神圣的，所行的事也是真正神圣的。天主发言，天主创造；如果发言的是真天主，那么创造的也是真天主；除非，在一方的言语中显示了真天主性的临在，在另一方的行动中显示出其缺失。因此，在天主圣子身上，我们看到了真正的天主本性。祂是天主，是创造者，是天主圣子，是全能者。这不仅仅是祂能随意而行，因为意愿常与能力相伴；而且祂也能做任何受命于祂的事。绝对的能力在于，其拥有者作为执行者，能完成其作为说话者能用言辞表达的一切。当无限的表达力与无限的执行力结合时，这种与命令之言相称的创造能力，便拥有天主的真本性。因此，天主圣子不是假天主，不是嗣子式的天主，也不是因名号恩赐的天主，而是真天主。如果陈述那些反对祂真天主性的论据，将毫无所得；祂拥有天主的名号与本性是决定性的证明。万有是藉着祂造的，祂是天主。世界的创造如此告诉我关于祂的一切。祂是天主，与天主之名称配；是真天主，与真天主之能力相配。天主的能力在创造之言中启示给我们；天主的能力也在创造行动中彰显。现在我再次问你，你根据什么权威，在宣认圣父与圣子时，否认祂——其名号揭示祂的能力，其证明祂拥有这名号——的真天主性？\par

6. 读者须牢记，我并非因健忘或对己见缺乏信心而讳言常见的反对意见。因为那常被引用的经文\BibleQuote{圣父比我大}及其同类段落，我十分熟悉，且已备好我的诠释，使它们成为圣子真天主性的见证。然而，依论敌进攻的顺序作答，最合我意；这样，我们敬虔的努力便能紧随他们亵渎计划的进程，一旦见他们偏离至不虔敬的异端，我们便能立即抹去错误的轨迹。为此，我们将福音作者和宗徒们的见证推迟到著作的末尾，眼下依据法律和先知书与亵渎者交战，用他们试图欺骗我们的那些经文本身，来封住他们基于误解和欺诈的歪曲论证。证明真理的可靠方法，是揭露针对它所提出的异议的谬误；若骗子的谎言自身被转化为真理的证据，他的耻辱便彻底了。事实上，人类的普遍经验已学会，虚假与真理互不相容，无法调和或连贯；按其本性，它们属于永远敌对、永远不能结合或一致的对立面。\par

7. 既然如此，我问：在\BibleQuote{让我们按我们的肖像和模样造人}这句话中，如何能在真天主与假天主之间作出区分？这些话表达了一个意义，意义是思想的结果；思想由真理所推动。让我们由话语追溯到其意义，从意义了解思想，借思想达到其下的真理。你的问题是：对祂所说的\BibleQuote{让我们按我们的肖像和模样造人}这句话所指向的那一位，说话者是否认为祂不是真的？因为这些话无疑表达了说话者的情感和思想。在说\BibleQuote{让我们造}时，祂清楚地指明了一位与祂自己并不不和、并非外来的或无能的存在，而是一位具有权能去执行祂所说之事的存在。祂自己的话向我们保证，这正是我们必须理解这些话语所说意思的方式。\par

8. 为更充分地向我们保证在圣子的本性和工作中所显示的真天主性，祂在我所引的话中表达了祂的意思，表明祂的思想是由祂对其所说\BibleQuote{按我们的肖像和模样}的那一位的真天主性所提示的。那位真天主对祂说\BibleQuote{按我们的肖像和模样}，祂怎会被虚假地称为天主？\Emph{我们的}一词与孤立、意图或本性上的差异都不相容。按字面严格意义理解，人是按他们的共同肖像被造的。现在，真与假之间不可能有共同之处。说话的天主是对天主说话；人是按圣父和圣子的肖像被造的。两者在名字上和本性上都是一体。人按之被造的是唯一的肖像。我尚未到讨论此事的时候；日后我将解释人是按天主圣父和天主圣子的什么肖像被造的。目前我们坚持这个问题：那位真天主对其说\BibleQuote{让我们按我们的肖像和模样造人}的那一位，是不是真天主？你若能，便将这一肖像中共同的真假成分分开；在你异端的疯狂中分裂那不可分割的。因为两者乃是一体，按其一体的肖像和模样，人成了那惟一的复制品。\par

9. 但现在让我们继续诵读这段圣经，以显示真理的一致性如何不受这些不诚实的反对所影响。接下来的话是：\BibleQuote{天主于是照自己的肖像造了人；就是照天主的肖像造了他}。肖像是共同的；天主照天主的肖像造了人。我要问那否认天主圣子是真天主的人，他以为天主是照着谁的肖像造了人？他必须始终牢记，万有都是藉着圣子而有的；异端的机巧不可为自己目的将此段扭曲为圣父的行动。因此，如果人是藉着天主圣子照着天主圣父的肖像被造的，他也是照着圣子的肖像被造的；因为所有人都承认，\BibleQuote{按我们的肖像和模样}这句话是对圣子说的。因此，祂的真天主性既被圣言明确肯定，又在神圣行动中显示出来；所以，是天主将人塑造成天主的肖像，祂显示自己为天主，而且是真天主。因为祂共同拥有天主的肖像，证明祂是真天主，而祂的创造行动则显示祂是天主圣子。\par

10. 何等疯狂的丧失理智之魂！何等鲁莽的亵渎之盲目的蛮勇！你听到天主和天主；你听到\BibleQuote{我们的肖像}。为何暗示一个是真天主，而另一个不是？为何区分本性上的天主和名义上的天主？为何以维护信仰为借口，你却摧毁信仰？为何竭力将一位天主、一位真天主的启示歪曲成否认天主是独一而真实的？我还不会用福音作者和先知们清晰的言语来扼杀你疯狂的企图，这些言语中圣父和圣子并非一个位格，而是本性为一，且各自为真天主。眼下，仅法律就足以消灭你。法律是否曾谈到一位真天主和一位非真天主？法律是否曾以天主之名——那是对他们本性的真实表达——来称呼两者之一？它谈到天主和天主；它也谈到天主是独一的。不，它所做的还不止这些。它通过两者共同肖像的确凿证据，表明他们是真天主和真天主。它首先以他们严格的天主之名来称呼他们；然后，将真天主性共同归于两者。因为当人，他们的受造物，按两者的肖像被造时，健全理性迫使我们得出结论：两者中的每一位都是真天主。\par

11. 但让我们再次踏上教诲的旅途，重温天主神圣法律所教导的课程。天主的天使对哈加尔说话；而这同一天使就是天主。但或许祂是天使，意味着祂不是真天主；因为这一称号似乎指示一种较低的本性；当名称指向本性上的差异时，人们便认为真实平等必定缺失。前一卷已揭露了这反对意见的空洞；天使的称号告诉我们祂的职务，而非祂的本性。我有先知性证据为此解释；\BibleQuote{祂使祂的天使成为神风，使祂的仆役成为火焰}。那火焰是祂的仆役；那来临的神风是祂的天使。这些形象显示了祂的使者（即天使）及祂的仆役的本性与能力。这神风是一位天使，那火焰是天主的仆役。它们的本性使之适于使者或仆役的职务。因此，法律——更确切地说，天主通过法律——为指示天主圣子是一个位格，却与圣父共享同一名称，便称祂为天使，即天主的使者。\OriginalTerm{使者}{Messenger}这一称号证明祂有祂自己的职务；而祂被称为天主，则证明祂的本性是真天主。但这一顺序——先是天使，然后是天主——是在启示的秩序中，而非在祂自身内。因为我们以最严格的方式承认圣父与圣子，彼此平等，以至于独生子因祂的出生，从非受生之父获得真天主性。这启示将祂们显示为派遣者与被派遣者，不过是圣父与圣子的另一种表达；既不否定圣子的真天主性，也不取消祂作为出生权所拥有的天主性。因为无人能怀疑圣子因出生而先天性地分有祂根源的本性，以至于从独一者中生出一个不可分的合一，因为独一者出自独一者。\par

12. 信德燃烧着炽热的热忱；沉默的负担难以忍受，我的思绪迫使我发言。在前一卷中，我已经偏离了既定的论证方法。我谴责异端者谈论独一天主时的亵渎含义，并解释了梅瑟论及天主与天主的段落。我因急切而虔诚的热忱，匆忙转向我们所持守的天主合一的正确信理。而如今，再次陷入另一探究的执着中，我任凭自己偏离了轨道；当我致力于圣子的真天主性时，我灵魂的热忱催促我过早地宣认真天主为圣父与圣子。但我们自己的信德必须在本论著中等待其适当的位置。此初步陈述乃是为读者之保障；日后将如此发展并解释，以致挫败反驳者的计谋。\par

13. 返回论证：这职务的称号并不表明本性上的差异，因为那作为天主天使的，就是天主。祂的真天主性之检验，在于祂的言语和行为是否属于天主。祂将依市玛耳繁衍为一个大民族，并应许许多多民族将奉他的名。请问，这在天使的能力范围内吗？若否，而这是天主的能力，那么你为何拒绝承认那拥有你所承认的天主真实能力者为真天主呢？因此，祂拥有天主本性真实而完美的权能。真实的天主，在祂为世界的救恩而启示自己的所有预像中，不是、也绝不可能是别的，唯有真天主。\par

14. 首先，我问：这些术语\OriginalTerm{真天主}{true God}与\OriginalTerm{非真天主}{not true God}是什么意思？若有人对我说\Quote{这是火，但不是真火；水，但不是真水}，我无法理解他话中的任何意义。同一类别的两个真实样本之间，能有什么本性上的差异？若某物是火，它必是真火；只要其本性保持不变，它就不能失去真实火的特性。剥夺水的本性，你就因此摧毁了它作为真水的特性；若它仍是水，那么它必然仍为真水。物体丧失本性的唯一方式是通过丧失存在；若它继续存在，它必定真是其自身。如果天主子是天主，那么祂就是真天主；如果祂不是真天主，那么祂在任何意义上都不是天主。如果祂没有那本性，那么祂无权拥有那名称；反之，如果那指示本性的名称是祂固有之权利，那么祂就不可能缺乏那本性最真实的意义。\par

15. 但或许有人会辩称，当天主的天使被称为天主时，祂是作为恩宠、通过义子身份而获得这名称，因此拥有的是名义上的、而非真实的天主性。如果当祂被称为天主的天使时，祂对自己的天主性作了不充分的启示，那么请判断，祂在低于天使的本性名下是否已充分彰显了祂的真天主性。因为一个人对亚巴郎说话，亚巴郎敬拜祂为天主。瘟疫般的异端者啊！亚巴郎承认祂是天主，你却否认。你亵渎了，还有什么希望获得应许给亚巴郎的祝福？他是外邦人之父，但不是为你；你不能离开你的重生，加入他的后裔之家，通过赐予他信德的福分。你不是从石头中给亚巴郎兴起的儿子；你是毒蛇的种类，是他信仰的仇敌。你不是天主的以色列，不是亚巴郎的承继者，因信德而成义；因为你不信天主，而亚巴郎因那信德而成义，并被立为外邦人之父——在那信德中，他敬拜了那他所信赖其话语的天主。那时，那位蒙福而忠信的族长所敬拜的确是天主；请注意，那位天主何等真实，在祂自己的话中，一切都是可能的。除了天主，还有谁没有什么是不可能的？而那位万事可能者，难道会缺乏真天主性吗？\par

16. 我进一步追问：这位毁灭索多玛和哈摩辣的天主是谁？因为\BibleQuote{主从主那里降下雨火}\newline{}\BibleRef{创 19:24}；这难道不是真主从真主那里降下吗？\newline{}你对这位主与主还有什么别的说法吗？\newline{}或者你对这些词语还有什么别的解释，除了在\Quote{主}与\Quote{主}中，他们的位格被区分开来？\newline{}请记住，你所明认的\Emph{唯一真天主}，你也曾明认祂是\Emph{唯一的公义审判者}。\newline{}现在请注意，那位从主那里降下雨火、不将义人与不义者一同击杀、并审判全地的主，\newline{}既是主，也是公义的审判者，且从主那里降下雨火。\newline{}面对这一切，我问你：你描述谁是唯一的公义审判者？\newline{}主从主那里降下雨火；你不会否认那从主那里降下雨火的是公义的审判者，\newline{}因为外邦人的祖先亚巴郎——但不是不信的外邦人的祖先——这样说道：\newline{}\BibleQuote{你决不能这样做，将义人与恶人一同杀死，那么义人与恶人便都一样了。\newline{}你，审判全地的，决不能施行这样的审判。}\newline{}\BibleRef{创 18:25}\newline{}所以，这位天主，公义的审判者，显然也是真天主。\newline{}亵渎者！你自己的谎言反驳了你。\newline{}我还没有引用福音关于天主审判者的见证；法律已经告诉我祂是审判者。\newline{}你必须剥夺圣子的审判权，然后才能剥夺祂真正的天主性。\newline{}你已庄严地明认，那唯一公义的审判者也是唯一真天主；\newline{}你自己的陈述迫使你承认：那公义的审判者也是真天主。\newline{}这位审判者是主，万物对祂都是可能的，祂是永恒福乐的应许者，义人和恶人的审判者。\newline{}祂是亚巴郎的天主，受到他的朝拜。\newline{}愚蠢的亵渎者，你无耻的舌辩必须发明一些新的谬论，才能证明祂不是真天主。\par

17. 祂仁慈而奥秘的自我启示，绝不与祂真实的天上本性相悖；\newline{}祂忠信的圣人总能穿透祂所取的外表，好让信德能看见祂。\newline{}法律的预象预示了福音的奥秘；\newline{}它们使圣祖得以看见并相信那后来宗徒所要注视和宣讲的。\newline{}因为，既然法律是未来之事的阴影，那所见到的阴影就是投下它的实体的真实轮廓。\newline{}天主曾以人的形态被看见、相信并朝拜，祂确实要在时期圆满时生而为人。\newline{}祂为了圣祖的肉眼而取了一个预象，这预象预示了未来的真理。\newline{}在那旧日，天主只是以人的形态被看见，并非为人所生；\newline{}到了预定时期，祂不仅被看见，而且真正为人所生。\newline{}人对祂所取、为使人能看见祂的这人形外貌的熟悉，是为预备他们迎接祂真正为人所生的时刻。\newline{}那时，阴影成了实体，外表成了真实，景象成了生命。\newline{}但天主始终不变，无论祂是在人形的外表下被看见，还是真正为人所生。\newline{}祂诞生后的自己与祂在景象中被看见的自己之间，有着完美的相似。\newline{}正如祂所诞生的那样，祂也曾显现；正如祂曾显现的那样，祂也诞生了。\newline{}但是，既然时机尚未到来，让我们将福音的记述与梅瑟先知的记述相比较，\newline{}就让我们继续在法律的书卷中走我们选定的道路。\newline{}今后我们将从福音中证明，真正的人子是真正的天主之子；\newline{}目前，我们从法律中证明：以人形显现给圣祖的，是真正的天主，即天主之子。\newline{}因为当一位以人形显现给亚巴郎时，祂被当作天主朝拜，并被宣告为审判者；\newline{}并且当主从主那里降下雨火时，毫无疑问，法律告诉我们主从主那里降下雨火，\newline{}是为了启示给我们父与子。\newline{}我们也不能有片刻认为，当圣祖明知而以天主朝拜子时，他竟看不清他所朝拜的是真天主。\par

18. 然而，不虔敬的不信者很难领会真信仰。\newline{}他们虔敬的能力从未因信仰而扩展，过于狭窄，无法接受真理的完整呈现。\newline{}因此，不信的灵魂不能理解天主所做的伟大工程：\newline{}降生成人，以完成人类的救恩；\newline{}在救恩的工程中，他们看不到天主的能力。\newline{}他们想到祂诞生的痛苦、婴孩的软弱、童年的成长、成年的成熟、身体的苦难以及随之而来的十字架，还有十字架上的死亡；\newline{}这一切都使他们看不见祂真正的天主性。\newline{}然而，祂为自己创造了所有这些能力，作为祂本性的增添；\newline{}这些能力是祂真实的天主性所不具有的。\newline{}因此，祂获得了这些能力而没有失去祂真实的天主性，\newline{}并且当祂成为人时，祂并没有停止是天主；\newline{}祂，永恒的天主，在时间中的某个时刻成了人。\newline{}他们看不到真天主的能力表现在：祂成了以前所不是的，却从未停止是祂以前的自己。\newline{}然而，若不是祂以自己全能的性体力量，在保持原有的同时，成为以前所不是的，\newline{}那么我们的软弱本性就不会被接纳。\newline{}异端是何等的盲目，世俗的智慧是何等的愚蠢，\newline{}竟看不见基督的耻辱是天主的德能，信仰的愚拙是天主的智慧！\newline{}所以，在你们眼中，基督不是天主，\newline{}因为那从永远就有的，竟然诞生了；那不变的，竟随着年岁增长；那不受苦的，竟然受难；那永活的，竟然死去；那死了的，竟然活着；\newline{}因为祂的全部历史都与自然的普遍进程相矛盾！\newline{}这岂不就是说：祂既是天主，便是全能的？\newline{}圣洁而可敬的福音，我还没有翻开你的书页，从其中证明基督耶稣在这些变化和苦难中是天主。\newline{}因为法律是福音的先驱，法律必须教导我们：当天主披上软弱时，祂并未失去祂的天主性。\newline{}法律的预象是我们对福音信仰奥秘的有力确证。\par

19. 请以你忠信的精神与我同在，圣洁而蒙福的雅各伯宗主教，以驳斥不信之蛇的毒嘶。愿你再次在与那人的角力中得胜，并因更强而再次恳求祂的祝福。为何祈求你本可从软弱对手那里强求的事？你强壮的手臂已征服了你祈求祝福的那位。你身体的胜利与你灵魂的谦卑形成鲜明对比，你的行为与你的思想截然不同。在你有力的掌握中，那人看似软弱无力；但在你眼中，这人是真实的天主，并非名义上而是本性上的天主。你所请求的不是义子天主的祝福，而是真实天主的祝福。你与人对峙，却与天主面对面相见。你肉眼所见与信德之眼所见大不相同。你感到祂是软弱的人；但你的灵魂得救，因为你看到祂身上的天主。当你角力时，你是雅各伯；如今你因所祈求祝福的信德而成为以色列。按肉体说，那人比你卑微，作为祂在肉身中受难的预像；但你却能在那软弱的肉身中认出天主，作为圣神内祝福的记号。眼目所见并不动摇你的信德；祂的软弱并未误导你忽视祂的祝福。虽然祂是人，但祂的人性并不妨碍祂是天主，祂的天主性并不妨碍祂是真实的天主；因为，既为天主，祂必真是真实的天主。\par

20. 法律在其进程中仍遵循福音奥秘的次序，而福音是法律的本体；法律的预像是宗徒所教导真理的忠实预期。在梦境的神视中，蒙福的雅各伯看见了天主；这是奥秘的启示，而非肉身的显现。因为向他显示了天使藉梯子降下，并升天，而天主安息于梯子之上；这神视按其解释，预言了他的梦境终将成为显明的真理。宗主教的话\Quote{天主的住所和上天之门}向我们展示了他神视的场景；然后，在对他所作所为的长篇叙述之后，叙述这样继续：\Quote{天主对雅各伯说：起来，上贝特耳去，住在那里；在那里给天主筑一座祭坛，就是当你躲避你哥哥厄撒乌时，曾显现给你的那位。}\BibleRef{创 35:1} 如果福音的信德通过天主圣子得以接近天主圣父，并且只有通过天主才能领悟天主，那么请向我们表明，这位要求人敬畏天主、安息于天上梯子之上的，在何种意义上不是真实的天主？什么本性的差异将两者分开，当两者都持有指示同一本性的同一名字时？所看见的是天主；也是那一位天主谈论所看见的天主。除非通过天主，不能领悟天主；同样，除非通过天主，天主不接受我们的敬拜。我们无法理解那一位必须被尊崇，除非另一位教导我们尊崇祂；我们无法知道那一位是天主，除非我们知道了另一位的天主性。奥秘的启示按既定的次序进行；我们藉由天主才被引入对天主的敬拜。当同一名字——表明同一本性——将圣父与圣子结合时，圣子怎能如此低于自身，以致不是真实的天主呢？\par

21. 人的判断不可对天主下论断。我们的本性并非如此，以致能凭自身力量提升自己以默观天上之事。我们必须从天主学习有关天主的思考；我们除了祂自己，别无知识之源。你可以尽你所能地接受世俗哲学的训练；你也可以过正直的生活。这一切都有助于你心智的满足，但并不能帮助你认识天主。梅瑟曾作为王后的义子被收养，并受过埃及人一切智慧的教育；此外，出于对本族的忠诚，他杀了埃及人为希伯来人报仇，但他仍不认识祝福他祖先的天主。因为当他因害怕自己的行为败露而离开埃及，在米德扬地做牧羊人时，他看到荆棘中的火焰，荆棘却没有烧毁。那时他听到了天主的声音，求问祂的名字，并领悟了祂的本性。这一切，除非通过天主自己，他无从得知。同样，我们论及天主时，必须局限于祂向我们的理智启示有关祂自己的那些言词。\par

22. 是天主的天使在荆棘的火焰中显现；也是天主在火焰中的荆棘里发言。祂显现为天使；这是祂的职分，而非祂的本性。表明祂本性的名字赐予你为\Quote{天主}；因为天主的天使就是天主。但也许祂不是真实的天主？那么，亚巴郎的天主、依撒格的天主、雅各伯的天主，难道不是真实的天主吗？因为从荆棘中发言的天使是他们的天主，永远如此。为免你暗示这名字只是祂义子的名号，是绝对的天主向梅瑟说话。祂的话是这样：\Quote{上主对梅瑟说：我是自有者；又说：你要这样对以色列子民说：那\InnerQuote{自有者}打发我到你们这里来。}\BibleRef{出 3:14} 天主的话语开始时是天使的发言，为要揭示在圣子内人类救恩的奥秘。接着祂显现为亚巴郎的天主、依撒格的天主、雅各伯的天主，使我们知道祂本性的名字。最后，是那\Quote{自有者}天主打发梅瑟到以色列那里去，使我们完全确信在绝对意义上祂是天主。\par

23. 亵渎的疯狂之虚妄愚昧还能杜撰出什么进一步的虚构？当所有族长们的知识都反驳你时，你是否仍执着于在纯洁的麦子中夜间撒稗子，注定要被焚烧？不仅如此：如果你相信梅瑟，你也会相信天主，天主子；除非你碰巧否认梅瑟所说的是祂。如果你想否认这一点，你必须听天主的话：——\BibleQuote{因为如果你们相信梅瑟，你们也会相信我，因为他论及我写了书。}\BibleRef{若 5:46} 是的，梅瑟将用整卷法律书反驳你，这部法律是藉天使经中保之手而颁布的。请问，那颁布法律者难道不是真天主吗？因为中保就是颁布者。难道梅瑟领百姓到山上去不是为了迎接天主吗？难道不是天主降临到山上吗？或者，那做这一切的，祂被称为天主，只是出于虚构或义子名分，而非出于本性的权利？听那号角的响声，火把的闪烁，如同从熔炉中冒出的烟云笼罩山顶，人在天主面前自觉无能的恐惧，百姓的承认——他们恳求梅瑟作他们的代言人，因为听到天主的声音他们就会死。在你看来，祂不是真天主吗？仅因害怕祂说话，以色列人就充满了死亡的恐惧？祂的声音是人的软弱所不能承受的？在你眼中祂不是天主，因为祂藉着人的软弱官能向你说话，好使你听见而生活？梅瑟上了山；在四十昼夜中，他获得了天上奥秘的知识，并按在那里向他启示的真理异象，安排了这一切。他与天主交谈，从与祂同在的那一位那里，领受了那无人能注视的荣耀的反照？祂那可朽坏的面容被转变成了与他同在的那一位不可接近之光的相似。这位天主是他所见证的，这位天主是他所谈论的；他召叫天主的天使们在异民的喜乐中来朝拜祂，并祈求那令祂喜悦的祝福降在若瑟的头上。面对这样的证据，谁敢说祂只有天主的名号，而否认祂的真天主性？\par

24. 我相信，这番长篇讨论已揭示出真理：在法律论及天主与天主、主与主的那些段落中，从未提出过任何合理的论证来支持一位是真天主而另一位不是真天主的区别。我已经证明，这些术语与二者在名号或本性上的差异不相容，并且我们可以用名号作为本性的试金石，以本性作为名号的线索。因此，我表明了天主性、权能、属性、名号都内在于法律称为天主的祂身上。我也表明了，法律逐渐揭示了福音的奥秘，通过显示天主在创造世界时顺从另一位天主的言语，在造人时创造了天主和天主的共同肖像，从而揭示了圣子是一个位格；此外，在对索多玛人的审判中，主是从主而来的审判者；在赐福和制定法律奥秘时，天主的天使就是天主。因此，为了支持救恩性的宣认，即天主始终在父和圣子的位格中显示自己，我们已表明法律如何通过严格使用天主的名称来教导真天主性；因为法律清楚地说他们是两位，但对他们每一位的真天主性不投下丝毫怀疑的阴影。\par

25. 现在是我们制止异端这一狡猾伎俩的时候了，他们藉此将法律中虔诚而敬畏天主的教导曲解为支持他们自己不虔敬的妄想。他们否认天主子，却先引用说：\BibleQuote{以色列，听！上主你的天主是唯一的}；然后，因为同名的身份会反驳他们的亵渎（因法律也谈到天主与天主），他们又援引先知的话\BibleQuote{他们要称颂你，真天主}，来证明这名号并非在真实意义上使用。他们争论说，这些话教导天主是唯一的，而天主、天主子只有名号而无本性；因此我们必须得出结论：真天主只有一个位格。但是，愚昧的人啊，你也许以为我们会反驳你的这些经文，从而否认有唯一真天主？我们当然不会按照你的理解来宣认而反驳它们。我们的信德接受它们，我们的理性接纳它们，我们的言语宣告它们。我们承认一位天主，而且是真天主。天主的名称对我们的宣认毫无危险，因为我们的宣认宣告在圣子的本性中有着唯一真天主。学习你自己的话的含义，承认唯一真天主，然后你才能忠实地宣认唯一而真实的天主。我们信仰的言语被你变成了亵渎的工具，你保全了声音却歪曲了意义。你穿着想象智慧的愚拙外衣，打着忠于真理的幌子，却是真理的毁灭者。你承认天主是唯一而真实的，目的却是要否认你所承认的真理。你的言语以不虔敬为虔敬，以虚假为真实。你对唯一真天主的宣讲导向了对祂的否认。因为你否认圣子是真实的天主，虽然你承认祂是天主，但只是名号上的天主，而非本性上的。如果祂的诞生是名号上的而非本性的，那么你否认祂对这名号的真实权利就情有可原了；但若祂是真实地作为天主而生，那么祂怎能不因祂的诞生而是真天主呢？否认事实，你就可以否认后果；如果你承认事实，祂又怎能不同于祂自身？没有一个存在物能改变自己的本质本性。关于祂的诞生，我稍后会讲；现在，我要用先知的话语来驳斥你们关于祂真天主性的亵渎谎言。但我要注意：在我们主张唯一真天主时，我既不给撒贝里乌异端留有任何余地，即认为父与子同一个位格，也不给你们从唯一真天主的合一中推导出的对圣子真天主性的诽谤留有任何余地。\par

26. 亵渎与智慧不兼容；哪里没有敬畏天主——智慧的开端，哪里就没有丝毫理智的光芒。异端者引用先知的话\Quote{\Emph{他们必称颂祢，真实的天主}}作为反对圣子天主性的证据，可见一斑。首先，我们看到愚昧在不信中阻碍了对预言前半部分的理解（或若理解了，则压制它）；再次，我们看到他们欺诈地插入了一个在经卷中找不到的小字。这种做法既愚蠢又不诚实，因为没有人会如此信任他们，以至于不经核对预言原文就接受他们的读法。因为那段经文并非如此写：\Quote{\Emph{他们必称颂祢，真实的天主}}，而是\Quote{\Emph{他们必称颂真实的天主}}。\EditorialNote{依撒意亚 65:16} \Quote{\Emph{祢，真实的天主}}与\Quote{\Emph{真实的天主}}之间并非毫无区别。若保留\Quote{\Emph{祢}}，第二人称代词表示正在对另一个说话；若省略\Quote{\Emph{祢}}，则\Quote{\Emph{真实的天主}}成为句子的宾语，即说话者本人。\par

27. 为确保我们对这段经文的解释完整而确凿，我全文引述这些话：——\Quote{\Emph{因此上主这样说：看，事奉我的必得吃喝，但你们却要饥饿；看，事奉我的必得畅饮，但你们却要干渴；看，事奉我的必欢喜快乐，但你们却要因心里忧愁而哀哭，因心灵痛苦而悲嚎。因为你们将把你们的名号留给我选民作为喜乐，但上主必要杀害你们。但我的仆人要被称为一个新名，这名字必在地上蒙福；他们必称颂真实的天主，在地上起誓的也必指着真实的天主起誓。}}任何偏离惯用表达方式都是有充分理由的，但新奇之处也为误解提供了机会。这里的问题是，既然先前关于天主已有如此多的预言，而\Quote{\Emph{天主}}这个名称单独不加修饰便足以表明神圣的威严和本性，为什么预言之神如今要藉依撒意亚预告，\Emph{真实}的天主将被称颂，人在地上要指着\Emph{真实}的天主起誓？首先，我们必须记住这段话是关于未来时代的。那么我问：在犹太人心中，那曾经被称颂、被指着起誓的，难道不是真实的天主吗？犹太人不知其奥秘的预表意义，因此不认识天主圣子，只简单地敬拜天主为天主，而非为父；因为若他们敬拜祂为父，也必敬拜圣子。所以他们称颂的是\Quote{\Emph{天主}}，指着起誓的也是\Quote{\Emph{天主}}。但先知证实，从今以后被称颂的是\Quote{\Emph{真实}的天主}；他称祂为\Quote{\Emph{真实的天主}}，因为祂降生成人的奥秘将使有些人对祂真实的天主性视而不见。当谬说将要被传扬时，真理必须被清楚宣告。现在让我们逐句回顾这段经文。\par

28. \Quote{\Emph{因此上主这样说：看，事奉我的必得吃喝，但你们却要饥饿；看，事奉我的必得畅饮，但你们却要干渴。}}请注意，一个子句中包含两种不同时态，为要教导关于两个不同时期的真理；\Quote{\Emph{事奉我的必得吃喝}}。现世的虔诚将获得未来的赏报，同样，现世的邪恶将遭受未来饥渴的惩罚。接着祂说：\Quote{\Emph{看，事奉我的必欢喜快乐，但你们却要因心里忧愁而哀哭，因心灵痛苦而悲嚎。}}这里与先前一样，启示了未来和现在。如今事奉祂的将欢喜快乐，而不事奉的则要在忧愁和痛苦中哀哭悲嚎。祂继续说：\Quote{\Emph{因为你们将把你们的名号留给我选民作为喜乐，但上主必要杀害你们。}}这些话是关于未来的，是对属血肉的以色列说的，它被讥讽将来不得不将其名号交给天主的选民。这名号是什么？当然是以色列；因为预言是对以色列说的。现在我问：今天什么是以色列？宗徒给出了答案：那些按精神而不按文字行事、遵行基督法律的人，就是天主的以色列。\BibleRef{罗 2:29}\par

29. 此外，我们必须得出结论：为何前述经文\Quote{\Emph{因此上主这样说}}之后紧接着\Quote{\Emph{但上主必要杀害你们}}，以及下一句\Quote{\Emph{但我的仆人要被称为一个新名，这名字必地上蒙福}}是什么意思。毫无疑问，\Quote{\Emph{因此上主这样说}}以及后来的\Quote{\Emph{但上主必要杀害你们}}都证明，那位说话并定意要杀害、要以那个新名赏报祂的仆人、以先知的言辞而被知晓、并将审判义人和恶人的，正是同一位主。因此，福音奥秘启示的其余部分消除了关于那位主是说话者也是杀害者的任何疑惑。它继续说：\Quote{\Emph{但我的仆人要被称为一个新名，这名字必地上蒙福。}}这里一切都是未来时。那么，这宗教的新名是什么？一个地上蒙福的名字？若在过去曾有\Quote{\Emph{基督徒}}这个名称的祝福，它就不是新名。但若这个我们对天主的虔敬圣名是新的，那么这授予我们信德的\Quote{\Emph{基督徒}}新称号，便是我们在世上所得的属天祝福。\par

30. 现在，有些话语与我们信德的内在确证完全和谐。祂说：\BibleQuote{他们必祝福真天主，凡在地上起誓的，必指真天主起誓}。的确，那些在天主服务中领受新名的人必祝福天主；而且他们所起誓所凭的天主是真天主。对于这位真天主是谁，人们要凭祂起誓、要祝福祂，藉着祂，一个崭新而蒙福的名字要赐给事奉祂的人，还有什么疑问呢？我这边，为反对异端亵渎性的歪曲，拥有教会信德清晰而明确的证据：即祢、基督、所赐的新名，以及祢因忠信事奉所赏报的蒙福称号。它确证祢是真天主。每一个信者的口，啊，基督，都宣告祢是天主。所有信者的信德都确证祢是真天主，承认、宣告、内在确证祢是真天主。\par

31. 因此，这段预言，连同其全部上下文，清楚地将祂描述为天主：既是我们因新名之事奉的祂，也是那新名在地上得蒙祝福所通过的那位。它告诉我们，谁蒙祝福为真天主，谁被起誓为真天主。这正是教会，在时期圆满时，对其主基督所做的信德的宣认。我们可以看出，预言的话语如何与真理完全一致，因为它们避免了插入第二人称代词。如果话语是\BibleQuote{祢，真天主}，那么它们可能被解释为对另一位说的。\BibleQuote{真天主}只能指说话者本人。这段文字，单独看，指明了它所指的那一位；前面的语词，与之相连，宣告了做出对天主宣认的那位说话者是谁。它们是这样的：——\BibleQuote{我向那不询问我的人显现，我从那不寻求我的人那里被找到。我对那不呼求我名的民族说：我在这里。我全天向一个不信和抗辩的民族伸出我的手}。难道还能有比这更彻底地暴露压制真理的不诚实企图，或更清晰地显示说话者就是真天主吗？请问，谁向那不询问祂的人显现，被那不求祂的人找到？什么民族先前不呼求祂的名？谁全天向一个不信和抗辩的民族伸出祂的手？将这些话与\BibleBook{}{申命纪}中那神圣之歌\BibleRef{申 32:21}相比，其中天主，因对不是神的神们发怒，激动不信者对那非民族和愚昧民族嫉妒。你自己来判断，谁使祂自己显于那不认识祂的人；谁，虽然一个民族属于祂自己，却成了外邦人的产业；谁在不信和抗辩的民族前伸出祂的手，将从前针对我们的判词的记录钉在十字架上\BibleRef{哥 2:14}。因为同一位圣神在先知中——我们现在所思考的——在同一个预言的过程中如此继续，该预言在论证上连贯，在话语上也连续：——\BibleQuote{但我的仆人要由一个新的名字称呼，这名在地上要被祝福，他们要祝福真天主，凡在地上起誓的要指真天主起誓}。\par

32. 如果异端以其愚昧和邪恶，试图引诱那些单纯和未受过教导的人远离真信仰，即这些话是指天主圣子而说的，而假称它们是天主圣父关于祂自己的言谈，那么它必听到外邦人的宗徒和导师对那谎言的判决。他解释所有这些预知都是指向主的苦难和福音信仰的时期，当他责备以色列的不信，即以色列不愿承认主成了肉身来到。他的话是：——\BibleQuote{因为凡呼求主名的，必得救。但他们怎能呼求他们所不信的祂呢？但他们怎能在没有听过祂的情况下信祂呢？没有传道者，他们怎能听见呢？除非他们被派遣，他们怎能传道呢？正如经上所记：那传报和平、传报美事者的脚是多么美丽啊。但并不是所有人都服从福音。因为依撒意亚说：主，谁信了我们的报告？所以信德来自听闻，听闻来自那话语。但我问：他们没有听过吗？是的，他们的声音传遍了全地，他们的话语到达了地极。但我问：以色列难道不知道吗？首先，梅瑟说：我要以那非民族激起你们的嫉妒，以愚昧的民族惹动你们的怒气。再者，依撒意亚放胆说：我向不寻求我的人显现，被不问我的人找到。但对于以色列，祂说了什么？我整天向一个不听命的民族伸出我的手}\BibleRef{罗 10:13}。你是谁，曾穿越层层诸天，不知道是在身内还是身外，竟能比这位更忠实地解释先知之言？你是谁，曾听见但不可道出天上隐秘之事不可言喻的奥秘，并以更大的确信宣告天主为启示而赐给你的知识？你是谁，曾被预定全然分享主在十字架上的苦难，并首先被提到乐园，从天主的圣经中汲取了比这蒙选的器皿更高贵的教训？如果有这样的人，难道他不知道这些都是真天主的作为和言语，是由祂自己真实而蒙选拣的宗徒向我们宣讲的，好让我们在其中认出祂是作者？\par

33. 但有人可能会辩称，宗徒借用这些先知的话时，并没有受到先知神恩的感动；他只是随意解释别人的话，虽然宗徒自己所说的一切无疑都是来自基督的启示，但他对依撒意亚话语的了解仅限于书本。我的回答是：在那段宣告的开头，其中说到真天主的仆人将祝福祂并指着祂起誓，我们读到先知这样的朝拜：\Quote{从永远我们未曾听见，我们的眼目也未曾看见天主，除了祢和祢将为等候祢仁慈的人所行的作为。} \EditorialNote{依撒意亚 64:4} 依撒意亚说他除了祂以外没有见过任何天主。因为他确实看见了天主的荣耀，他预言了天主从童贞女取肉身的奥秘。如果你在你的异端中不知道先知在那荣耀中看见的是天主独生子，请听福音者的话：\Quote{依撒意亚说了这些话，是因为他看见了祂的荣耀，并且指着祂说的。} \BibleRef{若 12:41} 宗徒、福音者、先知联合起来堵住你的反对。依撒意亚确实看见了天主；虽然经上写着\Quote{从来没有人见过天主，只有那在父怀里的独生子，祂将祂彰显出来}，但先知所见的是天主。他凝视着天主的荣耀，人们因如此荣耀赐予他先知的伟大而心生嫉妒。这就是犹太人对他判处死刑的原因。\par

34. 因此，那在父怀里的独生子将无人见过的天主告诉我们。要么反驳圣子这样告诉了我们的事实，要么就相信那位被看见的、向不认识祂的人显现、成为那不呼求祂的外邦人的天主、并向悖逆的百姓伸出双手的那一位。也要相信关于祂的这件事：侍奉祂的人将被称为新名，在地上人们将祝福祂并指着祂作为真天主起誓。预言告知，福音证实，宗徒解释，教会宣认：那被看见的是真天主；但没有人敢说圣父被看见过。然而，异端的疯狂已到了如此地步，以至于他们虽然声称承认这真理，实际上却否认了它。他们通过一种新奇且不敬神的诡计来否认它，即逃避真理，同时刻意假装坚持它。因为他们宣认独一天主是唯一真实、唯一公义、唯一智慧、唯一不变的、唯一不死的、唯一大能的，他们却给祂加上一个本质不同的圣子，不是由天主所生为天主，而是通过受造而被收养为子，拥有天主的名并非出于本性，而是作为由义子而得的称号；这样他们就不可避免地剥夺了圣子所有那些他们堆积在孤独威严的圣父身上的属性。\par

35. 异端扭曲的心智无法认识和宣认独一真天主；宣认所需的健全信德和理智与不信是不相容的。我们必须宣认父与子，然后才能领悟天主是独一而真实的。当我们认识了人类救恩的奥秘，这奥秘通过重生的能力在我们内得以完成，使我们进入父与子内的生命，那时我们才有希望洞察法律和先知的奥秘。对福音者和宗徒教导的不信神的无知，无法形成独一真天主的思想。我们将根据福音者和宗徒的教导，提出关于祂的正确道理，与真信徒的信德精确一致。我们将如此呈现祂：那出于父本质的独生子，在性体上不可分、不可分离，而非在位上不可分。我们将陈述天主是独一的，因为天主来自天主的性体。但我们也要以先知的话为根据，确立这完全合一天主的道理，并使它们成为福音结构的根基，证明只有一个天主，一个天主性体，因为天主独生子从未被单独列为第二位神。因为在我们这部论文的整卷书中，我们遵循了与前一卷相同的路线；在那一卷中证明圣子是神的方法，在此证明了祂是真天主。我相信对每一段经文的解释已经如此令人信服，以至于我们现在已经像先前证明祂的神性一样有效地显明了祂是真天主。本书的其余部分将致力于证明：那现在被承认为真天主的，绝不应被视为第二位神。我们反驳第二位神的概念将进一步确立合一；并且这真理将被显示为与圣子的位格存在不相矛盾，同时它维护了天主内性体的合一。\par

36. 我们探究的正确方法要求我们从他开始，天主藉着他首次向世界显现自己，即从梅瑟开始，天主独生子藉梅瑟的口如此宣告：\Quote{看，看，我是天主，除我以外没有别的神。} \BibleRef{申 32:39} 那不敬神的异端不能将这些话归于未生的天主父，这一点从经文的意义和我们已陈述过的宗徒的见证来看是清楚的，宗徒教导我们应将这整段话理解为由天主独生子所说。宗徒还指出\Quote{外邦人哪，你们当与祂的百姓一同欢乐}这些话是圣子的，并进一步引证了这话：\Quote{从叶瑟的根上必生出一位，兴起治理外邦；外邦人将要仰望祂。} 因此，梅瑟藉着\Quote{外邦人哪，你们当与祂的百姓一同欢乐}这话指明那说\Quote{除我以外没有别的神}的一位；而宗徒将同一段话引用到我们的主耶稣基督、天主独生子身上，祂按肉体从叶瑟的根兴起为王，外邦人的希望在于祂。因此，我们现在必须考虑这些话的含义，以便我们这些知道这些话是由祂所说的人，可以查明祂是在什么意义上说的。\par

37. 那真实、绝对而圆满的道理，构成我们的信德，就是宣认天主出自天主、天主在天主内，不是藉身体过程，而是藉天主的能力；不是藉性体注入性体，而是藉同一天主性体隐秘而强有力的运作：天主出自天主，不是藉分裂、伸展或流溢，而是藉一性体的运作，藉诞生而产生一个与自己同一的性体。这些事实将在下一卷中得到更充分的处理，那一卷将专用于阐述福音作者和宗徒的教导；目前，我们必须藉法律和先知来维护我们的主张和信仰。天主由之而生的性体必然与祂的根源同一。祂不能以非天主的方式存在，因为祂的来源惟独来自天主。祂的性体是同一的，并非指生者亦为受生者——因为那样无生者若受了生，就不再是祂自己——而是说受生者的本质包含一切在生者本质中所总括的元素，而生者是祂惟一的根源。因此，祂从独一者而来，祂的存在分享独一性，并非出于外因；祂内没有新元素，因为祂的生命来自那生活的；也没有缺乏元素，因为那生活的生了祂，使祂分享祂自己的生命。因此，在圣子的生发中，无形无变化的天主依照祂自己的性体，生了无形无变化的天主；这无形无变化的天主完美地从无形无变化的天主诞生，正如我们在天主出自天主的启示中所看见的，不涉及生者本质的减少。因此，天主独生子藉圣\Person{梅瑟}作证：\Quote{看，看，我是天主，除我以外没有别的神。}因为没有第二个天主性，所以祂以外不可能有别的神，因为祂是天主，但藉着祂性体的能力，天主也在祂内。因为祂是天主，天主在祂内，所以祂以外没有别的神；因为天主，除祂以外没有别的神性根源，就在祂内，因此在祂内不仅有祂自己的存在，还有那存在的本源。\par

38. 我们所宣认的这救恩信德，由预言之神所维系，它藉众多口舌同声发言，在漫长多变的世代中从未对启示真理作过模糊的见证。例如，我们得知藉梅瑟由天主独生子所说的话，为使我们更好受教，由预言之神通过那些先知伟人再次证实：\Quote{因为天主在你内，除你以外没有别的神。}让异端以绝望和狂怒的最大努力猛扑这名字与性体不可分地结合的宣告，并以其狂暴挣扎若能撕裂那在称号与事实上都完美的合一。天主在天主内，除祂以外没有别的神。异端者若能，将内在的天主与祂所在的祂分开，并按各自种类将那奥秘结合的成员分类。因为当祂说\Quote{天主在你内}时，祂教导天主圣父的真性体临在于天主圣子内；因为我们必须理解，那\BibleRef{出 3:14}的\Quote{我是自有者}的天主就在祂内。当祂加上\Quote{除你以外没有别的神}时，祂指出在祂以外没有别的神，因为天主住在祂自己内。而第三项断言\Quote{你是天主，我们却不认识}，为我们的教导指出了虔诚而有信德的灵魂必须作出何种宣认。当它得知了神圣诞生的奥秘，以及天使向\Person{若瑟}宣告的\OriginalTerm{厄玛奴耳}{Emmanuel}之名时，它必须呼喊：\Quote{你是天主，我们却不认识，以色列的天主，救主。}它必须承认天主性体在祂内的自立存在，因为天主在天主内，并且除了真实者以外没有任何别的神。因为，祂是天主，天主在祂内，任何其他神明的幻觉，不论何种，都必须放弃。这就是\Person{依撒意亚}先知的信息；他为圣父与圣子不可分割、不可分离的天主性作证。\par

39. 同样受默感的先知\Person{耶肋米亚}也教导，天主独生子与天主圣父是同一性体。祂的话是：\Quote{这是我们的天主，没有别的可以比拟祂；祂发现了全部知识的道路，并将它赐给祂的仆人\Person{雅各伯}和祂的爱人\Place{以色列}；以后祂在地上显现，在人间居住。}为什么试图将天主子变为第二位神？学习认识并宣认唯一天主。没有第二位神能与基督相比，因此能自称为天主。祂是天主出自天主，凭性体和诞生，因为祂天主性的来源是天主。而且，祂不是第二位神，因为没有别的可与祂比拟；祂内的真理不是别的，正是天主的真理。为什么在假装对天主的合一的敬拜中，将真与假、卑贱与真实、不相像与不相像联结在一起？圣父是天主，圣子是天主。天主在天主内；除祂以外没有别的神，也没有别的能与祂相比而成为神。如果你在这两位中认出天主的合一，而非孤立，你将有分于教会的信德，这信德在圣子内宣认圣父。但如果你对上天奥秘无知，坚持天主是独一以强调祂孤立的说辞，那么你便对天主的认识是外行，因为你否认天主在天主内。\par

\SourceRecord{Translated by E.W. Watson and L. Pullan.；Nicene and Post-Nicene Fathers, Second Series；Vol. 9.；Edited by Philip Schaff and Henry Wace.；Buffalo, NY: Christian Literature Publishing Co.,；1899.；<http://www.newadvent.org/fathers/330205.htm>.}

% structure: p0001 source=h1 target=section reason=page-title

\section{\WorkTitle{论天主圣三（第六卷）}}

\Person{希拉略}一开始便哀叹亚略异端的广泛传播；他对灵魂的爱驱使他与这异端作战，其阴险特性使之更加危险（第1至4节）。他在第5、6节中重复了第四卷中提到的同一亚略派信经。异端者通过谴责与自身不相符的谬误而获得正统的外观；这种谴责旨在使公教信仰被怀疑与这些谬误有牵连。因此，他必须推迟对新约的引用，直到他审查了这些谬误（第7、8节）。相应地，在第9至12节中，他依次解释了瓦伦廷、摩尼、撒伯流和希拉加斯的学说，并指出教会拒绝所有这些学说，如同她拒绝（第13节）亚略派在其信经中错误地归于她的学说一样。他们的目的是否认圣子与圣父同永恒并与祂同一性体（第14、15节）；但这种否认完全违背圣经，违背圣经是亵渎（第16、17节）。亚略派要把基督当作受造物（第18节），而在第19至21节中，\Person{希拉略}转向祂，激动地宣称祂确是真天主。然后他继续论证，从福音所记载的圣父和圣子的话语，证明基督是生而为子，而非被收纳为子（第22至25节）。这一点由福音对其行为的记载所证实（第26、27节），否则这些行为无法解释。通过讨论若望福音7:28-29和8:42（第28至31节），论证更加有力。基督真子身份进一步由宗徒们的信德证明，他们的确信随着知识增长而增强（第31至35节），尤其由圣伯多禄（第36至38节）、圣若望（第39至43节）和圣保禄（第44、45节）的见证所证明。拒绝如此重大的见证就是选择假基督而非基督（第46节）。此外，我们还有祂施行奇迹的对象、魔鬼、犹太人、在海上遇险的宗徒、十字架旁的百夫长的见证，证明基督确实是天主子（第47至52节）。\par

我充分了解这时代的危险与激情，才敢攻击这狂野而无神的异端，它宣称天主子是受造物。罗马帝国几乎每个行省中的众多教会，早已染上这致命学说的瘟疫；执着灌输的错误，虚假地声称是真理，已深深嵌入那些虚妄地以为自己忠于信仰的心灵。我深知，当对错误事业的热情受到人数众多的鼓励并得到普遍赞同的批准时，意志是多么难以彻底悔改。受迷惑的群众只能带着困难和危险接近。当人群迷途时，即使他们知道错了，也羞于回头。他们要求考虑其人数，并厚颜地命令将其愚昧视为智慧。他们假定其规模是其观点正确性的证据；因此，一个已被普遍相信的谎言被大胆地宣称已确立了其真理。\par

就我而言，不仅是我的圣召对我的要求，即作为主教我亏欠于教会勤勉宣讲福音的责任，驱动了我。我写作的热忱随着受这异端理论危害和束缚的人数增加而增长。想到许多可能得救的人，如果他们能认识对天主的正确信仰的奥秘，放弃人类愚昧的亵渎原则，离开异端者而投身于天主；如果他们能放弃捕鸟者设饵诱捕的诱饵，在自由和安全中翱翔，以基督为首领、先知为教师、宗徒为向导，并在对圣父和圣子的宣认中接受完美的信仰和确定的救恩，那将是极大的喜乐前景。如此，他们将要服从主的这话：\Emph{\BibleQuote{谁不尊敬子，就是不尊敬派遣祂来的父}}\BibleRef{若 5:23}，通过尊敬子来尊敬父。\par

因为近来一种致命的邪恶感染已蔓延至人类中，其蹂躏处处带来毁灭和死亡。城市及其居民突然被地震夷为平地，不断发生的战争带来可怕的屠杀，无法抵御的瘟疫造成广泛的死亡，这些都未曾像这异端在世界的进展那样造成如此致命的伤害。因为天主，众死者都为祂而活，只毁灭那些自取灭亡的人。对于那将要审判万民、其威严将仁慈地调节分配给无知错误的惩罚的祂，那些否认祂的人连审判也指望不上，只能指望否认。\par

4. 因为这种疯狂的异端确实否认；它借用我们宣认中的语句来否认真实信仰的奥秘，并将这些语句用于其不敬神的目的。这种谬误信仰的宣认，我在前一本书中已经引用过，开头如下：——我们宣认唯一的天主，独一非受造、独一永恒、独一\OriginalTerm{无始}{unoriginate}、独一真实、独一拥有不朽、独一良善、独一全能。因此，他们摆出我们自己宣认的开头语句，即\Quote{唯一的天主，独一非受造、独一\OriginalTerm{无始}{unoriginate}}，以便这貌似真理的表相能作为他们亵渎性添加的引子。因为，在大量同样不真诚地表达对圣子虔敬的词语之后，他们的宣认继续道：\Quote{天主完美的受造物，但并非像祂的其他受造物一样；祂的工艺，但并非像祂的其他作品一样。}接着，在间隔一段插入真确陈述以掩盖他们不敬神意图——即他们试图通过诡辩证明祂是从虚无中产生的——之后，他们又加上：\Quote{祂在万世之前受造并被确立，在祂出生之前并不存在。}最后，似乎他们错误教义的每一点——即祂既不应被视为圣子也不应被视为天主——都已坚不可摧地防御了攻击，他们继续道：——\Quote{至于诸如\Emph{出自祂}、\Emph{从母腹中}、\Emph{我出自父并来到}这些说法，如果被理解为表示圣父延伸出这一实体的一部分以及某种发展，那么按照他们的说法，圣父将具有\OriginalTerm{复合本性}{compound nature}、可分、可变、有形体；这样，就他们的言辞而言，无形体的天主将屈从于物质的属性。}但是，既然我们现在要再次覆盖整个领域，这次使用福音的语言作为武器对抗这最不敬神的异端，似乎最好在第六卷中重复整个异端文件，尽管我们在第四卷中已经提供了完整副本，以便我们的对手可以再次阅读它，并逐点与我们的回答比较，从而被福音作者和宗徒的清晰教导所迫，尽管不情愿和好争辩，也承认真理。异端的宣认如下：——\par

5. 我们宣认唯一的天主，独一非受造、独一永恒、独一\OriginalTerm{无始}{unoriginate}、独一真实、独一拥有不朽、独一良善、独一全能、创造者、命定者和安排者，不变、不改、正义、良善，法律和先知以及新约的天主。我们相信，这一位天主在万世之前生了\OriginalTerm{独生圣子}{Only-begotten Son}，藉着祂创造了世界和万物，祂不是貌似地，而是真实地生了祂，随从自己的旨意，使祂不变、不改、天主完美的\OriginalTerm{受造物}{Creature}，但并非像祂的其他受造物一样，祂的\OriginalTerm{工艺}{Handiwork}，但并非像祂的其他作品一样；不像\Person{瓦伦廷}所坚持的，圣子是圣父的\OriginalTerm{发展}{development}，也不像\Person{摩尼}所宣称的关于圣子的，是圣父的\OriginalTerm{同质}{consubstantial}部分，也不像\Person{撒伯里乌}，他从一个中造出两个，同时是圣子和圣父，也不像\Person{希拉克斯}，是光中之光，或一盏有两团火焰的灯，也不像祂先前存在而后出生，或被重新创造成为圣子——这种观念经常被您，有福的教宗，在教会中和弟兄的集会上公开谴责。但是，正如我们已肯定的，我们相信祂是由天主的旨意在时间和万世之前受造，并从圣父获得生命和存在，圣父赐予祂分享自己荣耀的完美。因为，当圣父将万有的产业赐予祂时，并未因此剥夺了祂那些无始的属性，祂是万有的根源。\par

6. \Quote{所以有三位：圣父、圣子和圣神。天主是万有的原因，完全无始，与一切分离；而圣子由圣父在时间之外产生，在万世之前受造并被确立，在祂出生之前并不存在，但在时间之外、万世之前出生，作为独一圣父的独一圣子而存在。因为祂既非永恒，也非与圣父\OriginalTerm{同永恒}{co-eternal}，也非与圣父\OriginalTerm{同非受造}{co-uncreate}，也没有与圣父并存的存在，像某些人主张的两个无始原理那样。但天主是在万有之前，作为不可分和万有的开始。因此祂也在圣子之前，正如我们从您公开讲道中所学到的。因此，既然祂从天主获得祂的存在、祂荣耀的完美和祂的生命，并被托付万有，所以天主是祂的根源。因为天主统治祂，作为祂的天主，因为祂在祂之前。至于诸如\Emph{出自祂}、\Emph{从母腹中}、\Emph{我出自父并来到}这些说法，如果被理解为表示圣父延伸出这一实体的一部分以及某种发展，那么按照他们的说法，圣父将具有\OriginalTerm{复合本性}{compound nature}、可分、可变、有形体；这样，就他们的言辞而言，无形体的天主将屈从于物质的属性。}\par

7. 谁能不在这里看到蛇的黏滑踪迹：蜷缩的毒蛇，力量集中准备扑击，将毒牙的致命武器隐藏在它的盘绕中？我们马上要把它展开并检查，暴露这隐藏头颅的毒液。因为他们的计划是先用某些正确的陈述留下印象，然后注入他们异端的毒液。他们对我们说话温和，以便暗中伤害我们。然而，在他们所有动听的表白中，我哪里也听不到天主的儿子被称为天主；我从未听到儿子身份归于圣子。他们说了很多关于祂拥有圣子之名的话，却从未提及祂拥有圣子的本性。这一点被隐藏了，使祂似乎甚至没有权利拥有这个名字。他们假装揭露其他异端，以掩盖他们自己是异端者的事实。他们大力主张只有一位、独一真实的天主，目的是剥去天主子真实而属位的天主性。\par

8. 因此，虽然我在前两卷中已根据法律和先知的教导证明：必须承认天主与天主、真天主与真天主、真天主圣父与真天主圣子，是因本性合一而非位格混淆而成为唯一的真天主，但为了完整地陈述信仰，我还必须援引福音作者和宗徒的教导。我必须据此证明：天主子，真天主，并非与圣父具有不同或异类的本性，而是通过真正的诞生拥有同一神性，同时具有独立的存在。事实上，我不认为有任何灵魂如此愚钝，竟会幻想：虽然我们认识天主的自我启示，却无法理解它们；或者以为，如果它们能被理解，人就不愿理解；又或者梦想人的理性能够对它们进行改进。然而，在我开始讨论这些救恩奥秘所包含的事实之前，我必须先降伏那些异端者用以指责其他异端名称的骄傲。我要揭露这种为他们自身不敬虔所披的巧妙外衣。我要表明：这种掩盖其教导致命性的手段，反而使其暴露和出卖自身，并且是一个广泛有效的警示，让人认清这种蜜毒的真实性质。\par

9. 例如，这些异端者主张天主子不是出自天主；天主并非出自并在天主的本性中由天主所生。为此，当他们庄严地见证\Quote{独一真天主}时，他们避免加上\Quote{圣父}。然后，为了通过否认真正的诞生来逃避承认圣父与圣子的唯一真神性，他们接着说：\Quote{并非如瓦伦提努所主张的那样，圣子是圣父的一种发展。}因此，他们试图通过称之为\Quote{发展}来诋毁天主出自天主的诞生，仿佛这是瓦伦提努异端的一种形式。因为瓦伦提努是污秽而愚蠢幻想的作者；除了首神之外，他还发明了一整个神祇家族和无数被称为永世的权能，并教导我们的主耶稣基督是通过一种隐秘的意志行动神秘地产生出来的发展。教会的信仰，福音作者和宗徒的信仰，对这种来自鲁莽无脑的梦想家头脑的想象发展一无所知。它不知道瓦伦提努的\Quote{深渊}和\Quote{沉默}以及三十个永世。它只知道独一天主圣父，万物出于祂，并独一耶稣基督，我们的主，万物藉着祂，祂是由天主所生的天主。但是，他们想到：祂作为天主从天主而生，既没有从祂根源的神性中取走任何东西，祂自己也不是作为非天主而生；祂成为天主，不是通过神性的新开始，而是通过从已有的天主诞生；而且，每一种诞生，就人的能力所能判断的，似乎都是一种发展，因此那诞生本身也可以被视为一种发展。这些考虑促使他们攻击瓦伦提努关于发展的异端，以此作为摧毁圣子真正诞生信仰的手段。因为日常生活的经验使世俗智慧认为，诞生和发展之间没有太大区别。人的心智迟钝而缓慢，难以把握天主的事物，需要经常提醒它我已多次陈述的原则：从人类经验中得出的类比并不能完全适用于神圣权能的奥秘；它们的唯一价值在于，这种与物质对象的比较使精神获得一种对天上事物的概念，使我们能藉着自然之梯上升到对天主威严的理解。但是，天主的诞生绝不能用人类诞生中的那种发展来判断。当一位从一位而生，天主从天主而生时，人类诞生的环境使我们能够理解这一事实；但是，以交合、受孕、时间和分娩为前提的诞生，绝不能向我们提供神圣方式的线索。当我们被告知天主从天主而生时，我们必须接受祂确实是出生的这一事实，并对此感到满足。然而，我们将在适当的地方讨论神圣诞生的真理，正如福音和宗徒所陈述的那样。我们现在的责任是揭露这异端者的巧妙计谋，这种攻击基督真正诞生的行为，是隐藏在攻击所谓发展的形式之下的。\par

10. 接着，在继续这种对信仰的欺诈性攻击时，他们的宣认这样写道：\Quote{也非像摩尼所宣称的，圣子是圣父的一个\OriginalTerm{同体}{consubstantial}部分。}他们已否认了圣子是发展出来的，以避免承认祂的诞生；如今他们又引入以摩尼之名标记的教义，即圣子是唯一天主实体的一部分，并否认它，以颠覆对出自天主的信仰。因为摩尼，那位法律和先知的狂暴敌人、魔鬼事业的热烈拥护者和太阳的盲目崇拜者，教导说：在童贞女胎中的那一位是唯一天主实体的一部分，而圣子必须理解为天主实体的某一片段，被切割下来，并在肉身上显现。于是他们便大肆利用这种异端，即在圣子的诞生中发生了唯一实体的分割，并以此作为手段来规避独生子的诞生教义，以及实体统一的名称本身。因为谈论由唯一实体分割而来的诞生是纯粹的亵渎，所以他们否认任何诞生；所有形式的诞生都被列入他们对摩尼式分割诞生的谴责之中。再者，他们也废除了实体统一，既废其名也废其实，因为异端者认为统一是可分割的；并且否认圣子是出自天主的天主，拒绝相信祂真正拥有天主性。这疯狂的异端为何要假装一种虚假的敬畏、一种无谓的焦虑？教会的信仰，正如这些疯狂的谬误提倡者提醒我们的那样，确实谴责摩尼，因为她不知道圣子是一个部分。她知道祂是全部出自全部的天主，是一体出自一体，不是被分割而是被生的。她确信天主诞生既不使生子者贫乏，也不使被生者低劣。如果这是教会自己的想象，那就以虚假智慧的愚昧来责备她；但如果她是从她的主那里学来的，那就承认被生者知道祂被生的方式。她是从独生天主那里学到这些真理的，即圣父与圣子原为一体，并且在天主性全部圆满居于圣子内。因此她憎恶将圣子视为唯一实体的一部分；因为她知道祂确是由天主所生，她便敬拜圣子为真正天主性的合法拥有者。但现在，让我们暂缓对这些个别指控的充分答复，并快速通过其余谴责。\par

11. 接下来是这样：\Quote{也非像撒伯里乌，他从一之中造出两个，即同时是圣子和圣父。}撒伯里乌持有这种观点，是出于对福音作者和宗徒启示的故意盲目。但我们在这里看到的，并非一个异端者诚实地谴责另一个异端者。促使他们责备撒伯里乌分割不可分位格的，正是希望不留任何圣父与圣子之间的结合点；这种分割并不导致第二位格的诞生，而是将一位格切为两部分，其中一部分进入童贞女的胎中。但我们承认诞生；我们拒绝将两个位格混淆于一个之中，同时我们仍坚持天主统一。也就是说，我们认为\Quote{出自天主的天主}意味着本质的统一；因为那一位，通过出自天主的真实诞生而成为天主，只能从天主那里获得祂的实体。既然祂继续从天主那里获得祂的存在，正如祂起初获得的那样，祂必定永远是真天主；因此二者是一体，因为出自天主的天主别无其他本质，只有天主性，也别无其他起源，只有天主起源。但这里之所以谴责这种亵渎性的撒伯里乌式两个位格混淆为一个，是因为他们想剥夺教会对一位天主中两个位格的真实信仰。但现在我必须审视这种扭曲才智的其余实例，以免被视为一个挑剔的、出于厌恶而非真正恐惧而评判诚实探寻者的法官。我将通过他们结束宣认的措辞，展示他们精心设计出的致命结论及其必然结果。\par

12. 他们的下一句是：\Quote{也不是像\Person{希耶拉加斯}所说的那样，光从光，或一盏灯有两道火焰，也不是好像祂先前已存在，而后才被生或受造成为子。}\Person{希耶拉加斯}无视独生子的诞生，完全不明白福音启示的意义，却谈论一盏灯的两道火焰。这对称的两道火焰，由同一盏灯碗内供给的油所滋养，正是他对圣父与圣子实体的比喻。好像那实体是独立于每一位之外的某种东西，如同灯里的油，与两道火焰不同，尽管火焰的存在依赖油；或者如同灯芯，整根是同一材料，两端燃烧，与火焰不同，却提供火焰并将它们连接在一起。这一切不过是人的狂妄的虚幻，它信赖自己而非天主以获得知识。但真信仰断言：天主由天主而生，如光从光，它倾注自身而不自我减损，给予所拥有的，却仍拥有所给予的。它断言，藉着祂的诞生，祂就是祂所是的，因为祂如何，祂就如何诞生；祂的诞生是现有生命的恩赐，这恩赐并未减少它所取用的源头；而圣父与圣子为一，因为祂所从出的那一位，与祂相似，而出生的那一位，既没有另一个源头，也没有另一种本性，因为祂是光从光。正是为了将人的信仰从这真实教义中引开，\Person{希耶拉加斯}的这盏灯或这火焰才被掷向那些承认光从光的人。因为这句话曾被用于异端意义，并且现今和早先都受过谴责，他们就想说服我们，认为这句话没有可以使用的真实意义。让异端立刻抛弃这些无根据的恐惧，别再借着用如此不诚实的方式赚来的热心名声，以教会信仰的保护者自居。因为我们不容许任何有形体的、无生命的事物在天主的属性中有一席之地；凡是天主的，都是完美的天主。祂内只有能力、生命、光、福乐、圣神。那本性不含任何愚钝的物质元素；因其不变，它内没有不一致。天主，因为祂是天主，是不变的；不变的天主生了天主。他们结合的联系，并不像两道火焰、同一盏灯的灯芯那样，是某种自身之外的东西。独生子由天主的诞生，不是空间上的延长，而是一个生；不是延伸，而是光从光。因为光与光的合一乃是本性的合一，并非无间断的连续。\par

13. 异端的巧思何等奇妙：\Quote{也不是好像祂先前已存在，而后才被生或受造成为子。}天主，既然祂由天主所生，当然并非从无而生，也不是从不存在之物而生。祂的诞生乃是永远生命本性的诞生。然而，虽然祂是天主，祂并不等同于先存的天主；天主由存在于祂之前的天主所生；在祂的诞生中，并且藉着祂的诞生，祂分享了祂源头的本性。如果我们说的是自己的话，这一切不过是不敬；但如果我们（正如我们将要证明的）天主亲自教导了我们如何说话，那么我们就有必要按照天主所启示的意义来承认这神圣的诞生。正是圣父与圣子中本性的合一，这活生生的诞生的不可言喻的奥迹，正是异端的疯狂企图从信仰中驱逐的，当它说：\Quote{也不是好像祂先前已存在，而后才被生，或受造成为子。}现在，谁愚蠢到以为圣父不再是祂自己；那先前已存在的同一者，后来才被生，或受造成为子？或者，以为天主消失了，祂的消失之后在诞生中出现，而事实上，那诞生正是其作者持续存在的证据？或者谁狂妄到以为子不经诞生就能存在？谁如此缺乏理性，竟说天主的诞生产生了别的什么，而不是天主被生？常存的天主并非被生，但天主从常存的天主被生；那诞生中所赋予的本性正是生者自身的本性。而天主藉着祂的诞生（这诞生是从天主到天主），因为祂是真正的诞生，所领受的不是新受造之物，而是天主永远拥有的、过去和现在都存在的事物。因此，被生的不是先存的天主；然而，天主被生，并开始存在，出于天主的属性并与天主的属性同在。由此可见，异端在这漫长的前奏中，是如何诡诈地引向这最亵神的教义。其目的是否认天主独生子，它以看似为真理辩护开始，进而断言基督不是从天主而生，而是从无而生，并且祂的诞生是由于天主从不存在之物受造的神圣计划。\par

14. 然后，经过一段旨在预备我们迎接后续内容的间隔之后，他们的异端发起了这样的攻击——\Quote{圣子是在时间之外被发出、在被造和建立于诸世界之前，祂在受生之前并不存在。}这\Quote{祂在受生之前并不存在}是异端用以自夸能达成两个目的的一种措辞：支持其亵渎，并为自己的教义在被质问时提供遮掩。支持其亵渎，因为如果祂在受生之前不存在，祂就不能与永恒的本源同质。祂必须从无中开始，如果除了与祂受生同时的能力外别无能力的话。遮掩其异端，因为如果这种说法被谴责，它便提供现成的回答。他们会说，已经存在的不可能受生；既然已存在，祂不可能通过经历受生的过程而进入存在，因为受生恰恰意味着受生者进入存在。愚昧和亵渎者！谁会在那未受生和永恒的天主身上梦想受生？我们怎能认为那\Emph{自有者}\BibleRef{出 3:14}的天主受生，而受生意味着受生的过程？你正试图否认天主独生子从天主圣父的受生，你的目的是通过这个\Quote{祂在受生之前并不存在}来逃避承认那真理：天主——天主子由此受生——永恒存在，而天主子正是从祂的永恒本性中通过受生而获得存在。那么，如果圣子是从天主受生，你必须承认祂的受生是那永恒本性的受生；不是预先存在之天主的受生，而是天主独生子从预先存在之天主的受生。\par

15. 但这个异端的狂热激愤如此强烈，以致它无法抑制自己的暴烈发作。为努力证明——符合其\Quote{祂在受生之前并不存在}的主张——圣子是从非存在中受生的，即祂并非通过真实而完美的受生从天主圣父受生成为天主子，它最终在愤怒和仇恨中将其宣认推升至可能亵渎的最高点：——\Quote{至于诸如\Emph{从祂而来}、\Emph{从母腹而出}、以及\Emph{我由圣父而来并来到了}这些短语，如果它们被理解为意指圣父延伸了那同一实体的一部分，并可以说是一种发展，那么按他们的说法，圣父将是复合的、可分的、可变的、有形体的本性；并且这样，就他们的话而言，无形体的天主将受制于物质属性。}为真信德辩护以驳斥异端谎言，确实是一项辛苦而困难的任务，如果我们必须追随他们的思想过程，直到他们沉入无神的深渊。但幸而对我们有益的是，正是思想的浅薄孕育了他们亵渎的热忱。因此，虽然反驳愚昧容易，但纠正愚昧之人却难，因为他既不会自己思考出正确的结论，也不会接受他人提供的正确结论。然而我深信，那些因虔诚的无知而并非因自负滋生的故意愚昧而被错误束缚的人，将会欢迎纠正。因为我们对真理的演示将提供令人信服的证明：异端无非是愚昧。\par

16. 你在你的愚昧中说了，并且今天仍在重复，不知道你的智慧是对天主的挑战，\Quote{至于诸如\Emph{从祂而来}、\Emph{从母腹而出}、以及\Emph{我由圣父而来并来到了}这些短语，}我问你，这些短语是天主的话，还是不是？它们当然是的；而且，既然是天主关于祂自己说的，我们必须完全照其所说接受它们。关于这些短语本身，以及每个短语的确切含义，我们将在适当的地方谈论。目前我只向每位读者的理智提出这个问题：当我们看到\Emph{从祂自身而来}时，我们是应当把它等同于\Quote{从别处而来，}还是\Quote{从无而来，}还是作为真理接受它？它不是\Quote{从别处而来，}因为它是\Emph{从祂自身而来}；也就是说，祂的天主性除天主外别无来源。它不是\Quote{从无而来，}因为它是\Emph{从祂自身而来}；这是对祂受生所出自之本性的宣告。它不是\Quote{祂自己，}而是\Emph{从祂自身而来}；这是声明他们作为圣父和圣子的关系。接下来，当启示\Emph{从母腹而出}被给出时，我问，我们能否可能相信祂是从无受生的，而祂受生的真理已清楚借用身体功能的术语表明？不是因为祂有身体的肢体，天主才用这些话记录圣子的受生：\Emph{在晓星之前，我从母腹中生了祢}。祂使用语言帮助我们的理解，以向我们保证祂的独生子是以不可言喻的方式从祂自己的真实天主性中受生的。祂的目的是通过用描述人性境况的言语来披戴神圣真理，从而教育人的能力以达到信德的知识。因此，当祂说\Emph{从母腹而出}时，祂教导我们，祂的独生子在神圣意义上受了生，并非通过从无中受造而进入存在。最后，当圣子说\Emph{我由圣父而来并来到了}时，祂是否留下了怀疑余地，即祂的天主性是否源自圣父？祂从圣父而来；也就是说，祂有一次受生，并且是圣父，而非别的，给了祂那次受生。祂作证，祂宣称从中而来的那位，是祂存在的本源。这一切的证明和解释将在后面给出。\par

17. 然而，让我们看看这些人究竟有何根据，竟敢自信地禁止我们接受天主关于祂自身的启示；这些启示的真实性是连他们也不否认的。人的愚昧和狂妄，还能对天主的自我启示投下比这更严重的侮辱吗？那就是以纠正的方式加以谴责。因为连怀疑和批评都不足以满足他们。还有什么比这种对天主的本质和权能进行亵渎的谈论和争辩更严重的呢？还有什么比妄自断言：如果圣子来自天主，那么天主就是可变的、有形的，因为祂已延伸或发展了自己的一部分成为祂的圣子？这种急于证明天主不变性的焦虑从何而来？我们明认圣子的受生，我们宣示独生子，因为天主就是这样教导我们的。而你们，为了从教会的信仰中铲除受生和独生子，却用一个不变的天主来对抗我们，声称祂的本性不能延伸或发展。我甚至可以举出世界上某些本性中的受生例子，来驳斥这种可悲的妄想，即认为每一次受生都必然是一种延伸。我本可以避免你们犯这样的错误，即认为一个存在者只能以生育者的损失为代价而产生，因为有许多生命传递的例子，并不需要身体上的交合。但是，当天主已经发言时，再用证据来论证便是亵渎；否认祂赐予我们信仰的权威，乃是极端的疯狂，因为我们的敬拜本身就是承认唯独祂能赐予我们生命。因为如果生命唯独来自祂，那么作为生命条件的信仰，岂不也应由祂来创立？如果我们认为祂关于祂自身的见证不可信，我们又怎能确信祂所赐予的生命呢？\par

18. 因为你们，最不虔敬的异端者，将圣子的受生归因于一个创造性的意志行为；你们说祂不是从天主所生，而是由创造者的选择所创造和存在的。而你们所解释的天主性合一，不允许祂是天主，因为既然天主保持为独一，圣子就不能在祂所进入的那个状态中保留祂的原始本性。你们说，祂通过创造而被赋予了一种与神性不同的实体，虽然从某种意义上说，祂是独生子，故优于天主的其他受造物和工程。你们说，祂被高举起来，为的是祂反过来可以执行托付给祂的任务，即举起受造的世界；但祂的受生并没有将神性赐予祂。按照你们的话，祂是受生的，意思是祂是从无中产生的。你们称祂为子，不是因为祂从天主所生，而是因为祂由天主所创造。因为你们想起，天主甚至认为圣人也配得上这个称号，你们认为圣子被赋予这一称号的方式，与\BibleQuote{我曾说过：你们是神，都是至高者的儿子}这句话完全一样；也就是说，祂是因赐予者的俯就而承受这名号，而非出于本性的权利。因此，在你们眼中，祂是义子，因称号而成为神，因恩宠而为独生子，因时间而为首生者，在各方面都是受造物，绝无神性。因为你们认为，祂的生育不是从天主而生的自然受生，而是一种受造实体生命的开端。\par

19. 全能的天主，现在我首先必须祈求祢宽恕我过度的愤慨，并允许我向祢陈辞；其次，祈求祢恩赐我这个尘土和灰烬般的人，却仍以忠诚的孝爱依附祢，在这场辩论中获得发言的自由。曾有一段时间，我这个可怜虫并不存在；在我的生命、意识和位格开始存在之前。我生命的存在全赖祢的仁慈；我毫不怀疑，祢出于美善，赐给我出生是为了我的益处，因为祢并不需要我，绝不会使我生命的开端成为邪恶的起点。然后，当祢将生命的气息吹入我内，并赋予我思维的能力时，祢借着祢的仆人梅瑟和先知们赐给我们的圣卷，教导我认识祢自己。我从中学到祢的启示，即我们不可敬拜祢为一位孤独的天主。因为这些书卷教导我，天主并非在本质上与祢不同，而是在奥秘的实体合一中与祢为一。我得知祢是在天主内的天主，不是通过混合或混淆，而是借着祢的本性，因为祢自身的神性居于那从祢而来的一位内。但完美受生的真确教义启示：祢——被内居者，和祢——内居者，并非同一位格；然而祢确实居于那从祢而来的一位内。而福音作者和宗徒们的声音不断重复这一教训，祢独生子的圣口所说的话也被记录下来，告诉我们：祢的圣子，独生子天主，从祢——无生者天主——而来，为了完成我救恩的奥迹，由童贞女降生为人；祢因着祂从祢而来的真实受生而居于祂内，祂也因着祂永恒受生中所赐的本性而居于祢内。\par

20. 这绝望的错误泥沼，祢将我投了进去，这究竟是什么呢？因为我已学到这一切，并且开始相信；这信德如此深植于我的心中，我既无力也无愿改变。为什么用关于祢自己的谎言来欺骗一个不幸的人，毁灭这可怜的灵魂和肉身？在红海被分开之后，梅瑟从山上下来时脸上的光辉欺骗了我。他在祢面前曾注视天上的一切奥秘，我相信了他由祢所默示的关于祢自己的话。还有达味，那按祢心意而找到的人，他把我引向了毁灭；撒罗满，被认为配得上天主智慧之恩赐的；依撒意亚，他看见了万军的上主并预言；耶肋米亚，在母腹中成形之前就被祝圣，成为列国先知，要被拔除和栽种；厄则克耳，复活奥秘的见证者；达内尔，蒙爱的人，知道时期；以及全体神圣的先知团队；还有玛窦，被选来宣讲整个福音奥秘，先是税吏，后成为宗徒；若望，主的密友，因此配得揭示天上最深的奥秘；蒙福的西满，在宣认奥秘之后被立为教会的基石，并得到了天国的钥匙；以及他所有因圣神而发言的同伴；保禄，特选的器皿，从迫害者变为宗徒，他作为活人留在深海之下\BibleRef{格后 11:25}，并升到三层天上，在殉道之前已在天堂，他的殉道是完美无瑕信德的完美祭献；所有人都欺骗了我。\par

21. 这些人教导了我所持守的道理，他们的教导如此深深地浸透了我，以致没有任何解毒剂能使我摆脱他们的影响。请宽恕我，全能的天主，我无力改变，愿意死在这信德中。这些亵渎的传播者——在我看来如此——是末世的产物，太新颖了，不能帮助我。要纠正我从祢所接受的信德，为时已晚。在我听到他们的名字之前，我已信靠祢，由祢重生，成为祢的，如今仍是。我知道祢是全能的；我并不指望祢向我启示那在祢与祢的独生子之间隐秘、不可言说的诞生奥秘。对祢而言，没有不可能的事，我毫不怀疑在生祢的圣子时，祢运用了祢全部的全能。怀疑这一点就是否认祢是全能的。因为我自己的出生教导我祢是良善的，所以我确信在生祢的独生子时，祢没有吝惜任何美物。我相信凡属于祢的，都属于祂；凡属于祂的，都属于祢。世界的创造足以向我证明祢是智慧的；我确信祢的智慧，与祢相似，必是由祢所生。祢是唯一天主，在我眼中千真万确；我永不相信在祂——由祢而来的天主——内，有丝毫不属于祢的。请在祂内审判我，如果通过祢的圣子，我过于信任法律、先知和宗徒是罪过的话。\par

22. 但这些狂言必须停止；揭露异端愚妄的修辞必须让位于构建论证的苦工。所以，我相信，那些他们中能得救的人，将转向福音作者和宗徒所教导的真信德，并承认祂是真正的天主子，不是因义子名分，而是因本性。因为我们的回应计划必须是：首先证明祂是天主子，因此完全拥有那构成其子性的天主性。因为我们正在考虑的异端的主要目的，就是否认我们的主耶稣基督是真天主并真正是天主子。许多证据确保我们确信，我们的主耶稣基督是，并启示为，天主独生子，真正是天主子。祂的圣父为此作证，祂自己宣称，宗徒们宣告，信徒相信，魔鬼承认，犹太人否认，外邦人在祂受难时认出了祂。天主的名称是以绝对主权赐予祂的，并非因为祂被接纳与其他拥有这称号者共享。基督的每一言行都超越那些被称为儿子者的能力；我们从祂生命中最突出的一切所学到的首要教训，就是祂是天主子，并且祂并非将圣子之名作为一个与广泛友人群体共享的头衔来持有。\par

23. 我不会以自己言语的混杂来削弱这真理的证据。让我们听圣父，当耶稣基督的圣洗完成时，祂多次论及祂的独生子所说的话，为使我们不至被祂有形可见的身体误导，而不认识祂是圣子。祂的话是：— \BibleQuote{这是我的爱子，我所喜悦的}。\BibleRef{玛 3:17} 这里的真理是以模糊的轮廓呈现的吗？宣告是以犹豫不决的声调发出的吗？由天使从圣神带来的童贞生子的应许，贤士的明星指引，在马槽中向祂致敬，由洗者若翰所证实的屈尊受洗者的尊威——这一切都被视为不足以见证祂的光荣。圣父亲自从天上说话，祂的话是：\BibleQuote{这是我的圣子}。这证据意味着什么？不是名号，而是代词？名号可以随意附加在名字上；代词却确凿地指明其所指的人。而我们在这里有\Emph{这个}和\Emph{我的}，最清晰的指示。请留意这些话语的真正含义和目的。你已经读到：\BibleQuote{我生了儿女，将他们抚养长大}；\EditorialNote{依撒意亚1:2}但你在那里并没有读到\BibleQuote{我的儿女}，因为祂是藉分散在外邦人中并从祂产业的人民中为自己生了那些儿女。并且免得我们以为\BibleQuote{圣子}这个名号不过是加给天主独生子的一个额外称号，藉以表示祂因义子之名在某种共同产业中有份，祂的真实本性却由那个给予确凿所有权意义的代词表达出来。我允许你，如果你愿意，将\BibleQuote{圣子}一词解释为基督是众多中的一个，如果你能举出一个例子，在关于那众多中的另一个时，曾说过\BibleQuote{这是我的圣子}。反之，如果\BibleQuote{这是我的圣子}是祂特有的称号，那么当圣父断言祂的所有权时，你为何控告祂提出没有根据的要求呢？当圣父说\BibleQuote{这是我的圣子}时，我们难道不能这样阐述祂的意思吗：\Quote{祂将\InnerQuote{儿子}的名号给了别人，但祂自己就是我自己的儿子；我藉义子之名将这名字赐给众多的人，但这圣子是我自己的。不要寻求另一位，免得你失去\InnerQuote{这就是祂}的信德。藉手势和声音，藉\InnerQuote{这个}、\InnerQuote{我的}和\InnerQuote{圣子}，我将祂宣告给你们。}如今还有什么合理的借口可以用于不信呢？正是这，而且正是这，才是圣父的声音所宣告的。祂愿意我们不致对那来受洗以满全一切义德者的本性一无所知；使我们藉天主的言语，认出那以可见的人形显现者是天主子，以完成我们救恩的奥迹。\par

24. 因为信徒的生命与这信德的宣认息息相关——因为除了确信耶稣基督——天主独生子——是天主子之外，没有通往永生的其他道路——宗徒们再一次听到天上重复同一讯息的声音，为加强这赐予生命的信仰，否定它就是死亡。当主身穿荣光，站在山上，梅瑟和厄里亚在祂身旁，还有被选为异象和声音之真理见证的教会三大柱石时，圣父就这样从天上说：— \BibleQuote{这是我的爱子，我所喜悦的，你们要听从祂！}\BibleRef{玛 17:5} 他们所看见的荣光不足以证明祂的尊威；那声音宣告：\BibleQuote{这是我的圣子}。宗徒们不能面对天主的荣光；凡俗的眼目在它面前变得昏暗。伯多禄、雅各伯和若望的信心丧失，他们俯伏在地，恐惧不已。但这来自圣父知识的庄严宣告，解除了他们的重负；祂被启示为祂父自己的真圣子。而且，除了\Emph{这个}和\Emph{我的}对祂真实圣子身份的见证之外，还说出了\BibleQuote{你们要听从祂}。这是圣父从天上所作的见证，以确认圣子在地上所作的见证；因为我们奉命听从祂。虽然圣父对圣子的这一认识消除了所有疑惑，但我们仍然奉命接受圣子的自我启示。当圣父的声音命令我们以听从祂来显示我们的服从时，我们被命令对圣子的话语寄予绝对的信任。因此，既然圣父在要我们听从圣子的这个讯息中显明了祂的旨意，就让我们听听圣子关于祂自己告诉了我们什么。\par

25. 我想不到有任何人会如此缺乏普通理智，以致在各福音中认出圣子因取了肉身而屈尊的宣认——例如祂屡次说的话：\Quote{父啊，光荣我吧！}，以及\BibleQuote{你们要看见人子}\BibleRef{玛 26:64}，和\BibleQuote{父比我大}\BibleRef{若 14:28}，以及更强烈的：\Quote{现在我的心灵极其烦乱}，甚至还有：\BibleQuote{我的天主，我的天主，你为什么舍弃了我？}\BibleRef{玛 27:46}，以及许多其他经文，我将在适当的时候谈论它们——然而，面对这些不断表达的谦卑，却指责祂因称天主为祂的父而僭越，例如祂说：\Quote{凡不是我天上的父所栽种的植物，必要连根拔起}，或\BibleQuote{你们不要使我父的殿宇成为商场}\BibleRef{若 2:16}。我想不到会有人愚蠢到认为祂一贯主张天主是祂的父，不是出于确定知识的真诚告白，而是一种大胆而无根据的主张。我们不能指责这种不断宣认的谦卑是蛮横地要求他人的权利，是染指不属于祂的东西，是攫取唯有天主才能行使的能力。同样，当祂称自己为子时，例如：\Quote{因为天主没有派遣子到世界上来审判世界，而是为叫世界藉着他而获救}，以及\Quote{你信天主子吗？}我们也不能指责祂犯了与称天主为祂的父相同的僭越。然而，如果我们仅仅允许耶稣基督享有义子的名称，这岂不是一种指控？当祂称天主为祂的父时，我们岂不是指控祂大胆地作无根据的主张？圣父的声音由天上说：\Quote{你们要听从祂}。我听见祂说：\BibleQuote{父啊，我感谢你}\BibleRef{若 11:41}，以及\Quote{你们说我亵渎了，因为我说\InnerQuote{我是天主子}吗？}如果我不能相信这些名号，并假定它们意味着它们所声称的，那么我该如何信任和理解呢？没有给出其他含义的暗示。圣父由天上作证：\Quote{这是我的圣子}；圣子则谈到\Quote{我父的殿宇}和\Quote{我父}。当信仰在问题中被要求时，那个名号的宣认给予救恩：\Quote{你信天主子吗？}代词\Emph{我的}表示随后的名词属于说话者。我问你们这些异端者，有什么权利作相反的假设？你们违背了圣父的话和圣子的断言；你们掏空了语言的意义，扭曲了天主的言语，使它们不能承受的含义。只有你们担当这无耻亵渎的罪责，即天主关于祂自己说了谎。\par

26. 因此，虽然唯有真诚相信这些名号具有真实意义——即当我们读到\Quote{这是我的圣子}和\Quote{我父}时，这些话确实指明了所谈及和所对话的位格——才能使它们变得可理解，然而，为避免有人以为\Emph{子}和\Emph{父}只是义子与尊贵的称号，让我们看看圣子自己赋予祂的圣子之名什么属性。祂说：\BibleQuote{我父将一切交给了我；除了父外，没有人认识子；除了子和子所愿意启示的人外，也没有人认识父。}\BibleRef{玛 11:27} 我们正在谈论的话，即\Quote{这是我的圣子}和\Quote{我父}，是否与\Quote{除了父外，没有人认识子；除了子和子所愿意启示的人外，也没有人认识父？}一致？因为只有藉着相互作证，子才能透过父被认识，父也透过子被认识。我们听见来自天上的声音；我们也听见圣子的话。我们不认识子的藉口，与不认识父的藉口一样少。一切都被交付给祂；这个\Emph{一切}没有例外。如果祂们拥有同等的权能；如果祂们共享同等的、对我们隐藏的相互认识；如果父和子这些名号表达了祂们之间的关系，那么，我请问，祂们岂不在事实上正如在名号上一样，是同一全能的行使者，被笼罩在同一不可测透的奥秘中吗？天主说话不是为了欺骗。圣父的父职和圣子的子职是字面上的真理。现在来看事实如何印证这些名号所启示的真理。\par

27. 子这样说：\BibleQuote{因为父交给我要完成的工程，就是我所行的这些工程，为我作证：是父派遣了我；那派遣我的父也为我作了证。} \BibleRef{若 5:36} 天主独生者不仅以名字，也以能力证明他的子性；他所行的工程证明他是父所派遣的。请问，这些工程证明的事实是什么？就是他受了派遣。\Quote{受了派遣}用作证明他子性的服从和父的权威：因为除了那由父所派遣的以外，无人能行他所行的工程。然而，他的工程证据未能使不信的人相信父派遣了他。因为他接着说：\BibleQuote{那派遣我的父也为我作了证；你们从来没有听过他的声音，也没有见过他的形状。} 父关于他的见证是什么？翻阅福音书，审查其内容。请给我们读出父所作的其他见证，除了我们已经听过的以外：\BibleQuote{这是我的爱子，我所喜悦的，你是我的儿子。} 若翰听了这些话，并不需要它们，因为他已知道真理。父说话是为了教导我们。但这还不止。若翰在旷野中得到了这一启示；宗徒们也不应被拒绝同样的保证。这保证以同样的言语临到他们，但加上了若翰没有领受的部分。若翰从母胎中已是先知，并不需要\Quote{你们要听从他}这条命令。是的；我将听从他，而且只听从他和他的宗徒，宗徒是为教导我而听的。即使圣经没有包含父为子所作的进一步见证，只说他就是子，为了确认真理，我也有他所作的父的工程为证据。这现代的诽谤是什么？说他名字是义子的恩赐，天主性是谎言，称号是虚假？我们有父为他子性的见证；子藉着与父同等的工作，见证自己与父的平等。为何如此盲目，看不见他显然拥有真实子性，他既宣称又显示出来？基督拥有子的名号，不是因为天主父屈尊恩赐；也不是因圣德赢得此称号，如同许多人因在信仰宣认中忍受苦难而赢得。这样的子性不是本有的；而是天主藉着一项与自己相称的恩惠赐予称号。但由\Emph{这个}、\Emph{我的}和\Emph{听从他}所指明的，与其他的在性质上不同。那是真实、实在、真正的子性。\par

28. 而且，事实上，子从未为自己要求低于父所赐予他的这称号的内涵。父的话\BibleQuote{这是我的儿子}启示他的本性；接下来的\BibleQuote{你们要听从他}是召唤我们聆听他从天堂带来的奥秘和信仰，以便学习：如果我们想要得救，我们的宣认必须肖似他的教导。同样，子亲自用自己的话教导我们，他确实是生和确实来——\BibleQuote{你们不认识我，也不知道我从哪里来，因为我不是由自己来的，但那派遣我的是真实的，你们不认识他，我却认识他，因为我来自他，是他派遣了我。} \BibleRef{若 7:28} 没有人认识父；子常向我们保证这一点。他说除了他自己以外无人认识父，原因是他来自父。请问，他来自父，是由于受造行动的结果，还是由于真实的生？如果是受造行动，那么一切受造物都来自天主。那么，当子说他认识父是因为他来自父时，为何它们都不认识父？如果他受造而非受生，我们就会看到他与其他来自天主的存有相似。既然按照这假设，一切都来自天主，为何他不像其他存有那样对父无知？但如果对父的认识是他特有的，他来自父，那么\Quote{他来自父}这一事实岂不也是他特有的吗？也就是说，他岂不是必须是从天主性所生的真子吗？因为他唯独认识天主的原因是他唯独来自天主。你看，一种特有的认识，是由于一种特有的生。你承认，他来自父，不是因创造的能力，而是通过真实的生；而这就是为何唯有他认识父，父为一切其他来自他的存有所不认识。\par

29. 但祂立即补充说：\BibleQuote{因为我出于祂，祂派遣了我}，以杜绝异端妄自假定祂出于天主是从祂降世之时开始的。福音对于奥秘的启示依逻辑顺序进行；首先祂被生，然后祂被派遣。同样，在之前的声明中，我们被告知无知：首先是关于祂是谁，然后是祂从何而来。因为\BibleQuote{我出于祂，祂派遣了我}这句话包含两个独立的陈述，同样\BibleQuote{你们既不认识我，也不知道我从哪里来}这句话也是如此。每一个人都是由肉身所生，然而普遍的觉知岂不使每个人都出于天主？那么基督何以断言祂自己或祂存在的根源不为人知？祂只能通过将祂的直接父系归于存在的最终根源来这样做；而一旦如此，祂就能藉着他们不认识祂是天主子这一事实，来证明他们对天主的无知。让这可怜幻想的受害者们思想这句话：\BibleQuote{你们既不认识我，也不知道我从哪里来}。他们争辩说，万物皆出于无；他们不容许有任何例外。他们甚至胆敢将天主独生子曲解为从无而生。我们如何解释亵渎者对基督以及基督根源的无知？正如圣经所说，他们不知道祂从何而来，这一事实本身就表明了祂所来自的那不可知的根源。如果我们能说某物是从无中产生的，那么我们就不是不知道它的起源；我们知道它是从无中被造的，而这是一种确定的知识。那来临的并非祂自己存有的创造者；但派遣祂的那位是真实的，是亵渎者所不认识的。正是祂派遣了祂；而他们不知道祂是派遣者。因此，被派遣者出于派遣者；出于他们不认识的那位作为祂的根源。他们之所以不知道基督是谁，是因为他们不知道祂是从谁而来。凡否认祂是受生者的人，都不能承认圣子；凡认为祂是从无而来的人，都不能理解祂是受生的。事实上，祂远非从无而被造，以至于异端者无法说出祂从何而来。\par

30. 那些将天主的名字与天主的本性分离的人是完全无知的；他们无知，且安于无知。但请他们听一听圣子对不信者因缺乏这种认识而施加的斥责，那时犹太人说天主是他们的父亲：\BibleQuote{如果天主是你们的父亲，你们必然会爱我；因为我出于天主并来了；我不是凭自己来的，而是祂派遣了我}。\BibleRef{若 8:42}天主子在此没有责备那些将宣认祂是真天主、天主子与他们自称是天主子女结合起来的虔敬信心。祂所责备的是犹太人胆敢宣称天主为父的傲慢，而同时他们并不爱祂——圣子：\BibleQuote{如果天主是你们的父亲，你们必然会爱我；因为我出于天主}。所有因信德而以天主为父的人，都是通过那同一信德以祂为父，即我们宣认耶稣基督是天主子。但是，在一种涵盖所有圣人集合的意义上宣认祂是圣子；实际上说祂是众天主子之一——那有什么信德可言？其余的一切，尽管是软弱受造物，不也是在这种意义上的儿子吗？如果祂作为众多儿子中的一员，只有儿子的名份而没有儿子的本性，那么那宣认耶稣基督是天主子的信德的高超之处在哪里？这种不信不爱基督；这些真理的曲解者宣称天主为父，是对信德的嘲弄。如果祂是他们的父亲，他们就会爱基督，因为祂出于天主。现在我必须探究这\Quote{出于天主}的意思。祂的\Quote{出}显然不同于祂的\Quote{来}，因为在这段经文中两者并列提到：\BibleQuote{我出于天主并来了}。为了阐明\BibleQuote{我出于天主}和\BibleQuote{我来了}各自的意思，祂立即补充说：\BibleQuote{我不是凭自己来的，而是祂派遣了我}。祂告诉我们祂不是自己存在的根源，用这句话：\BibleQuote{我不是凭自己来的}。在这句话中，祂告诉我们祂再次出于天主，并由祂所派遣。但当祂告诉我们，那些称天主为父的人必须爱祂自己，因为祂出于天主，祂就把祂的受生作为他们爱的原因。\Quote{出}将我们的思想带回那无形的受生，因为正是藉着对那由祂而生的基督的爱，我们才能获得虔敬地称天主为父的权利。因为当圣子说：\BibleQuote{恨我的，也恨我的父}\BibleRef{若 15:23}，这个\InnerQuote{我的}是宣称一种与圣父的关系，是无人能共享的。另一方面，祂谴责那称天主为父却不爱圣子的人，认为这是滥用圣父的名号；因为恨祂——圣子的人，也必恨圣父，除了爱圣子的人以外，没有人能孝爱圣父。因为祂给予爱圣子的唯一理由就是祂出于圣父。因此，圣子出于圣父，不是因祂的降世，而是因祂的受生；只有那些相信圣子出于圣父的人，才能爱圣父。\par

31. 对此，主的话作证——\BibleQuote{我必不向你们说：我要为你们求父，因为父自己爱你们，因为你们爱了我，并相信我是从天主出来的，且从父来到这世界。} \BibleRef{若 16:26} 关于圣子的一种完备信德，即接受并热爱他从天主出来这个真理，无需他的中介便能亲近圣父。承认圣子是生于天主并从天主派遣而来，为他赢得了圣父的直接垂听和爱。因此，关于他的出生与来临的叙述必须按最严格、最字面的意义来理解。他说：\BibleQuote{我是从天主出来的}，传达了他的本性正是由他的出生所赐予的；因为除了天主，谁能从天主出来呢？也就是说，谁能因他而开始存在呢？然后他接着说：\BibleQuote{且从父来到这世界}。为使我们确信这\Quote{从天主出来}是指从圣父出生，他告诉我们他从圣父来到这世界。后一句话指他的降生成人，前一句指他的本性。此外，他先记录了他从天主出来的事实，然后是他从圣父来的事实，这禁止我们将\Quote{出来}与\Quote{来临}等同。从圣父来和从天主出来不是同义的；它们可以解释为\OriginalTerm{出生}{Birth}和\OriginalTerm{临在}{Presence}，其含义像这些词一样不同。一件事是从天主出来，通过出生进入实体存在；另一件事是从圣父来到这世界，完成我们救恩的奥迹。\par

32. 在我心中所安排的辩护顺序中，这似乎是证明第三点的最合适的地方：宗徒们相信我们的主耶稣基督是天主子，不仅是名义上的，而且是本性上的，不是因义子关系而是因出生。诚然，还有许许多多、极其重要的天主独生子关于他自己的话尚未提及，其中他神圣出生的真理被如此清楚地阐明，足以平息任何反对的窃窃私语。然而，由于用大量证据加重读者思想的负担是不智的，并且已经提供了充足的证据证明他出生的真实性，我将保留他其余的话语，留待我们探究的后续阶段。因为我们如此安排了论证的进程：现在，在听了圣父的见证和子的自我启示之后，我们要受宗徒们的信德教导，认识真正的、我们必须承认的真正生于天主的圣子。我们必须看看他们能否在主的话\BibleQuote{我是从天主出来的}中找到任何其他的含义，除了在他内有一个神圣本性的出生。\par

33. 在许多隐晦的言词之后，他用比喻讲论，他们已知他是\Person{梅瑟}和先知所预言的基督，\Person{纳塔乃耳}曾承认他是天主子、以色列的君王，他本人也曾责备\Person{斐理伯}关于圣父的问题，因为他没有通过他所行的工程看出圣父在他内、他在圣父内；在已经多次教导他们他是从圣父派遣来的之后，仍然，直到他们听见他声称他是从天主出来的，他们才用福音中紧接着的话语承认——\BibleQuote{他的门徒们向他说：看，现在你明明地讲论，并不说什么比喻了。现在我们知道你一切都知道，也不需要有人问你；由此我们相信你是从天主出来的。} \BibleRef{若 16:29} 这个措辞\Quote{是天主出来的}有什么如此奇妙之处呢？你们这些神圣而蒙福的人，为了你们信德的赏报，已经领受了天国钥匙，以及在天上地下捆绑和释放的权柄，你们看见我们的主耶稣基督、天主子行了如此伟大、如此真正神圣的工程，而你们却声称，直到他首先告诉你们他是从天主出来的，你们才获得真理的知识？然而你们看见在婚宴中水变成了酒；一个本性变成另一个本性，无论是通过变化、发展或创造。你们的手曾将五个饼分给那大群人吃，当所有人都吃饱时，你们发现需要十二个篮子来装饼的碎片；一小块物质，在缓解饥饿的过程中，变成了同一种性别的巨大物质。你们看见枯萎的手恢复了柔软，哑巴的舌头被松开说话，跛子的脚快跑，瞎子的眼睛得见光明，死人恢复生命。已发臭的\Person{拉匝禄}因一句话而站起来。他从坟墓中被召唤，立刻出来，话语与成就之间没有停顿。他站在你们面前，是一个活人，而空气中仍带着死亡的气味传到你们的鼻孔。我不提其他他大能、神圣能力的运用。然而，尽管如此，直到你们听见他说\Quote{我是从天主出来的}，你们才明白他是谁，是从天上派遣来的？这是第一次没有用比喻向你们说真理吗？这是第一次他本性的能力向你们显明他是从天主出来的吗？尽管他沉默地察看你们意志的意图，不需要问你们任何事，好像他不知道，他无所不知？因为这一切，以天主的能力和本性所做的一切，都证明了他必须是从天主出来的。\par

34. 通过这事，圣宗徒们并不理解祂是从天主出来的，意思是祂被派遣。因为他们常听见祂在之前的讲论中承认祂是奉派遣的；但现在他们听见的是明确的宣告：祂是从天主出来的。这开了他们的眼，使他们从祂的作为中认出祂的天主性。祂是从天主出来的这一事实向他们显明了祂的真天主性，因此他们说：\BibleQuote{现在我们晓得，祢知道一切事，也不需要谁问祢；因此我们相信祢是从天主出来的。}他们相信祂是从天主出来的原因，是祂既能又实行天主的作为。他们对祂天主性的完全确信，是由于他们知道——不是祂从天主而来，而是祂从天主出来。因此我们发现，正是这现在第一次听到的真理，坚定了他们的信德。主曾作了两个声明：\BibleQuote{我从天主出来}，和\BibleQuote{我从圣父来到这个世界}。其中一个，\BibleQuote{我从圣父来到这个世界}，他们常听见，并不引起诧异。但他们的答复显明，他们现在相信并理解了另一个，即\BibleQuote{我从天主出来。}他们的回答：\BibleQuote{因此我们相信祢是从天主出来的}，正是对这一个的回应，且只对这个；他们没有加上\Quote{并从圣父来到这个世界。}前一个声明以信德的宣告被欢迎，后一个却被默然放过。这宣认是被一个新真理的突然呈现逼出来的，这真理说服了他们的理性，迫使他们承认自己的确信。他们已知祂像天主一样能行一切事，但祂的诞生——这就是那全能的原因——并未被启示。他们知道祂是从天主被派遣的，却不知道祂是从天主出来的。如今，终于由这讲话所教导，明白了圣子不可言喻且完美的诞生，他们宣认祂对他们讲话没有用比喻。\par

35. 因为天主并非由天主通过人类分娩的通常过程而生；这并非一个存有通过自然力的作用从另一个存有出现。那诞生是纯洁、完美、无玷的；事实上，我们不如称之为发出，而非诞生。因为这是从一而来的一；没有分割、分离、减少、流出、延展、或遭受变化，而是从活的本性产生活的本性。这是天主从天主发出，而非被选来承担天主之名的受造物。祂的存在并非从无开始，而是从永恒者发出；这发出正当地被称为诞生，尽管称它为开始是错误的。因为天主从天主发出，与一个新实体的产生完全不同。而且，虽然我们对这不可言喻的真理的理解无法用言辞界定，但圣子的教导，当祂向我们启示祂是从天主发出的，就赋予这真理以确信信德的确凿性。\par

36. 相信天主子仅是名义上的儿子，而非本性上的儿子，并非福音和宗徒们的信仰。如果这只是一个称号，祂仅凭义子名分而拥有这称号；如果祂并非因出自天主而为子，那么，请问，有福的西满巴尔约纳曾向祂宣认说：\BibleQuote{祢是默西亚，永生天主之子}\BibleRef{玛 16:16}？这是从何而来？难道是因为祂藉重生的圣事，与全人类分享了成为天主子女的能力吗？如果基督只是以这种名义的方式成为天主子，那么父在天上不是藉血肉，而是藉启示传给伯多禄的，又是什么呢？他做出一项普遍适用的声明，能得什么赞美？陈述一个众所周知的事实，又有什么功劳可言？如果祂是义子，那么伯多禄的宣认之福在哪里？他的宣认对那儿子表示敬意，而按那种情况，那儿子并不比诸圣中的任何一位更有资格。宗徒的信德穿透了人类推理所不能及的领域。他无疑常听到：\BibleQuote{谁接纳你们，就是接纳我；谁接纳我，就是接纳那派遣我来的}\BibleRef{玛 10:40}。因此，他深知基督已被派遣；他曾听见祂——他知道祂是被派遣的那一位——作出声明：\Quote{我父将一切交给了我；除了父，没有人认识子；除了子，也没有人认识父。}那么，这真理是什么？是父如今向伯多禄启示的，并受到有福宣认的赞美？不可能是因为\Quote{父}和\Quote{子}的名词对他来说是新奇的；他常听到它们。然而，他说出了人类的口舌从未说出的话：\Quote{祢是默西亚，永生天主之子}。因为，虽然基督在肉躯内居住时，曾自称是天主子，但如今宗徒的信德才第一次在祂身上认出了天主性的临在。伯多禄受赞美，不仅因为他表达敬拜，更因为他辨认出这奥妙的真理；因为他不只宣认基督，还宣认了基督是天主子。显然，如果他只说\Quote{祢是默西亚}，不再多言，也足以表达敬意。但如果伯多禄只欢呼祂为默西亚，而不宣认祂是天主子，那宣认就是空洞的。因此，他的话\Quote{祢是}宣告了那关于祂的断言严格而真实地符合祂的本性。接下来，父的宣告\Quote{这是我的子}启示伯多禄，他必须宣认\Quote{祢是天主子}；因为在\Quote{这是}这话中，启示者天主指出祂来，而回应\Quote{祢是}，乃是信者对真理的欢迎。这就是宣认的磐石，教会建立在其上。但血肉的感知能力不能达成对这项真理的辨认和宣认。这是一个由天主启示的奥秘：基督不仅要被命名，而且要被相信为天主子。难道启示给伯多禄的只是天主性之名，而非天主性本身吗？如果是名字，他常从主那里听到主自称天主子。那么，他宣告这名字又有什么荣誉呢？不；不是天主性，而是本性，因为名字早已被反复宣告过。\par

37. 这信德就是教会的根基；凭这信德，阴间的门不能战胜她。这就是拥有天国钥匙的信德。凡这信德在地上所释放或束缚的，在天上也必被释放或束缚。这信德是父藉启示所赐的恩赐；甚至包括这样的知识：我们不可想象一个假的基督，一个从无中受造的受造物，而必须宣认祂是天主子，真正拥有天主性。是何等亵渎的疯狂和可怜的愚昧，竟不理会那有福的殉道者伯多禄本人可敬的年龄和信德！父曾为他祈祷，使他的信德在诱惑中不至失落；他曾两次重复被要求的对天主的爱的宣告，并因第三次重复这个问题而忧伤，好像他的虔诚是可疑而动摇的；然而，因为第三次考验清除他的软弱，他获得酬报，三次听到主的委托：\Quote{牧放我的羊}。当所有宗徒都沉默时，他独自从父的启示认出了天主子，并因宣认他幸福的信德而赢得了超越人性软弱的光荣的卓越地位。我们研究他的话，迫使我们得出什么结论？他宣认基督是天主子；而你，这新宗徒职的谎言之主教，把你们现代的观念强加给我们，认为基督是一个受造物，从无中受造。这是何等的暴力，如此歪曲这些荣耀的话语？他之所以有福，正是因为他宣认了天主子。这是父的启示，这是教会的根基，这是她恒久的保证。因此，她拥有天国的钥匙，因此，在天上有审判，在地上也有审判。藉着启示，伯多禄学到了从世界开始以来隐藏的奥秘，宣扬了信德，公布了天主性，宣认了天主子。谁若否认这一切真理，宣认基督是一个受造物，他必须先否认伯多禄的宗徒职、他的信德、他的真福、他的主教职、他的殉道。而当他做完这一切时，他必须知道他已与基督隔绝；因为正是藉着宣认祂，伯多禄赢得了这些荣耀。\par

38. 你这今日的悲惨异端者，你以为伯多禄若说：\Quote{祢是基督，天主的完美受造物，祂的化工，虽超越祂所有其他工程。祢的起源是从无而来，因着唯独善的天主的良善，祢被赐予了义子的名号，尽管事实上祢并非由天主所生}，他如今会更蒙福吗？你想，这样的话会得到什么回答？当同一伯多禄因预告苦难而说：\Quote{主，请祢远离我，这事绝不会临到祢身上}，却受斥责说：\Quote{撒殚，退到我后面去！你是我的绊脚石}\BibleRef{玛 16:22}。然而伯多禄可以以其人性的无知来减轻他的过犯，因为那时圣父尚未启示苦难的全部奥秘；但仅仅是信德上的缺陷就受到了如此严厉的谴责。那么，为什么圣父没有向伯多禄启示你那真实的宣信，即那在一位义子受造物上的信德呢？我想，天主必定是吝于让他知道真理；祂想把它延迟到一个更晚的时代，并把它作为一件新奇事物留给你们现代的宣讲者。是的，如果天国的钥匙改变了，你们的信德也可以改变。如果那教会（阴间的门不能战胜她）改变了，你们的信德也可以改变。如果有一个新的宗徒职位，在天上捆绑和释放地上所捆绑和释放的，你们的信德也可以改变。如果有另一位天主子基督（除了真基督之外）被宣讲，你们的信德也可以改变。但如果那宣认基督为天主子的信德——并且唯独那信德——在伯多禄身上领受了一切累积的祝福，那么，那宣告祂为一个从无中造出的受造物的信德，就必然不持有教会的钥匙，也与宗徒的信德和能力无分。它既不是教会的信德，也不是基督的信德。\par

39. 因此，让我们列举每一个由一位宗徒作出信德宣言的例子。他们所有人，当他们宣认天主子时，宣认祂不是一位名义上的、义子的儿子，而是因拥有天主本性而为子。他们从未将祂贬低到受造物的层次，而是将祂从天主真实出生的荣耀归于祂。让若望对我们说话吧，他正等待主的来临，正如他所是的；若望，他被留下并指定了一个隐藏在天主计划中的命运，因为他被告知他不会死，而只是要存留。让他用他熟悉的声音对我们说：\Quote{从来没有人看见过天主，只有那在父怀里的独生子，给我们详述过}\BibleRef{若 1:18}。在他看来，似乎\Quote{子}的名号不足以足够清晰地表明祂真实的天主性，除非他通过指出基督与所有其他人的不同来为基督特有的尊威提供外部支持。因此，他不仅称祂为子，还进一步加上了\Emph{独生子}的称号，从而剪断了这想像的义子论的最后一根支柱。因为祂是独生子这一事实，是祂享有子名份的明证。\par

40. 我将\Quote{在父怀里}这句话的思考推迟到更合适的地方。我目前探究的是\Emph{独生子}的意义，以及这意义可能对我们的要求。首先让我们看看，这个字是否如你所主张的那样，意谓天主的完美受造物；\Emph{独生子}等同于完美，\Emph{子}是受造物的同义字。但若望描述独生子是天主，而非完美的受造物。他的话：\Quote{在父怀里}，表明他预见到了这些亵渎性的称号；而且，他确曾听见他的主说：\BibleQuote{天主竟这样爱了世界，甚至赐下了自己的独生子，使凡信祂的人不至丧亡，反而获得永生}。爱世界的天主赐下祂的独生子，作为祂爱的显明标记。如果祂爱的证据是这样的：祂把受造物赐给受造物，将一个世俗的存在为世界的缘故赐予，将一个从无中兴起者为同样从无中兴起的对象的救赎而赐予，那么这廉价而微小的牺牲，是对祂赐予我们的恩宠的一种贫乏保证。珍贵的礼物是爱的证据：牺牲的伟大是爱的伟大的证明。爱世界的天主赐下的不是义子，而是祂自己的、祂的独生子。这里有个人性的关切、真实的儿子身份、真诚；而非受造、义子或伪装。这就是祂的爱和感情的证明：祂赐下了祂自己的、祂的独生子。\par

41. 我现在不诉诸任何赐给子的称号；当选择因令人尴尬的丰富而延迟时，并无损失。我现在的论点是：成功的结果意谓着充分的原因；每一个有效的成就背后必须有清晰而令人信服的动机。因此，圣史不得不揭示他写作的动机。让我们看看他承认的目的是什么：\Quote{这些事记录下来，为叫你们信耶稣是默西亚，天主子}\BibleRef{若 20:31}。他提出写福音的唯一理由，是叫所有人都可以信耶稣是默西亚，天主子。如果相信祂是默西亚就足以得救，为什么他又加上\Quote{天主子}呢？但如果真信德恰恰是相信基督不仅是默西亚，而且是天主子基督，那么\Quote{子}的名号就不是作为因义子而有的惯例附加在基督身上，因为它是救恩所必需的。既然如此，救恩就在于宣认这名号，那么这名号岂不必须表达真理？如果这名号表达真理，祂又怎能被称为受造物呢？使我们得救的，不是对受造物的宣认，而是对子的宣认。\par

42. 因此，相信耶稣基督是天主圣子，便是真实的救恩，是真诚信德的悦纳服务。因为我们内心对天主圣父的爱，唯有藉着对圣子的信德才能拥有。让我们听祂藉着宗徒书信向我们所说的话：\Quote{\Emph{凡爱圣父的，也爱那由祂所生的}}。\BibleRef{若一 5:1}请问，这\Quote{由祂所生}是什么意思？难道是受祂所造吗？难道福音作者说祂是由天主所生，而异端者更正确地教导说祂是受造的吗？让我们都来听听这位异端导师的真实面目。经上记载：\Quote{\Emph{凡否认圣父及圣子的，便是假基督}}。现在，你这受造物的拥护者，从虚无中编织出新基督的人，你要怎么办呢？如果你坚持你的主张，请听等待你的称号。或者你以为，你仍可将圣父和圣子描述为造物主和受造物，却藉着巧妙的语言模糊，不被认作假基督？如果你的信仰承认一个真实的圣父和一个真实的圣子，那么我就是诽谤者，用你不应得的恶名攻击你。但如果在你的信仰中，基督的一切属性都是虚假和名义上的，而非祂本有的，那么请从宗徒那里学习你这种信仰的正确描述，并听听那相信圣子的真实信仰是什么。接下来的话是：\Quote{\Emph{凡否认圣子的，就没有圣父；凡明认圣子的，也有圣父也有圣子}}。否认圣子的，就没有圣父；明认并拥有圣子的，也有圣父。这里哪有义子之名的空间？难道每一句话不是在论及天主性吗？请学习这性体是如何全然临在的。\par

43. 若望这样说道：\Quote{\Emph{因为我们知道天主子已来临，并为我们降生成人，受苦，从死者中复活，将我们收归自己，并赐给我们智慧，使我们认识那真实的，并得以在祂真实的圣子耶稣基督内。祂是真实的，是永生，又是我们的复活}}。智慧注定有邪恶的结局，缺乏天主之神，注定拥有假基督的精神和名字，对天主子来成就我们救恩的奥迹视而不见，并因那盲目而不配看见那至高知识的光明！因为这智慧断言耶稣基督并非天主的真子，而是祂的一个受造物，借着义子之名承受了天主的称号。这隐藏知识的黑暗神谕是从哪里学来的？这是谁的探求，使我们得到了这当代的伟大发现？难道是你曾依偎在主怀中？是你，在爱的亲密交往中，祂向你揭示了奥迹？是你独自跟随祂直到十字架下？当祂吩咐你接受玛利亚为你的母亲时，祂是否将这秘密作为祂对你特别之爱的标志教导了你？或者你跑向坟墓，甚至比伯多禄更快到达，并因此在那里获得了这知识？还是在天使的群集、无人能开启的封印书卷、天上各种征兆的影响、永恆歌队未知的歌声中，你的向导羔羊向你启示了这虔诚的教义：圣父不是父，圣子不是子，既没有性体，也没有真理？因为你将这一切都变成了谎言。宗徒藉着所赐予他的那极卓越的知识，称天主子为真实的。你却断言祂是受造，宣告祂是义子，否认祂是所生。既然真实的天主子是我们的永生和复活，那么在他眼中，若祂不是真实的，就没有永生，也没有复活。这就是主所爱的门徒若望所教导的。\par

44. 那位曾为迫害者，后来归化成为宗徒和蒙选器皿的，也传达了同样的信息。他的哪篇讲论没有预设对圣子的明认？他的哪封书信不是以明认那奥迹的真理开篇？当他说：\Quote{\Emph{我们藉着祂圣子的死，得以与天主和好}} \BibleRef{罗 5:10}，又说：\Quote{\Emph{天主派遣了自己的圣子，成为罪恶肉身的模样}} \BibleRef{罗 8:3}，以及\Quote{\Emph{天主是信实的，你们原是蒙祂召叫，与祂的圣子共融}} \BibleRef{格前 1:9}，难道还给异端者的曲解留有任何缝隙吗？\Quote{祂的圣子，天主子}——我们读到的是这样，但丝毫没有提到祂的义子身份或天主的受造物。这名称表达了性体：祂是天主圣子，因此这子嗣关系是真实的。宗徒的明认肯定了这关系的真实性。我看不出圣子的天主性还能被更完全地陈述出来。那位蒙选器皿已用坚定无弱的声音宣告：基督就是那一位的圣子，我们相信那一位是圣父。外邦人的导师、基督的宗徒，在呈现这一教义时，没有给我们留下任何不确定、任何错误的余地。在义子的问题上，他的表述非常清楚；这是指那些因信德而得以成为义子并被称为义子的人。用他自己的话说：\BibleQuote{\Emph{因为凡受天主圣神引导的，都是天主的子女。其实你们并没有领受奴仆的神，仍存恐惧；而是领受了义子的神，藉着祂我们呼叫：\InnerQuote{阿爸，父啊！}}} \BibleRef{罗 8:14}这是赐给我们这些相信之人的名字，藉着重生的圣事；我们信德的明认为我们赢得了这义子身份。因为我们顺从天主圣神而作的工，给了我们天主的子女的称号。\Quote{阿爸，父啊}是我们发出的呼喊，而不是我们本性的表达。因为我们的本性并不受那声音的赞颂所影响。天主被称为父是一回事，祂是祂圣子的父是另一回事。\par

45. 但现在让我们学习宗徒所持的有关天主圣子的信德是什么。因为虽然在他所发表的关于教会道理的众多言论中，没有一篇不提及圣父而同时宣认圣子的，然而，为了以人类语言所能达到的最大确定性来展示那个名称所传达的关系的真理，他这样说：— \Quote{\Emph{什么？如果天主为我们，谁能反对我们？祂没有怜惜自己的子，反而为我们将祂交出？}} 难道\Emph{子}这个称号，当祂被明确称为天主的自己的子时，还能有任何可能通过义子身份而获得吗？因为宗徒为了显明天主对我们的爱，使用了一种比较，让我们能够估计那爱有多大，当他说天主没有怜惜的是祂自己的子时。他暗示，这并不是为那些祂打算收养的人而牺牲一个义子，为受造物而牺牲一个受造物，而是为陌生人牺牲祂的子，为那些祂愿意赐予儿子名分的人牺牲祂自己的子。要找出这个术语的全部含义，以便你理解那爱的程度。思考一下\Emph{自己的}的意义；注意它暗示的圣子身份的确实性。因为宗徒现在描述祂是天主自己的子；之前他常常称祂为天主子或天主子。虽然许多抄本因为翻译者缺乏理解，在此处读作\Quote{祂的子}而非\Quote{祂自己的子}，但原文希腊文——宗徒写作所用的语言——更准确地翻译为\Quote{祂自己的}而非\Quote{祂的}。尽管普通读者可能看不出\Quote{祂自己的}和\Quote{祂的}之间有多大区别，但宗徒在他所有其他陈述中一直说\Quote{祂的子}（在希腊文中是\OriginalTerm{祂自己的子}{\GreekText{τὸν} \GreekText{ἑαυτοῦ} \GreekText{υἱόν}}）。\par

\SourceRecord{Translated by E.W. Watson and L. Pullan.；Nicene and Post-Nicene Fathers, Second Series；Vol. 9.；Edited by Philip Schaff and Henry Wace.；Buffalo, NY: Christian Literature Publishing Co.,；1899.；<http://www.newadvent.org/fathers/330206.htm>.}

% structure: p0001 source=h1 target=section reason=page-title

\section{\WorkTitle{论天主圣三}（第七卷）}

第七卷。亚略派善于隐藏其含义；他们使用圣经术语却赋予非圣经的含义（§1）。他们已被驳斥，因为基督是真且同永恒的圣子；现在希拉里进而证明基督的真天主性，这从逻辑上与祂的真圣子身份不可分割（§2）。但危险在于，在攻击一种异端时，他可能使用会认可其他异端的语言（§3）。然而真理是唯一的，而异端是多样的。它们中的每一个都可以被信任来拆毁其他异端，而没有一个能建立自己的案例。他通过撒伯流、亚略和福提诺互相毁灭的论证来说明这一点（§§5-7）。基督被证明是天主，因为圣经中给了祂\Emph{天主}这个名号：\BibleQuote{圣言是天主}（§§8,9）。这个名号是在严格意义上属于祂，而不是任何派生含义（§§10,11）。然而圣父和圣子不是两位，而是一位天主（§13）。作为天主子，祂有天主的本性，因此是天主（§§14-17），然而与圣父不是同一个位格（§18）。再者，祂在工程中所显示的能力，证明了祂的天主性（§19），正如所有审判都被圣父赐给了祂这一事实（§20）。基督自己的话显明了真理（§21）。亚略派对圣经的明显含义视而不见，比犹太人更亵渎；基督对后者的答复也回应了前者的异议（§§22-24）。祂宣告了与圣父的合一（§25），并以祂的工程为证明（§26）。圣父在圣子内，圣子在圣父内（§27）：这通过生理特性从父母传递给孩子以及从火焰到火焰的传递来说明（§§28-30）。事实上，唯有公教（的解释）才是对圣经言辞的理性说明（§§31,32）。再者（§§33-38），通往圣父的路是通过圣子，而认识圣子就是认识圣父。若非圣子在与圣父同样的意义上即是天主，这是不可能的。因此，撒伯流和亚略的敌对教义被驳斥；既非一个位格，也非两位天主（§§39,40）。基督召叫我们相信真理，而信德不仅是可能的，而且是合理的（§41）。\par

1. 这是我们针对现代异端之狂诞所著论文的第七卷。按顺序，它必须紧随前卷；按重要性，作为正确信德之奥迹的阐述，它超越并优于前卷。我深知我们所攀登的福音教导之路是何等艰难崎岖。对自己无能的认识所引发的恐惧将我拉回，但信德的热情催逼我前进；异端的攻击使我热血沸腾，对无知者的危险激起我的怜悯。我害怕发言，却也不能沉默。双重的畏惧压制了我的灵：说话或许会，沉默或许会，使我犯下背弃真理的罪过。因为这狡诈的异端已经用扭曲才智的奇妙手段将自己围护起来。首先是虔诚的外表；其次是精心选择以麻痹坦诚听众疑心的言辞；又次是将其观点迎合世俗哲学；最后，通过假装解释神圣方法，使人远离明显的真理。他们大声宣认天主的合一，是对信德的欺诈模仿；他们声称基督是天主子，是玩弄言辞以迷惑听众；他们说祂在生之前并不存在，是为了争取世上哲学家的支持；他们宣认天主是无形不变，通过谬误逻辑的展示，导致否认天主生于天主。他们用我们自己的论证来反对我们；教会的信德被用作摧毁自身的工具。他们已使我们陷入同等的危险困境，无论我们是与他们推理还是保持沉默。因为他们利用我们允许他们某些主张不受挑战的事实，作为支持我们所反对的那些主张的论证。\par

2. 我们记得，在前几卷中曾敦促读者研读那整个亵渎的宣言，并留意它如何自始至终贯穿着一个目的，即传播我们的主耶稣基督既非天主、亦非天主子这个信念。它的作者们争辩说，祂因一种义子身份而被允许使用天主和子的名称，尽管天主性和子性并非祂本性所固有。他们利用天主是永恒不变和非物质体这一本身真实的事实，作为反对子从祂出生的论据。他们珍视天主圣父是独一的那真理，只是为了作为反对我们信仰基督天主性的武器；他们辩称，一个非物质体的本性不能被合理地设想为产生另一个，并且我们信独一天主与宣认来自天主的天主是不一致的。但我们早期的几卷书已经通过诉诸律法和先知驳斥并挫败了他们的这个论点。我们的辩护一步一步地跟随着他们攻击的进程。我们陈述了来自天主的天主，同时宣认了独一的真天主；表明这种信仰表述既不会因为将单独一位位格归于独一的真天主而违背真理，也不会因断言存在第二位神祇而增加信仰。因为我们既不宣认一位孤立的天主，也不宣认两位天主。因此，既不否认天主是独一，也不坚持祂是孤独的，我们持守真理的正路。每一个神圣位格都在统一体中，但没有一位位格是独一的天主。其次，我们的目的是要通过福音作者和宗徒的证据来证明这奥迹不可辩驳的真理，我们的首要职责是使读者认识天主子那真实存在且真实出生的本性；证明祂没有外在于天主的起源，不是从虚无中被创造，而是从天主所生的圣子。这是上一卷书所引证据已毋庸置疑地确立了的真理。关于祂因义子身份而拥有子之名的主张已被驳倒，祂以真正的出生显明为真子。我们现在的任务是从福音中证明，因为祂是真子，所以祂也是真天主。因为除非祂是真子，祂不能是真天主；除非祂是真天主，祂也不能是真子。\par

3. 没有什么比濒临危险的感觉更折磨人性了。如果未知或未曾预料的灾祸降临到我们身上，我们或许需要怜悯，然而我们始终无忧无虑；没有焦虑的重担压迫我们。但心中充满种种可能麻烦的人，已在恐惧中受着折磨。我现在冒险出海，是一个习惯遭遇海难的水手，是一个经验知道有神圣匪徒潜伏在森林中的旅行者，是一个探知非洲沙漠有蝎子、毒蛇和蛇怪危险的探险者。我对当前危险的认识和恐惧，没有一刻的缓解。每一个异端者都在监视，留意从我口中吐出的每一个字。我的整个论证过程充斥着埋伏、陷阱和罗网。我不是抱怨道路的艰难或陡峭；我正跟随宗徒的足迹，而不是选择自己的路。我的麻烦是不断的危险、不断的恐惧，唯恐误入某个埋伏，跌进某个坑穴，被某张网缠住。我的目的是宣讲天主的合一，按律法、先知和宗徒的意义。\Person{撒伯里乌}就在眼前，急切地想以残酷的善意欢迎我，凭借这合一，把我吞没在他的毁灭中。如果我抵抗他，否认天主按撒伯里乌的意义是独一，一个新的异端就准备接纳我，指出我教导两位天主的存在。再者，如果我试图讲述天主子如何从玛利亚出生，\Person{福提诺}，我们时代的厄俾翁派，就会迅速把这真理的断言扭曲为对他谎言的确认。我不需要提及其他异端，只提一个；全世界都知道他们与教会无关。这是一个常被谴责、常被拒绝的异端，却仍在蚕食我们的生命。\Place{迦拉达}培育了一大群否认天主合一的邪恶断言者。\Place{亚历山大}几乎在整个世界播撒了她的否认——其实是对两位天主学说的肯定。\Place{潘诺尼亚}支持她那恶毒的学说，即耶稣基督唯一的出生是从童贞女。而教会，被这些相互竞争的信仰所困扰，面临通过真理被引向拒绝真理的危险。有些学说为了邪恶的目的而被强加给她，根据它们被使用的方式，要么支持信仰，要么推翻信仰。例如，我们作为真信徒，不能断言天主是独一，如果我们意指祂是孤独的；因为信仰一位孤独的天主就否定了子的天主性。另一方面，如果我们如实地断言子是天主，我们就会像他们天真想象的那样，有放弃天主是独一这一真理的危险。我们两方面都处于危险之中；我们可以否认合一或坚持孤立。但这危险对\BibleQuote{世上的愚拙事物}\BibleRef{格前 1:27}毫无威慑。我们的对手瞎了眼，没有看到祂宣称自己不孤独与合一是一致的，虽然祂是独一却并非孤独。\par

4. 但我相信，教会借其教义之光，将如此启迪世间虚妄的智慧，以致即使它不接受信德的奥迹，也会承认，在我们与异端者的争战中，我们而非他们，才是那奥迹的真正代表。因为真理的力量是伟大的；它不仅是自己充足的见证，越是受到攻击，就越发显明；每日所承受的冲击，仅使其固有的稳固性增加。教会的特质是：当她受鞭笞时，她得胜；当她受论理攻击时，她证明自己正确；当她被支持者遗弃时，她坚持阵地。她愿所有人都留在她身边，在她怀中；若由她决定，没有一人会变得不配留在这位尊贵母亲的庇荫之下，没有一人会被驱逐或任由离开她平静的憩息之所。但当异端者离弃她，或她驱逐他们时，她所承受的损失——因不能拯救他们——被一种更大的确信所补偿：唯独她能赐予福乐。这是一个被各个敌对异端的狂热激情带到最明显凸显的真理。教会由主所立，由祂的宗徒所建立，是为众人而一；但分裂教派的疯狂愚行已将她们与她分离。显而易见，这些关于信德的分歧源于扭曲的理智，它歪曲圣经的话语以符合自己的意见，而不是将自己的意见调整到圣经的话语。于是，在彼此毁灭的错误的碰撞中，教会不仅凭自己的教导，也凭其对手的教导而显明。她们全都排列起来反对她；而她独自屹立这一事实本身，就是对她们那邪恶妄想的充分回答。异端的大军集结反对她；她们中每一个都能打败其他一切，却没有一个能为自己赢得胜利。唯一的胜利是教会对她们全体举行的凯旋。每一个异端都使用一些已在另一情形下被教会谴责所粉碎的武器来进攻其敌人。她们之间毫无联合点，她们自相残杀的结果是信德的巩固。\par

5. 撒伯流扫除了圣子的诞生，然后宣讲天主的合一；但他并不怀疑那在人性基督内行动的大能本性就是天主。他对所启示的圣子身份的奥迹闭目不见；所行的奇事在他看来如此奇妙，以致他无法相信行这些事的那位能经历真实的诞生。当他听到\BibleQuote{\Emph{看见了我，就是看见了父}}\BibleRef{若 14:9}时，他妄下结论，认为圣父与圣子在本性上是不可分且无区别的同一，因为他未能看到，诞生之启示正是祂们本性合一向我们显示的方式。因为圣父在圣子中被看见，这证明圣子的天主性，而非否认祂的诞生。因此，我们对于祂们每一位的认识，是以对另一位的认识为条件的，因为祂们之间在本性上没有差别，且在这层意义上祂们原是一体，虔敬地研究其中任何一位的特质，便会让我们对圣父和圣子的本性有真正的洞见。因为，确实，祂虽具有天主的形体，但在自我启示中，必定恰好以天主的形体的样貌呈现给我们。\BibleRef{斐 2:6}再者，这种乖戾疯狂的妄想从\BibleQuote{\Emph{我与父原是一体}}\BibleRef{若 10:30}中得到了进一步的鼓励。从同一本性中的合一，他们不虔敬地推断出位格的混淆；他们的解释——说这句话意指单个的权能——与经文的要旨相悖。因为\BibleQuote{\Emph{我与父原是一体}}并不指示独一的天主。连词\Emph{与}明确表示所指的不仅是一个位格；而\Emph{原是一体}需要复数的主语。此外，\Emph{一体}并不与诞生相冲突。其意思是，两个位格共有同一本性。\Emph{一体}与差异不相容；\Emph{原是一体}与同一性不相容。\par

6. 让我们将现代的异端与沙贝里同样妄诞的迷惑并列；让他们各自尽量辩护。新异端者会引用这句经文：\Quote{父比我大}。他们忽视天主降生的奥迹，以及天主自己空虚、取了肉身的奥迹，却因祂宣称父是更大的，就论断祂的本性是低劣的。他们会反对沙贝里，声称基督是子，既然是子，就比父小，并需要祈求恢复祂的光荣，且怕死，也确实死了。作为回应，沙贝里将列举祂的行为来证明祂的天主性；而我们的新异端为了避免承认基督真正的子身份，会衷心同意天主是独一的，沙贝里则会坚决肯定同一信条，意思是没有任何子存在。一方强调子的行动；另一方则主张在那行动中天主彰显自己。一方将证明合一性，另一方则否定同一性。沙贝里会这样辩护：\Quote{所行的事，除非是天主性，别无其他本性能够成就。罪过得赦免，病人得痊愈，瘸子奔跑，瞎子看见，死人复活。只有天主有这能力。\InnerQuote{我与父原是一体}这句话只能出自自知之明；父之外的本性不可能说出这样的话。那么为什么暗示有另一个实体，并迫使我相信有第二位天主？这些工作唯独属于天主；独一的天主行了这些事。} 他的对手同样怀着恶毒的仇恨来反对信仰，会论证子在本质上与天主父不像：\Quote{你对自己救恩的奥迹一无所知。你必须相信一位子，万物藉祂而造，人由祂塑造，祂藉天使颁赐法律，由玛利亚所生，由父派遣，被钉十字架，死亡并被埋葬，又从死者中复活，坐在天主的右边，是审判生者死者的那一位。我们必须向祂复活，我们必须宣认祂，我们必须赢得祂国度中的位置。} 教会的这两个敌人，每一个都在为教会作战。沙贝里藉着在祂工作中显露的天主性之见证，显示基督是天主；沙贝里的对手则依据启示信仰的证据，宣认基督是天主子。\par

7. 再者，我们的信仰赢得何等光荣的胜利，在那胜利中，厄俾翁——换言之，福提诺——既得胜又失败！他责备沙贝里否认天主子是人，而他自己又不得不忍受亚略狂热者的指责，因为他看不出这个人就是天主子。反对沙贝里时，他诉诸福音，以及其中关于玛利亚之子的见证；亚略却剥夺了他的这个盟友，证明福音所宣示的基督，远不止是玛利亚之子。沙贝里否认有天主子；福提诺则高举人为子来反对他。福提诺不愿听任何在万世之前出生的子；亚略则否认天主子唯一的出生是祂的人性出生。让他们尽情地互相击败吧，因为他们每一次得胜都伴随着失败。我们现在的对手在天主子的天主性问题上受到诘问；沙贝里在天主子被启示的存在问题上受到诘问；福提诺因否认子在万世之前出生，被证明是无知或说谎。与此同时，教会，其信仰基于福音作者和宗徒的教导，坚持反对沙贝里，宣称子存在；反对亚略，宣称祂本性是天主；反对福提诺，宣称祂创造了宇宙。她越发确信自己的信仰，因为他们无法联合起来反驳它。因为沙贝里指出基督的工作，以证明行这些事者的天主性，虽然他不知子就是这些事的作者。亚略派承认祂子的名号，虽然他们不承认天主真实的本性住在祂内。福提诺维护祂的人性，虽然在维护时他忘记了基督在万世之前已经作为天主出生。因此，在他们各自的肯定和否定中，每个异端在某些点上，无论辩护或攻击，都有正确之处；而他们冲突的结果，是我们所宣认的真理得以更加彰显。\par

8. 我觉得必须腾出一点篇幅来指出这一点。这不是出于喜好铺陈，而是为了起到警示作用。首先，我想揭露这群异端模糊而混乱的特性，他们的相互争斗，如我们所看到的，反对我们有利。其次，在与现代异端亵渎性教义的争战中——也就是说，在我宣告天主父与天主子都是天主的任务中，换言之，父与子在名号上、在本性上、在它们所拥有的天主性种类上都是一体——我希望保护自己免受任何可能对我的指控，无论是指控我主张两位天主，还是主张一位孤寂而隔离的神祇。因为在我所阐述的天主父与天主子中，既无法觉察出位格的混淆，也无法在我阐述祂们共同的本性时，看出两者天主性之间的任何差异。在前一卷书中，我已经藉着福音的见证，充分反驳了那些否认天主子藉着从天主真实的出生而存在的人；我现在的责任是表明，那在本性真实上是天主子者，在本性真实上也是天主。但这项证明绝不能堕落为对一位孤寂天主或第二位天主的致命宣认。它应显明天主是独一的，却不孤独；但在小心避免使祂孤独的错误时，也不要陷入否认祂合一的错误。\par

9. 因此我们共有这关于我们的主耶稣基督天主性的不同确证：——祂的名字、祂的出生、祂的本性、祂的能力、祂自己的宣称。就名字而言，我认为不可能有任何怀疑。经上记载：\BibleQuote{在起初已有圣言，圣言与天主同在，圣言就是天主} \BibleRef{若 1:1}。有什么理由可以怀疑祂不是祂名字所指示的那位？难道这个名字不是清楚地描述了祂的本性吗？若一个声明被反驳，必然有某种理由。请问，在这件事上，有什么理由否认祂是天主？这名字被赋予祂，清楚而分明，且没有任何不协调的附加语可能引起怀疑。我们读到，那成了血肉的圣言，正是天主。这里没有任何漏洞可以让人猜测祂是作为恩惠而接受这名字，或自己取用了这名字，从而拥有一个并非祂本性的名义上的天主性。\par

10. 请考虑其他记载中这名字被作为恩惠赐予或自行取用的实例。曾对梅瑟说：\BibleQuote{我使你成为\Emph{法郎的}神} \BibleRef{出 7:1}。这附加语\Emph{给法郎}难道不解释了这一称号吗？难道天主将天主性传授给了梅瑟？祂难道不是使梅瑟在法郎眼中成为一个神，而法郎在梅瑟的蛇吞食了魔术的蛇并变回棍杖时要恐惧战栗，当他驱回他所召来的毒蝇，当他用召来冰雹的同一能力止住冰雹，并用带来蝗虫的同一力量使蝗虫离去；当他在所行的奇事中，术士们看到了天主的手指？正是在这意义上，梅瑟被立为法郎的神；他被惧怕，被恳求，他惩罚，他医治。被立为神是一回事，而成为神是另一回事。祂成了法郎的神；祂并不拥有天主所由构成的哪本性及那名号。我想起另一处那名字作为称号被赐予的例子：经上记载：\Emph{\Quote{我曾说过}：你们是神}。但这显然是恩惠的赐予。\Emph{我曾说过}证明这不是一个定义，而只是选择这样说话的一位所作的描述。定义给我们关于被定义对象的知识；描述取决于说话者的任意意愿。当说话者明显地授予一个称号时，那称号只起源于说话者的话语，而非事物本身。这称号不是表达其本性及种类的名字。\par

11. 但在这种情况下，圣言确确实实是天主；天主性的本质存在于圣言内，这本质在圣言的名字中表达出来。因为\Quote{圣言}这名字内在于天主子，作为祂奥秘出生的结果，同样还有\Emph{智慧}和\Emph{能力}这些名字。这些连同祂因真实出生而拥有的实体，被召唤存在而为天主子；然而，由于它们是天主本性的要素，它们仍然以不减少的幅度内在于祂内，尽管它们从祂而生为祂的儿子。因为，如我们经常所说，我们所宣讲的奥秘是一位存在非由分裂而由出生的儿子。祂不是被割下的一片故而残缺，而是所生的后裔，因此完美；因为出生不减少生者，且为被生者提供了完美的可能性。因此，那些实体属性的称呼被应用于天主独生子，因为当祂藉着出生进入存在时，正是它们构成了祂的完美；并且尽管它们并没有因此离弃圣父，在祂内，由于祂本性不变，它们永远存在。例如，圣言就是天主独生子，然而非受生的圣父从未无祂的圣言。并不是说圣子的本性是发出的声音。祂是天主出自天主，藉着真实出生而自立；是天主自己的儿子，由圣父所生，在本性上与祂无区别，因此不可分离。这是祂的\Quote{圣言}称号所要教导我们的功课。同样地，基督是天主的智慧和能力；不是说祂如人们常以为的那样，是圣父能力或思想的内在活动，而是说祂的本质，藉着出生拥有真实实体的存在，被这些内在力的名字所指示。因为一个因出生而自有其存在的对象，不能被视为一种属性；属性必然内在于某个存在者，不能有独立存在。但正是为了避免我们得出结论说圣子与祂圣父的天主性相异，祂，从永远的天主圣父所生的独生子，作为天主出生而有自己的实体存在，才在这些属性的名字下向我们启示了自己，而圣父，祂从祂而生，没有遭受任何减少。因此，祂既是天主，就只是天主。因为当我听到\BibleQuote{圣言就是天主}这句话时，它们不仅告诉我圣子被称为天主；它们向我的理智启示祂是天主。在先前那些例子中，梅瑟被称为神，其他人也被称为神，那不过是增添一个名字作为称号。这里陈述了一个坚实的本质真理：\BibleQuote{圣言是天主}。这个\Quote{是}指示的不是偶然的称号，而是永恒的现实，祂存在的一个永久要素，祂本性中内在的性格。\par

12. 现在，让我们看看宗徒\Person{多默}的宣认，当他呼喊\Quote{\Emph{我的主，我的天主！}}时，是否与福音作者的这一断言相符。我们看到，他称他所宣认为天主的那一位为\Quote{我的天主}。\Person{多默}无疑熟悉主的那句话：\Quote{\BibleQuote{以色列，请听！上主你的天主是唯一的上主}}。那么，一位宗徒的信德怎会对这一首要诫命如此疏忽，以至于宣认基督为天主，而永生又取决于对天主唯一性的宣认呢？这是因为，在复活的光照下，信德的一切奥秘已向宗徒显现。他曾多次听到\Quote{\BibleQuote{我与父原是一体}}、\Quote{\BibleQuote{凡父所有的，都是我的}}以及\Quote{\BibleQuote{我在父内，父在我内}}这样的话；现在他可以宣认，天主的名字表达了基督的本性，而无损于信德。在不违背对唯一天主、圣父的忠诚下，他的孝爱如今能视天主子为天主，因为他相信，凡在圣子本性内的一切，都确实与圣父有同一本性。他不再需要害怕这样的宣认是宣告第二位天主，是对天主性统一的背叛；因为那圆满的天主性所诞生的并不是第二位天主。因此，\Person{多默}是在对福音奥秘的完全认识下，宣认了他的主和他的天主。这不是一个尊称，而是对本性的宣认。他相信基督在实体和能力上是天主。而主反过来表明，这一朝拜举动不仅是敬畏的表达，更是信德的表达，当祂说：\Quote{\BibleQuote{因为你看见了我，才相信吗？那些没有看见而相信的，才是有福的！}}。因为\Person{多默}是先看见后才相信。但或许你会问：\Person{多默}相信了什么？毫无疑问，正如他的话\Quote{我的主，我的天主！}所表达的。除了天主性之外，没有本性能凭自身力量从死亡复活到生命；而正是这个事实——基督是天主——被\Person{多默}以确定信德的信心所宣认。那么，我们岂能梦想祂的天主之名并非实在的本体呢？当那名已被基于确凿证据的信德所宣告时？当然，一位顺服于祂圣父、不寻求自己旨意而寻求那派遣祂者的旨意、不寻求自己光荣而寻求那派遣祂者的光荣的圣子，会拒绝接受包含这样一个破坏祂所教导的天主统一性的名字的朝拜。然而，事实上，祂确认了这位相信的宗徒对奥秘真理的宣告；祂接受那属于圣父本性的名字作为祂自己的。并且祂教导说，那些虽然没有看见祂从死者中复活，但因复活的保证而相信祂是天主的人是有福的。\par

13. 因此，表达祂本性的名字证实了我们信德宣认的真理。因为，指示任何一个单一本体的名字，也指出同类的另一个本体；在此情形中，并非两个本体，而是同一类的一个本体。因为天主子是天主；这是祂名字所表达的真理。这一个名字并不包含两位天主；因为天主这个名字是同一不可分的天主性的名字。既然圣父是天主，圣子是天主，而那唯独属于天主性的名字内在于两者，因此两者是同一的。因为圣子虽通过从天主性的诞生而存在，却在祂的名字中保持了统一；而圣子的诞生并不迫使忠信的信徒承认两位天主，因为我们的宣认宣告圣父和圣子在本体和名字上都是同一的。因此，天主子因着祂的诞生而拥有天主之名。我们论证的第二步原是要表明，祂之所以是天主，是由于祂的诞生。我仍需引用宗徒们的证据，表明天主之名被用于祂是确切意义上的；但目前我打算继续探究福音的用语。\par

14. 首先我要问，藉由出生，什么能破坏祂天主性的新元素能被导入圣子的本性中？普遍理性拒绝接受一个存有藉由出生过程变得与其本源本性不同的假设；尽管我们承认，从不同类的父母可能生育出一个兼具两者本性却又与两者不同的后裔。这在驯养和野生的野兽中很常见。但即使在这种情况下，也没有真正的创新；新性质已经存在，隐藏在两种不同的父母本性中，只是通过结合而发展。它们共同后裔的出生并不是后裔与父母不同的原因。差异是来自它们的各种不同性的恩赐，这些不同性被接收并结合在一个框架中。当即便在身体差异的传递和接收上也是如此时，断言天主\OriginalTerm{独生子}{Only-begotten}是由天主出生而本性低于祂自己，岂非一种疯狂？因为出生是生命传递者真实本性的功能；没有那真实本性的临在和行动，就不可能有出生。这一切激烈争论的目标是要证明：天主子不是出生，而是受造；神性不是祂的起源，祂在祂的\OriginalTerm{位格生存}{personal subsistence}中并不拥有那本性，而是从无中抽取一种与神性不同类的本性。他们发怒是因为祂说：\BibleQuote{那由肉生的属于肉，那由神生的属于神。}\BibleRef{若 3:6}因为天主是神，显然在从祂所生的一位中，不可能有任何与祂所从生的神相异或不同的东西。因此，天主的出生使祂成为完美的天主。由此也清楚可见，我们不能说祂开始存在，只能说祂被生。因为开始与出生在某种意义上是不同的。开始存在的事物要么是从无中产生，要么是从一种状态发展为另一种状态，不再是以前的样子；例如，黄金由土形成，固体融为液体，冷变为热，白变红，水生出活动物，无生命物变成有生命。与此相反，天主子不是从无开始成为天主，而是作为天主被生；祂在神性之前并没有另一种存在。因此，那被生而为天主者，其天主性既无开始，也无发展至它。祂的出生为祂保留了那出身的本性；天主子在祂的独特存在中，就是天主所是的，而非其他。\par

15. 再者，凡对此事存疑者，可从犹太人那里获得对基督本性的准确认识；或者更恰当地说，从福音中得知祂确实被生，那里记载着：\BibleQuote{因此犹太人越发想要杀祂，因祂不但破坏了安息日，并且称天主为自己的父，使自己与天主平等。}\BibleRef{若 5:18}这段经文与大多数不同，没有给出犹太人所说的话，而是宗徒解释他们想要杀主的原因。我们看到，这些亵渎者的邪恶不能以误解为借口；因为宗徒证明基督的真实本性已完全启示给他们。他们能谈论祂的出生：——\BibleQuote{祂说天主是祂的父，使自己与天主平等。}当祂通过称天主为自己的父而宣示祂本性的平等时，难道这不显然是本性从本性的出生吗？现在显然，平等在于平等者之间没有差异。同样显然，出生的结果必然是一种在圣子与圣父之间没有差异的本性。这是真正平等的唯一可能起源；出生只能使一种与其起源平等的本性存在。但再次，我们不能认为平等存在于混乱之中，正如不能认为平等存在于差异之中一样。因此，作为肖像\BibleRef{希 1:3}的平等，与孤立和多样性不相容；因为平等不能与差异共存，也不能在孤独中存在。\par

16. 现在，虽然我们已经发现圣经的意义，按我们所理解的，与普通理性的结论一致，两者都同意平等不能与多样性或孤立共存，但我们仍须从主自己的实际话语中为我们的主张寻求新的支持。因为只有这样，我们才能抑制那任意的解释自由，这些大胆诋毁信德的人甚至敢于对主庄严的自我启示吹毛求疵。祂对犹太人的回答是这样的：\BibleQuote{子不能由自己做什么，唯有看见父做什么；因为父所作的事，子也照样作。父爱子，凡自己所作的都指示给祂；并且还要把比这些更大的工作指示给祂，叫你们惊奇。父怎样叫死人起来，使他们活着，子也照样随自己的意愿使人活着。父不审判任何人，但把审判权全交给了子，为叫众人尊敬子如同尊敬父。不尊敬子的，就是不尊敬派遣祂来的父。}\BibleRef{若 5:19}我的论证进程，按我心中所塑造的，要求逐一处理辩论的各个要点；既然我们受教知悉我们的主耶稣基督、天主子，在名、出生、本性、能力、自我启示上都是天主，那么我们信德的证明应当按此顺序逐一确立每个点。但祂的出生阻碍了这样的处理；因为对出生的考虑包括对祂的名、本性、能力和自我启示的考虑。因为祂的出生包含这一切，而这一切都属于祂，因为祂是被生的。因此，关于祂出生的论证已经进行到如此地步，以至于我们不可能再将这些其他问题留待以后逐一分开讨论。\par

17. 犹太人想要杀害主的主要理由，是祂称天主为祂的父，把自己等同于天主；因此，祂在回答中谴责了他们邪恶的激情，并以我们信德全部奥迹的阐述形式表达出来。因为就在此前，祂治好了瘫子，而他们判定祂因违反安息日而该死，祂曾说：\Quote{\Emph{我父继续工作，我也工作}}。他们因祂使用父的名称而将自身提升到天主的地位，嫉妒被激怒到极点。现在祂要肯定祂的受生，并揭示祂本性的能力，于是祂说：\Quote{\Emph{我实实在在告诉你们：子不能由自己做什么，只有看见父做什么，才做}}。祂回答的开头几句话针对的是犹太人的邪恶热忱，这热忱甚至驱使他们想要杀害祂。针对违反安息日的指控，祂说：\Quote{\Emph{我父继续工作，我也工作}}。祂希望他们明白，祂的实践由神圣权威所证明；祂同样用这些话教导他们，祂的工作必须被视为父的工作，父在祂内工作祂所做的一切。此外，为了平息他们因祂称天主为父而激起的嫉妒，祂说出了这些话：\Quote{\Emph{我实实在在告诉你们：子不能由自己做什么，只有看见父做什么，才做}}。免得祂使自己等同于天主——拥有天主圣子的名号和本性——这一事实使人不再信从祂已经受生的真理，祂说子只能看见父做什么才做。接着，为了确认我们宣认父与子时真理的救恩和谐，祂展示了祂因受生而拥有的本性；这本性的行动能力不是来自于为特定行为而陆续给予的力量，而是来自知识。祂表明这知识不是通过父以任何有形的工作作为模式，让子模仿父先前所做的事；而是通过神圣本性的运作，祂分享了神圣本性的自立，换言之，祂作为子由父所生。祂告诉他们，因为天主的权力和本性自觉地在祂内存在，祂不可能做任何祂没有看见父做的事；既然天主独生子是在父的德能中执行祂的作为，祂行动的自由范围与祂对天主父本性强力的知识范围一致；这本性是祂不可分割的，并因祂的受生而合法拥有。因为天主不是以有形的方式观看，而是因祂的本性拥有全能的目光。\par

18. 接下来的话是：\Quote{\Emph{因为凡祂——父——所做的，子也同样做}}。这\Emph{同样}是加进来指示祂的受生；\Emph{凡}和\Emph{相同}指示祂本性真实的神性。\Emph{凡}和\Emph{相同}使得祂不可能有任何与父的行动不同或超出父的行动。因此，那本性有能力做一切与父相同之事的那一位，与父包含在同一本性之内。但是，与此对比，当我们读到所有这些相同的事是由子\Emph{同样}做时，工作与另一位相似这一点推翻了那种认为祂是孤立工作者的假设。因此，父所做的相同之事全都由子同样地做。这里我们有了祂真实受生的清楚证据，同时也令人信服地证明了我们信德的奥迹，它建基于天主本性的一体，承认父和子内存在不可分的神性。因为子做与父相同的事，并且同样地做；在祂以相似方式行动的同时，祂做相同的事。两个真理结合在一个命题中：祂的工作做得同样证明了祂的受生；它们相同的工作证明了祂的本性。\par

19. 因此，在我们主的答复中所包含的渐进启示，与教会信德宣认中对真理的渐进陈述是一致的。两者都不分割本性，都宣认受生。因为基督接着说的话是：\BibleQuote{圣父爱圣子，并将自己所做的一切事指示给祂；还要将比这些更大的工程指示给祂，使你们惊奇。因为圣父怎样使死人复活并赐予生命，圣子也同样随己意赐予生命。}天主工作方式的这个启示，除教导真正的受生——相信一个由实在的天主、祂的圣父所生的实在的圣子——之外，还能有其他目的吗？唯一的另一种解释是，天主的独生子如此无知，以至于需要在此指示中所传达的教导；但这一建议的鲁莽亵渎使得这种替代解释不可能。因为祂既知道一切所教导的事，就不需要教导。因此，在\BibleQuote{圣父爱圣子，并将自己所做的一切事指示给祂}这句话之后，我们接着被告知，这所有的指示都是为了我们信仰的教导；使圣父和圣子在我们的宣认中享有同等的份额，并且我们因此得救，避免以为圣子的知识不完全的错觉，因为圣父将祂所做的一切都指示给圣子了。为此，祂接着说：\BibleQuote{祂还要将比这些更大的工程指示给祂，使你们惊奇。因为圣父怎样使死人复活并赐予生命，圣子也同样随己意赐予生命。}我们看到圣子完全知道圣父将来要指示给祂的未来的工程。祂知道自己将被告知，如何效法圣父的榜样，使死人得生命。因为祂说圣父将指示给圣子一些事，这些事将使他们惊奇；并立即告诉他们这些事是什么：\BibleQuote{因为圣父怎样使死人复活并赐予生命，圣子也同样随己意赐予生命。}权能是相等的，因为本性是同一的。工程的指示并非为弥补祂的无知，而是为加强我们的信德。它并没有给圣子带来未知事物的知识，而是使我们有信心宣认祂的受生，因为向我们保证圣父已将祂所能做的一切工程都指示给了圣子。这神圣讲论中的措辞是经过极其审慎地挑选的，以免任何含糊的语言暗示两者之间的本性差异。基督说圣父的工程被指示给祂，而不是说为了能使祂完成这些工程，赐给了祂一种强大的本性。祂借此想向我们揭示，这个指示是祂受生过程的一个实质性部分，因为与受生同时，由于圣父的爱，祂被赋予了圣父愿意祂做的工程的知识。此外，为了免使我们因这个指示的宣告而被引导去认为圣子的本性是无知的，因而与圣父不同，祂明确指出祂已经知道将要指示给祂的事。事实上，祂远不需要先例的权威来使自己行动，以致祂随己意赐予生命。意愿意味着一个自由的本性，它存在于拥有选择能力的权能中，在全能的喜乐运用中。\par

20. 其次，为了避免似乎随己意赐予生命并不属于真正受生者的权能，而只是非受生全能的特权，祂急忙补充说：\BibleQuote{圣父不审判任何人，但已将一切审判交给了圣子。}一切审判都已交给祂的陈述，教导了祂的受生和祂的圣子身份；因为只有与圣父完全合一的本性才能拥有一切；而圣子除了藉恩赐外，不能拥有任何事物。但一切审判都已赐给祂，因为祂随己意赐予生命。现在我们不能认为审判从圣父那里被夺走了，尽管祂不执行审判；因为圣子全部的审判权能都来自圣父，是来自圣父的恩赐。并且，审判赐给圣子的原因并未隐瞒，因为接下来的话是：\BibleQuote{但祂已将一切审判交给了圣子，为叫众人尊敬圣子如同尊敬圣父。不尊敬圣子的，就是不尊敬派遣祂来的圣父。}怀疑或不敬的否认还有什么借口呢？审判之恩赐的原因，是使圣子得到与圣父同样的尊敬；因此，侮辱圣子的人也犯了侮辱圣父的罪。经过这一证明，我们怎能想象祂因受生而得的本性不同于圣父，当祂在工程、权能、尊荣以及分派给不信者的惩罚上，与圣父相等呢？因此，整个神圣答复无非是对祂受生奥秘的揭示。圣父与圣子之间唯一正确或可能做出的区分，就是后者是受生的；但受生的意义是祂与圣父成为一体。\par

21. 因此，圣父工作到如今，圣子也工作。在圣父和圣子中，你们有表达它们彼此关系的名字。也要注意，是天主性，即天主借以工作的本性，在这里工作。请记住，免得你们陷入错误，以为这里描述的是两种不同本性的运作；关于那瞎子，曾如此说：\BibleQuote{但为叫天主的工程在他身上显扬出来，我应当做那派遣我来的者的工程。}\BibleRef{若 9:3} 你们看见，在他身上，圣子所行的工程是圣父的工程；圣子的工程是天主的工程。我们所考虑的话语的其余部分也论及工程；但我现在的辩护只关心将全部工程归于两者，并指出他们在工作方式上是一致的，因为圣子从事于圣父向来所做的工程。这一事实所含的认可——因他神圣的受生，圣父在他所做的一切事上与他一同工作——将拯救我们免于认为安息日的主在安息日工作是不对的。他的子身份不受影响，因为他的天主性与圣父的没有混淆，也没有否定；他的天主性不受影响，因为他的天主本性未受触及。他们的合一不受影响，因为没有显明差异来分离他们；他们的合一也不是以否定他们各自存在的方式呈现的。首先承认圣子的子身份；\BibleQuote{子不能由自己做什么，除非他看见父做什么。} 这里他的受生是明显的；因此，在他看见它被做成之前，他不能由自己做什么。他不能是非受生的，因为他不能由自己做什么；他没有发起的能力，因此他必然是受生的。但他能看见圣父的工程这一事实证明他拥有属于天主性的有意识拥有者的理解。其次，注意他确实拥有这真实的天主性——\BibleQuote{凡他所做的事，子也同样做。} 现在我们已经看见他被赋予那本性的能力，注意这如何导致合一，一个本性如何居住在两者之中——\BibleQuote{为叫众人尊敬子，如同尊敬父一样。} 然后，免得对这合一的思考使你陷入一位孤独自足的天主的错觉，请留心这些话语中所显示的信德的奥迹：\BibleQuote{那不尊敬子的，就不尊敬派遣他的父。} 异端的愤怒和诡计可能尽其所能；我们的立场是不可攻破的。他是子，因为他不能由自己做什么；他是天主，因为凡父所做的，他也同样做；两者是一，因为他在尊荣上与父平等，并做完全相同的工程；他不是父，因为他被派遣。这一受生的道理所包含的奥秘真理是多么丰富！它包含他的名、他的本性、他的能力、他的自我启示；因为在他受生中赋予他的一切，必然包含在那从他受生而来的本性中。在他的本性中，没有引入任何与他本源不同种类的实体元素，因为从单一本性而出的本性，必定与其母本性完全同一。合一是不包含不和谐元素、与其自身同类的；通过受生构成的合一不能是孤独的；因为孤独只能有一个居住者，而通过受生构成的合一却意味着两者的结合。\par

22. 此外，让祂自己的天主性言语为祂作证。祂说：\BibleQuote{属于我的羊听我的声音，我也认识他们，他们跟随我；我赐给他们永生，他们永不丧亡，谁也不能从我手里把他们夺去。我圣父赐给我的，比一切更大，谁也不能从我圣父手里把他们夺去。我与圣父原是一体。}\BibleRef{若 10:27} 什么倦怠能如此彻底地钝化我们悟性的锋刃，以至这如此明确的陈述在我们看来片刻晦涩？什么骄傲的诡辩能如此戏弄人的驯良，以至那些从这些话语中学会了天主是什么的人，竟在他们面前启示了天主性的祂身上，不承认天主？异端要么应提出其他福音来支持自己的教义；要么，若我们现有的福音是唯一教导关于天主的文献，他们为何不信所受的教导？若它们是知识的唯一来源，为何不从它们中汲取信德，如同汲取知识一样？然而现在我们发现，他们的信德是公然对抗知识而持守的；因此，这信德不是植根于知识，而是植根于罪恶；是对公开承认的真理的一种大胆不敬的信德，而非敬畏的谦卑。天主独生子，如我们所看到的，完全确信自己的性体，以最精确的语言启示了祂出生的奥秘。祂启示了那虽是无可言喻的奥秘，使我们可以相信并宣认它；使我们能理解祂是受生的，并相信祂拥有天主的性体，与圣父是一体，其一体之意在于：天主并非独一，圣子也不是圣父的别名，而是祂确实是圣子。首先，祂向我们保证祂天主性体的权能，论到祂的羊说：\BibleQuote{谁也不能从我手里把他们夺去。}这是自觉能力的宣告，是对自由且不可抗拒的力量的承认，绝不容许任何人从祂手中夺走祂的羊。不仅如此；祂不仅拥有天主的性体，更要我们知道，这性体是因祂由天主所生而属于祂的，因此祂接着说：\BibleQuote{我圣父赐给我的，比一切更大。}祂毫不隐瞒自己由圣父所生，因为祂从圣父所领受的，祂说比一切更大。那领受者，在祂出生时领受了，而不是在出生之后；然而这来自另一位，因为祂是领受者。但那从另一位领受这恩赐的，禁止我们设想祂自己与那一位在种类上不同，也不是与赐恩者永恒地具有相同性体，祂说：\BibleQuote{谁也不能从我圣父手里把他们夺去。}没有人能从祂手里把他们夺去，因为祂从圣父那里领受了比一切更大的事物。那么，这矛盾的断言是什么意思——谁也不能从圣父手里把他们夺去？圣子的手从圣父那里领受了他们，圣父的手将他们赐给了圣子：说什么不能从圣子手里夺走的也不能从圣父手里夺走，这是什么意思？请听，如果你愿意知道：\BibleQuote{我与圣父原是一体。}圣子的手就是圣父的手。因为天主性体在通过生而传递时，并不会变质或不再相同；这种相同也不妨碍我们对生的信德，因为在那生中，没有异己的成分被接纳进祂的性体。在此祂提到圣子的手，就是圣父的手，好使你借着身体的比拟，学习存在于两者中的唯一天主性体的能力；因为圣父的性体和能力就在圣子内。最后，为使你在生这奥秘的真理中，能辨别天主性体真正不可区分的合一，才说了这话：\BibleQuote{我与圣父原是一体。}说这话，为叫我们在这合一中看不到差异或孤立；因为祂们是两位，然而并没有第二个性体借着那真实的生和受生而存在。\par

23. 若我理解正确，这些疯狂灵魂中仍存留着同样的欲望，尽管他们实现它的机会已经丧失。他们苦涩的心灵仍然怀有对邪恶的渴望，却再无法指望满足。主在天上的宝座，而异端的狂怒不能再像犹太人那样将祂拖上十字架。但不信的精神是一样的，虽然现在它以否认祂的天主性的形式出现。他们向祂的话语挑战，虽然不能否认祂曾说过。他们用亵渎发泄仇恨；用谩骂代替石头。若可能，他们甚至会将祂从宝座上拉下来，再次钉死。当犹太人因基督教导的新奇而愤怒时，我们读到：\BibleQuote{犹太人因此拿起石头要砸死祂。祂回答他们说：我从父显示了许多善行给你们，为了哪一件你们要砸死我呢？犹太人回答祂说：为了善行我们不砸死祢，而是为了亵渎，因为祢是人，却把祢自己当作天主。}\BibleRef{若 10:31} 异端者啊，我提醒你，要在其中认出你自己的行为，你自己的话。你要确信你是他们的同伙，因为你已以他们的不信为榜样。正是在那句话，\Quote{我与父原是一体}，犹太人拿起了石头。他们对那救恩奥迹启示的不敬恼怒，甚至促使他们企图杀害。现在没有人你可以投石；但你的罪过是否比他们小？意愿相同，尽管被祂天上的宝座所阻。不，你比犹太人更不敬。他向身体投石，你向圣神投石；他以为是向人，你是向天主；他向地上的一位过客，你向坐在威严宝座上的那一位；他向一位他不认识者，你向你所明认的那一位；他向可死的基督，你向宇宙的审判者。犹太人说：\Quote{是人}；你说：\Quote{是受造物}。你和他同声呼喊：\Quote{把祢自己当作天主}，带着同样的亵渎傲慢。你否认祂是由天主所生的天主；你否认祂是出于真实诞生的子；你否认祂的话，\Quote{我与父原是一体}，包含两者同一本性的宣称。你以一位现代的、陌生的、异样的天主顶替祂，使祂成为与父不同类的天主，或者根本不是天主，因为不是由天主生育而存在。\par

24. 那话\Quote{我与父原是一体}所含的奥迹激起你的愤怒。犹太人回答说：\Quote{祢是人，却把祢自己当作天主}；你的亵渎与他匹敌：——\Quote{祢是受造物，却把祢自己当作天主。}你实际上说：\Quote{祢不是由生育而生的子，祢不是真天主；祢是受造物，超越所有其他受造物。但祢并非生而为天主，因为我拒绝相信无形体的天主生育了祢的本性。祢和父不是一体。不仅如此。祢不是子，祢不像天主，祢不是天主。}主对犹太人有回答；这回答甚至比应对他们的亵渎更适用于你的亵渎：——\BibleQuote{在你们的法律上不是记载着：\InnerQuote{我说过：你们是神}么？如果法律上称那些承受天主话的人为神，而经书是不能废弃的，那么，父所祝圣并派遣到世界上来的，因为说过：\InnerQuote{我是天主子}，你们就说：你说亵渎的话么？假使我不作我父的工作，你们就不必信我；但若是我作了，你们纵然不肯信我，至少要信这些工作，如此你们必定认出父在我内，我在父内。}\BibleRef{若 10:34} 这回答的内容是受对祂亵渎攻击的内容所决定的。指控是祂是人，却把自己当作天主。他们对这一主张的证明是祂自己的声明，\Quote{我与父原是一体}。因此祂着手证明，那属于祂本性的天主性，使祂有权宣称祂和父是一体。祂首先揭露了这样的指控——即祂虽然是人，却把自己当作天主——的荒谬和傲慢。法律把这一称号赐给了圣人；天主的圣言——不可上诉——认可了这一名称的公开使用。那么，父所祝圣并派遣到世界上来的那一位，自称天主之子，有什么亵渎可言？天主的圣言不可改变的记录，已经为法律所指定的人确认了这一称号。因此，当法律把\Quote{神}的称号给予那些明明白白是人的人时，那指控祂是人却把自己当作天主的控告就瓦解了。而且，如果其他人可以使用这个名称而无亵渎，那么父所祝圣的那个人使用它显然也无亵渎——注意，在整段论证中祂称自己为人，因为天主子也是人子——因为祂超越其余的人，而那些人自称神却毫无不敬。祂超越他们，因为祂被祝圣为子，正如蒙福的保禄所说，他教导我们这一圣化：——\BibleQuote{这福音是天主先前借自己的先知在圣经上所预许的，是论及他的儿子，我们的主耶稣基督，他按肉身是生于达味的后裔，按至圣的神性，由于他从死者中复活，被立为具有大能的天主之子。}\BibleRef{罗 1:2} 因此，关于祂把自己当作天主的亵渎指控，就归于无有。因为天主的圣言已将这名称赐给许多人；而祂，被圣父祝圣并派遣的那一位，只不过宣告自己是天主子而已。\par

25. 我认为，毫无疑问，\BibleQuote{\Emph{我与父是一体}}这话是论及祂由受生而得的本性。犹太人曾斥责祂，因为祂作为人，却藉这些话使自己成为天主。祂答复的过程证明，在这\BibleQuote{\Emph{我与父是一体}}中，祂确实是自称天主子，首先是名称上，然后是本性上，最后是因受生。因为\Emph{我}和\Emph{父}是实体存在的名称；\Emph{一体}是宣示他们的本性，即本质上是同一的；\Emph{是}禁止我们将两者混淆；\Emph{是一体}在禁止混淆的同时，教导两者的一体是受生的结果。现在，这一切真理都从那名称——天主子——中引出，这名称是祂被父祝圣后，赐给自己的；祂对这名称的权利，因祂宣称\BibleQuote{\Emph{我与父是一体}}而得以确认。因为受生不能将任何本性赋予子嗣，除了那产生子嗣的父母的本性之外。\par

26. 天主的独生子再次用祂自己的话语，为我们总结了整个启示的信仰奥秘。当祂答复了那指责——祂作为人，却使自己成为天主——之后，祂决意表明祂的话\BibleQuote{\Emph{我与父是一体}}是一个清楚而必然的结论；因此祂这样继续论证：\BibleQuote{\Emph{你们说我亵渎了，因为我说了：我是天主子。我若不做我父的工作，你们就不要信我；但若做了，你们纵然不信我，也要信这些工作，好使你们知道并确信：父在我内，我在父内}}。在此之后，仍坚持其道路的异端有意识地陷入绝望，恣意妄为；不信的断言是蓄意的无耻。坚持己见者以愚昧为荣，且死于信德，因为反驳此话不是无知，而是疯狂。主曾说过\BibleQuote{\Emph{我与父是一体}}；祂所启示的受生奥秘，就是父与子在本性上的合一。再者，当祂因自称为天主本性而被控告时，祂提出理由来辩护：\BibleQuote{\Emph{我若不做我父的工作，你们就不要信我}}。除非祂做父的工作，我们不应相信祂自称是天主子。由此我们看出，祂的受生并未赋予祂任何新奇或异类的本性，因为做父的工作正是我们必须相信祂是子的理由。在被证明是天主子的证据乃是祂做属于父的本性工作时，这里还有什么收养、或使用名称的许可、或否认祂是由天主本性所生呢？受造物没有与天主相等或相似的，任何外在于祂的本性在能力上都不可与祂相比；唯有那从祂自己而生的子，我们可以不亵渎地比拟并相等祂。任何外在于祂自己的存在若被证明像祂且在能力上与祂相等，那么，天主在接纳一个合作伙伴分享祂的宝座时，就丧失了祂的卓越。天主不再是独一的，因为第二个与祂不可区分者出现了。另一方面，使祂自己的真子与祂平等并无侮辱。因为那像祂的是祂自己的；那与祂相比的是从祂自己而生的；那能做祂自己工作的能力并非外在于祂。不仅如此，这实际上是祂荣耀的增高，因为祂生出了全能，却没有将那全能的本性从自己割裂。子做父的工作，并基于此要求我们相信祂是天主子。这不是单纯的傲慢之词，因为祂将之建立于祂的工作上，并让我们审视它们。祂亲自作证这些工作不是祂自己的，而是祂父的。祂不愿我们的思想被这些事迹的光彩所分心，从而忽略祂受生的证据。因为犹太人不能洞悉祂所取的身体——即由玛利亚所生的人性——的奥秘，而认出天主子，祂便求助于祂的事迹来证实祂对此名称的权利——\BibleQuote{\Emph{但若我做了这些工作，而你们又不信我，就相信这些工作吧}}。首先，祂不愿他们相信祂是天主子，除非有祂所做的天主的工作为证。其次，如果祂做了这些工作，却因身体上的谦卑而似乎不配承受神圣的名称，祂要求他们相信这些工作。为何祂按人性出生的奥秘会阻碍我们认出祂作为天主的出生呢？因为那由天主所生者，借着祂所取了的人性，完成了每一项神圣的任务。如果我们不信那人，却因工作的缘故，当祂告诉我们自己是天主子时，就让我们相信这些工作吧，因为它们无疑是天主的工作，并由天主子显明地完成了。因为天主子因受生而拥有天主的一切；因此子的工作就是父的工作，因为祂的受生并未将祂排除在那作为祂根源且祂常住的本性之外，也因祂本身拥有那本性，祂的存在永恒归于此。\par

27. 因此，圣子既行圣父的工程，又要求我们：若不信祂，至少因祂的工程而信，祂必告诉我们，我们应从工程相信什么要点。祂在接下来的话中告诉我们：——\Emph{\BibleQuote{但我若作了，你们纵然不肯信我，至少该信这些工程，好使你们知道并确知：父在我内，我在父内}}。这同一真理也包含在\Emph{\BibleQuote{我是天主子}}和\Emph{\BibleQuote{我与父原是一体}}中。这是祂由\OriginalTerm{出生}{birth}而得的\OriginalTerm{本性}{nature}；这是得救信德的奥秘：我们既不可分割合一，也不可将本性同出生分开，而必须宣认\OriginalTerm{生活的天主}{living God}确实从生活的天主所生。天主，祂是\OriginalTerm{生命}{life}，不是由各种无生命的零件组合而成的存有；祂是\OriginalTerm{德能}{power}，不是由软弱的元素凑成；祂是\OriginalTerm{光明}{light}，与黑暗的阴影毫无掺和；祂是\OriginalTerm{神}{Spirit}，不能与任何不协调之物相谐。祂内在的一切都是一；神是什么，就是光明、德能和生命；生命是什么，就是光明、德能和神。那说\Emph{\BibleQuote{我是，且不改变}}\BibleRef{拉 3:6}的一位，既不能在细节上改变，也不能在种类上转化。因为这些我提及的属性，并非附着于祂的不同部分，而是在整个生活的天主的存有中全然完美地相遇并合一。祂是生活的天主，是生活的天主本性的永恒德能；那由祂所生的一位，按照祂所启示的奥秘真理，不可能不是生活的。因为当祂说：\Emph{\BibleQuote{如同生活的父派遣了我，我因父而生活}}\BibleRef{若 6:57}，祂教导我们，祂是由于生活的父而在自己内有生命。此外，当祂说：\Emph{\BibleQuote{因为父怎样在自己内有生命，也赐给子同样在自己内有生命}}，祂作证，生命至圆满的程度，是祂从生活的天主得来的恩赐。如今，如果生活的子是从生活的父所生，那出生并未使一个新本性产生存在。当生活者由生活者所生时，没有新事物产生存在；因为生命并非从无中寻求而受生；而生命从生命受生，由于本性的合一和那完美、不可言说的出生事件，必定常住在祂内生活，并在自己内拥有生活者的生命。\par

28. 我记得，在本论著的开头，我曾警告说：人类的类比不完全对应于神圣的原型，然而，通过将后者与可见的形象相比较，我们的理智确实得到了真实（即使不完全）的启发。现在我诉诸人类在\OriginalTerm{出生}{birth}方面的经验：子女存在的根源是否不留在父母之内？因为，即使无生命和卑贱的物质（它启动\OriginalTerm{生命}{life}的开端）从一方父母传入另一方，但父母各自保持着自然的力量。他们产生了一个与自己\OriginalTerm{本性}{nature}合一的本性，因此生育者与受生者的存在紧密相连；受生者通过生育者传递且未丧失的力量而接受出生，便寓居于那生育者之内。作为对人类出生中发生的事情的陈述，这也许足够了。但作为与天主独生子的完美出生的类比，它是不足的；因为人类在软弱中出生，从两个不同本性的结合而生，并由无生命物质的组合维持生命。此外，人类并非立刻进入其指定的生活实施，也从未完全活出那种生活，总是被众多易朽坏和不自觉抛弃的肢体所累。相反，在天主内，神圣生命得以最圆满地活出，因为天主是生命；从生命只能生出真正可生活的事物。祂的出生并非以流出方式，但出于\OriginalTerm{德能}{power}的行动。因此，既然天主的生命在其强度上是完美的，既然由祂所生者在德能上是完美的，天主具有生育的德能，但不承受改变。祂的本性能够增加，而非减少，因为祂继续存在于祂所生的圣子内，并分享那圣子的生命；祂在出生时赐给了圣子一个与自己相似且不可分离的本性。那圣子，生活者从生活者所生，并不因祂的出生事件而与那生了祂的本性分离。\par

29. 另一个对信德含义有所启发的类比是\OriginalTerm{火}{fire}：火在自己内含有火，且寓居于火中。火含有\OriginalTerm{光明}{light}的亮度、作为其本质属性的\OriginalTerm{热}{heat}、以及通过燃烧摧毁火焰闪烁不定的特性。然而它始终是火，在所有这些表现中只有一个\OriginalTerm{本性}{nature}。它的弱点在于它的存在依赖于可燃物质，并且它随它所赖以存在的物质而消亡。与火的比较，在某种程度上使我们得以洞察天主那无可比拟的本性；它帮助我们相信天主的属性，我们发现这些属性在某种程度上存在于一个尘世元素中。那么我问：从火中产生的火，是否存在任何\OriginalTerm{分离}{division}或分割？当火焰从另一火焰点燃时，原初的本性是否与派生的本性切断，以致不再寓居其中？它岂不反而如同通过出生那样，以一种增加的方式跟随并居住在第二火焰中？因为从第一火焰的本性中没有切下任何部分，却有了光从光而来。第一火焰岂不活在第二火焰中，而第二火焰存在——虽非通过分割——归于第一火焰？第二火焰岂不仍然居住在第一火焰中，它未曾与之切断；它从那里出去，保持与它本性所属的实体的合一？当从实体上不可能通过分割从光得到光，并且从逻辑上不可能在本性上区分它们时，那二者岂不是一？\par

30. 我再说一次，这些比喻只可用作理解信德的辅助，决不可作为与天主尊威比较的标准。我们的方法是借有形之物为线索，以领悟无形之事。敬畏与理性都容许我们使用这样的帮助，因为在天主自己的见证中我们也看到祂使用了它们，但我们并不奢望找到与天主性体完全平行的东西。然而，朴实信友的内心曾因异端者疯狂的反对而困扰，他们宣称：接受一个需要借助有形比喻才能理解的关于天主的道理是错误的。因此，依照我们已引述的主的话：\BibleQuote{那由肉而生的，是肉；那由圣神而生的，却是圣神。}\BibleRef{若 2:6}，我们认为应该（既然天主是圣神）在论证中给这些比喻留一席之地。这样做可以避免天主因使用这些比喻而被指责为欺骗了我们：正如我们已表明的，祂受造物本性的这类比喻使我们能把握天主自我启示的意义。\par

31. 我们看到永生之父的永生之子，那位出自天主的，天主，如何在以下话语中揭示了天主性体不可分解的同一与一体，以及祂诞生的奥秘：\Quote{我与父原是一体。}因为这话中看似的自傲招致了对祂的偏见，祂更清楚地表明祂是在拥有天主性的自觉中说话，说道：\Quote{你们说我亵渎了，因为我曾说：我是天主子。}由此显示祂与天主性体的同一源于从天主而生。然后，为了用明确的断言来巩固他们对祂诞生的信德，同时防止他们以为这诞生包含性体的差异，祂以这话结束论证：\Quote{你们当因这些工作相信：父在我内，我在父内。}祂所启示的诞生，是否显示出祂的天主性并非本性所拥有、并非自有权利？两者互在其中；圣子的诞生仅来自圣父；没有异己或不同的性体被提升到天主性而作为天主存在。出自天主的天主，永恒存在，祂的天主性除了天主外别无来源。你若找到机会，就将两位天主带入教会的信仰；你尽可能将圣子与圣父分离，但既然你承认诞生，父仍在子内，子仍在父内；这不是因为流溢的互换，而是因为活的本性圆满的诞生。因此你无法将天主父与天主子相加，视其为两位天主，因为祂们两位是同一至尊。你也不能混淆祂们，因为祂们两位不是一个位格。所以宗徒信德拒绝两位天主；因为它不知有两位父或两位子。在宣认父时，它宣认子；它相信子，因它相信父。因为父的名字包含子，因为没有儿子，谁也不能为父。子存在的证据证明曾有父，因为子只能由父而生。当我们宣认天主是独一至尊时，我们否认祂是单一的；因为子是父的圆满，子存在是由于父。但诞生并未改变天主性体；在父和子内，这性体都按自己的种类真实存在。对这一性体一体的正确表达，就是宣认祂们虽因诞生与受生而为二位，却是一位天主，而非一个位格。\par

32. 让那不愿宣认一体却能宣讲一位天主的人去宣讲两位天主吧；让那否认有位格互居（因本性的力量和给予与接受的诞生奥迹）的人去宣告天主是孤独的吧。让那不知父与子的合一是一项启示真理的人去给两位分别指派不同的性体吧。让异端者抹去圣子自我启示的话：\Quote{我在父内，父在我内；}然后他们才能断言有两位天主，或一位孤独的天主。关于他们拥有同一性体，我们所得的启示中并无暗示有一个以上性体。出自天主的天主并不使天主倍增为二；诞生排除了孤独天主的假定。再者，因为祂们互相内在，形成了合一；祂们互相内在，这由一位出自一位的事实所证明。因为那一位在生那位时，并未赐予任何不属于祂自己的东西；而那一位在被生时，从那一位只领受了属于那一位的东西。因此，宗徒信德在宣认父时，宣认祂为一位天主；在宣认子时，宣认祂为一位天主；因为同一的天主性体就在两者内，且因为父是天主，子是天主，天主的同一名称表达了两者的性体，\Quote{一位天主}这词指明两者。出自天主的天主，或在天主内的天主，并不意味有两位天主，因为天主（一位出自一位）永恒地持有同一的天主性体与同一的天主名称；天主也不缩减为一个位格，因为一位和一位永不会在孤独中。\par

33. 主并没有使这伟大奥迹所传达的教训处于疑惑或隐晦之中；祂没有遗弃我们在朦胧的不确定性中迷失方向。请听祂向祂的宗徒启示这信德的全备知识时所说的话：——\BibleQuote{我是道路、真理、生命；除非经过我，谁也不能到圣父那里去。你们若认识我，也就认识我的圣父；从现在起，你们已认识祂，并且已看见祂。斐理伯对祂说：\InnerQuote{主！请把圣父显示给我们，这对我们便足够了。} 耶稣对他说：\InnerQuote{斐理伯！这么长久的时间，我与你们在一起，你还不认识我吗？看见我的，就是看见了圣父；你怎么说：把圣父显示给我们呢？你不信我在圣父内，圣父在我内吗？我对你们所说的话，不是凭我自己说的；而是住在我的圣父，祂做祂的工程。你们当信我，我在圣父内，圣父在我内；不然，至少因那些工程而信。}}。\BibleRef{若 14:6} 祂是道路，引领我们不走入歧途或荒芜之地；祂是真理，不以谎言嘲讽我们；祂是生命，不将我们出卖于导致死亡的幻觉。祂亲自选用了这些悦耳的名号，以指明祂为我们得救所指定方法。作为道路，祂将引导我们到达真理；真理将使我们立足于生命。因此，认识祂所启示的获得这生命的奥秘方式，对我们至关重要。\BibleQuote{除非经过我，谁也不能到圣父那里去。} 到圣父那里的道路是通过圣子。现在我们必须探究，这是否要通过服从祂教导的历程，还是通过信德相信祂的天主性。因为可以设想，我们到圣父那里的道路可能是通过坚守圣子的教导，而不是通过相信圣父的天主性住在圣子内。因此，接下来让我们寻求在此所赐予的教导的真义。因为我们不是通过坚持预想的意见，而是通过研究话语的力量，来拥有这信德。\par

34. 紧接着上述引文的，是这些话：\BibleQuote{你们若认识我，也就认识我的圣父。} 他们所注视的是那人耶稣基督。对祂的认识怎能成为对圣父的认识呢？因为宗徒们看见祂带着属于祂的人类本性的面貌；但天主并不受身体和血肉的累赘，并且不能被居住在我们软弱有肉身的身体内的人所认识。答复来自主，祂断言在祂所奥秘地取了的肉身之下，祂圣父的本性住在祂内。祂按适当顺序陈述事实如下：\BibleQuote{你们若认识我，也就认识我的圣父；从现在起，你们已认识祂，并且已看见祂。} 祂在看见的时间和认识的时间之间作了区分。祂说从现在起，他们将认识他们早已看见的那一位；并且从这启示的时刻起，他们将拥有那本性的知识——他们长久以来在祂内注视的那本性。\par

35. 但这些话语的新奇声音惊动了宗徒斐理伯。一个人在他们眼前；这人自称是天主子，并宣告当他们认识祂时，他们也会认识圣父。祂告诉他们，他们已经看见了圣父，并且因为他们看见了祂，他们以后要认识祂。这真理对软弱的人类而言过于宽广；他们的信德在这些悖论面前失效了。基督说圣父已经看见了，且现在将被认识；并且，尽管是看见，这就是认识。祂说，如果圣子被认识了，圣父也被认识了；并且尽管圣子通过可见和可听的感官传达了关于祂自身的知识，而圣父的本性与他们所认识的可见之人的本性完全不同，无法从他们所见过的那一位本性的知识中得知。祂也多次作证说没有人见过圣父。因此斐理伯以宗徒的忠诚和信心，迸发出祈求：\BibleQuote{主！请将圣父显示给我们，这对我们便足够了。} 他并非在玩弄信德；这只是一个出于无知的错误。因为主曾说过圣父已经看见了，且从此应当被认识；但宗徒并没有理解祂已经被看见。因此他没有否认圣父已经被看见，而是请求看见祂。他并非祈求圣父被展示给他肉体的凝视，而是祈求获得某种启示，以光照他关于已被看见的圣父。因为他曾在人的面貌下看见圣子，但无法理解他如何因此看见了圣父。他在祈祷\BibleQuote{主！请将圣父显示给我们}中加上\Quote{这对我们便足够了}，清楚表明他所渴望的是对圣父的精神性而非肉体性的神视。他并未拒绝相信主的话语，而是请求对他的心智给予光照，以使他能够相信；因为主已说话这一事实对宗徒来说就是他有责任相信的确凿证据。促使他请求显示圣父的考虑是：圣子曾说过，祂已被看见，并应因曾被看见而被认识。在这祈祷中，并无假想那已被看见者现在应被显明出来。\par

36. 因此主这样回答\Person{斐理伯}——\Quote{我与你们在一起这么久了，你们还不认识我吗，\Person{斐理伯}？}祂责斥宗徒对祂认识不足；因为此前祂曾说过，当祂被认识时，圣父也就被认识了。但这抱怨说他们这么久还不认识祂是什么意思？这就是说，如果他们认识祂，就必须在祂身上认出属于祂圣父本性的天主性。\par

37. 因此主责备他们虽然这么久行了这些事，却不认识祂，并以回答他们祈求祂显示圣父的话说：\Quote{看见我的，就是看见了圣父}。祂不是指肉身的显现，不是肉眼的感知，而是指祂曾说过的那眼——\Quote{你们不是说：还有四个月才到收获期吗？看，我给你们说：举起你们的眼，观看田地，因为已经发白，可以收割了。}\BibleRef{若 4:35}。一年的季节、发白的田地这些比喻都与地上的、可见的景象不相符。祂是叫他们举起理智的眼目，默想最后丰收的福乐。同样地，祂现在的话\Quote{看见我的，就是看见了圣父}。祂从\Person{童贞玛利亚}所生所得的肉身，不能将天主的肖像和相似显示给他们。祂所穿戴的人性相貌，对于灵性看见无形无像的天主没有任何帮助。但天主在基督里被那些藉着祂神性本性的德能认出基督是天主子的人所认识；而对天主圣子的认识就产生对天主圣父的认识。因为圣子作为肖像，其意义是与圣父在同一类中，而又指明圣父是祂的根源。其他用金属、颜料或其他材料通过各种技艺制作的肖像，再现了所代表物体的外观。然而，没有生命的复制品能和有生命的原版相提并论吗？绘画、雕刻或铸造的肖像能与它们所模仿的本性相比吗？圣子作为圣父的肖像并非如此；祂是那生活者的生活肖像。由圣父所生的圣子，其本性丝毫不异于圣父；因为祂的本性没有不同，祂就拥有与自己相同的本性的权能。祂是肖像这一事实，证明天主圣父是独生子的诞生之源，而圣子自己被启示为无形天主的相似和肖像。因此，这相似与天主性结合，是不可磨灭地属于祂的，因为那本性的权能是不可转让地属于祂自己的。\par

38. 这就是这段经文的意思：\Quote{我与你们在一起这么久了，你们还不认识我吗，\Person{斐理伯}？看见我的，就是看见了圣父。你怎么说：把圣父显示给我们呢？你不信我在圣父内，圣父在我内吗？} 我们人类在谈论神圣之事时，只能通过天主圣言来推理。其他属于天主性的一切，都是黑暗而困难、危险而模糊的。如果有人提议用天主所提供的以外的词汇来表达已知之事，他必然要么显露自己的无知，要么使读者的思想完全困惑。主被要求显示圣父时，祂说：\Quote{看见我的，就是看见了圣父}。想要改变这话的是敌基督；否认这话的是犹太人；不知道这话的是外邦人。如果我们遇到困难，让我们归咎于自己的理性；如果天主的宣告似乎有模糊之处，让我们假定这是我们的信德不足所致。这些话精确地陈述了天主并非独居，而又在天主性内没有差异。因为圣父在圣子内被看见，如果祂是独居的，或者如果祂不像圣子，这都不可能。圣父是通过圣子被看见的；而圣子所启示的奥秘就是他们是同一的天主，但不是同一位格。对于主的话\Quote{看见我的，就是看见了圣父}，你还能赋予什么别的含义？这并非同一；连词\Quote{也}表明圣父是除了圣子之外被命名的。这些话\Quote{圣父也}与一个孤立、单一的一位格的概念不相容。唯一可能的结论就是圣父通过圣子成为可见的，因为他们是一体且在本性上相似。并且，为了使我们在这方面的信德不留任何疑惑，主接着说：\Quote{你怎么说：把圣父显示给我们呢？} 圣父已经在圣子内被看见；那么人怎么还不认识圣父？有什么需要显示圣父呢？\par

39. 再者，生者与受生者的合一，彰显于本性相同与种类真实一体，证明圣父在其真正本性中被看见。这由主接下来的话显示：\Quote{\Emph{你们不信我在父内，父在我内吗？}}\BibleRef{若 14:10}除子所用的这些话外，无法以其他言词陈述这一事实：圣父与圣子，在本性上相似，是不可分离的。子，祂是道路、真理和生命，并没有以某种名字与形貌的戏剧性变换来欺骗我们，当祂穿戴人性时，称自己为天主子。祂并没有虚假地隐瞒祂是天主父的事实；祂不是一个隐藏面貌于面具之下的单一实体，好让我们以为有两者同在。祂不是一个孤存者，时而伪装为自己的子，时而又称自己为父；用变化的名字来装饰一个不变的本性。远离此的是这些话的坦白诚实。父是父，子是子。但这些名字及其所代表的现实，对天主性并无创新，毫无矛盾，毫无异质。因为天主性，忠于自身，持续为自身；从天主而来的即是天主。神性受生既不带来削减，也不带来差异于天主性，因为子生于并与一个在天主性内且相似于它的本性共同存在，而圣父没有寻求任何异质元素掺杂于祂独生子的本性，而是以祂自己的一切赋予祂，且不损害给予者。因此，子并不缺乏天主性，因为祂是天主，从天主而来，而非来自其他；祂与天主并无差异，而是实在不外乎天主，因为从天主受生的即是子，且唯有子，而天主性在作为子受生时并未丧失其天主性。如此，圣父在圣子内，圣子在圣父内，天主在天主内。这并非由于两种和谐但不同种类的存在物的结合，也非由于较广大的实体包容较狭小的力量；因为物质的属性使得包含他物者不可能也被包容。这是由活生生的本性生自活生生的本性。本体保持不变，受生并不引起天主性的恶化；从天主生的天主不会是别的，只是天主。这里没有创新，没有疏离，没有分裂。相信圣父与圣子是两位天主是罪，断言圣父与圣子是一位孤存天主是亵圣，否认基于种类相同的从天主而来的天主的合一则是亵渎。\par

40. 免得那些信仰合乎福音的人将这奥迹视为模糊不清，主已以此顺序阐释了它——\Quote{\Emph{你们不信我在父内，父在我内吗？我对你们所说的话，不是凭自己说的，而是住在我内的父，祂做祂的作为。}}\BibleRef{若 14:10}除了这些话，过去或现在还有何言词可以宣告圣父与圣子拥有天主性，同时又不掩蔽子的受生？当祂说：\Quote{\Emph{我对你们所说的话，不是凭自己说的}}，祂既没有抹杀自己的位格，也没有否认子的身份，也没有隐瞒父的天主性在祂内的临在。当祂论及自己——这由代名词\Emph{我}证明——祂是作为常存于天主性本体而说；当祂说不是凭自己说的，祂为那在祂内从天主父而来的天主受生作证。祂在本性的合一中与父不可分离、不可区分；因为祂说话，却非凭自己说话。那说话的，虽非凭自己说话，必然而存在，因祂说话；而因祂非凭自己说话，祂表明祂的话不是祂自己的。因为祂接着说了：\Quote{\Emph{但住在我内的父，祂做祂的作为}}。父住在子内，证明父并非孤立独存；父通过子做工，证明子并非异类或外来的。不能只有一位格，因为祂非凭自己说话；反之，当一位通过另一位的口说话时，他们不能分离或分裂。这些话是祂们合一奥迹的启示。再者，祂们二者彼此并非有别，因为因着祂们内在的本性，彼此在对方内；而祂们是一，因为那位说话的，非凭自己说话，而那位非凭自己说话的，却确实说话。然后，在教导了父在祂内说话并作工之后，子确立了这完全的合一作为我们信德的准则：\Quote{\Emph{但住在我内的父，祂做祂的作为。你们当信我，我在父内，父也在我内；否则，因那些作为而信。}}\BibleRef{若 14:10-11}父在子内作工；但子也作祂父的作为。\par

41. 为此，为避免我们相信并说，圣父借着祂自己的全能能力在圣子内行事，而非借着圣子因生而得的神性本质作为其天赋权益，基督说：\BibleQuote{你们当信我，我在父内，父也在我内}。\BibleQuote{你们当信我}是什么意思？显然，这指向前文\BibleQuote{求你将父显示给我们}。他们的信德——那要求父被显示的信德——因这命令他们相信而得坚固。祂不满足于说\BibleQuote{看见我的，就是看见了父}。祂更进一步，扩展我们的认识，使我们在圣子内默观圣父，同时牢记圣子在圣父内。如此，祂要我们避免想象二者彼此流出的错误，而是教导他们通过赐予和接受的生育在同一个本性内合一。主愿意我们按祂的话语接受祂，以免因祂屈尊取了人性而使我们对信德的把握动摇。如果祂的肉身、祂的身体、祂的受难似乎使祂的天主性变得可疑，那么至少让我们因着那些工作的证据而相信：天主在天主内，天主出自天主，祂们是同一的。因为藉着本性的能力，每一位都在另一位内。圣父没有失去任何属于祂的，因为那在圣子内；圣子从圣父得到祂完整的子职。受造的身体本性并非如此被造，以致它们能相互包含对方，或拥有一个永存本性的完美合一。在它们的情况下，一个独生子不可能永恒地存在，与祂圣父的真天主性不可分离。然而，这是天主独生子的特有属性，这是在其真生育的奥秘中所启示的信德，这是圣神能力的工作：对基督而言，存在与在天主内是同一回事；并且这在天主内的存在不是一个事物在另一个之内的临在，如同一个身体在另一身体内，而是基督的生命与\OriginalTerm{自立}{subsistentia}是这样的：祂是在自立的圣父内，并在祂内，然而拥有祂自己的自立。因为每一位都如此自立，而不与另一位分离，因为他们通过赐予和接受的生育而成为两位，因此只有一个天主本性存在。这就是以下话语的意义：\BibleQuote{我与父原是一体}，以及\BibleQuote{看见我的，就是看见了父}，以及\BibleQuote{我在父内，父在我内}。它们告诉我们：所生的圣子与圣父既不相同也不低等；祂作为天主之子凭着生育的权利拥有天主本性，因此除了天主别无其他，这是在父与子独一天主性的奥秘启示中所传达的至高真理。因此，关于独生子受生的教义并不犯\OriginalTerm{二位神论}{ditheismus}之罪，因为天主子在受生进入天主性时，在祂自身内显示了祂生父的天主本性。\par

\SourceRecord{Translated by E.W. Watson and L. Pullan.；Nicene and Post-Nicene Fathers, Second Series；Vol. 9.；Edited by Philip Schaff and Henry Wace.；Buffalo, NY: Christian Literature Publishing Co.,；1899.；<http://www.newadvent.org/fathers/330207.htm>.}

% structure: p0001 source=h1 target=section reason=page-title

\section{论天主圣三（卷八）}

卷八。虔诚对主教是必要的，但面对像\Person{希拉利}时代那样猖獗的异端，他同样需要知识和辩证技巧；因为异端者在热忱上超过正统者，并且善于为不警惕的推理者设下陷阱（§§1-3）。他在§4中主张，到目前为止他已确立了自己的论点；现在在§5中，他转向亚略派对\BibleQuote{我与父原是一体}的解释，即认为他们在意志上是一体，而非在本质上是一体。这种错误的荒谬性通过比较基督徒在基督内的合一得到表明（§§7-9）；这种合一被公认是本质上的合一，却并不比圣父与圣子的合一更自然，后者是它的原型（§10）。的确，\BibleQuote{我与父原是一体}这句话并不适合仅仅表达意志的和谐（§11）。这种本质合一的恩赐，若赐予者自身不拥有它，就无法像现在这样通过降生成人和圣体圣事赐予基督徒；也就是说，除非圣父和圣子是一个天主（§§12-14）。事实上，我们通过基督的中介与圣父有完美的结合；这是本质上的统一，是永久的居住；这向我们保证了圣父内住于圣子、圣子内住于圣父，以及基督不是受造物、在意志上与圣父为一，而是圣子，在本质上与祂为一（§§15-18）。再者（§§19-21），圣神的使命是共同来自圣父和圣子；祂有时被称为圣父的神，有时被称为圣子的神，这进一步证明了圣父与圣子的本质合一。\Person{希拉利}现在查考圣经论及圣神的各种含义（§§22-25）。有时这个称号赋予圣父，有时赋予圣子，在这两种情况下都是为了使我们免于对天主的物质化观念。但严格来说，它也被用于护慰者，如在五旬节日那样。圣神居住在基督徒内；但这位神，无论称为天主之神、基督之神或真理之神，都由圣父所发并受圣子所遣，只是同一圣神。因此，天主性是唯一的，并且那在天主性内的诸位的本性也是唯一的（§§26,27）。他接着指出（§28），亚略派在敬拜基督的同时又称祂为受造物，这是自相矛盾的；因为他们因此落在法律的诅咒之下，并丧失了圣神。再者（§§29-34），天主所赋予的权能和恩宠被不加区分地描述为天主性内此一位格或彼一位格的恩赐。因此，圣子作为赐予者，必须与同样也是赐予者的圣父为一，并与圣神为一。\BibleQuote{只有一个天主，一个主}（§35）；如果我们否认圣子是神，我们也必须否认圣父是主；这是荒谬的。他们是一个天主，拥有一个圣神，但并非一个位格（§36）。圣保禄明确地说基督是\BibleQuote{万有之上的天主}；这一表述必须像宗徒的全部教导一样，具有公教的意义，并且与亚略主义不相容（§§37-39）。因此，亚略主义的支持者远离信仰（§40）。在重申真理之后（§41），\Person{希拉利}开始从圣子被圣父所印证的事实推衍圣子的神性（§§42-45）。这印证使祂成为圣父的对应者，从而成为祂的肖像，尽管是在仆人的形象中。如果祂在降生成人之后尚且是天主的肖像，那么在祂屈尊之前的程度就更高了（§46）。在§47中，他再次否认这一教导将圣父和圣子缩减为一个位格；然后（§§48-50）详细说明基督是天主肖像的意义。这意味着他们具有同一本质和同一能力，并且圣子是首生者，万有都是借祂受造。但创造以及和好都是圣父与圣子的共同工程（§51）。基督不可能比祂已经所做的更明确地陈述祂与圣父的合一；对这一真理的认可是真教会的试金石（§52）。异端对赋予生命的基督与从祂获得生命的受造宇宙之间的根本区别视而不见（§53）。在祂内住着\BibleQuote{整个神性的圆满有形有体地}。内住者和被内住者都是位格，然而是一个天主；并且整个天主性居于每一位中（§§54-56）。\par

1. 蒙福的宗徒保禄在制定选立主教的格式，并以他的训诲创建了一种全新类型的教会成员时，用以下的话教导我们在他身上成全的所有德行之总和：——\BibleQuote{持守那合乎信德道理的话，使他能劝勉纯正的道理，并驳斥那些反对的人。因为有许多不受约束的人，空谈家和欺骗者。}\BibleRef{铎 1:9}因为这样他指出，秩序和道德的要素只有在同时不缺乏认识和保存信仰所需的能力时，才对司铎职务的良好服务有益，因为一个人并非仅凭无辜的生活或仅凭讲道的知识就能成为良好且有用的司铎。因为一个无辜的侍奉者倘若未经教导，则只对自己有益；而受过教导的人，若无辜，则其教导无所支撑。因为宗徒的话并非仅以诚实和正直的诫命使人适应此世生活，另一方面也并非仅仅为解释法律而将一位会堂的经师培养成讲道的专家：而是宗徒在培育一位教会的领袖，以最大德行的完美成全而完善，使其生活由教导所装饰，其教导由生活所装饰。因此，他给弟铎，即他的话所指向的人，提供了关于完美实践宗教的如下训诫：——\BibleQuote{你在一切事上，要显出善行的模范；在教导上要持重、纯全、无可指摘的话，使那反对的人因无从说我们有什么可耻或邪恶的话而感到羞愧。}\BibleRef{铎 2:7}这位外邦人的导师和教会特选的博士，因意识到那在他内发言并居住的基督，深知毒害言语的感染会蔓延，而有害道理的腐败会狂暴地攻击信德之言的纯全形式，将其邪恶信条的毒素注入灵魂深处，并以根深蒂固的祸患潜行蔓延。正是关于这些人，他说：\BibleQuote{他们的话如同毒疮一样蔓延}\BibleRef{弟后 2:17}，以隐秘而隐蔽的传染侵袭心灵的健康，使健康被其侵蚀。为此，他希望在主教那里有纯全的话的教导，在信德中有良好的良心，并有劝勉的专长以抵抗邪恶虚假的狂言。因为有许多人假冒信德，却不服从信德，反倒为自己建立一种信仰，而非接受那赐予的信仰，他们因人性虚妄的思念而自大，只知自己所愿知的，不愿知那真实的；因为有时知道我们所不喜欢的事物是真正智慧的特征。然而这种意志之智慧伴随着愚昧的讲道，因为愚昧地学得的，必然愚昧地宣讲。但愚昧的宣讲对听众是何等大的祸害，当他们因智慧的幻觉而被导向愚昧的意见！因此，宗徒如此描述他们：\BibleQuote{有许多不受约束的人，空谈家和欺骗者。}\BibleRef{铎 1:9}因此，我们必须出声反对傲慢的邪恶、自夸的傲慢和诱惑的自夸——是的，我们必须通过我们道理的纯正、信德的真理、宣讲的真诚来反对这些事，使我们拥有真理的纯洁和纯正道理的真理。\par

2. 我刚刚提及宗徒此话的原因是这样：心术不正、伪善、毫无希望且言语恶毒的人，将抨击他们的必要性强加于我，因为他们在宗教的外衣下，将致命的道理、传染性的思想和腐败的欲望灌输给听众的单纯心灵。他们这样做时完全不顾宗徒教导的真义，以致圣父不再是圣父，圣子不再是圣子，信德不再是信德。在抵抗他们狂野的谬误时，我们已将答复的进路延伸到如此之远，先是从法律中证明天主与天主不同，且真实的天主在真实的天主内，然后从福音作者和宗徒的教导中证明独生天主的完美和真实的诞生；最后，我们在论证的适当过程中指出，天主子是真实的天主，且与圣父的本性相同，以使教会的信仰既不应承认天主是独一的，也不应承认有两位天主。因为天主的诞生不容许天主是孤独的，而完美的诞生也不容许将不同的本性归于两位天主。如今，在驳斥他们虚妄的空谈时，我们有一个双重目标：第一，我们教导圣洁、完美和纯全的道理，并且我们的讲论不应迷失于任何旁路和弯曲之道，挣扎于迂回曲折的隧道，似乎是在寻找真理而不是宣示真理。第二个目标是，我们应使所有人确信他们虚妄而欺骗性意见的那些狡诈论点的愚昧和荒谬，这些论点适合于一种诱人真理的貌似合理的外观。因为对我们来说，仅仅指出什么是善是不够的，除非通过驳斥其对立面，这些善被理解为绝对的善。\par

3. 然而，正如善人智者的本性及努力在于全力预备自己，以确保某些珍贵望德的实在性或机遇，免得他们的预备在某些方面不及他们所期盼的——同样，那些充满异端疯狂之狂热的人，也将他们最大的焦虑放在用尽一切不敬的机巧来反对虔诚信德的真理，以便对抗虔诚之人，建立自己的不敬；他们以自己的无望来超越我们生命的望德，并且花更多心思在虚假教导上，超过我们花在真实教导上的心思。因为他们针对我们信德的虔诚断言，精心设计了他们不敬谬信的反对意见，例如首先问我们是否相信独一天主，其次问基督是否也是天主，最后问圣父是否大于圣子，为的是当他们听到我们承认天主是独一时，就能利用我们的回答来证明基督不能是天主。因为他们并不询问圣子是否是天主；他们询问关于基督的问题，其全部愿望是证明祂不是圣子，从而通过陷害那些单纯信德的人，借着对独一天主的信仰，使他们从相信基督为天主转向，理由是如果基督也必须被承认为天主，那么天主就不再是独一了。此外，他们以何等世俗智慧的诡诈争辩说：\Quote{如果天主是独一的，那么无论那另一位被证明是谁，他都不会是天主。}因为如果另有他神，祂就不再是独一了，因为本性不允许在有另一位之处只有独一，也不允许在只有独一之处另有另一位。随后，当他们用这阴险论证的狡猾诡计误导了那些愿意相信和听从的人之后，他们便应用这一命题（仿佛现在能以更容易的方式建立它）：基督是天主，与其说在本质上，不如说在名称上，因为祂里面的这个一般名称不会在任何人心中摧毁对独一天主那唯一的真实信仰；并且他们争辩说，由此圣父大于圣子，因为本性不同，既然只有一位天主，圣父因其本性的本质特征而更大；而另一位仅仅被称为圣子，实际上却是因圣父的旨意而存在的受造物，因为祂小于圣父；并且祂不是天主，因为天主是独一的，不接纳另一位天主，因为那较小的必然在本质上与那较大的不同。此外，他们在试图为天主治定定律时何等愚昧，他们主张没有诞生能从一个单独存在者发出，因为在整个宇宙中，诞生来自于两个的结合；而且，不变的天主不能从自身赐予一个被生者以诞生，因为不变者不能增加，一个孤独单独的存在者的本性也不能在其自身内含有生育的属性。\par

4. 相反，我们通过属灵教导，达到了福音作者和宗徒的信德，并借着宣认圣父和圣子，追随永福的望德，而且从律法中证明了天主与天主的奥迹，未曾逾越我们对独一天主信德的界限，也没有未能宣认基督是天主，我们采用了来自福音书的答复方法，即我们宣告\OriginalTerm{独生天主}{Only-begotten God}从天主圣父的真正诞生，因为藉此，祂既是\OriginalTerm{真天主}{very God}，又不远离独一真天主的本质，因此祂的天主性既不能被否认，祂自己也不能被描述为另一位天主，因为诞生使祂成为天主，而祂内里来自天主的独一天主的本质并没有将祂分离为另一位天主。虽然我们人的理性引领我们得出这个结论，即不同本质的名称不可能在同一本质中相遇而合一，除非各自的本质在种类上没有差别；然而，我们觉得最好从我们主明确的言论中证明这一点，主在多次表明我们信德与望德的天主是独一的之后，为肯定独一天主的奥迹，在宣告并证明自己的天主性时说：\BibleQuote{我与父原是一体；}\BibleQuote{你们若认识我，也就认识我的父；}\BibleQuote{看见我的，就是看见了父；}\BibleQuote{你们当信我，我在父内，父也在我内；若不然，至少因那些工作而信。}祂在\Quote{父}之名中表明了自己的诞生，并宣告在认识祂时，父就被认识。祂承认本性的合一，因为那些看见祂的人看见了父。祂证明自己与父不可分，因为祂住在父内，父也住在祂内。祂拥有自知之信，要求人因祂能力的工作而相信祂的话。因此，在这圆满诞生的至福信德中，一切错误——无论是两位天主的错误还是独一天主的错误——都被消除了，因为他们在本质上是一体，却不是同一人；而那位与存有者不是同一人的，却又与祂毫无差别，以致两者是一体天主。\par

5. 如今，异端者既不能否认这些事，因为它们是如此清晰地陈述和理解，但他们仍以最愚昧邪恶的谎言曲解它们，以便随后否认它们。因为基督的话：\BibleQuote{我与父原是一体}，他们试图解释为仅指意志的一致，使得其中只有意志的合一，而非本性的合一——就是说，他们之所以为一，不是由于存有的本质，而是由于意志的同一。他们引用\BibleBook{}{宗徒大事录}中的经文来支持他们的观点：\BibleQuote{众信友的心身都合而为一}\BibleRef{宗 4:32}，以证明多种灵魂和心志可以通过纯粹的意志一致而合为一个心一个灵魂。或者，他们也引用\BibleBook{格林多}{前书}中的话：\BibleQuote{所以栽种的和浇灌的，原是一体}\BibleRef{格前 3:8}，以表明，既然祂们在为我们得救的工作中和在一个奥秘的启示中是一体，祂们的合一便是意志的合一。或者，他们又引用吾主为将信祂的万民得救所作的祈祷：\BibleQuote{我不但为他们祈祷，而且也为那些因他们的话而信从我的人祈祷，愿众人都合而为一。父啊！愿他们在我们内合而为一，就如祢在我内，我在祢内}\BibleRef{若 17:20}，以表明，既然人不能，可以说，被融合回天主或自身合并成一个无分别的团块，这种合一必然源于意志的合一，同时所有人行天主所喜悦的事，并彼此在思想的和睦中联合，如此使他们成为一的不是本性，而是意志。\newline{}\par

6. 谁不认识天主，谁就显然不认识智慧。既然基督就是智慧，那么不认识或憎恶基督的人必定处于智慧之外。例如，那些主张荣耀之主、宇宙君王和独生子天主是天主的受造物而非祂的圣子的人，他们除了这些愚昧的谎言之外，还在辩护其谬误时展现出更加愚昧的狡诈。因为即使暂且搁置那存在于天主父与天主圣子之间的合一的本有特性，他们也能被自己所引用的经文本身驳倒。\newline{}\par

7. 至于那些心灵合一的人，我问：他们是以信德而合一的吗？是的，确实，藉着信德，因为藉着信德，众人的心灵才能合一。我再问：信德是一个还是有两个？无疑只有一个，这有宗徒本人的权威为证，他宣告同一信德、同一主、同一圣洗、同一天主\BibleRef{弗 4:4}。那么，如果众人是因信德，即因同一信德的本质而合一的，你怎么不理解那些因同一信德的本质而成为一的人身上的本性的合一呢？因为所有人都重生到天真无邪、不死、认识天主和信德望德之中。如果这些事物因其自身不能彼此有别，因为同一天主、同一望德、同一主和同一重生的圣洗；如果这些事物是藉一致而非藉本性而成为一，那么请你也将意志的合一归于那些重生进入其中的人。然而，如果他们已重生进入同一生命与永生的本性，那么，既然他们的心灵合一，意志的合一就不足以解释那些因重生进入同一本性而成为一的人的情况。\newline{}\par

8. 这些不是我们提出的臆测，也不是我们虚伪地拼凑这些以歪曲词义来欺骗听者；而是我们紧守健全教导的规范，认识并宣讲真实的事。因为宗徒写信给迦拉达人时，表明这信友的合一源于圣事的本质：\BibleQuote{凡受洗归于基督的，就是穿上了基督；不再分犹太人或希腊人，不再分为奴的或自主的，不再分男或女，因为你们众人在基督耶稣内已是一个了。}\BibleRef{迦 3:27}这些人在如此大的种族、地位、性别差异中成为一——这是出于意志的一致，还是出于圣事的合一？因为他们领受了同一洗礼，都穿上了同一基督。那么，在此，心志的一致又有什么用呢？他们因同一圣洗的本性穿上了同一基督而成为一。\newline{}\par

9. 或者，栽种的和浇灌的原是一体，他们岂不就是因为自身在同一圣洗中重生，而成为同一重生圣洗的职事吗？他们不是做同样的事吗？他们不是在本一者内而成为一吗？所以，那些因同一事物而成为一的人，不仅因意志，而且也因本性而成为一，因为他们自己已成为同一事物，且是同一事物和同一权能的职事。\newline{}\par

10. 然而，愚人的反驳总是证明他们的愚蠢，因为对于他们以不智或歪曲的理解所设计出来反对真理的过错，真理本身却屹立不摇，那么反对真理的事物必然被视为虚假和愚昧。因为异端者试图用\BibleQuote{我与父原是一体}\BibleRef{若 10:30}的话来欺骗他人，不让他们承认在父与子内神性的合一与相似本质，而只是一种源于彼此相爱和意志的一致——我说，这些异端者甚至从吾主的话中引出了一个合一的实例，正如我们以上所显示的：\BibleQuote{愿众人都合而为一，父啊！就如祢在我内，我在祢内，使他们也在我们内合而为一。}凡是不信这些福音应许的人，就是处于应许之外，并且因邪恶理解的罪责而失去了所有简单的望德。因为不知道你所信的，与其说是可原谅的，不如说是一种奖赏，信德最大的服务就是盼望你所未见的事。但是，最彻底邪恶的疯狂，要么是不相信那些被理解的事，要么是腐蚀了自己所信的涵义。\newline{}\par

11. 然而，尽管人的邪恶能够扭曲其理智能力，但言语仍然保持其意义。主向圣父祈祷，使那些信祂的人合而为一，并如同祂在圣父内，圣父在祂内一样，使众人都在他们内合而为一。你们为何在这里引入思想的一致，为何引入因意志一致而产生的灵魂与心灵合一？因为如果使众人合而为一的是意志，主不会缺少合适的言辞，祂本可以这样祈祷——圣父，如同我们在意志上合而为一，愿他们也因意志合而为一，使我们众人通过一致而合而为一。或者，难道那身为圣言者竟不了解言语的含义？那身为真理者竟不懂得讲真理？那身为智慧者竟在愚妄的言谈中迷失？那身为德能者竟被软弱所困，以致不能说出祂要让人理解的话？祂已清楚讲述了福音信德的真实而纯正的奥迹。祂不仅说了使我们理解，也教导了使我们相信，说：\Emph{\Quote{愿他们众人合而为一，就如祢、圣父，在我内，我在祢内，使他们也在我们内。}}因为首先是为那些人所祈祷，经上说：\Emph{\Quote{愿他们众人合而为一。}}然后，合一之增进藉着合一的模型被展示出来，当祂说：\Emph{\Quote{就如祢、圣父，在我内，我在祢内，使他们也在我们内}}，使如同圣父在圣子内，圣子在圣父内，藉着这合一的模型，众人都在圣父和圣子内合而为一。\par

12. 但是，因为只有圣父和圣子才本性上合而为一：天主出于天主，独生者出于非受生者，只能存在于祂根源的本性中；所以那受生者应当存在于祂出生的实体中，而出生所拥有的神性真理与那所从出的神性真理既不相同也不同；因为主藉着在随后的全部论述中肯定这完全合一的本质，已使我们的信仰毫无疑惑。因为紧接着的话是：\Emph{\Quote{愿世界相信祢派遣了我。}}\BibleRef{若 17:21}这样，世界要相信圣子是由圣父派遣的，因为所有信祂的人都将在圣父和圣子内合而为一。而他们如何如此，我们很快就得知——\Emph{\Quote{祢赐给我的光荣，我已赐给了他们。}}现在我问：光荣是否等同于意志？因为意志是心灵的感受，而光荣是本性的装饰或美化。所以，正是从圣父那里接受的光荣，圣子赐给了所有信祂的人，而肯定不是意志。若是意志被赐予，信德便没有赏报，因为意志的必然性加于我们，也会把信德强加给我们。然而祂已展示了接受的光荣的赐予所产生的效果：\Emph{\Quote{愿他们合而为一，如同我们合而为一一样。}}\BibleRef{若 17:22}因此，接受的光荣被赐予是为了这个目的：使众人合而为一。所以如今众人在光荣中合而为一，因为所赐的光荣正是所接受的；并且它被赐予没有任何其他原因，只为使众人合而为一。既然众人藉着赐给圣子并由圣子赐予信徒的光荣合而为一，我问：圣子怎能具有与圣父不同的光荣？因为圣子的光荣把一切信徒都带入圣父光荣的合一中。此时，人的望德之言或许有些过分，但并不违反信德；因为对此存有望德虽是冒失，但若不相信它便是罪，因为我们希望和信德的同一作者。我们将在适当的地方以更清楚、更详尽的方式处理这个问题。然而与此同时，从我们目前的论证中可以轻易看出，我们的这个望德既非虚妄也非冒失。所以，藉着接受并赐予的光荣，众人合而为一。我持守信德，并承认合一的原因，但我还不明白所赐的光荣如何使众人合而为一。\par

13. 如今主没有让祂忠信追随者的心灵存疑，反而解释了祂本性运作的方式，说：\Emph{\Quote{愿他们合而为一，如同我们合而为一：我在他们内，祢在我内，使他们得以完全合而为一。}}现在我问那些主张圣父与圣子之间意志合一的人：今天基督是在我们内是因着本性的真实，还是因着意志的同意？因为如果圣言确实成了血肉，而我们确实从主那里作为食粮领受那成了血肉的圣言，我们岂不必须相信祂自然地居留于我们内？祂生而为人，已取了我们血肉的本性，这本性如今与祂不可分离，并将祂自己血肉的本性与永恒神性的本性，在藉着祂的血肉通传于我们的圣事中结合起来。因为这样我们众人都合而为一：圣父在基督内，基督在我们内。所以，凡否认圣父自然地在基督内的人，必须先否认他自己自然地在基督内，或基督在他内，因为圣父在基督内、基督在我们内，使我们在他内合而为一。因此，如果基督确实取了我们身体的血肉，那由玛利亚所生的人确实是基督，而我们确实以奥迹方式领受祂身体的血肉——（并因此我们将合而为一，因为圣父在祂内，祂在我们内）——那么，既然通过圣事所接受的本性的特殊属性是完全合一的圣事，如何还能维持意志的合一呢？\par

14. 谈论天主之事所用的言语，不可带有纯粹人性和世俗的意义，也不可凭借任何强横妄断的宣讲，从天上言语的正直中，歪曲出异端乖戾的解释。让我们阅读所写下的，理解我们所读的，然后成就完备信德的要求。因为关于基督的性体在我们内之事实，除非由祂教导，我们的话便是愚昧不敬的。因为祂亲自说：\Quote{\Emph{我的肉真是可食的，我的血真是可喝的。那吃我的肉、喝我的血的人，就住在我内，我也住在他内}}。\BibleRef{若 6:55}关于这血肉的真实性，毫无怀疑的余地。因为如今，从主自己的宣告和我们自己的信德来看，这确实是肉，确实是血。这些被吃喝时，便使我们既在基督内，基督也在我们内。难道这不真实吗？然而那些否认耶稣基督是真天主的人，大可认为它是假的。因此，祂藉着肉在我们内，我们在祂内，而同时我们自己也与祂一同在天主内。\par

15. 至于我们如何藉着赐予我们的血肉圣事而在祂内，祂亲自作证说：\Quote{\Emph{世界不再看见我，你们却要看见我，因为我活着，你们也要活着；因为我在父内，你们在我内，我也在你们内}}。如果祂只想表示意志的合一，为什么祂在成全合一的过程中设定了一种等级和顺序呢？除非因为祂藉着神性之本性在父内，而我们则相反藉着祂在身体内的诞生而在祂内，祂愿意我们相信祂藉着圣事之奥秘而在我们内。如此，通过中保得以教导一种完美的合一：我们住在祂内，祂住在父内，而祂住在父内时也住在我们内；这样我们就能达到与父的合一，因为在那位由诞生而自然地住在父内的祂内，我们也自然地居住，同时祂也自然地住在我们内。\par

16. 再者，这合一在我们内是多么自然，祂亲自如此作证：\Quote{\Emph{那吃我的肉、喝我的血的人，就住在我内，我也住在他内}}。\BibleRef{若 6:56}因为没有人能住在祂内，除非祂自己住在那人内，因为祂所取的血肉唯独是那些领受了祂血肉之人的血肉。如今，祂在这完美合一之圣事以前已教导说：\Quote{\Emph{就如生活的父派遣了我，我因父而生活，同样，那吃我肉的人，也要因我而生活}}。因此，祂因父而生活，正如祂因父而生活一样，我们也要因祂的血肉而生活。因为一切类比都是为塑造我们的理解，使我们能借助摆在面前的类比来把握我们所探讨的主题。这就是我们生命的原因：我们有基督藉着血肉居住在我们血肉之身内，而我们将因祂而生活，正如祂因父而生活一样。那么，如果我们按血肉自然地因祂而生活——即分有了祂血肉的本性——难道祂不也按神性自然地拥有父在祂内，因为祂自己因父而生活吗？祂因父而生活，因为祂的诞生并未在祂内植入一种异类或不同的本性，因为祂的存有本身来自父，却未因本性的任何不似而隔离于父，因着诞生在本性的大能内，祂内拥有父。\par

17. 我之所以详述这些事实，是因为异端者错误地断言父与子之间的结合只是意志上的，并利用我们与天主结合的例子，仿佛我们仅仅通过服从和虔诚的意志而与子结合，又通过子与父结合，并未藉着身体和血的圣事获赐任何共融的自然真实。然而，由子藉着圣子按血肉住在我们内而赐予我们的荣耀，使我们与祂在身体上且不可分地结合，这荣耀命令我们宣讲这真实而自然合一的奥秘。\par

18. 因此，我们已对我们狂暴反对者的愚妄作了答复，仅仅是为了证明他们虚假的空洞，从而防止他们因虚妄愚蠢的言论之错误而误导不慎之人。但福音的信德并不需要我们的答复。主曾为我们祈求与天主的合一，但天主保守自己的合一，并住在其中。他们合一并非藉着天主的神秘安排，而是藉着本性的诞生，因为天主在从自己生出祂时并未失去任何东西。他们是一，因为从祂手里不能被夺去的东西，也不能从父手里被夺去\BibleRef{若 10:28}；因为当祂被认识时，父也被认识；当祂被看见时，父也被看见；因为祂所说的，父作为住在祂内者所说的；因为在祂的工作中，父工作；因为祂在父内，父在祂内。这并非来自受造，而是来自诞生；不是因意志而成，而是因能力；不是思想的协议在说话，而是本性；因为受造与诞生并非同一，如同愿意与能够不同；同意与住在也不同。\par

19. 这样，我们并不否认圣父与圣子之间的同心合意——因为异端者惯于散布这样的谎言，说我等因不单凭和谐作为合一的纽带，便宣称他们彼此不合。然而，请他们听听我等如何不否认这样的同心合意。圣父与圣子在本性、尊荣、权能上都是一体，同一本性不会意愿相反之事。再者，请他们听听圣子关於祂与圣父本性合一的证言，因为祂说：\BibleQuote{当那护慰者来到时，就是我从圣父那里要派遣给你们的，就是从圣父所出的真理之神，祂要为我作证。} 护慰者要来，圣子要从圣父那里派遣祂，而祂就是从圣父所出的真理之神。让整个异端追随者团体激发他们才智最敏锐的能力吧；让他们现在寻找他们能向无知者散布的谎言，并宣告圣子从圣父那里派遣的事物到底是什么。那派遣者在其所派遣的事物中显示其权能。至于祂从圣父那里所派遣的，我们应如何看它，是领受的、还是被派出的、或是受生的？因为祂说\Quote{从圣父那里派遣}这句话，必然暗示这些派遣方式中的一种。并且祂将从圣父那里派遣那从圣父所出的真理之神；因此祂不可能是领受者，因为显明祂是派遣者。剩下的只是要确定我们在这点上的确信，即我们是相信\OriginalTerm{共存的存在}{coexistens}的\OriginalTerm{外出}{egressus}，还是\OriginalTerm{受生的存在}{genitus}的\OriginalTerm{出自}{processio}。\par

20. 目前我暂不揭露他们思辨的放肆，有些人在主张护慰者圣神来自圣父或来自圣子。因为我们的主在这事上没有留下不确定之处，因为说过同样的话之后，祂这样说：\BibleQuote{我还有许多事要告诉你们，然而你们现在不能担负。当那真理之神来到时，祂要引导你们进入一切真理，因为祂不是凭自己说话，而是凡祂所听到的，祂就说出来，并且要把将来的事宣报给你们。祂要光荣我，因为祂要从我领受，并宣报给你们。凡圣父所有的都是我的；为此我说：祂要从我领受，并宣报给你们。}\BibleRef{若 16:12} 因此祂是从圣子领受，祂既被圣子派遣，又从圣父而出。现在我问：从圣子领受与从圣父出自是一回事吗？但如果有人认为从圣子领受与从圣父出自有区别，那么从圣子领受与从圣父领受必然被视为同一回事。因为我们的主亲自说：\BibleQuote{因为祂要从我领受，并宣报给你们。凡圣父所有的都是我的；为此我说：祂要从我领受，并宣报给你们。} 祂所要领受的——无论是能力、卓越，还是教导——圣子已说必须从祂领受，而祂又指明同样的事必须从圣父领受。因为当祂说凡圣父所有的都是祂的，并因此宣告那必须从祂自己的领受时，祂也教导：从圣父领受的实际上也是从祂自己领受，因为凡圣父所有的都是祂的。这样的合一不允许区别，也不造成从谁那里领受的区别，那由圣父赐予的被描述为圣子赐予。难道这里也只是提出意志的合一吗？凡圣父所有的都是圣子的，凡圣子所有的都是圣父的。因为祂亲自说：\BibleQuote{凡我所有的都是你的，你所有的都是我的。} 现在还不是表明祂为何如此说的地方，因为\BibleQuote{因为祂要从我领受：}指向后来某个时刻，即启示祂将要领受的时刻。无论如何，既然凡圣父所有的都是祂的，祂现在说祂要从自己领受。你如果能割裂本性的合一，并引入某种必然的不相似，使圣子不在本性合一中存在。因为真理之神从圣父而出，并由圣子从圣父那里派遣。凡圣父所有的都是圣子的；因此，那将要被派遣者无论领受什么，祂都要从圣子领受，因为凡圣父所有的都是圣子的。本性在各方面保持其法则，因为两者是一体，同一神性因生和诞生而被指明存在于两者之中；因为圣子肯定：真理之神将从圣父领受的那事物要由祂自己赐予。所以，异端者的悖逆不得被允许放纵不虔信仰的无法无天，拒绝承认主的这话——因为凡圣父所有的都是祂的，所以真理之神将从祂领受——应归于本性的合一。\par

让我们聆听那位被选的器皿和外邦人的导师，当时他已因罗马人认识真理而称赞了他们的信德。因为他愿意教导关于圣父和圣子本性的合一，就这样说：\newline{}\BibleQuote{但你们不属于血肉，而属于圣神，只要天主的圣神住在你们内。但谁若没有基督的圣神，就不属于祂。但如果基督在你们内，身体固然因罪而死，但圣神因正义而成为生命。但如果那使基督从死者中复活者的圣神住在你们内，那使基督从死者中复活的，也必因那住在你们内的圣神，使你们必死的身体复活。}\BibleRef{罗 8:9}\newline{}如果天主的圣神住在我们内，我们都是属神的。但这天主的圣神也就是基督的圣神，并且虽然基督的圣神在我们内，但那使基督从死者中复活者的圣神也在我们内，而那使基督从死者中复活的，也必因那住在我们内的圣神，使我们的必死身体复活。因此，我们乃是因那住在我们内的基督的圣神，借着那使基督从死者中复活者而复活。既然那使基督从死者中复活者的圣神住在我们内，而基督的圣神也在我们内，那么，在我们内的圣神不能不是天主的圣神。所以，异端者啊，请把基督的圣神与天主的圣神分开，把从死者中复活的基督的圣神与那使基督从死者中复活的天主的圣神分开；虽然住在我们内的基督的圣神就是天主的圣神，而那位从死者中复活的基督的圣神也是那使基督从死者中复活的天主的圣神。\par

现在我问你，你认为在\OriginalTerm{天主的圣神}{Spirit of God}中所指的是一种本性，还是一种属于本性的属性？因为本性并不等同于属于它的东西，就如一个人不等同于属于人的东西，火也不等同于属于火本身的东西，同样，天主也不等同于属于天主的东西。\par

因为我确实知道，天主子在\Emph{天主的圣神}这一称号下被启示出来，为使我们能理解父在祂内的临在，并且\Emph{天主的圣神}这个术语可用来指称两者中的任一，这不仅由先知们的权威，也由圣史们的权威所表明，当它说：\BibleQuote{主的圣神在我身上，因此祂给我傅了油。}\BibleRef{路 4:18} 又说：\BibleQuote{看，我的仆人，我所拣选的，我的爱者，我心灵所喜悦的；我要把我的圣神放在祂身上。}\BibleRef{玛 12:18} 而当主自己为祂作证时：\BibleQuote{但如果我藉天主的圣神驱逐魔鬼，那么天主的国就来到了你们中间。} 因为这几段经文似乎毫无疑问地指示父或子，同时却彰显了本性的卓越。\par

因为我认为，\OriginalTerm{天主的圣神}{Spirit of God}这个说法被用于两者中的每一位，以免我们认为圣子是纯属肉体地临在于圣父内，或圣父亲临在于圣子内，也就是说，以免我们认为天主固定在某个位置，而除此之外不存在于别处。因为一个人或任何像他的东西，当他在一个地方时，不能在另一个地方，因为在某个地方的东西被限定于它所处的地方：他的本性不允许他当存在于某个位置时无处不在。但天主是一个活的力量，拥有无限能力，无所不在，无处缺席，并通过祂自己的东西显示祂整个自己，表明祂自己的东西无非就是祂自己，以至于在祂们所在之处，祂可以被理解为祂自己。然而，我们不可认为，按肉体方式，当祂在一个地方时，祂就不再无所不在，因为通过祂自己的事物，祂仍然存在于所有地方，而祂的事物正是祂自己。如今，这些话是为了使我们理解\Quote{本性}的含义。\par

现在我认为，应当清楚理解的是，天主圣父由\Quote{天主的圣神}所指代，因为我们的主耶稣基督宣称主的圣神在祂身上，因为祂给祂傅油并派遣祂宣讲福音。因为在祂内，父本性的卓越得以彰显，显明圣子分享祂的本性，即使是当祂通过这属灵傅油的奥迹生于血肉时，因为在祂受洗时所确认的诞生之后，这关于祂内在本性之子的提示，随着从天而来的声音作证而被听到：\BibleQuote{你是我的儿子；我今天生了你。} 因为连祂自己也不能被理解为安息在自己身上，或从天上来到自己这里，或把儿子的称号加给自己：但这整个显示是为了我们的信德，为使我们能在完全而真实的诞生奥迹下，认识到本性的合一存在于那也开始成为人的圣子内。这样我们便发现，在\Quote{天主的圣神}中父被指定；但我们明白，圣子也同样被指出，当祂说：\BibleQuote{但如果我藉天主的圣神驱逐魔鬼，那么天主的国就来到了你们中间。} 也就是说，它清楚地表明，祂凭祂本性的能力驱逐魔鬼，而魔鬼除非藉天主的圣神是不能被驱逐的。\Quote{天主的圣神}这个短语也指护慰者圣神，这不仅由先知们的见证，也由宗徒们的见证，当它说：\BibleQuote{这正是借着先知所说过的：在最后的日子，天主说，我要将我的圣神倾注在一切血肉之上，你们的儿子和女儿要讲预言。}\BibleRef{宗 2:16} 而且我们知道，这全部预言都在宗徒们身上实现了，当圣神倾注之后，他们都用外邦人的各种语言说话。\par

26. 如今我们不得不将这事陈述出来，是为了达到这样一个目的：无论异端者的欺骗转向何方，它仍能被福音真理的界限和范围所约束。因为基督住在我们内，而哪里住着基督，哪里就住着天主。当基督之神住在我们内时，这居留并不意味着有别于天主之神的任何其他神灵住在我们内。然而，如果我们理解基督藉着圣神住在我们内，那么我们必须承认这圣神既是天主的，也是基督的。既然那本性作为单一实体的本性住在我们内，我们就必须认为圣子的本性与其父相同，因为那既是基督之神又是天主之神的圣神被证明是一个同一本性的存有。那么，我现在请问，他们怎能不是本性上同一的呢？真理之神从父发出，由子派遣并从子领受。但父所有的一切，都是子的；因此那从子领受的，是天主之神，同时也是基督之神。这神灵是子本性的存有，但这同一存有也是父本性的存有。他是那从死者中复活了基督者的神灵；但这无非是那被复活的基督的神灵。基督与天主的本性必须在某方面有所差异以致不能相同，如果可以证明那出于天主的灵同时也并非基督之灵的话。\par

27. 但是，异端者，当你狂乱咆哮，被你致命教义的神灵驱策时，宗徒抓住并约束了你，为我们建立了基督作为信德的根基，他也深知我们主的那句话：\BibleQuote{\Emph{若有人爱我，他也会遵守我的话；我父也必爱他，并且我们要到他那里去，与他同住。}}\BibleRef{若 14:23} 因为藉此他证实了，当基督之神住在我们内时，天主之神也住在我们内，而那从死者中复生者的神灵与那使他从死者中复生者的神灵并无差别。因为他们来并住在我们内：我请问，他们是作为联合在一起的异乡者前来并居留，还是在本性的合一中？然而，外邦人的导师争辩说，在信者内出现的并不是两个神灵——天主之神和基督之神——而是基督之神，也就是天主之神。这不是一个共同居留，而是一个独一的居留：然而这居留带有一种共同居留的神秘表象，因为并非有两个神灵内住，也不是内住的一个与另一个不同。因为在我们内有天主之神，也有基督之神；当基督之神在我们内时，天主之神也在我们内。所以，既然那出于天主的也出于基督，那出于基督的也出于天主，那么基督不可能与天主的任何不同。因此，基督就是天主，与天主同一神灵。\par

28. 现在宗徒断言福音中的那些话，\BibleQuote{\Emph{我与父原是一体}}\BibleRef{若 10:30}，意味着本性的合一，而非一个孤寂单一的存有，正如他写给格林多人的信中所说：\BibleQuote{\Emph{所以我告诉你们，没有人在天主之神内说耶稣是可咒骂的。}}\BibleRef{格前 12:3} 异端者，你现在觉察到，你在什么神灵中称基督为受造物吗？因为那些侍奉受造物超过侍奉造物主的是受诅咒的——你在肯定基督为受造物时，要知道你是什么，因为你很清楚侍奉受造物是可咒骂的。并请注意接下来的话：\BibleQuote{\Emph{除非在圣神内，也没有人能说耶稣是主。}}\BibleRef{格前 12:3} 你觉察到，当你否认基督所固有的东西时，你缺少了什么吗？如果你认为基督因他的天主性而是主，你就有圣神。但如果他仅仅是因一个义子的名号而是主，你就缺少圣神，而被一个错误的神灵所鼓动：因为除非在圣神内，没有人能说耶稣是主。但是当你说他是受造物而非天主时，尽管你称他为主，你仍然没有说他是主。因为对你来说，他像是同类中的一个，藉着一个普通的名号，而非本性，才是主。但请从保禄那里了解他的本性。\par

29. 因为宗徒接着说：\BibleQuote{\Emph{神恩虽有区别，却是同一圣神；职分虽有区别，却是同一主；功效虽有区别，却是同一天主，在一切人身上行一切事。但圣神显示在每人身上，是为了益处。}} 在上面这段经文中，我们看见一个四重的陈述：在神恩的区别中是同一圣神，在职分的区别中是同一主，在功效的区别中是同一天主，而在赐予那有益的事时，有圣神的显示。并且为使那有益的事的赐予在圣神的显示中被认出，他继续道：\BibleQuote{\Emph{有人藉着圣神领受了智慧的话；另有人因着同一圣神领受了知识的话；另有人因着同一圣神领受了信德；另有人因着同一圣神领受了治病的恩惠；另有人能行异能；另有人能作先知；另有人能辨别神灵；另有人能说各种语言；另有人能解释语言。}}\par

30. 确实，我们称为第四个陈述的，即圣神在赐予有益之事时的显现，其含义是清晰的。因为宗徒曾列举了这些有益的恩赐，圣神的显现正是通过这些恩赐发生。在这些不同的活动中，那一恩赐被明确无误地彰显出来，主曾对宗徒们论及这恩赐，祂教导他们\BibleQuote{不要离开耶路撒冷，但要等候父的恩许，就是你们从我所听见的；因为若翰固然以水施洗，但你们却要因圣神受洗，你们过不多几日也必要领受}\BibleRef{宗 1:4}。又说：\BibleQuote{但当圣神降临于你们身上时，你们要得着能力，并要在我耶路撒冷、犹太全地、撒玛黎雅，直到地极，作我的证人}\BibleRef{宗 1:8}。祂吩咐他们等候父的恩许，就是他们从祂口中听见的。我们可以确定，这里指的是父的同一恩许。因此，正是通过这些奇迹性的运作，圣神的显现发生了。因为圣神的恩赐是显明的：当智慧发出言语，生命之言被听见，且藉由天主赐予的洞察力获得知识时，免得我们像野兽一样因对天主的无知而未能认识我们生命的作者；或藉由对天主的信德，免得因不信天主的福音而置身于祂的福音之外；或藉由治病的恩赐，藉由疾病得痊愈来为赐予这些恩赐的天主之恩宠作证；或藉由奇迹的运作，使我们所行的被理解为天主的能力；或藉由先知之恩，使我们藉由对教义的理解得以知道是被天主所教导的；或藉由辨别神恩，使我们能分辨是否有人以圣神或邪灵说话；或藉由各种语言之恩，使说语言的恩赐作为圣神恩赐的标记被赐予；或藉由语言的解释之恩，使听者的信德不致因无知而受损，因为语言的解释者将语言解释给那些不懂的人。因此，在这一切分施给各人、各人以之获益的事上，有圣神的显现，圣神的恩赐借着这些加于各人的奇妙益处而显明。\par

31. 有福的宗徒保禄在揭示这些最难以理解的天上奥秘的秘密时，保留了清晰的表述和措辞谨慎的告诫，以表明这些不同的恩赐是通过圣神并在圣神内赐予的（因为通过圣神和\Quote{在圣神内}赐予是不同的事），因为一项在圣神内实行的恩赐的授予仍然是通过圣神赐予的。但他将恩赐的多样性总结如下：\BibleQuote{然而这一切都是同一且同一圣神所运作，按照祂的愿意分给各人}\BibleRef{格前 12:11}。现在，因此我问：运作这一切、按照祂的意愿分给各人的那位圣神，是通过祂还是\Quote{在祂内}进行恩赐的分施？但如果有人胆敢说所指的是同一位，宗徒将驳斥这一错误意见，因为他在上面说：\BibleQuote{运作的分别也有，却是同一的天主，在一切人身上运作一切}\BibleRef{格前 12:6}。所以，有一位分施者，另有一位是在其内分施者。然而要知道，运作这一切的总是天主，但如此运作，即基督运作，而子在祂的运作中履行父的工程。如果你在圣神内承认耶稣为主，要理解宗徒书信中那三重指示的力量：因为在恩赐的多样性中，是同一圣神；在职务的多样性中，是同一主；在运作的多样性中，是同一天主；而且，运作这一切、按照祂的愿意分给各人的，是同一圣神。要尽力理解，如果你能：在职务分施中的主，在运作分施中的天主，就是同一圣神，祂既运作又按照祂的愿意分施；因为在恩赐的分施中，有一位圣神，同一圣神运作并分施。\par

32. 如果这位同一神性的圣神，通过生之奥迹在神和主中为一的圣神，不令你满意，那么请向我指出，运作并分施这些不同恩赐给我们的是哪一位圣神，以及祂是在哪一位圣神内如此行。但你必须向我指明与我们的信德相符的事，因为宗徒向我们指明应当理解谁，他说：\BibleQuote{就如身体只是一个，却有许多肢体，但身体的许多肢体只是一个身体，基督也是这样}\BibleRef{格前 12:12}。他断言恩赐的多样性来自一位主耶稣基督，祂是一切的身体。因为在他使人知道在职务中的主，又使人知道在运作中的天主之后，他仍表明同一圣神运作并分施这一切，为建立同一身体而分施祂恩慈恩赐的各种样式。\par

33. 或许我们以为宗徒没有持守合一的原则，因为他说：\BibleQuote{服务有多种，却是同一的主；工作有多种，却是同一的天主}。\BibleRef{格前 12:5} 因此，因为他将服务归于主，将工作归于天主，他似乎不明白在服务和工作中是同一位存在。你当明白，这些服务的肢体也是工作的肢体，当他说：\BibleQuote{你们是基督的身体，各自都是肢体。天主在教会中设立了一些人，第一是宗徒}，在他们身上有智慧之言；\BibleQuote{第二是先知}，在他们身上有知识的恩赐；\BibleQuote{第三是教师}，在他们身上有信德的道理；\BibleQuote{其次是异能}，其中包括\BibleQuote{治病的恩赐、助人的恩赐、治理的恩赐}借由先知，以及说或译各种方言的恩赐。显然这些是教会的服务和工作之代理人，基督的身体由他们组成；天主安排了他们。但或许你会声称他们不是由基督安排的，因为是天主安排了他们。但你将听到宗徒自己说：\BibleQuote{恩宠是按照基督恩赐的尺度赐予我们每人。}又说：\BibleQuote{那下降的，正是那升到诸天之上，以充满万有的。他赐予一些人作宗徒，一些人作先知，一些人作传福音者，一些人作牧者和教师，为成全圣徒，为服务的工作。}难道服务的恩赐不是属于基督，同时也是天主的恩赐吗？\par

34. 但若不虔敬自认为，因为他（宗徒）说：\BibleQuote{同一的主和同一的天主}\BibleRef{格前 12:5}，他们不是在性体上合一，我将用你认为更强有力的论证来支持这种解释。因为同一位宗徒说：\BibleQuote{但我们只有一位天主，就是父，万物出于祂，我们也在祂内；并有一位主耶稣基督，万物借着祂，我们也借着祂。}又说：\BibleQuote{一个主，一个信德，一个圣洗，一个天主和众人的父，祂超越众人，贯乎众人，也在我们众人内。}\BibleRef{弗 4:5}借此\BibleQuote{一位天主}和\BibleQuote{一位主}的话，似乎只有天主，作为一位天主，才拥有天主性；因为\InnerQuote{一}的属性不容许与他人共享。实在，这些神恩何等稀少、何等难求！在赐予这样有用的恩赐时，圣神的彰显何等真实！恩宠分赐的次序被这样设立，是有理由的：首先应是智慧之言；因为确实如此：\BibleQuote{除非在圣神内，没有人能说耶稣是主}\BibleRef{格前 12:3}，因为若非借着这智慧之言，基督不能被理解为是主；接着应是理解之言，使我们能带着理解讲论我们所知道的，并能认识智慧之言；第三项恩赐应是信德，因为那些首要和更高的恩宠若没有我们相信祂是天主，就成了无用的恩赐。因此，按宗徒这伟大而崇高的言论的真实意义，没有异端者拥有智慧之言、知识之言或宗教的信德，因为蓄意的邪恶，由于无法理解，缺乏对圣言的知识和真诚的信德。因为没有人能说他所不知道的；也不能相信他无法言说的；因此，当宗徒宣讲一位天主时——他原是法律中的归化者，被召叫接受基督的福音——他达到了完美信德的宣认。为了避免一句看似毫无防备的简单陈述可能给异端者提供机会，借着一位天主的宣讲否认圣子的出生，宗徒在提出一位天主的同时，用这些话指明祂的独特属性：\BibleQuote{一位天主父，万物出于祂，我们也在祂内}，使那位是天主者也被承认为父。此后，因为仅仅相信一位天主父不足以致得救恩，他补充说：\BibleQuote{并有一位主耶稣基督，万物借着祂，我们也借着祂}，表明拯救信德的纯正全在于宣讲一位天主和一位主，好使我们相信一位天主父和一位主耶稣基督。因为他十分清楚我们的主曾说过：\BibleQuote{这是我父的旨意，凡看见子，并信祂的，必得永生。}\BibleRef{若 6:40}但在确定教会信仰的次序，并将我们的信仰建立在父和子之上时，他用\BibleQuote{一位天主和一位主}的话，说出了那不可分割、不可解散的合一与信德的奥迹。\par

35. 首先，不自量力的异端者啊，你在宗徒借以说话的圣神内无分，当知你的愚昧。如果你误用对一位天主的宣认，否认基督的天主性，理由是既然存在一位天主，祂就必须被视为独一的，而\InnerQuote{一}是那位\InnerQuote{独一者}的特性与特有属性——那么，你将如何解释耶稣基督是一位主这句话呢？因为，如你所主张，唯有父是天主，已使基督没有成为天主的可能；那么按你的逻辑，基督是一位主，也必须使天主没有成为主的可能，因为你坚持认为\InnerQuote{一}必须是那位\InnerQuote{独一者}的本质属性。因此，如果你否认那位唯一的主基督也是天主，你就必须否认那位唯一的天主父也是主。那么，如果天主不是主，天主的伟大还剩下什么？如果主不是天主，主的权能又是什么？因为正是这种伟大或权能使那为主者成为天主，使那为天主者成为主。\par

36. 宗徒维护了主所说的：\BibleQuote{我与圣父原是一体}　\BibleRef{若 10:30} 的真正意义，在肯定两者是一的同时，表明了二者的一并非如同单一存有之独一性，而是在圣神的合一中；因为独一天主圣父和独一主基督，既然每一者都是主也是天主，在我们的信经中尚不能承认两个天主或两个主。因此每一者都是一，虽是单一却并非独一。除非用宗徒的话，我们无法表达信德的奥秘。因为\Quote{只有一个天主，只有一个主}，这一事实证明了在主内同时有天主的主权，在主内同时有天主性。你无法维持一个位格的合一，从而使天主成为单一性的；也不能分割圣神，从而阻止二者成为一体。在独一天主和独一主中，你也不能分离权能，使那为主者不再是天主，使那为天主者不再为主。因为宗徒在宣告这些名号时已谨慎地不宣讲两个天主或两个主。为此他采用了一种教导方法：在独一主基督中同时阐明独一天主，在独一天主圣父中同时阐明独一主。为了避免我们陷入天主是孤立的亵渎，从而摧毁独生天主的出生，他承认了圣父和基督两者。\par

37. 除非极端绝望的疯狂胆敢走得太远，以致既然宗徒称基督为主，便无人应承认祂是主以外的任何东西，并且因祂具有主的属性就没有真正的天主性。但保禄深知基督是天主，因为他说：\BibleQuote{圣祖们也是他们的，并且按血统，基督也出于他们，祂是在万有之上的天主。}　\BibleRef{罗 9:5} 这里并非受造物被算为天主；相反，受造物的天主是万有之上的天主。\par

38. 那万有之上的天主也是与圣父不可分的圣神，这一点也要从我们现在讨论的宗徒这话中学习。因为他承认独一天主圣父，万物从祂而来，以及独一主耶稣基督，万物藉祂而成；请问，他说万物从天主而来，万物藉基督而成，意图指出何种区别？祂，万物从祂而来并藉祂而成，怎能被视为具有与自身分离的性体和圣神？因为万物藉圣子从无中产生，宗徒将它们归于天主圣父：\BibleQuote{万物从祂而来}，但也归于圣子：\BibleQuote{万物藉祂而成}。我在这里找不到区别，因为每一者行使同一权能。如果关于宇宙的存续，确切地说明受造物从天主而来就足够了，那么还需要说明从天主而来的事物是藉基督而来的吗？除非\Quote{藉基督}和\Quote{从天主}是同一回事。但正如主和天主这两个称谓已归属两者，以至于每个名号都属于两者，同样地这里\Quote{从祂而来}和\Quote{藉祂而来}也指两者，这是为了显示两者的合一，而非表明天主的单一性。宗徒的话没有为邪恶的错误提供空间，他的信德也并非过于高深以致无法谨慎陈述。因为他用那些特别恰当的词语保守了自己不被理解为两个天主或一个孤立的天主：因为他既拒绝位格的单一性，又不分割天主性的合一。因为这\BibleQuote{万物从祂而来，万物藉祂而成}，虽然没有确立一个独占威能的孤立神性，却必须表明一位效率无别的作者，因为\Quote{万物从祂而来}和\Quote{万物藉祂而成}必指向一位相同性体的作者从事同一工作。此外，他确证两者都适当地具有同一性体。因为在宣布天主的丰盛、智慧、知识的深度，声明祂不可测度的判断奥秘，承认我们无法追查其道路的无知之后，他仍运用了人信德的行为，对那天上不可探寻和不可测度的奥秘之深度表达了敬意：\BibleQuote{因为万有都出于祂、藉着祂、并在祂内；愿光荣归于祂，至于永世。阿们。}　\BibleRef{罗 11:36} 他用这来指明同一性体，因为同一性体的工作不能不是这样。\par

39. 因为他特别地将万物从祂而来归于天主，并将万物藉祂而成当作基督的特殊属性，而现在\Quote{万有都出于祂、藉着祂、并在祂内}是天主的光荣；既然天主圣神与基督圣神是同一的，或者说，在主的管理和天主的运作中，同一圣神运作并分施，那些属性属于同一者的就不能不是一；因为在同一主圣子和同一天主圣父内，同一圣神在同一个圣神中分施并成就一切。这位圣人是何等配得上高超和天上奥秘的知识，被收纳并拣选分享天主的隐秘事，对不可言说的事保持应有的沉默，真正的基督宗徒！他以清晰教导的宣告抑制了人任性妄为的想象，承认独一天主圣父和独一主耶稣基督，使人既不能宣讲两个天主，也不能宣讲一个孤立的天主；虽然那不是一个位格者不能倍增为两个天主，另一方面，那非两个天主的也不能被理解为单一的位格；同时，天主启示为圣父证明了基督的真出生。\par

现在，伸出你们颤抖咝咝的舌头吧，你们这些异端的毒蛇，无论是你\Person{萨贝里乌斯}，还是你\Person{福提努斯}，或是那些现在宣扬独生的天主是一个受造物的人。凡否认圣子的，必要听到一位天主圣父，因为既然父亲只有有了儿子才成为父亲，父这个名必然蕴含圣子的存在。再者，凡剥夺圣子同一本性之合一的，要承认一位主耶稣基督。因为除非借着圣神的合一，祂是同一主，否则就没有空间让天主圣父为主。再者，凡认为圣子是因时间及降生成人而成为子的，要学习：万有是借着祂而有的，我们也是借着祂而有的，在时间之前，祂那无始的无限性就在创造万有。同时，让他再读一遍：只有一个蒙召的望德，一个圣洗，一个信德；如果此后他还反对宗徒的宣讲，他因编造自己私意的新奇道理而被诅咒，既不是蒙召，也不是受洗，更不是相信的；因为在一体天主圣父和一位主耶稣基督内，存在着一个信德、一个望德和一个圣洗。任何异端的道理都不能夸耀自己在属于独一天主、独一主、独一望德、独一圣洗和独一信德的真理中有任何位置。\par

所以，这独一的信德就是：借着不可分本性的合一——不是混杂而是不可分离，不是混合而是同一，不是连结而是共存，不是不完全而是完美的——承认父在子内，子在父内。因为有出生而非分离，有子而非收养；祂是天主，而非受造物。祂也不是另一种类的天主，而是父与子为一：因为本性并未因出生而改变，以致与源头的特性相异。因此，当宗徒宣告对他而言只有一位天主圣父和一位主基督时，他持守了子居住在父内、父居住在子内的信德，因为在基督主内也有天主，在天主父内也有主，二者即是那一位天主，二者也是那一位主，因为理性表明：在天主内若没有成为主，必有某种缺陷；在主内若没有成为天主，也必有缺陷。因此，既然二者为一，二者均被每一名称所蕴涵，且二者均在合一之外不存在，宗徒在其教导中没有超出福音的宣讲，基督在保禄内说话时，也与祂曾在世上以有形身体居住时所讲的话无异。\par

因为主在福音中说过：\BibleQuote{不要为那必坏的食物劳碌，要为那存留到永生的食物劳碌，就是人子将要赐给你们的：因为父，即天主，已印证了祂。他们遂问祂说：我们该做什么，才算做天主的工作呢？耶稣回答说：天主要求于你们的工作，就是信祂所派遣的那位。}\BibleRef{若 6:27}在启示祂降生成人和天主性的奥迹时，我们的主也讲出了我们信德和望德的教导：我们应为食物劳碌，不是那必坏的，而是那存留到永远的；我们应记住这永生的食物是由人子赐给我们的；我们应知道人子是被天主圣父所印证的；我们应知道这是天主的工作，即信祂所派遣的那位。父所派遣的是谁？就是父所印证的那位。父所印证的是谁？实在是人子，就是赐予永生食物的那一位。再者，祂赐给谁？只赐给那些为不坏食物劳碌的人。因此，为这食物劳碌同时也是天主的工作，就是信祂所派遣的那位。但这话是由人子说出的。人子如何赐予永生的食物？凡不知道赐予永生食物的人子已被天主圣父印证的人，就不明白他自己救恩的奥迹。在此，我现在要问：我们应如何理解\Quote{人子已被天主圣父印证}？\par

首先，我们应当认识：天主说话不是为祂自己，而是为我们，并且祂如此调整祂言语的表达，使我们软弱的本性能够领会和理解。因为当犹太人斥责祂因自称是天主子而自视为与天主平等时，祂曾回答说：祂所做的一切都是父所做的，并且祂从父领受了全部审判；此外，祂应受尊敬如同父一样。在这一切中，祂既已先自称为子，便在尊荣、权能和本性上使自己与父平等。之后祂又说：如同父在自己内有生命，同样赐给了子在自己内有生命，其中祂指明：借着出生的奥迹，祂拥有同一本性的合一。因为当祂说祂拥有父所拥有的，意思就是祂拥有父自身。因为天主不像人一样是由复合而成，在祂内拥有者和被拥有者之间没有种类差异；凡祂所是都是生命，一个圆满、绝对、无限的本性，不是由不同元素组成，而是整个被同一生命所渗透。既然这生命是在被拥有的同样方式下被赐予，虽然被赐予的事实显明了领受者的出生，却并不导致种类差异，因为所赐的生命就是如同所拥有的那样。\par

44. 因此，在如此多样而精确地启示了圣父的本性在祂内临在之后，祂接着说：\Quote{因为祂是天主圣父所印证的}。\BibleRef{若 6:27} 印鉴的本性在于展现其所雕刻的全部形式，而由它印出的印迹也在各方面再现它；既然它接收了所印的一切，它也完全在自身内显示那印在它身上的那位。然而，这种比拟不足以说明天主的诞生，因为在印鉴中，有材质、本性的差异以及印压的行为，藉此，较强本性的形像被印在较柔顺的事物上。但独生子天主——祂也因我们救恩的奥迹而成为人子——为了向我们指出圣父本性的相似在祂内，祂说自己是由天主所印证的；因为人子将要赐予永生的食物，好使我们由此在祂内认识赐予永恒食物的能力，因为祂在自身内拥有那印证祂的天主圣父之形式的全部圆满：如此，天主所印证的应在自身内只展示那印证它的天主的形式。这些事确实是主向犹太人所说的，他们因不信而无法接受祂的话。\par

45. 但在我们内，福音的宣讲者藉着那以祂发言的基督之神，灌输了对祂这本性的认识，祂说：\Quote{祂虽具有天主的形体，却没有以自己与天主同等为应当把持不舍的，反而空虚了自己，取了奴仆的形体}。因为天主所印证的，除了天主的形体之外，不能是别的；而那在天主形体中被印证的，也必须同时在其内反映出天主所拥有的一切。为此，宗徒教导说，天主所印证的那位，是寓居于天主形体中的天主。因为当他将要论及所取和所生的身体的奥迹时，他说：\Quote{祂没有以自己与天主同等为应当把持不舍的，反而空虚了自己，取了奴仆的形体}。就祂在天主的形体中而言，因天主在祂身上的印证，祂仍然是天主。但既然祂要取了奴仆的形体，并服从至死，不以与天主同等为把持不舍，祂便因服从而空虚了自己，取了奴仆的样式。祂空虚了自己那与天主同等的形体——并非祂视与天主同等为掠夺——虽然祂具有天主的形体，与天主同等，并作为天主被天主所印证。\par

46. 在此，我问：那位作为天主而寓居在天主形体中的天主，是否是一种不同种类的天主？正如我们在印鉴上所感知的那样，关于印压的形像和被印压的形像，例如钢印印在铅上或宝石印在蜡上，都会塑造刻在其中的图形或印上突起的图案。但如果有人如此愚蠢无知，竟认为天主所塑造为天主的、属于祂自己的那位，不是天主，并且那位具有天主体形的，在祂降生成人和谦卑的奥迹之后，藉着服从至死，甚至死在十字架上而得以成全，在各方面除了天主之外还是别的，那么他将听到天上、地上和地底下的一切，以及万口的承认：耶稣是在天主圣父的光荣中。如果，当祂的形体成了奴仆的形体时，祂仍留在如此的光荣中，那么我问，当祂在天主的形体中时，祂是怎样留下的？基督之神岂不曾在诸神的本性中——因为这就是\Quote{在天主的光荣中}的意思——当基督作为耶稣，即作为人出生时，存在在天主圣父的光荣中吗？\par

47. 在一切事上，蒙福的宗徒保持福音信德不变的教导。主耶稣基督被宣认为天主，如此：宗徒的信德既不因称祂为不同种类的天主而堕入承认两位天主的境地，也不因使天主圣子与圣父不可分离而为独一且孤单的天主的邪恶教义留下空隙。因为当他说\Quote{在天主的形体中}和\Quote{在圣父的光荣中}时，宗徒既未教导他们彼此不同，也不允许我们想祂不存在。因为那位在天主形体中的，既不会最终成为另一位天主，也不会丧失祂的天主性：祂不能与天主的形体分离，因为祂存在于其中；并且那位在天主形体中的，不是非天主。正如那位在天主光荣中的，不能是别的，只能是一位天主；并且，既然祂是在天主光荣中的天主，就不能被宣认为另一位与真天主不同的神，因为由于祂在天主光荣中这一事实，祂自然地从那在其光荣中的祂那获得了神性的特质。\par

48. 但同一信德不会因宣讲的多样性而停止成为信德。福音作者曾教导说，我们的主说：\Quote{谁看见了我，就是看见了父}。\BibleRef{若 14:9} 但是外邦人的导师保禄是否忘记或隐瞒了主的话的意义，当他说：\Quote{祂是不可见的天主的肖像}？\BibleRef{哥 1:15} 我问：祂是不可见的天主的可见肖像吗？无限的天主也能以有限形体的相似来呈现吗？因为肖像必须重复那所肖似之物的形式。然而，那些愿意在圣子内有不同本性的人，让他们决定他们希望圣子成为怎样的不可见天主的肖像。是可见的形体肖像，可被凝视，并以人的步态和动作从一处移到另一处吗？不，但让他们记得：根据福音和先知，基督是神，天主是神。如果他们将这位基督之神限制在形状和形体的范围内，那么这样有形体的基督将不是不可见天主的肖像，有限的界限也不能代表那无限的。\par

49. 但目前，主并没有让我们怀疑：\Quote{看见了我，就是看见了父；}宗徒也没有沉默祂的本性：\Quote{祂是不可见的天主的肖像。}因为主曾说：\BibleQuote{我若不行我父的作为，你们不要信我} \BibleRef{若 10:37}，教导他们藉着祂行父的作为，在祂内看见父；好使他们通过感知祂本性的德能，理解所感知之德能的本性。因此，宗徒宣告这就是天主的肖像，说：\BibleQuote{祂是不可见的天主的肖像，一切受造物的首生者；因为在祂内万有是在天上和地上被造的，有形与无形之物，无论宝座、宰制、率领、权能，都是藉着祂并在祂内受造的；祂万有之先，万有在祂内得以存立。祂是身体——教会的头，祂是元始，由死者中首生者，为使祂在万有之上居首位。因为父乐意使一切圆满居住在祂内，并藉着祂使万有与祂和好。} \BibleRef{哥 1:15} 因此，藉着这些作为的德能，祂是天主的肖像。因为制造不可见之物的造物者，并不被祂本性中的任何必然所迫而成为不可见之天主的可见肖像。并且，免得祂被视为形状的相似而非本性的相似，祂被称为不可见之天主的相似，好使我们通过祂行使神圣本性之德能（而非不可见的属性）而理解，那本性是在祂内。\par

50. 因此祂是一切受造物的首生者，因为在祂内万有被造。并且，免得有人敢将万有在祂内的创造归于任何其他者，他说：\BibleQuote{万有都是藉着祂并在祂内受造的，祂万有之先，万有在祂内得以存立。} 这确实描述了受造物的起源。但是关于祂取我们身体的那一安排，他补充说：\BibleQuote{祂是身体——教会的头：祂是元始，由死者中首生者：为使祂在万有之上居首位。因为父乐意使一切圆满居住在祂内，并藉着祂使万有与祂和好。} 宗徒将属灵奥迹与其物质效果联系起来。因为那不可见之天主的肖像，祂自己是祂身体——教会的头；而那一切受造物的首生者，同时也是元始、由死者中首生者：为使祂在万有之上居首位，为我们成为身体，同时祂也是天主的肖像，因为祂既是受造物的首生者，同时也是为永生的首生者；以致属灵的事物，既在首生者内受造，就得以存立；同样，所有属人的事物也因祂，在由死者中首生者内得以重生进入永生。因为祂自己就是元始，作为子，因此是肖像，并且因为是肖像，故属天主。此外，祂是一切受造物的首生者，在自己内拥有宇宙的起源：祂又是祂身体——教会的头，并由死者中首生者，使祂在万有之上居首位。并且因为万有为祂而存立，在祂内天主性圆满乐意居住，因为万有在祂内藉着祂与祂和好，万有也是藉着祂在祂内受造的。\par

51. 你现在明白成为天主的肖像是什么意思吗？它意味着万有在祂内藉着祂受造。既然万有在祂内受造，你当明白，那祂是祂肖像的那一位，也在祂内创造万有。并且，既然在祂内受造的万有也是藉着祂受造，你就承认，在那位是肖像者内，存在着那祂是祂肖像者的本性。因为祂藉着祂自己创造在祂内受造的事物，正如万有藉着祂自己在祂内和好一样。既然它们在祂内和好，你在祂内承认父合一的本性，在祂内使万有与自己（父）和好。既然万有藉着祂和好，你当看见祂在祂自己内使万有与父和好，这些万有是祂藉着自己和好的。因为同一宗徒说：\BibleQuote{但万有出于天主，祂藉着基督使我们与祂自己和好，并将和好的职务赐给了我们：就是说，天主在基督内使世界与自己和好。} \BibleRef{格后 5:18} 请将整个福音信德的奥迹与此比较。因为当耶稣被看见时，那被看见者；在祂的作为中工作，在祂的话语中发言，也在祂的和好中和好。因此之故，在祂内并藉着祂有和好，因为父藉着相似的本性居住在祂内，藉着和好并藉着祂并在祂内，使世界与自己重新和好。\par

52. 因此，圣父顾念人的软弱，并没有将信德置于空泛而不确定的言论中。因为尽管主的话语本身就足以迫使人接受，祂却仍藉着解释其含义的启示来光照我们的理智，好叫我们藉着那本身是合一之原因的事物，学习认识祂的话：\BibleQuote{我与父原是一体}\BibleRef{若 10:30}。因为，当祂说圣父在祂的话中说话，藉祂的作为而工作，藉祂的审判而审判，在祂的显现中被看见，藉祂的和好而和好，并住在祂内，而祂又住在圣父内时——我问，为了适应我们理性的能力，使我们相信祂们的合一，祂还能在教导中使用比这些更恰当的话语吗？这些话藉着诞生的真理与本性的合一，宣告了无论圣子做什么、说什么，圣父都在圣子内说了、做了。这并非论及一种与祂自己不同的本性，或藉受造而加给圣父的本性，或藉圣父的割分而生入神性中的本性；而是表明那一位因完美诞生而生为完美之圣父的独一性，祂对自己的本性如此确信而说：\Quote{我在圣父内，圣父在我内}，又说：\Quote{凡圣父所有的都是我的}。因为那一位，在其作为、言语和显现中，圣父作为、言语并被看见，在祂内毫无欠缺天主性。他们并非两位天主，在祂们的作为、言语和显现中装扮成合一的外表。祂也不是一位孤独的天主，在祂的作为、言语和圣父的显现中，祂自己作为天主而工作、说话并被看见。教会理解这事。犹太会堂不信，哲学不知：那由独一者而来之独一者，由整体而来之整体，圣父与圣子，既没有因祂的诞生剥夺了圣父的圆满，也没有因诞生而未能拥有这同一圆满在自己内。凡陷入这种不信的愚昧中的人，或是犹太人的门徒，或是外邦人的门徒。\par

53. 现在，为使你理解主所说的话，当祂说\BibleQuote{凡圣父所有的都是我的}\BibleRef{若 16:15}时，要学习宗徒的教导和信德，他曾说：\BibleQuote{你们要小心，免得有人以哲学，以虚伪的妄言，按照人的传授，按照世俗的原理，而不是按照基督，把你们勾引了去；因为在他内，住有整个圆满的天主性，以身体的方式。}\BibleRef{哥 2:8}那属于世界、沾有人的教导、并被哲学所左右的人，就是那不认识基督是真天主、不在祂内认出天主性之圆满的人。人的心智只认识它所理解的，世间的信德能力是有限的，因为它按照物质元素的法则判断，认为只有它所能看见或做到的才是可能的。因为世界的元素是从无中产生的，但基督的存续并非始于非存在，祂也从未开始存在，而是从起初就取了永恒的开始。世界的元素或是没有生命的，或是从无生命的状态进入了生命，但基督是生命，从永生的天主而生为永生的天主。世界的元素虽由天主建立，但它们不是天主；基督是天主出自天主，祂本身完全就是天主所是的一切。世界的元素，既然在世界内，就不可能超出其状态而不再在世界内，但基督，藉着奥迹让圣父住在自己内，祂自己却在圣父内。宇宙的元素，从其自身产生具有相似生命的受造物，确实通过其身体功能的行使，从自己的身体赋予它们生命的开端，但它们自身并不以活生生的形式存在于其子嗣中；而在基督内，整个圆满的天主性却以身体的方式临在。\par

54. 现在我问，那住于祂内的圆满究竟是谁的天主性？如果不是圣父的，那么你这误导人宣讲独一真神的人，又向我强推哪一位别的神，说祂的天主性圆满地住在基督内？但如果是圣父的，请问我这圆满如何以身体的方式住于祂内？你若认为圣父以身体的方式住在圣子内，那么圣父住在圣子内时，祂就不会单独存在。另一方面，——而这一点更为真实——如果以身体的方式住于祂内的天主性，在祂内显示了出自圣父的真正本性，因为圣父在祂内，不是藉着屈尊或意愿，而是藉着诞生而真实地、完全地、以身体的方式住于祂内，就祂所是而言；并且，在天主性内没有差异或不相似，除非那在基督内以身体方式住着的，是按照天主性的圆满；那么，为何还要追随人的学说？为何还要依附空洞妄言的教训？为何还要谈论\Quote{同意}或\Quote{意志的和谐}？或\Quote{受造物}？天主性的圆满以身体的方式住在基督内。\par

55. 宗徒在此坚守了他信德的准绳，教导说天主性的圆满以身体的方式住在基督内；这是为了使信德的教导不致堕落为不敬的位格合一之宣认，或使罪恶的狂怒爆发为相信两种不同本性的异端。因为以身体方式住在基督内的天主性，既非孤独的，也非可分离的；因为以身体形式的圆满，不接纳另一身体圆满的任何分拆，而住于内的天主性，也不能被视为天主性的居所。基督是那样构成的，使得天主性的圆满以身体的方式住在祂内，并且这圆满必须被认为在本质上与基督为一。抓住每一个可供你诡辩的机会，磨尖你亵渎才智的锋芒。至少指出那想象中是谁的天主性的圆满以身体的方式住在基督内。因为祂是基督，并且天主性的圆满以身体的方式住在祂内。\par

56. 你若想知道什么叫做有形地居住，就应明白什么是发言者发言、被见者被见、工作者工作，在天主内的天主，整个的整个，一个的一；这样你就了解天主性圆满地具有形体的意思。也请记住，宗徒并未缄默于这个问题：那圆满地有形居住在基督内的天主性是谁的。因为他说：\BibleQuote{\Emph{其实，自创世以来，他那看不见的事，借着所造之物，得以清楚看见，即他永恒的大能和神性。}}\BibleRef{罗 1:20}所以，正是他的天主性有形地居住在基督内，不是局部地，而是完整地，不是零碎地，而是圆满地；如此居住，使两者成为一体，且如此一体，以致作为天主的这一位与作为天主的另一位并无不同：两者同样神圣，如同完美的生育产生了完美的天主。这生育因其完美而存在，因为天主性的圆满有形地居住在出自天主的天主内。\par

\SourceRecord{Translated by E.W. Watson and L. Pullan.；Nicene and Post-Nicene Fathers, Second Series；Vol. 9.；Edited by Philip Schaff and Henry Wace.；Buffalo, NY: Christian Literature Publishing Co.,；1899.；<http://www.newadvent.org/fathers/330208.htm>.}

% structure: p0001 source=h1 target=section reason=page-title

\section{\WorkTitle{论天主圣三}{（卷九）}}

卷九。在总结（§\,1）已得出的成果之后，依拉略在§\,2回到某些亚略异端的论证经文，并警告他的读者：他们的生命取决于在基督内认出真天主和真人，因为这双重本性使祂成为中保（§\,3）。普遍的类比和我们对上升至天主内生命的能力的意识，使我们确信祂内的这两个本性，祂使这上升成为可能（§\,4）。但异端抓住了降生成人的基督所说的话，这些话适合于祂作为人的谦卑，却将它们归于祂先前的状态；因此他们使祂否认祂的真天主性。但祂在降生成人之前、在世生活期间以及回到荣耀之后的言论，必须仔细区分（§§\,5,6）。依拉略现在考察基督降生成人的目的和成就，并表明祂为人类的工作是一个神圣的工作，是祂为我们完成的，正因为祂自始至终既是天主也是人，祂内的两个本性不可分离（§§\,7‑14）。在从基督在世生活的总体考察得出这一结论后，他以此光照检查亚略异端从孤立词语中得出的论证。他们断言基督拒绝被称为\Quote{善}或\Quote{师傅}。祂并未拒绝这两个称号，却宣称两者只属于天主（§§\,15‑18）。而且，祂本不能再与圣父更紧密地联合，同时却保持祂位格的分别（§\,19）。圣父亲自为圣子作证；犹太人的罪和损失在于：他们看见基督所行的圣父的工程，却没有在祂内看见圣子（§§\,20,21）。基督的尊荣和荣耀与天主的不可分离（§§\,22,23）。经师正确地宣认了天主的唯一性，但仍在天国之外，因为他不信基督是天主（§§\,24‑27）。其次，从\BibleQuote{永生就是：认识祢，唯一的真天主，和祢所派遣的耶稣基督}这些话中得出的亚略异端论证，通过与相关段落的比较而被驳斥（§§\,28‑35）。因为，如果圣父是唯一的真天主，那么圣子也必须是唯一的真天主（§\,36）。那为圣父和圣子所共有的神性本性不受任何限制，而永恒的受生不能被任何受造物的类比所说明（§\,37）。基督取了人性，由于圣父的本性并未分享此人性，统一性在一定程度上受损。但人性已在基督内被提升至天主；这只能是因为祂在神性内与圣父的统一性是完美的。否则，基督所取的肉躯就不能进入天主的荣耀（§\,38）。圣父和圣子只有一个荣耀；圣子在降生成人中寻求的，并非圣言的荣耀，而是肉躯的荣耀（§§\,39,40）。圣父和圣子的荣耀是一个；在那统一性中，圣子既赐予荣耀，也接受荣耀（§§\,41,42），而这为两者所共有的荣耀，证明了神性本性也为两者所共有（§\,42）。再者，亚略异端引用\Quote{圣子不能凭自己做什么}这些话，依拉略通过语境考察表明这支持公教的事业（§§\,43‑46）。圣子行圣父的工程，不是出于强制作为等级较低者，而是因为他们是一体。祂的意志是自由的，却因本性的统一而与圣父的意志完全和谐（§§\,47‑50）。亚略异端还引用\Quote{圣父比我大}这经文。事实上，圣父更大，首先因为祂是非受生的，其次因为圣子屈尊就卑取了人的状态，却未丧失祂的天主性（§\,51）。但圣父在本性上并不大于圣子，圣子是圣父的肖像；或者，生者更大，而圣子作为受生者并不小于圣父，因为虽然受生，圣子却没有存在之始（§§\,52‑57）。其次，基于\BibleRef{谷 13:32}的对无知的指控，因此也是与全知的天主在本性上的差异的指控，被驳斥（§§\,58‑62），这既通过圣经明确的陈述，也通过考虑天主的属性。在旧约中，天主有时被说成是知道、有时被说成是不知道特定事实，这只是比喻意义上的用法（§§\,63,64）。基督也是如此；祂的无知不过是一种明智而仁慈的知识隐藏（§§\,65‑67）。然而，亚略异端虽然承认基督优于人，知晓人类的一切秘密，却断言祂不能穿透天主的奥秘（§\,68）。但基督明确宣称祂能够且确实如此，因为两者彼此内在于对方，并在对方内被反映（§\,69）。无知只能是隐藏。只有圣父知道，即祂除了圣子以外没有告诉任何人；圣子不知道，即祂不愿揭示祂的知识（§§\,70,71）。天主是无限的；因此知识也是无限的。圣父和圣子的本性是一，所以圣子不可能不知道圣父所知道的。正如意志上，他们在知识上也是一（§§\,72‑74）。宗徒们通过在复活后重复他们的提问，表明他们意识到祂的无知意味着保留。而基督这次没有提到无知，尽管祂保留了他们所问的知识（§\,75）。\par

1. 在上一卷书中，我们论及了圣父天主与圣子天主不可区分的本性，并证明了\BibleQuote{我与父原是一体}\BibleRef{若 10:30}这句话所证明的并非一位孤独的天主，而是藉着圣子的诞生而未被分裂的天主性的合一：因为天主只能由天主所生，而那由天主所生为天主的，必然拥有天主所是的一切。我们虽非详尽无遗，却足以阐明我们的意思，回顾了主与宗徒们教导圣父与圣子本性及能力不可分离的言论；继而我们来到宗徒教导中的一段，他说：\BibleQuote{你们要谨慎，恐怕有人用哲学和虚妄的欺骗，按照人的传统，按照世事的原理，而不是按照基督，把你们掳去；因为在祂内，圆满地住有整个天主性}\BibleRef{哥 2:8}。我们指出，这里\BibleQuote{在祂内，圆满地住有整个天主性}这句话，证明祂是由圣父的本性而来的真实而完美的天主，既未将祂与圣父分离，也未将祂与圣父视同一体。一方面，我们得知，既然无形体的天主有形体地居住在祂内，那么作为由天主所生的天主的圣子，在本质上与圣父合而为一；另一方面，如果天主居住在基督内，这便证明了祂所居住的那位个人的基督的诞生。因此，在我看来，我们已经足够有力地回应了那些人的不敬，他们主的话——如\BibleQuote{谁看见了我，就是看见了父}\BibleRef{若 14:9}，或\BibleQuote{父在我内，我在父内}，或\BibleQuote{我与父原是一体}，或\BibleQuote{凡父所有的，都是我的}——解释为意志的合一或一致。这些假教师不敢否认话本身，却披着虔诚的外衣，曲解话的意义。例如，诚然，当本性合一时，意志的一致也不能否认；但为了排除那由诞生而来的合一，他们只声称一种相互和谐的关系。然而，蒙福的宗徒在多次明确陈述真实真理之后，以\BibleQuote{在基督内，圆满地住有整个天主性}这句话打断了他们鲁莽亵渎的主张，因为藉着无形体的天主有形体地居住在基督内，便教导了他们本性严格的合一。因此，圣子并非独居，而是父住在祂内，这并非言语上的，而是真实的真理；而且不仅是居住，也是工作和说话；不仅是工作和说话，也在祂内彰显自己。藉着诞生奥迹，圣子的能力是父的能力，祂的权柄是父的权柄，祂的本性是父的本性。藉着诞生，圣子拥有父的本性：作为父的肖像，祂从父那里再现父内所有的一切，因为祂是父的实体也是肖像，因为完美的诞生产生完美的肖像，而\BibleQuote{整个天主性有形体地居住在祂内}表明祂本性的真实性。\par

2. 这一切确实是如此：那本性上是由天主所生的天主，必然拥有祂本源的本性，即天主所拥有的本性，而一个活生生的本性的不可区分合一，不能因活生生的本性的诞生而分裂。然而，异端者却藉着福音信仰中救恩的宣认，偷偷潜入以颠覆真理：因为他们强行将自己的解释加诸于具有其他含义和意图的话语上，剥夺了圣子本性的合一。因此，为了否认天主子，他们引用祂自己的话的权威：\BibleQuote{你为什么称我为善？除了天主一个外，没有谁是善的。}他们说，这些话宣告了天主的独一性：因此，任何其他分享天主之名者，不能拥有天主的本性，因为天主是唯一的。并且从祂的话\BibleQuote{永生就是：认识你，唯一的真天主}\BibleRef{若 17:3}，他们试图建立一种理论，即基督被称为天主仅仅是称号，而非作为真实的天主。此外，为将祂排除在真天主的特有本性之外，他们引用\BibleQuote{子不能由自己做什么，惟有看见父做什么}\BibleRef{若 5:19}。他们还使用经文\BibleQuote{父比我大}。最后，当他们重复\BibleQuote{至于那日子和那时辰，没有人知道，连天上的天使也不知道，连子也不知道，惟有父知道}的话，好像这是对祂神性主张的绝对放弃时，他们夸口说已经推翻了教会的信仰。他们说，诞生不能将因无知而贬低的本性提升到平等。父的全知与子的无知揭示了神性中的不同，因为天主必须无所不知，而无知者不能与全知者相比。他们既不理性地理解所有这些段落，也不按它们的不同场合加以区分，也不在福音奥迹的光照下领悟，也不按话语的严格意义去领会，因此以粗鲁无知的鲁莽攻击基督的神性，引用单句孤立的言论去捕捉粗心者的耳朵，而隐瞒解释它们的后续部分或促使它们的发生情境，尽管话语的意义必须在其前后文脉中寻求。\par

3. 我们稍后将亲自用福音和书信的话语来解释这些经文。但我们首先认为，提醒我们共同信仰的成员是正当的：对\OriginalTerm{永恒者}{the Eternal}的认识，是在那赐予永生的同一宣认中呈现的。那不认识基督耶稣既是真天主又是真人的人，不能、也不可能认识他自己的生命。否认基督耶稣是\OriginalTerm{神性的天主}{God the Spirit}，或否认祂是我们身体的血肉，两者同样危险：\BibleQuote{凡在人前承认我的，我在我天上的父前也必承认他；凡在人前否认我的，我在我天上的父前也必否认他。}\BibleRef{玛 10:32} 如此，成了血肉的圣言说了；如此，耶稣基督——荣耀的主，为了教会的救恩在自己位格内成为中保，并在那人与天主之间的中保奥秘中，祂自己是同一的位格，既是人又是天主——教导了。因为祂，为了那中保而联合了两个性体，是每个性体的圆满实在；祂在每个性体中都存留，不缺少任何一个；祂不因成为人而不再是天主，也不因永远是天主而不再是真人。这是人类真福的真实信仰：同时宣讲天主性和人性，宣认圣言和血肉，既不因祂是人而忘记祂是天主，也不因祂是圣言而忽略血肉。\par

4. 这与我们对自然的经验相悖：祂作为人而生却仍是天主；但这与我们期望的脉络相符，因为当较高的性体出生为较低的性体时，较低者也被出生为较高者是可信的。实际上，根据自然的法则和习惯，我们期望的作用甚至预见了神圣的奥秘。因为在每一个出生的东西中，本性具有增长的能力，但没有减少的能力。看树木、庄稼、牲畜。看看人本身，理性的拥有者。他总是通过成长而扩展，不会因减少而收缩；他也从未失去他已成长成的自我。他确实因年老而衰弱，或因死亡而被剪除；他因时间的流逝而经历变化，或到达生命结构所分配给他的终点，但他没有能力停止成为他之所是；我的意思是，他不能通过从旧的自我减少而造出一个新的自我，即，从老人变成婴儿。因此，永恒增长的必要性，它是由自然法则强加给我们的本性的，使我们有充分理由期望它被提升到更高的本性，因为它的增长是符合本性的，而它的减少是违反本性的。唯有天主能够成为不同于以前的事物，却又不停止祂一直是的事物；祂能够收缩在子宫、摇篮和婴儿期的界限内，却又不离开天主的大能。这是一个奥秘，不是为祂自己，而是为我们。我们本性的被取用对天主不是提升，而是祂愿意降卑自己却是我们的提升，因为祂没有放弃祂的天主性，而是将天主性赐予了人。\par

5. 因此，\OriginalTerm{独生天主}{Only-begotten God}，当祂由童贞女作为人出生，并在时期圆满时，将要亲自将人性提升到天主性时，始终保持着这种福音教导的形式。也就是说，祂教导要相信祂是天主子，并劝勉要宣讲祂是人子；人说出并做了所有属于天主的事；天主说出并做了所有属于人的事。然而，祂从未说话而不通过这些话语的双重方面既表明祂的人性又表明祂的天主性。尽管祂宣告独一的天主圣父，祂却藉着祂受生的真理宣告自己是在独一天主的本性内。然而，在祂作为子的职分和作为人的状况中，祂使自己服从于天主圣父，因为一切受生的都必须回溯到它的作者，并且所有血肉之躯都必须在天主面前承认自己的软弱。在这里，异端者得以有机会欺骗单纯和无知的人。这些在祂的人性角色中说出的话语，他们错误地归于祂神性本性的软弱；并且因为祂在所有话语中是同一个位格，他们声称祂总是谈论祂的全部自我。\par

6. 我们不否认所有保存下来的祂的言论都指向祂的本性。但是，如果耶稣基督既是人又是天主，既不是当祂成为人时才初次成为天主，也不是那时就不再是天主，也不是在祂成为人之后，祂在天主内是小于完美的人和完美的天主，那么祂话语的奥秘必定与祂本性的奥秘相同。当我们按照所指示的时间，将祂的天主性与人性分开时，那么也让我们将祂作为天主的言语与作为人的言语分开；当我们同时宣认祂是天主和人时，让我们同时区分祂作为天主的言语和祂作为人的言语；在祂的人性和天主性之后，我们再次认出祂整个的人性完全属于天主的时刻时，让我们将那时所启示的一切都归于那个时刻。祂在成为人之前是天主，这是一回事；祂是人又是天主，这是另一回事；在祂成为人和天主之后，祂是完美的人和完美的天主，这又是一回事。那么，不要将时间和本性混淆在\OriginalTerm{救恩计划}{dispensation}的奥秘中，因为根据祂不同本性的属性，祂必须相对于祂人性的奥秘以不同的方式谈论祂自己：在祂出生之前是一种方式，在祂仍要死的时候是另一种方式，而在祂作为永恒者时又是一种方式。\par

7. 因此，为了我们，耶稣基督保留了所有这些属性，并生于我们的身体成为人，按照我们本性的方式说话，毫不隐藏神性属于祂自己的本性。在祂的诞生、受难和死亡中，祂经历了我们本性的一切境遇，但祂都凭借祂自己的力量承受了它们。祂自己是祂诞生的原因，祂愿意承受祂所不能承受的苦难，祂死了，虽然祂永远活着。然而，天主这样做不仅是通过人，因为祂是由自己而生，祂出于自己的自由意志受苦，并自己死去。祂也作为人做了这一切，因为祂确实诞生、受苦并死亡。这些是上天密旨的奥迹，在世界受造之前就已预定。天主独生子要按自己的意志成为人，而人将永远住在天主内。天主将按自己的意志受苦，以免魔鬼的恶意，在人性的软弱中运作，在我们内确立罪的法律，因为天主已承担了我们的软弱。天主将按自己的意志死去，以免任何权势，在永生天主将自身约束于死亡的法律之后，能抬起头来反对祂，或施展祂在其中所创造的自然力量。因此，天主诞生是为将我们纳入祂内，受苦是为使我们成义，死亡是为报复我们；因为我们的人性永远住在祂内，我们人性的软弱与祂的力量结合，邪恶与不义的精神权势在我们血肉的胜利中被制服，因为天主透过血肉死了。\par

8. 宗徒知道这个奥迹，并亲自从主那里接受了信德的知识，他不会忘记，世界、人类或哲学都无法包含祂，因为他写道：\BibleQuote{\Emph{你们要谨慎，免得有人用哲学和虚妄的欺骗，照人的传统，照世俗的原理，而不是照耶稣基督，把你们掳去，因为在天主内，有形地住有整个神格的圆满，你们也是在祂内得以圆满，祂是一切率领者与掌权者的元首。}}\BibleRef{哥 2:8} 宣布在基督内有形地住有整个神格的圆满之后，紧接着就是我们被收纳的奥迹，即\BibleQuote{\Emph{你们在祂内得以圆满}}这句话。因为神格的圆满在祂内，我们也在祂内得以圆满。宗徒不仅说\BibleQuote{\Emph{你们得以圆满}}，而是说\BibleQuote{\Emph{你们在祂内得以圆满}}；因为所有藉着信德的希望重生而得到永生的人，无论现在或将来，都已在基督的身体内；将来他们不再是在祂内得以圆满，而是在自己内得以圆满，就是在宗徒所说的时期：\BibleQuote{\Emph{祂要改变我们卑微的身体，使它相似祂荣耀的身体。}}\BibleRef{斐 3:21} 因此，现在我们是在祂内得以圆满，也就是说，藉着祂血肉的被取纳，因为神格的圆满有形地住在祂内。我们在祂内的这个希望，在祂内并不是轻微的权威。我们在祂内的圆满构成了祂对一切权势的元首地位和主权，正如经上所记载：\BibleQuote{\Emph{致使天主耶稣的名字，一切在天上的、地上的和地下的，无不屈膝叩拜；一切唇舌无不明认耶稣基督是主，以光荣天主圣父。}} 耶稣将在天主圣父的光荣中被明认，祂生于人，但如今不再住在我们身体的软弱中，而是在天主的光荣中。一切唇舌都将明认此事。但是，虽然天上地上的一切都要向祂屈膝，但祂在此作为一切率领者与掌权者的元首，正是因为整个宇宙都要屈膝顺服于祂，我们在祂内得以圆满，祂藉着有形地住在祂内的神格的圆满，将在天主圣父的光荣中被明认。\par

9. 但在宣布基督本性的奥迹和我们被收纳，即神格的圆满住在基督内，以及我们因祂降生成人而在祂内得以圆满之后，宗徒继续人类救恩的分施，写道：\BibleQuote{\Emph{你们在祂内也受了割损，不是人手所做的割损，而是脱去血肉之身，属于基督的割损；你们在圣洗中与祂一同埋葬，也在其中藉着信德，依靠那使祂从死者中复活的天主的工作，与祂一同复活。}}\BibleRef{哥 2:11} 我们所受的割损不是肉身的割损，而是基督的割损，也就是说，我们重生成为新人；因为与祂一同埋葬在祂的圣洗中，我们必须死于旧人，因为圣洗的重生具有复活的力量。基督的割损不是指除去包皮，而是指完全与祂同死，并藉着那死亡从此完全为祂而活。因为我们是在祂内藉着信德，依靠那使祂从死者中复活的天主而复活；因此我们必须相信天主，基督是因祂的工作而从死者中复活，因为我们的信德在基督内并与基督一同复活。\par

10. 于是所取人性的全部奥迹得以完成：\Quote{你们从前因着过犯和肉身未受割礼而死亡，祂却使你们与基督一同生活，赦免了我们的一切过犯，涂抹了那以规条所写、反对我们、与我们为敌的债卷，并将它撤去，钉在十字架上；祂脱去自己的肉身，将众权势公开示众，并在祂身上凯旋了。}属世的人不能领受宗徒的信德，也没有任何言语能像宗徒的言语那样解释其含义。天主使基督从死者中复活；基督是天主性圆满有形有体地居住的那一位。但祂也使我们与祂一同复活，赦免了我们的罪，涂抹了那因先前所立的规条而反对我们的罪恶法律的债卷，将它撤去，钉在祂的十字架上，借着死亡的法律脱下自己的肉身，将权势公开示众，并在祂身上凯旋了。关于这些权势、祂如何在身上凯旋、如何公开示众、所涂抹的债卷以及所赐的生命，我们已经谈论过。然而，谁能理解或表达这一奥迹？天主的作为使基督从死者中复活；同一天主的作为使我们与基督一同复活，赦免我们的罪，涂抹债卷，钉在十字架上；祂脱下自己的肉身，公开示众权势，并在祂身上凯旋了。我们有天主的作为，使基督从死者中复活；我们也有基督亲自在祂身上成就天主在祂内所成就的事，因为正是基督死去，从自己身上脱下肉身。因此，你要紧握由天主从死者中复活的基督的人性，也要紧握在尚未死亡时就成就我们救恩的基督的天主性。天主在基督内工作，但正是基督脱下祂的肉身而死。是基督死了，也是基督在死前以天主的德能工作；然而，正是天主的作为使死去的基督复活，而除了基督自己——那位在死前工作并脱下肉身而死的那一位——以外，没有别人使基督从死者中复活。\par

11. 你已经明白宗徒信德的奥迹了吗？你以为你已经认识基督了吗？那么请告诉我：是谁从自己身上脱下肉身？那被脱下的肉身又是什么？我看出宗徒表达了两个概念：被脱下的肉身和脱下肉身的那一位；然后我听到，基督由天主的作为从死者中复活。如果基督是从死者中复活的那一位，而天主是使他复活的那一位，那么请问是谁脱下肉身？是谁使基督从死者中复活，又使我们与祂一同复活？如果死去的基督与被脱下的肉身不是同一的，那么请告诉我被脱下的肉身的名字，并向我解释那脱下肉身者的本性。我发现，那位从死者中复活的基督天主，正是从祂自己身上脱下肉身的那一位；而那肉身，正是从死者中复活的基督；然后我看到祂将率领者与掌权者公开示众，并在祂身上凯旋。你明白这\Quote{在祂身上凯旋}吗？你领悟到被脱下的肉身和脱下肉身的那一位并非彼此不同吗？祂在自己身上凯旋，即在祂从自己身上脱下的那肉身中凯旋。你看到如此宣告了祂的人性和祂的天主性：死亡归属于人，肉身的复活归属于天主，尽管死去者和使死者复活者并非两位，而是同一位位格。被脱下的肉身是死去的基督；使基督从死者中复活的那一位，正是从自己身上脱下肉身的同一基督。要在祂复活的德能中看见祂的神性，并在祂的死亡中认出祂人性的安排。虽然每一功能由其本性完成，但请记住，死去和复活者是同一位，基督耶稣。\par

12. 我记得宗徒常提及天主圣父使基督从死者中复活；但他并不自相矛盾，也不与福音的信德相悖，因为主自己说：\Quote{圣父爱我，因为我舍掉我的性命，为再取回来。没有人夺去我的性命，而是我甘心情愿舍掉它。我有权舍掉它，也有权再取回来。这是我由我父接受的命令。}\BibleRef{若 10:17} 再者，当有人要求祂显一个关于祂自己的征兆，以便他们相信祂时，祂论到祂身体的圣殿说：\Quote{你们拆毁这座圣殿，三天之内，我要把它重建起来。}祂借着再取灵魂的权能和重建圣殿的权力，宣告自己是天主，并宣告复活是祂自己的工程；然而祂将一切都归于祂父命令的权柄。这与宗徒的意思并不矛盾，当他宣告基督是\Quote{天主的德能和天主的智慧}\BibleRef{格前 1:24}时，他将祂一切工程的壮丽都归于圣父的光荣：因为基督所做的一切，正是天主的德能和智慧所做的一切；而天主的德能和智慧所做的一切，无疑就是天主自己做的，因为基督正是天主的德能和智慧。所以，基督是由天主的作为从死者中复活的；因为祂自己以与天主不可区分的天主性，成就了天主圣父的工程。我们对复活的信仰，建基于那位使基督从死者中复活的天主。\par

13. 蒙福的宗徒所强调的，正是基督位格双重层面的宣讲。他指出基督人性的软弱，以及祂的神性德能与本性。因此，他致书给格林多人说：\BibleQuote{祂固然由于软弱而被钉在十字架上，却由于天主的德能而生活。}\BibleRef{格后 13:4}，将祂的死归于人性的软弱，将祂的生归于神性的德能；又致书给罗马人说：\BibleQuote{祂死，是死于罪恶，仅仅一次；祂活，是活于天主。你们也要这样看自己是死于罪恶，在基督耶稣内活于天主的人。}\BibleRef{罗 6:10}，将祂的死归于罪恶，即归于我们的身体，将祂的生归于天主，因为天主的本性就是生活。因此，他说，我们应当死于自己的身体，以便在基督耶稣内活于天主。基督在取了我们罪恶的身体之后，如今完全活于天主，将祂与我们共有的本性结合于神圣不朽的共融之中。\par

14. 我不得不简略地停留于此，以免我们忘记我们讨论的主耶稣基督是具有两种本性的位格，因为那原处于\OriginalTerm{天主的形态}{form of God}中的，取了\OriginalTerm{仆人的形态}{form of a servant}，在那种形态中服从至死。死亡的服从与天主的形态毫不相干，正如天主的形态并非内在于仆人的形态。然而，藉着福音安排的奥迹，同一位主既处于仆人的形态，又处于天主的形态，尽管取仆人的形态与处于天主的形态并非同一回事；那处于天主的形态中的，若不\OriginalTerm{空虚自己}{emptying Himself}，就不能取仆人的形态，因为两种形态的结合是不协调的。然而，空虚自己并取仆人的形态的，并不是另一位不同的人。取任何事物都不能加于不存在者，因为只有存在者才能取。因此，形态的空虚并不意味本性的消灭：祂空虚了自己，但并未失去自我；祂取了新的形态，但仍然是祂所是的。同样，无论是空虚还是取，祂都是同一位；因此，存在着一个奥迹：祂空虚了自己，取了仆人的形态，但祂并未终结以致在空虚自己时不再存在，也未在取形态时成为非存在。空虚本身促成了仆人形态的取得，但并不阻止那处于天主的形态中的基督继续是基督，因为确实是基督取了仆人的形态。当祂空虚自己而成为人基督，同时继续是神性基督时，祂身体形态的改变和另一本性在祂身体中的接纳，并未终止祂永恒神性的本性，因为祂在改变形态时和接纳我们的本性时，都是同一位基督。\par

15. 我们现在已阐述了诸奥迹的安排，异端者正是藉此欺骗某些无知的人，将基督通过祂所取的人性所说所做的一切归之于神性中的软弱，而将只适用于仆人形态的归之于天主的形态。那么，让我们继续详细回答他们的说法。我们总能安全地区分这两类言论，因为唯一真实的信仰在于宣认耶稣基督是圣言和肉身，即天主和人。异端者认为必须否认我们的主耶稣基督凭其本性是神性的，因为祂曾说：\BibleQuote{你为什么称我为善？除了天主一个外，没有谁是善的。}现在，一个令人满意的答案必须与所问之事直接相关，因为只有那样才能对所提的问题给出回答。那么，首先，我要问这些曲解者：\Quote{你们认为主不乐意被称为善吗？}难道祂宁愿被称为恶吗？正如祂的话语似乎暗示的：\Quote{你为什么称我为善？}我想没有人会如此无理，将祂归为承认恶，祂曾说过：\BibleQuote{凡劳苦和负重担的，你们都到我这里来，我要使你们安息。你们背起我的轭，跟我学吧！因为我是良善心谦的：这样你们必要找得你们灵魂的安息。因为我的轭是柔和的，我的担子是轻松的。}祂说自己是良善心谦的：我们能相信祂因被称为善而生气吗？这两者是不一致的。那见证自己良善的，不会拒绝善的名。显然，祂并非因被称为善而生气：如果我们不能相信祂因被称为善而恼怒，那么我们必须问，人们对祂说了什么，以致祂发怒。\par

16. 那么，让我们看看问话者除了称祂为善之外，还如何称呼祂。他说：\BibleQuote{善师，我该行什么善？}在\Quote{善}的称号之外，又加上了\Quote{师}的称号。如果基督并非因被称为善而责备，那么必定是因为祂被称为\Quote{善师}。此外，祂责备的方式表明，祂所恼怒的是问话者的不信，而不是师的名称或善的名称。一个年轻人，他以遵守法律为满足，却不知道法律的终向\BibleRef{罗 10:4}就是基督；他以为靠行为成义，没有觉察到基督是来寻找\BibleQuote{以色列家迷失的羊}，并来到那些相信法律不能藉着信德的成义而拯救的人中；他问法律的主、独生子天主，好像祂只是一个教授普遍规条和法律著作的老师。但是，主厌恶这种不虔不信的宣告，即把他当作一个法律教师来称呼，便回答说：\Quote{你为什么称我为善？}并为了指出我们如何能认识祂并称祂为善，祂又加上：\Quote{除了天主一个外，没有谁是善的。}这并不是拒绝善的名字，只要是以天主的身份而称祂。\par

17. 然后，作为祂对\Quote{善师}这名号感到不悦的证明——因为这不相信的人以人的身份称呼祂——祂回答那虚荣的少年，以及他夸耀自己已遵守了法律的狂言：\BibleQuote{你缺少一样：去，变卖你所有的一切，施舍给穷人，你将有宝藏在天上；然后来，跟随我。}在许诺天上宝藏时，祂并不避讳\Quote{善}的称号；在引领通往真福的道路时，祂也不拒绝被视为\Quote{师}。但祂斥责那不信，这不信从善唯独属于天主的教导中得出对祂属世的看法。为表明祂既是善又是天主，祂行使善的职分，开启天上的宝藏，并自荐为向导。祂拒绝所有作为人所给予祂的敬意，但祂并不否认自己向天主所献的；因为在祂承认唯一的天主是善的那一刻，祂的言行都带有唯一天主的权能、良善和本性。\par

18. 祂并非避讳\Quote{善}的称号，或拒绝师的职分，而是不悦于那在祂身上只看见身体和血肉的不信，这可以从祂言语的不同得到证明：当宗徒们承认祂为他们的师父时，祂说：\BibleQuote{你们称呼我师父和主，你们说得对，因为我就是。}\BibleRef{若 13:13}；在另一场合，祂说：\BibleQuote{不要被称为师父，因为你们的师父只有一位，就是基督。}\BibleRef{玛 23:10}。从信众——祂是他们的师父——祂以赞赏的话接受这称号，但在这里，当祂不被承认为主和基督时，祂拒绝了\Quote{善师}的名称，并宣称只有唯一的天主是善的，但祂并未将自己与天主区分，因为祂称自己为主、基督和天上宝藏的向导。\par

19. 主始终维护教会信德的这一定义，即教导有唯一的天主圣父，但祂并不将自己与唯一天主的奥迹分离，因为祂借着祂由生而来的本性，宣称自己既非第二位天主，也非唯一的天主。既然唯一天主的本性在祂内，祂就不能是异于祂的天主；祂的出生要求，既然祂是圣子，就应当有完满的子职。因此，祂既不能与天主分离，也不能与天主相混。所以祂使用精心选择的言辞，凡祂为圣父所要求的，祂都以谦逊的语言表明也同样适用于自己。以这命令为例：\BibleQuote{你们要信赖天主，也要信赖我}\BibleRef{若 14:1}。祂在尊荣上与天主相同；请问，祂如何能与祂的本性分离？祂说\BibleQuote{也要信赖我}，正如祂说\BibleQuote{要信赖天主}。\Quote{在我内}这词不正是表明祂的本性吗？将两种本性分开，但你也必须将两种信仰分开。如果我们不信基督的情况下信赖天主就是生命，那么就将基督的名号和天主性特质剥去。但如果圆满的生命只赐给那些既信赖天主，又信赖基督的人，就让细心的读者深思这句话的含义：\Quote{你们要信赖天主，也要信赖我}，因为这些话将对祂的信赖与对天主的信赖结合起来，便将祂的本性与天主的本性结合了。祂首先命令我们当信赖天主，但又加上我们也要信赖祂的命令；这暗示祂就是天主，因为凡信赖天主的也必须信赖祂。然而祂排除了相反宗教的合一的暗示，因为\Quote{要信赖天主，也要信赖我}这一劝诫，禁止我们认为祂是独自孤独的。\par

20. 在许多，甚至可说在祂几乎所有的言论中，祂都提供了这一奥秘的解说，从未将祂自己与天主的合一性分离，当祂宣认天主圣父时，也从未将天主描绘为单独孤立的，当祂将祂自己置于与祂的合一之中时。但祂更为清晰地教导祂的合一与祂的诞生之奥秘，莫过于当祂说：\BibleQuote{但我有比若翰更大的证词，因为圣父交给我要我完成的工作，就是我正在做的这些工作，为我作证：圣父派遣了我。派遣我的圣父，祂亲自为我作了证。你们从未听过祂的声音，也从未见过祂的形状，并且你们没有祂的话存在你们内，因为祂所派遣的那位，你们不相信。} \BibleRef{若 5:36} 圣父如何能真被说成是为圣子作了证，而祂自己既未被看见，声音也未被听见？然而我记得有声音从天上说：\BibleQuote{这是我的爱子，我所喜悦的，你们要听祂。} 如何能说他们没有听见天主的声音，因为他们所听见的声音自己声称那是圣父的声音？但或许耶路撒冷的居民没有听到若翰在旷野的孤独中所听到的。那么，我们必须问：\Quote{圣父如何在耶路撒冷作证？} 这已不再是给若翰的证词，他听到了从天上来的声音，而是比若翰更大的证词。祂继续说明那证词是什么：\BibleQuote{圣父交给我要我完成的工作，就是我正在做的这些工作，为我作证：圣父派遣了我。} 我们必须承认这证词的权威，因为除了被圣父派遣的圣子，没有人能做这样的工作。因此，祂的工作就是祂的证词。但接下来呢？\BibleQuote{派遣我的圣父，祂亲自为我作了证。你们从未听过祂的声音，也从未见过祂的形状，并且你们没有祂的话存在你们内。} 他们因未曾知道圣父的证词，而圣父在他们中间从未被听见或被看见，且祂的话没有存在他们内，难道是无可指责的吗？不，因为他们不能辩解说祂的证词对他们隐藏了；正如基督所说，祂工作的证词就是圣父为祂所作的证词。祂的工作为祂作证，证明祂是由圣父所派遣的；但这些工作的证词就是圣父的证词；因此，既然圣子的工作就是圣父的证词，那么必然得出结论：在基督内运行的那同一性体，正是圣父借着祂为圣子作证的那性体。所以，行这些工作的基督，和借着它们作证的圣父，被揭示为借着诞生拥有一个不可分离的性体，因为基督的运作本身被指明就是天主为祂所作的证词。\par

21. 因此，他们不能因未能认出这证词而被宣告无罪；因为基督的工作就是圣父为祂所作的证词。他们也不能借口未曾听见作证者的声音，也未曾看见祂的形状，并且祂的话没有存在他们内，而申辩对证词的无知。因为紧接着这些话：\BibleQuote{你们从未听过祂的声音，也从未见过祂的形状，并且你们没有祂的话存在你们内，} 祂指出了为什么声音未被听见、形状未被看见、话未存在他们内，尽管圣父为祂作了证：\BibleQuote{因为祂所派遣的那位，你们不相信；} 也就是说，如果他们相信了祂，他们就会听见天主的声音，看见天主的形状，并且天主的话会在他们内，因为借着祂们性体的合一，圣父在圣子内被听见、显示并拥有。祂难道不也是圣父的表达吗，既然祂是从圣父那里被派遣来的？当祂说，圣父为祂作证，却未被听见、未被看见、未被领悟，因为他们不相信圣父所派遣的那时，祂难道用任何性体的差异将祂自己与圣父区别开来吗？因此，独生的天主并不将祂自己与天主分离，当祂宣认天主圣父时；而是借着\Quote{圣父}这个词宣布祂与天主的关系，并将祂自己包括在应归于天主的尊荣之中。\par

22. 因为，就在这同一篇谈话中，祂宣告祂的工作为祂作证，证明祂是由圣父所派遣的，并断言圣父为祂作证，证明祂是由圣父所派遣的，祂说：\BibleQuote{那独一之天主的荣耀，你们并不寻求}。然而，这并不是一个简单的陈述，而没有任何预备使人相信祂与圣父的合一。请听祂在前面所说的话：\BibleQuote{你们不肯到我这里来，为要得生命。我不接受从人来的荣耀。但我知道你们心里没有天主的爱。我因我圣父的名而来，你们并不接待我；若有人因自己的名而来，你们倒要接待他。你们既然彼此寻求荣耀，却不寻求那独一之天主的荣耀，怎么能相信呢？}\BibleRef{若 5:40} 祂蔑视人的荣耀，因为荣耀应从天主那里寻求。彼此接受荣耀是不信者的标记：因为人能给人什么荣耀呢？祂说祂知道他们心里没有天主的爱，并指出原因，就是他们不接待祂因圣父的名而来。\Quote{因祂圣父的名而来：} 这除了意味着\Quote{因天主的名而来}之外，还有什么意思呢？岂不是因为他们拒绝了因天主的名而来的那位，所以天主的爱不在他们里面吗？这不是暗示祂具有天主的性体，当祂说\BibleQuote{你们不肯到我这里来，为要得生命}时吗？请听祂在同一篇谈话中论及自己所说的话：\BibleQuote{我实实在在告诉你们，时候将到，现在就是了，死人要听见天主子的声音，而凡听见的，就必生活。}\BibleRef{若 5:25} 祂因圣父的名而来：这就是说，祂自己并不是圣父，却与圣父有同一的天主性；因为作为子和天主，祂因圣父的名而来是自然的。然后，另一个人若因同一的名而来，他们将要接待他：但那人是从之期望荣耀，并要回报荣耀的，尽管他将假装是因圣父的名而来的。无疑地，这预指了敌基督，妄用圣父的名而自夸。他们将荣耀他，也将从他得荣耀：但那独一之天主的荣耀，他们却不寻求。\par

23. 祂说，他们心里没有天主的爱，因为他们拒绝了因圣父的名而来的祂，却接受了另一个因同一的名而来的人，并且彼此接受荣耀，而忽视了那独一真天主的荣耀。难道可能认为祂将自己与独一真天主的荣耀分开，当祂给出他们不寻求独一真天主的荣耀的原因，是因为他们接受敌基督，而不肯接受祂自己吗？拒绝祂就是忽视独一真天主的荣耀；那么，祂的荣耀岂不就是独一真天主的荣耀，如果坚定地接受祂就等于寻求独一真天主的荣耀？这同一篇谈话就是我们的见证：因为在它的开始我们读到：\BibleQuote{为叫众人尊敬子，如同尊敬圣父一样。不尊敬子的，就是不尊敬派遣祂的圣父。}只有相同性体的事物才在荣耀上平等；荣耀的平等表明被荣耀者之间没有分离。但与生子的启示相结合的是对相等荣耀的要求。既然子应受尊敬如同圣父一样，而他们不寻求那独一真天主的荣耀，祂就没有被排除在独一真天主的荣耀之外，因为祂的荣耀与天主的荣耀是同一且相同的：正如\BibleQuote{不尊敬子的，也不尊敬圣父}一样，所以谁不寻求独一真天主的荣耀，也就不会寻求基督的荣耀。因此，基督的荣耀与天主的荣耀不可分离。当拉匝禄生病的消息传给祂时，祂的话说明了圣父与圣子在荣耀上的完全等同：\BibleQuote{这病不至于死，而是为天主的荣耀，为叫人子因此受到光荣。}拉匝禄为天主的荣耀而死，为叫天主子通过他受到光荣。难道还有疑问吗？天主子 的荣耀就是天主的荣耀，因为拉匝禄的死荣耀了天主，同时也荣耀了天主子。因此，基督通过祂的诞生被宣告与天主圣父在性体上是同一的，因为拉匝禄的病是为天主的荣耀，同时信德的奥迹没有被破坏，因为天主子要通过拉匝禄受到光荣。天主子应被视为天主，但祂仍然必须被承认是天主子：因为通过拉匝禄荣耀天主，天主子也被荣耀了。\par

24. 因着天主性体的奥迹，我们被禁止将永生圣子的诞生与祂永生的圣父分开。天主子并未遭受如此类的变化，以至圣父本性的真理不在祂内存留。因为即使在唯独承认唯一天主时，祂似乎藉着\Quote{独一}一词放弃了天主性体，然而，祂仍在不破坏唯一天主信仰的前提下，将自己置身于圣父本性的合一之内。因此，当经师问祂，法律中哪条诫命最大时，祂回答说：\BibleQuote{以色列，你要听！上主我们的天主是唯一的主；你当全心、全灵、全意、全力爱上主你的天主。这是第一条诫命。第二条与此相似：你当爱近人如你自己。再没有别的诫命比这两条更大的了}。他们认为，当祂宣告首要诫命为\BibleQuote{以色列，你要听！上主我们的天主是唯一的主}，并且甚至在第二条诫命中也没有将祂自己作为崇拜对象，因为法律命令我们爱近人，如同命令我们信唯一的天主，这样祂就使自己脱离了唯一天主的本性与崇拜。我们也不可忽略经师的回答：\BibleQuote{老师，你说得对：天主是唯一的，除祂以外再没有别的；并且全心、全灵、全力爱祂，并爱近人如自己，这远超过一切全燔祭和牺牲}。\BibleRef{谷 12:52}经师的回答似乎与主的话语一致，因为祂也宣认对唯一天主最深切、最内在的爱，并承认爱近人如同爱自己一样真实，将爱天主和爱近人置于一切牺牲的燔祭之上。但让我们看看接下来发生了什么。\par

25. \BibleQuote{耶稣见他回答得明智，便对他说：\Quote{你离天主的国不远了}}。这般适度的称赞有何含义？信唯一的天主，并全心、全灵、全力、全意爱祂，且爱近人如自己；如果这信仰使人对天主的国得以完善，为何那位经师不是已经在其中，而是\Emph{离天国不远}？祂以另一种语气将天国赐给那些给赤身者穿衣、给饥饿者食物、给口渴者水喝、探望病人和囚犯的人：\BibleQuote{你们蒙我父降福的，来承受自创世以来为你们预备的国度吧} \BibleRef{玛 25:34}；或奖赏神贫的人：\BibleQuote{神贫的人是有福的，因为天国是他们的}。他们的得着是完美的，他们的拥有是全备的，为他们预备的国度之继承是确定的。但这青年的宣认是否不及他们？他的责任理想将爱近人提升到爱己的层面；为达至善行完美，他还缺少什么？偶尔行善并乐意帮助，并非完美的爱；但完美的爱履行了全部的爱德责任，那时人不再亏负近人任何债，而是给予近人如同给予自己一样多。但这位经师被排斥在完美之外，因为他不知道已经完成的奥迹。他确实因自己的信仰宣认而受到主的称赞，听到回答说他不离天国不远，但他并未实际获得那有福的希望。他的道路虽无知，却是有利的：他将爱天主置于一切之上，并将对近人的爱德置于与爱己同等的地位。当他将爱天主甚至置于爱近人之上时，他打破了燔祭和牺牲的律法；那离福音的奥迹不远了。\par

26. 我们也从主自己的话中看出，为何祂说\Quote{你离天国不远}，而非\Quote{你将在天国之内}。随后接着说：\BibleQuote{从此以后，没有人敢再问祂什么。耶稣在殿里教训人时，说：\InnerQuote{经师们怎么说基督是达味之子呢？达味自己藉着圣神说：\InnerQuote{上主对我主说：你坐在我右边，等我把你的仇敌放在你的脚下。}达味自己称他为主，他怎麼又是达味之子呢？}} \BibleRef{咏 109:1} \BibleRef{谷 12:34} 当经师宣认唯一的天主，祂当受万有之上的爱时，他离天主的国不远了。但他自己所陈述的法律恰恰责备他，因为他错过了法律的奥迹，不知道主基督，天主子，凭着祂诞生的本性被包含在唯一天主的宣认中。按照法律，宣认唯一天主似乎没有给天主子在唯一主的奥迹中留有余地；因此祂问经师，当达味称基督为他的主时，经师怎能称基督为达味之子，因为按自然秩序，如此伟大圣祖之子竟也是他的主，这是不合常理的。祂要引导这位只从肉身和出自达味之女玛利亚的诞生来看待祂的经师记住，按祂的神性而言，祂是达味的主，而非他的儿子；如此，\BibleQuote{以色列，你要听！上主我们的天主是唯一的主}这些话并未将基督与唯一主的奥迹分开，因为如此伟大的圣祖和先知称祂为自己的主，如同那在晨星之前由主所生的圣子。祂并未超越法律，也并未忘记无人可被宣称为主，而是不违背法律的信仰，教导说祂是主，因为祂藉着从无形天主体质而来的自然诞生之奥迹而存在。祂是独一者，由独一者所生，唯一主的本性使祂按本性成为主。\par

27. 还有什么可怀疑的余地？主亲自宣告法律的首要诫命是宣认并爱慕唯一的主，祂借先知的话证明自己是主，而非凭自己的言语，却总是指明祂是主，因为祂是天主子。藉着诞生，祂安住于唯一天主的奥迹中，因为那传递天主本性的诞生，并非生出另一位本性不同的天主；又因为此生育是真实的，故圣父并未被剥夺主的身份，圣子生而并不次于主。圣父保持其权威，圣子获得其本性。天主圣父是唯一的主，但独生子天主，即主，并不与那一位分离，因为祂从唯一的主获得其为主的本性。这样，基督藉法律教导只有一位主；又藉先知们的见证证明祂自己也是主。\par

28. 但愿福音的信德常藉不敬神者鲁莽的争辩而得利，用攻击的武器自卫，并以那为毁灭而准备的武器取胜，证明同一圣神的话是同一信德的道理！因为基督正是所宣讲的那一位，即真天主，且安住于唯一真天主的荣耀中。正如祂似乎否认自己为主时，却仍从法律宣告自己是主；同样，在福音中，即使当祂似乎承认相反的说法时，祂仍然证明自己是真天主。为避免承认基督是真天主，异端者援引祂所说的：\BibleQuote{永生就是：认识你，唯一的真天主，和你所派遣的耶稣基督。}\BibleRef{若 17:3}当祂说\BibleQuote{你，唯一的真天主}时，他们以为祂因\Quote{独一}的限制而将自己排除于天主的真实性之外；因为唯一的真天主只能被理解为一位独一的天主。宗徒信仰确实不容许我们相信两位真天主，因为凡与唯一天主本性相异之物，都不能与那本性的真理并列；若在唯一真天主的本性之外，另有一位藉诞生未得同样本性的另一类真天主存在，则在那唯一天主的实在中就存有不止一位天主。\par

29. 但正是藉着这些言词，祂明明白白地宣告自己在唯一真天主的本性中是真天主。为理解这一点，让我们从祂先前说过的话来作答，尽管上下文一直连贯至这些言词。这样我们可以逐步建立信德，并使我们的自由信心最终安歇于我们论证的顶峰，即基督的真天主性。首先有祂言词的奥迹：\BibleQuote{看见我的，就是看见了父}；以及\BibleQuote{你不信我在父内，父在我内吗？我对你们所说的话，不是凭自己说的；而是住在我内的父，祂自己做祂的工。你们当信我在父内，父在我内；否则，因那些事本身而信吧。}\BibleRef{若 14:9}在这充满深邃奥迹的训诲结束时，门徒们回答说：\BibleQuote{现在我们晓得你知道一切，不需要有人问你；我们因此信你是从天主而来的。}他们藉着祂所行使的神圣能力，在祂身上感知到天主的本性；因为知道一切并洞察人心是属于圣子，而非单纯的天主使者。因此他们承认祂是从天主来的，因为神圣本性的能力在祂内。\par

30. 主称赞他们的理解，并回答说祂不是被派遣，而是从天主出来的；用\Quote{出来}一词表示祂从无形体天主诞生的伟大事实。祂早已用同样的言辞宣告了这诞生，当祂说：\BibleQuote{你们爱我，并相信我是从父出来的，并从父来到这世界。}祂是从父来到这世界的，因为祂是从天主出来的。为指明\Quote{出来}表示祂的诞生，祂补充说祂是从父来的；既然祂是从天主出来的——因为祂是从父来的——那么这\Quote{出来}随后紧跟着圣父的名字的宣认，就纯粹只是诞生。于是，向理解这\Quote{出来}奥迹的宗徒们，祂继续说：\BibleQuote{现在你们信了。看，时候要来，且已来到，你们要各自分散，撇下我独自一人；但我不是独自一人，因为父与我同在。}祂要表明这\Quote{出来}并非与天主圣父分离，而是诞生，藉着诞生，祂持续在圣父的本性中，因此祂补充说祂不是独自一人，而是圣父与祂同在；即在能力和本性合一中，因为圣父住在祂内，在祂的言语中发言，在祂的工程中工作。最后，为表明整个训诲的理由，祂补充说：\BibleQuote{我对你们讲论了这些事，是为叫你们在我内得到平安。在世界上你们要受苦难；但放心，我已经战胜了世界。}\BibleRef{若 16:33}祂向门徒讲论这些事，是为叫他们安住在祂内的平安中，不被关于信德辩论的热烈纷争所撕裂。祂被独自撇下，却并非独自一人，因为祂是从天主出来的，那从祂出来的天主仍然在祂内。因此，祂吩咐他们，当他们在世上受困苦时，要期待祂的应许，因为既然祂是从天主出来的，而天主仍在祂内，祂已战胜了世界。\par

31. 最后，为了用言语表达整个奥秘，祂举目向天，说：\BibleQuote{父啊，时候到了，求祢光荣祢的子，使子也光荣祢。正如祢赐给祂权柄掌管凡有血肉的人，叫祂将永生赐给一切祢所赐给祂的人。} 你称祂软弱，因为祂求受光荣？也可如此说，如果祂求受光荣不是为了借此使光荣祂的那一位得光荣。关于获得和给予光荣，我们在另一本书中已经谈过，再重复这个问题是多余的。但至少我们确定：祂祈求光荣，是为了使圣父因赐予光荣而受光荣。但或许祂是软弱的，因为祂 \Emph{接收} 权柄掌管凡有血肉的人？的确，接收权柄可能是一种软弱的标记，如果祂不能将永生赐给那些祂所接收的人。然而，接收这一事实本身被用来证明本性的低劣。如果基督不是像 \OriginalTerm{非生者}{Unbegotten} 那样由出生而为真天主，那或许如此。但如果接收权柄恰恰就是出生——藉此祂接受了祂所有的一切——那么这一恩赐并不贬低被生者，因为它使祂完全而彻底地成为天主。天主非生者使天主 \OriginalTerm{独生者}{Only-begotten} 由神圣福乐的完美出生而出：因此，是圣父作为出生之源的奥秘，但圣子因真实的出生而成为其源头的完美肖似，这并不是贬低。赐予掌管一切有血肉之人的权柄，并藉此使永生赐给一切有血肉的，这就预设了赐予者的父性和接收者的天主性：因为赐予表明一位是父，而在接收赐予永生的权柄中，另一位仍是天主圣子。所以，一切权柄对天主子而言都是本性的、与生俱来的；虽然它是被赐予的，但这并不使他与他的源头分离，因为所赐予的是他源头的属性，即赐予永生的权柄，将可朽的变为不朽的。父赐予了一切，子接受了一切；从祂的话可以清楚地看出：\BibleQuote{凡父所有的一切，都是我的。}\BibleRef{若 16:15} 祂在这里并非谈论受造物的类别和物质变化的过程，而是向我们揭示那荣福而完美的天主性的光荣，并教导我们说，天主在这里显明为他属性的总和：他的权能、永恒、眷顾、权威；不是要我们认为他拥有这些作为外在于他的东西，而是藉着这些属性，他自己以部分可被我们感官理解的方式表达出来。因此，独生者教导说，祂拥有父所有的一切，并且圣神将从祂领受：正如祂说：\BibleQuote{凡父所有的一切，都是我的；为此我说：祂要从我领受。} 凡父所有的一切，都是子所有的，被交付并接受了；但这些恩赐并不贬低祂的天主性，因为它们使祂拥有与父相同的属性。\par

32. 这些是祂增进对自己认识的步骤。祂教导自己是从父而来的，宣告父与祂同在，并证明祂已战胜了世界。祂将要被父光荣，也将光荣父；祂将运用所接受的权柄，赐予一切有血肉的永生。然后，倾听那总结这一整系列的最高点：\BibleQuote{永生就是：认识祢，唯一的真天主，和祢所派遣的耶稣基督。} 异端者，请学习宣认——如果你不能相信——这赐予永生的信仰。如果你能，就将基督与天主分开、子与父分开、万有之上者与真天主分开、那一位与唯一者分开——如果你说永生就是相信唯一的真天主而撇开耶稣基督。但如果宣认唯一的真天主（与基督分离）没有永生，那么请问，基督怎能从我们的信仰中与真天主分离，既然祂在我们的救恩中是不可分离的？\par

33. 我知道对疑难问题的详尽解答并不讨读者喜欢，但如果我允许自己暂时推迟全面真理的阐述，而用福音的这些话与异端者搏斗，或许对信仰有利。你听到主的话：\BibleQuote{永生就是：认识祢，唯一的真天主，和祢所派遣的耶稣基督。} 请问，是什么让你觉得基督不是真天主？没有进一步的指示告诉你应如何思想基督。只有\Quote{耶稣基督}，没有祂通常自称的\Quote{人子}，没有祂常宣告的\Quote{天主子}，没有祂多次重复（令许多人跌倒）的\BibleQuote{从天上降下的生活的食糧}\BibleRef{若 6:51}。祂说：\BibleQuote{祢，唯一的真天主，和祢所派遣的耶稣基督}，省略了祂所有常用的名号和头衔，无论是本性的还是取用的。因此，如果宣认唯一的真天主和耶稣基督赐给我们永生，那么毫无疑问，耶稣基督这个名字在此具有天主的全部意义。\par

34. 但也许由于说\Quote{祢是唯一的}，基督就将自己与天主的关系与合一割裂了。是的，但祂在说\Quote{祢，唯一的真天主}之后，不是立即接着说\Quote{和祢所派遣的耶稣基督}吗？我诉诸读者的理智：当我们奉命信仰父——唯一的真天主——以及耶稣基督时，我们必须相信基督是什么？或者，也许，如果父是唯一的真天主，就没有基督为天主的余地了？这是可能的，如果因为有一个天主父，基督就不是那唯一的主。天主父是一个，这并不妨碍基督仍是唯一的主；同样，父的唯一真天主性并不妨碍基督仍是真天主：因为我们只有相信基督以及唯一的真天主，才能获得永生。\par

35. 来吧，异端者！你那愚妄的教义会教导我们如何相信基督？基督，祂分施永生，祂被圣父所光荣，也光荣圣父，祂战胜了世界，祂虽被离弃，却非独处，因圣父与祂同在，祂出于天主，并从圣父而来。祂带着如此神圣的能力而生：你又允许祂保有天主的本性与实在吗？我们信从惟一真天主圣父是徒然的，除非我们也相信祂所派遣的那位，即耶稣基督。你为何犹豫不决？告诉我们，基督当被宣认为何者？你否认了所记载的，那么剩下的，岂不是相信那未记载的吗？啊，不幸的固执！啊，抗拒真理的谎言！基督在与惟一真天主圣父的信德和宣认中是合一的：请问，否认祂是真天主，称祂为受造物，这是何等信德？因为若没有基督，信从惟一真天主便不是信德。但你心胸狭隘，异端者，无法领受圣神。你躲避属天言语的含义；被错误之毒蛇咬伤，你忘了若要获得永生，基督必须在惟一真天主的信德中被宣认为真天主。\par

36. 然而，教会的信德在宣认惟一真天主圣父的同时，也宣认基督。它不宣认基督是脱离圣父这一惟一真天主的真天主；也不宣认圣父是脱离基督的惟一真天主。它宣认基督是真天主，因为它宣认圣父是惟一真天主。因此，天主圣父是惟一真天主这一事实，也构成了基督是真天主。独生子天主因祂的自然诞生并未遭受本性的改变：祂按神圣源起的本性，生于永生天主的永生天主，藉那本性的真理，不能与惟一真天主分离。因此，从真正的神圣本性中产生了必然的结果：真正的神性所出者必是真生，且独一天主不能从自身产生出次等天主。天主的奥迹既不在单纯性中，也不在多元性中：因为既没有另一位天主从天主带着祂本性的特质而出，天主也不作为独一位格而存在，因为圣子的真生教导我们宣认祂为圣父。因此，受生的天主并未失去祂本性的特质：祂拥有那位的本性能力，祂藉自然诞生将那位本性保留在自己内。祂内的神性并未改变或退化，因为倘若祂的诞生带来了任何缺陷，那将更合理地归咎于那本源本性，即未能将自身的属性植入其后裔。这改变不会贬低那藉诞生进入新实体的圣子，而是贬低圣父，因祂在圣子的诞生中未能维持其本性的恒常，反而产生了某种外在且与自己相异之物。\par

37. 然而，正如我们常说的，人类观念的不足并不对应天主圣父与天主圣子合一中的不足：仿佛有延伸、连续或流溢，如同泉源涌出溪流，或树木在树干上支撑树枝，或火放出热量进入空间。在这些例子中，我们看到了扩展而没有任何分离：各部分相连，不独立存在，但热在火中，枝在树中，溪在泉中。因此，唯有事物本身独立存在；一个不转入另一个，因为树与枝是同一的，正如火与热、泉与溪一样。但独生子天主是完美而不可言喻的诞生而自存的天主，是无始无终天主的真苗裔，无形体本性的无形体后裔，永生真天主的永生真天主，本性不可与天主分离的天主。诞生的事实并不使他成为不同本性的天主，也不使产生祂实体的生成改变其本性的种类。\par

38. 再论祂所取的血肉之救恩计划，并藉着服从，祂空虚了自己天主的形象，基督，生为人者，为自己取得了新的本性，并非因失去能力或本性，而是因形态的改变。祂空虚了自己天主的形象，取了奴仆的形体，当祂降生时。但圣父的本性，祂与之有本性合一，并未因这血肉的摄取而受影响；然而基督，虽然存留在其本性的能力中，但就这暂时变化中所取的人性而言，祂连同天主的形象也失去了与神圣本性的合一。但降生成人可归结为：整个圣子，即祂的人性连同神性，因圣父恩宠的许可而继续留在圣父本性的合一中，并保留了不仅神圣本性的能力，而且那本性自身。因为要达成的目的是使人成为天主。但所取的人性绝不能以任何方式留在天主的合一中，除非藉着与天主的合一，它达到了与天主本性的合一。于是，既然天主圣言是在天主的本性内，那么成了血肉的圣言也将反过来留在天主的本性内。因此，如果血肉与圣言的荣耀结合，人耶稣基督就能留在天主圣父的荣耀中，而成了血肉的圣言也能恢复圣父本性的合一，即使就其人性而言，因为所取的血肉已获得了圣言的荣耀。因此，圣父必须使圣言重回祂的合一，使祂本性的后裔再次回到自身内受光荣：因为合一曾因新的救恩计划而受损，唯有当圣父将那由圣子所取的血肉与祂一同光荣时，才能恢复如初的完美。\par

39. 为此，主已很好地准备了他们的心智以理解这一信仰，接着祂引用了这些话：\BibleQuote{永恒的生命就是认识祢，唯一的真天主，和祢所派遣的耶稣基督}，以及祂降生成人时所显示的服从：\BibleQuote{我在地上已光荣了祢，完成了祢交给我的工作}。\BibleRef{若 17:3}然后，为使我们知道祂服从的报酬和整个天主计划的隐秘目的，祂继续说道：\BibleQuote{现在，父啊，求祢以祢自己光荣我，就是我在世界之前与祢所有的光荣}。难道有人否认基督仍居于天主性内，或相信祂与唯一真天主是可分离、可区别的？让他告诉我们这段祈祷的意义。\BibleQuote{现在，父啊，求祢以祢自己光荣我}。父为何应以自己光荣祂？这些话的含义是什么？从它们的含义推导出什么？父既不需要光荣，也没有\OriginalTerm{空虚了自己}{emptied Himself}光荣的形态。祂怎能以自己并以上帝与世界之先与祂所有的光荣来光荣子？\Emph{祂与祂共有的？}又是什么意思？基督并没有说：\Quote{我在世界之先与祢同在时所拥有的光荣}，而是说\BibleQuote{我与祢共有的光荣}。\Emph{我与祢同在时}会意味着\Quote{我居住在祢旁边}；但\BibleQuote{我与祢共有的}则教导祂本性的奥秘。此外，\BibleQuote{求祢以祢自己光荣我}与\Quote{光荣我}不同。祂不仅仅求被光荣，使祂拥有某种自己的特殊光荣，而是求被父以自己光荣。父要以自己光荣祂，使祂如同从前一样居于与父的合一中，因为因着\OriginalTerm{降生成人的服从}{obedience of the Incarnation}，与父光荣的合一离开了祂。这意味着光荣化应使祂恢复到那本性中，就是祂因\OriginalTerm{神圣出生的奥秘}{Mysterium divinae nativitatis}与之结合的本性；使祂被父以自己光荣；使祂重新拥有从前与父同在的一切；使\OriginalTerm{仆人形态}{form of a servant}的取用不会使祂疏离\OriginalTerm{天主性形态}{form of God}的本性，反而天主应在自己内光荣仆人形态，使它永远成为天主的形态，因为此前居于天主形态中的祂，如今居于仆人形态中。既然仆人形态要在天主形态中被光荣，那么它就要在祂内被光荣，仆人的形态要因祂的形态而受尊荣。\par

40. 但主的这些话并非新事，也非首次在福音的教导中被证实，因为当犹大出去出卖祂时，祂因希望实现而充满高贵的喜乐，以此见证了天主父以自己光荣子的奥秘。祂因自己的目的即将完全实现而充满喜乐地说：\BibleQuote{现在人子受到了光荣，天主也在祂内受了光荣。如果天主在祂内受了光荣，祂在自己内光荣了祂，并立刻光荣了祂}。\BibleRef{若 13:31}我们这些灵魂背负着泥土的身体、心智被罪恶的污秽意识污染玷污的人，怎能如此自大，去判断祂神圣的要求？我们怎能自以为是地批评祂属天的本性，以不圣洁和亵渎的争辩反抗天主？主以最简明的言语宣告了福音的信仰，并按照我们本性软弱所容许的程度，使祂的言论适合我们的理解，同时没有说出任何有损祂本身尊威的话。我想，祂开头话的意义是不可怀疑的：\BibleQuote{现在人子受到了光荣}；也就是说，祂所获得的一切光荣不是为\OriginalTerm{圣言}{the Word}，而是为\OriginalTerm{祂的肉身}{His flesh}：不是为\OriginalTerm{祂天主性的出生}{the birth of His Godhead}，而是为\OriginalTerm{祂降生于世界的属人本性的安排}{the dispensation of His manhood}。那么，请问，接下来\BibleQuote{天主在祂内受了光荣}的意思是什么？我听到天主在祂内受了光荣；但根据你的解释，异端者，我不知道那是什么意思。天主在祂内受光荣，就是在人子内受光荣：那么请告诉我，人子与天主子是同一个吗？既然人子不是一个，天主子也不是另一个，而是那位是天主子的祂自己也是人子，那么请问，在这位既为人子又为天主子的祂内被光荣的天主是谁？\par

41. 因此，天主在人子——也是天主子——内受到光荣。我们且来看，所添加的第三句：\Emph{如果天主在祂内受到光荣，天主也在祂自己内光荣了祂}，这是什么意思？请问，这是什么奥秘？天主，在光荣的人子内，在祂自己内光荣了一位受到光荣的天主！天主的光荣在于人子，而天主的光荣也在于人子的光荣。天主在祂自己内光荣，但人并非通过自己受到光荣。再者，那位在人间受到光荣的天主，虽然接受了光荣，但祂自己仍然是天主。然而，由于在光荣人子时，光荣的天主在祂自己内光荣了天主，我认出基督本性的光荣被纳入了那光荣祂本性的本性的光荣中。天主并不光荣自己，但祂在祂自己内光荣了在人间受到光荣的天主。而这句\Quote{在祂自己内光荣}，尽管不是祂自己的光荣，却意味着祂将那受光荣的本性纳入了祂自己本性的光荣中，因为那位光荣在人间受到光荣的天主的天主，在祂自己内光荣了祂，祂证明祂所光荣的天主是在祂自己内，因为祂在祂自己内光荣了祂。来吧，无论你是何方异端者，提出你那曲折教义中无法解决的反对意见吧；尽管它们自相缠绕，但随你如何罗列，我们都不会陷入它们的陷阱。人子受到光荣；天主在祂内受到光荣；天主在祂自己内光荣了那位在人间受到光荣的祂。人子受光荣，与天主在人子内受光荣，或天主在祂自己内光荣那在人间受光荣的祂，并不相同。用你不信神的信仰术语，说出你所谓天主在人子内受光荣是什么意思。这必定要么是基督在肉身内受光荣，要么是圣父在基督内受光荣。如果是基督，基督显然是天主，在肉身内受光荣。如果是圣父，我们便面对合一的奥秘，因为圣父在圣子内受光荣。因此，你若承认是基督，即使不情愿，你也宣认祂是天主；你若理解为天主圣父，你便不能否认天主圣父在基督内的本性。关于光荣的人子和在祂内受光荣的天主，说到此已足够。但当我们思考天主在祂自己内光荣那在人子内受光荣的天主时，请问，你那亵渎的教义能有什么漏洞逃避宣认基督按其本性真实是天主？天主在祂自己内光荣了那生而为人的基督；那么基督是在祂之外，而祂在祂自己内光荣祂吗？祂在祂自己内将那先前与祂同在的光荣归还给基督，如今祂所取的仆人之形又重归天主的形像，在人间受光荣的天主在祂自己内受到光荣；在祂空虚自己的安排之前，祂已在天主自身内，如今祂与天主自身联合，既在仆人之形内，也在属于祂出生的本性内。因为祂的出生并未使祂成为一个新异本性的天主，而是藉着出生，祂成为自然圣父的自然圣子。在祂人性出生后，当祂在祂的人性内受光荣时，祂再次以自己本性的光荣发光；圣父在祂自己内光荣祂，当祂被纳入圣父本性的光荣中时，这光荣是祂在安排中所空虚自己的。\par

42. 宗徒信德的话是对你鲁莽疯狂亵渎的屏障，禁止你把思想的自由变成放纵，误入歧途。他说：\Emph{众口要宣认耶稣基督是主，以光荣天主圣父}。圣父在祂自己内光荣了祂，因此祂必须在圣父的光荣内被宣认。如果祂应在圣父的光荣内被宣认，而圣父在祂自己内光荣了祂，那么祂岂不显然就是祂圣父所是的一切，因为圣父在祂自己内光荣了祂，而祂要在圣父的光荣内被宣认？如今祂不仅在天主的光荣内，而且在天主圣父的光荣内。圣父光荣祂，不是用外在的光荣，而是在祂自己内。圣父将祂带回那属于祂自己的、祂先前与祂同在的光荣中，从而在祂自己内并与祂自己一同光荣了祂。因此，这宣认与基督不可分离，即使是在祂人性卑微的时候，正如祂所说：\Quote{永生就是：认识祢，唯一的真天主，和祢所派遣的耶稣基督}\BibleRef{若 17:3}；因为首先，没有耶稣基督，宣认天主圣父就没有永生；其次，基督在圣父内受光荣。永生正是：认识唯一的真天主和祂所派遣的耶稣基督；如果你能通过相信天主而不信祂而获得生命，那就否认基督是真天主吧。至于天主圣父是唯一的真天主这一真理，除非基督的光荣完全在唯一的真天主圣父内，否则这于天主基督而言就是假的。因为如果圣父在祂自己内光荣祂，而圣父是唯一的真天主，那么基督就不在唯一的真天主之外，因为圣父——唯一的真天主——在祂自己内光荣了被提升到天主光荣中的基督。由于祂被唯一的真天主在祂自己内光荣，祂就不与那唯一的真天主分离，因为祂被真天主在祂自己内——唯一的天主——所光荣。\par

43. 但或许无神的不信者会对虔诚的信徒提出这样的主张：我们无法理解对真天主的软弱自白，比如：\Emph{\BibleQuote{我实实在在告诉你们：子不能由自己做什么，他看见父做什么，才能做什么。}}\BibleRef{若 5:19} 倘若犹太人的双重愤怒没有要求双重的回答，那么子不能由自己做什么，除非他看见父做什么，这确实就是软弱的自白。但基督在同一句话中回应了犹太人的双重指控，他们指控祂违反了安息日，并且因称天主为自己的父而把自己与天主同等。那么，你是否认为，通过将注意力集中在祂答复的形式上，就能将其内容抽离？我们已在另一本书中讨论过这段经文；然而，由于信仰的阐述因重复而得益而非受损，既然此刻要求我们如此，就让我们再次思考这些话。\par

44. 请听这答复的必要性是如何产生的：\Emph{\BibleQuote{因此犹太人迫害耶稣，并设法要杀害祂，因为祂在安息日做了这些事。}}\BibleRef{若 5:16} 他们向祂发怒，以致想要杀害祂，因为祂在安息日做了祂的工。但让我们也看看主答复了什么：\Emph{\BibleQuote{我父到现在一直工作，我也工作。}}告诉我们，异端者，父的工作是什么？既然借着子、在子内，一切有形无形的都是受造的？你这比福音还聪明的人，无疑从某个其他秘密的知识源头获得了关于父工作的知识，为将祂启示给我们。但父在子内工作，正如子自己所说的：\Emph{\BibleQuote{我对你们所说的话，不是凭我自己讲的，而是住在我内的父，祂做祂的工。}}你理解这些话的意思吗：\Emph{\BibleQuote{我父到现在一直工作？}}祂这样说，是为了让我们在祂身上认识父本性的能力，在安息日使用那拥有此能力的本性去工作。父在祂内工作，当祂工作时；无疑，祂与父的工作一同工作，因此祂说：\Emph{\BibleQuote{我父到现在一直工作}}，以便祂言语和行为此时的工可以被视为父的本性在祂内的工作。这个\Emph{到现在一直工作}将时间与说话的时刻等同起来，因此我们必须认为祂是指父正在做的那份工，因为它暗示了父在祂说话的时刻就在工作。并且，唯恐信仰仅限于对父的认识，而失去了永生的希望，祂立刻补充说：\Emph{\BibleQuote{我也工作}}，也就是说，父到现在所做的一切，子也照样做。这样祂就阐释了整个信仰；因为现在的工属于现在的时间；如果父工作，子也工作，那么它们之间就没有合为一人的结合。但旁观者的愤怒现在加倍了。请听接下来所说的：\Emph{\BibleQuote{因此犹太人越发想要杀害祂，因为祂不仅违反了安息日，而且称天主为自己的父，使自己与天主平等。}}\BibleRef{若 5:18} 请允许我在这里重复，根据圣史的判断和人类的普遍认可，子与父的本性是平等的；而这种平等只能通过本性的同一性才能存在。受生者不能从自身源头以外获得其所是，所生者也不能与生它者相异，因为只有从它那里它才成为它所是的。那么，让我们看主对这双重的愤怒答复了什么：\Emph{\BibleQuote{我实实在在告诉你们：子不能由自己做什么，他看见父做什么，才能做什么；凡父所做的，子也照样做。}}\par

45. 除非我们视这些话为祂陈述的有机部分，否则我们以任意而悖信的诠释强加于它们，便是对它们施以暴力。但如果祂的回答指的是他们愤怒的理由，那么我们的信仰便正确地表达了祂要教导的，而恶人的乖谬则失去其亵渎妄想所依仗的凭据。让我们看看，这个回答是否适用于安息日做工的指控。\BibleQuote{\Emph{圣子不能由自己做什么，除非看见圣父做什么}}。祂刚刚说过：\BibleQuote{\Emph{我父一直工作至今，我也工作}}。如果因圣父的本性在祂内所具有的权威，祂所做的每一件事，都是与在祂内的圣父同工，而圣父甚至在安息日也一直工作，那么那以圣父工作的权威为自己辩护的圣子，便被宣告无罪。因为\Quote{\Emph{不能做什么}}这些词，所指的不是能力，而是权威；祂不能由自己做什么，除非看见了什么。然而，看见并不赋予做事的能力，因此祂并非软弱，如果祂没有看见就不能做什么，但祂的权威被表明取决于看见。再者，\Quote{\Emph{除非祂看见了}}这些词，表示由看见而来的意识，正如祂对宗徒们说：\BibleQuote{\Emph{我告诉你们：举目观看田地，它们已经发白了，可以收割了}}。\BibleRef{若 4:35} 祂意识到圣父的本性住在祂内并在祂工作时在祂内工作，为防止有人以为安息日的主犯了安息日，祂宣布：\BibleQuote{\Emph{圣子不能由自己做什么，除非看见圣父做什么}}。这样，祂表明祂的每个行动都源于祂对内在工作的本性的意识；当祂在安息日工作时，圣父即使在安息日也一直工作。然而，在接下来的话中，祂提到了他们愤怒的第二个原因：\BibleQuote{\Emph{凡祂所做的，圣子也照样做}}。那么，凡圣父所做的，圣子也照样做，这是假的吗？天主子承认圣父的能力和工作与祂自己有区别吗？祂是否不愿要求与同等能力和本性相配的同等敬拜？如果祂如此，那么轻视祂的软弱，将祂从与圣父本性的同等中降格吧。但祂自己稍后又说：\BibleQuote{\Emph{为使所有人都尊敬圣子，如同尊敬圣父；不尊敬圣子的，就不尊敬派遣祂来的圣父}}。如果你能，在两者同样受尊敬时找出祂的卑下；在两者以同样能力工作时，证出祂的软弱。\par

46. 你为什么歪曲回答的背景，以贬损祂的天主性？针对安息日做工，祂回答说自己不能由自己做什么，除非看见圣父做什么：为表明祂的平等，祂宣称凡圣父所做的，祂也照样做。请用祂关于安息日的回答来加强你软弱的指控，如果你能反驳\BibleQuote{\Emph{凡圣父所做的，圣子也照样做}}的话。但如果\Quote{凡所做的}包括了毫无例外的一切事物，那么当圣父所做的事没有一件是祂不能做的事时，祂哪里被发现有软弱呢？当为祂和为圣父要求同样的一份尊荣时，祂的平等主张被任何软弱的插曲驳倒了吗？如果两者在行动上有相同的能力，并且在敬拜中要求同样的崇敬，我无法理解还能存在什么卑下的羞辱，因为圣父和圣子拥有相同行动的能力和同等的尊荣。\par

47. 虽然我们已经按照事实本身来解释这段经文，但为证明主的话\BibleQuote{\Emph{圣子不能由自己做什么，除非看见圣父做什么}}不仅不支持对祂本性的这种不敬的贬低，反而见证了祂有意识地拥有圣父的本性，并且因圣父的权威祂在安息日工作，让我们向他们表明，我们可以引出另一句主的话，与此问题相关：\BibleQuote{\Emph{我不由自己做什么，而是照圣父所教导我的，我说这些话。派遣我的与我同在；祂没有留下我独自一人，因为我常做祂所喜悦的事}}。\BibleRef{若 8:28} 你感受到\BibleQuote{\Emph{圣子不能做什么，除非看见圣父做什么}}这些词所暗示的是什么吗？或者，在\BibleQuote{\Emph{我不能由自己做什么，祂没有留下我独自一人，因为我常做祂所喜悦的事}}这句话中，包含了何等奥秘？祂不由自己做什么，因为圣父住在祂内；你能使这一点与圣父不离开祂（因为祂做喜悦的事）这一事实协调吗？异端者，你的解释使这两个陈述之间产生了矛盾：即祂不由自己做什么，除非受住在祂内的圣父教导；以及圣父住在祂内，因为祂常做圣父所喜悦的事。因为如果圣父住在祂内意味着祂不由自己做什么，那么祂如何能通过常做圣父所喜悦的事而配得圣父住在祂内呢？不是由自己而做的事，不算功劳。相反，如果圣父自己住在圣子内是这些事的作者，那么圣子的行为又如何能取悦圣父呢？不虔敬的人，你陷入了严重的困境；信仰的全副武装的虔敬已将你包围。圣子要么是一位行动者，要么不是。如果祂不是行动者，祂的行为如何能取悦？如果祂是行动者，那么由\Quote{\Emph{不是由自己}}而成的事，在何种意义上算是祂自己的？一方面，祂必须做了那些喜悦的事；另一方面，做了却不出于自己，就不是功劳。\par

48. 然而，我的对手，祂们本性的统一是如此之深，以至于每一位的单独行动都意味著双方的共同行动，而双方的共同活动又意味著每一位的单独活动。设想圣子行动，而圣父藉著祂行动。祂并非从自己行动，因为我们必须解释圣父如何住在祂内。祂以自己位格行动，因为按照祂为圣子的出生，祂自己做了中悦父之事。祂若\Emph{不从自己}行动，就会证明祂软弱，除非祂这样做使得祂所做的中悦圣父。但如果祂所行的作为以及祂所中悦的事不属于祂自己，而只是被住在祂内的圣父推动去行动，那么祂就不在神圣本性的统一中。因此，圣父住在祂内，教导祂，而圣子在行动时，不从自己行动；另一方面，圣子虽不从自己行动，却亲自行动，因为祂所做的中悦圣父。这样，在行动中保持了祂们本性的统一，因为一位虽亲自行动，却不从自己行动，而另一位虽未行动，却是主动的。\par

49. 请将此与那句话联系起来——你抓住那句话来支持软弱的指控——\BibleQuote{凡父赐给我的人，必到我这里来；到我这里来的，我决不把他赶出去；因为我从天上降下来，不是为行自己的旨意，而是为行那派遣我来的父的旨意。}\BibleRef{若 6:37} 但也许你说，圣子没有意志自由：祂本性的软弱使祂受必然性的支配，祂被剥夺自由意志，并被置于必然性之下，好使祂不拒绝那些被赐给祂并从父那里来的人。主并不满足于用祂的行动来证明统一之奥秘——祂不拒绝那些被赐给祂的人，也不寻求行自己的旨意而代替那派遣祂者的旨意——当犹太人重复了\Emph{派遣我来的}这话之后开始抱怨时，祂用这话确认了我们的解释：\BibleQuote{凡从父听见而学习的，都到我这里来。这不是说有人看见过父，惟独那从天主来的，祂看见过父。我实实在在告诉你们，信我的人有永生。} 现在，请先告诉我：父在哪里被听见？祂又在哪里教导祂的听众？除了那从天主来的以外，没有人见过父：有人曾听见那无人见过的那位吗？凡从父听见的，就来到子那里；凡听见子的教导的，就听见了父本性的教导，因为祂的属性在子内启示出来。因此，当我们听见子教导时，必须明白我们在听父的教导。没有人见过父，然而那来到子面前的，从父听见并学习来参加：可见，父藉著子的话语教导，虽然不被任何人看见，却藉著子的显现对我们说话，因为子因祂完美的出生，拥有父本性的一切属性。因此，独生天主渴望为父的权威作证，同时又灌输祂与父本性的合一，祂不赶走那些由父赐给祂的人，也不行自己的旨意代替那派遣祂者的旨意：这并不是说祂不意愿祂所做的，或祂教导时不被听见；而是为了在那不可区分的一性面目下启示那派遣祂者和祂自己这被派遣者，祂表明祂所意愿、所说、所做的一切，都是父的意愿和作为。\par

50. 但是，祂用这话充分证明了祂的意志是自由的：\BibleQuote{就如父唤起死者并赋予生命，同样，子也随意赋予生命。}\BibleRef{若 5:21} 当指明父与子在能力和尊荣上平等时，圣子意志的自由便显明出来；当表明祂们的合一时，祂对父旨意的顺服便得到指示，因为父所意愿的，子就去做。但\Quote{去做}比\Quote{服从一个意愿}更多：后者意味着外在的必然性，而行另一个的意愿则需要与祂合一，是一种意志行动。在遵行父的旨意时，子教导说，由于本性的同一，祂的意愿在本性上与父的意愿相同，因为祂所做的一切都是父的意愿。子显然意愿父所意愿的一切，因为同一本性的意愿不能彼此不同。父的意愿在以下话语中启示出来：\BibleQuote{这是我父的旨意：凡看见子而信祂的，必得永生，并且在末日我要叫他复活。} 现在请听，子的意愿是否与父的意愿不一致，祂说：\BibleQuote{父啊，你所赐给我的人，我愿我在那里，他们也在那里。} 这里毫无疑问子意愿：因为父意愿那些信子的人得永生，而子意愿信者与祂同在。因为与基督同住不就是永生吗？祂说这话时，岂不赐给信祂的人一切福乐的圆满吗？\BibleQuote{除了父，没有人认识子；除了子及子所愿意启示的人，也没有人认识父。}\BibleRef{玛 11:27} 当祂愿意将父奥秘的知识传授给我们时，难道祂没有意志自由吗？祂的意志岂非如此自由，以至于祂能随己意将祂自己和祂父的知识赐给祂愿意的人？这样，父与子显然共同拥有一个因出生而共同、因合一而共同的本性：因为子意志自由，但祂自愿所作的，乃是父旨意的一行动。\par

51. 凡尚未掌握信仰之明确真理者，显然不能理解其奥秘；因为他没有福音的道理，便与福音的望德无分。我们必须宣认：圣父在圣子内，圣子在圣父内，因本性合一、因德能之力、因尊荣相等，作为生者与受生者。但或许你会说，主亲自的见证与此宣告相反，因为祂说：\Emph{\BibleQuote{父比我大}}。\BibleRef{若 14:28} 异端者啊，这就是你亵渎的武器吗？这就是你狂怒的兵器吗？你岂忘记了：教会不承认两个无始者，也不宣认两个父？你岂忘记了中保的降生成人，以及与之相属的诞生、马槽、童年、苦难、十字架和死亡？当你重生时，岂不曾宣认天主圣子、由玛利亚所生？如果那为此等事物所证的天主圣子说：\Emph{\BibleQuote{父比我大}}，难道你不知道，为你的救恩而降生成人是\InnerQuote{空虚了天主的形像}，而不受这种人性状况影响的圣父，却永恒地安居于自己不朽本性的真福中，并未取我们血肉？我们宣认：独生天主既保持于天主的形像内，也保持于天主的本性内，但我们并不立即将祂带有仆役形像的合一，再次吸收到天主性合一的实体中。我们也不教导说：父在子内，如同祂以形体方式进入子内；而是说：由父所生、与父同类的本性，自然拥有那生了它的本性；而这本性，既存留于生了它的本性的形像中，便取了人性和软弱的本性。基督拥有属于祂本性的一切；但天主的形像已由祂而出，因为\InnerQuote{祂以空虚自己}，取了仆役的形像。天主性并未消失，仍存在于祂内，只是已承受了尘世诞生的谦卑，并在所取的谦卑形态中施展其本有的德能。如此，由天主所生的天主，在仆役形像下显现为人，但在奇迹中却行天主之事：祂既是天主，因祂的作为证明；又是人，因祂显于人的形态。\par

52. 因此，在我们上面所论述的言论中，祂已为祂与父本性的合一作了见证：\Emph{\BibleQuote{看见我的，就是看见了父}}\BibleRef{若 14:9} \Emph{\BibleQuote{父在我内，我在父内}}。这两处经文完全一致，因为两个位格本性相等：看见了子，就等于看见了父；一位存于一位之内，表明他们不可分离。并且，为避免他们误解，好像他们看见祂的身体时，就在祂内看见了父，祂又加上说：\BibleQuote{你们该相信我：我在父内，父也在我内；否则至少该因那些事业而相信}。祂的德能属于祂的本性，而祂的作为就是这德能的运用；因此，在德能的运用中，他们可以在祂内认出与父本性的合一。凡在祂本性的德能中认出祂是天主的人，就会逐步认识天主圣父——存在于那大能的本性中。与父平等的圣子，以自己的作为显示了父可以在祂内被看见：为使我们在子内感知到与父相似的德能本性时，得以知道在父与子之间没有本性的区别。\par

53. 于是，独生天主在即将完成祂在肉躯中的工程、并圆满那取了仆役形像的奥迹之前，为坚固我们的信德，如此说：\Emph{\BibleQuote{你们听见我给你们说过：我去，但我还要回到你们这里来。如果你们爱我，就该喜欢我往父那里去，因为父比我大}}。祂在这同一篇言论的前部，已经全面展开了关于祂天主性的教导；那么，我们能否凭此宣认，就剥夺圣子那由祂真实诞生所成全的平等？或者，对独生天主而言，无始天主是祂的父，这岂是屈辱？因为祂由无始者所独生，便赋予了祂独生的本性。祂不是自己存在的根源，也不是在自己不存在时从虚无中生了自己；而是作为活的本性、来自活的本性，祂拥有那本性的德能，并宣告那本性的权威：祂为自己的尊荣作证，并在其尊荣中为所受诞生的恩宠作证。祂向父献上服从派遣祂者意志的顺服，但顺服的谦卑并不取消祂本性的合一：祂服从至死，但死后，祂超越一切名号\BibleRef{斐 2:8}。\par

54. 但若因祂脱去了天主的形体而怀疑祂的平等，我们就因忽视祂所取的谦卑的奥秘而羞辱了祂。祂人性的出生给祂带来了一个新的本性，祂的形体在祂的谦卑中改变了，借着取了奴仆的形体，但现在赐予名字恢复了祂形体的平等。问你自己这是什么，即被赐予的是什么。如果恩赐是有关天主的，那么对接受本性的赐予并不损害赐予本性的天主性。再者，经文\BibleQuote{并赐给祂那名字}在赐予中包含一个奥秘，但赐予名字并不使其成为另一个名字。赐给耶稣，使祂\BibleQuote{天上、地上和地底下的一切，无不屈膝叩拜，并一切舌头都承认耶稣是主，以光荣天主圣父}。这尊荣赐给祂，使祂以光荣天主圣父被承认。你听见祂说\BibleQuote{父比我大}吗？也要认识祂，关于祂因顺服的赏报而说\BibleQuote{并赐给祂那超乎万名之上的名字}；要听那说\BibleQuote{我与父原是一体；看见我的，就是看见了父；我在父内，父在我内}的祂。思考赐给祂的承认的尊荣，即耶稣是主，以光荣天主圣父。那么，父何时比子大？无疑，当祂赐给祂超乎万名之上的名字时。另一方面，何时子与父是一体？无疑，当一切舌头都承认耶稣是主以光荣天主圣父时。那么，如果父因赐予的权柄而更大，子是否因接受的承认而更小？赐予者更大；但接受者并不更小，因为祂被赐予与赐予者成为一体。如果耶稣没有被赐予以光荣天主圣父被承认，祂就比父小。但如果祂被赐予在那光荣中，即父所在的光荣中，我们在赐予的特权中看到赐予者更大，在恩赐的承认中看到二者是一体。因此，父比子大：因为显然更大的是那使另一位成为祂自己所是的一切，借着出生的奥秘将祂那未生的本性的肖像分施给子，从祂自身生出子到祂自己的形体中，并再次从奴仆的形体恢复到天主的形体，祂的作为是：基督，按圣神以光荣天主圣父作为天主而生，但现在按肉身死去的耶稣基督，应再次以光荣天主圣父作为天主。因此，当基督说祂要往父那里去时，祂揭示了为什么如果他们爱祂，他们应当喜乐的原因，因为父比祂大。\par

55. 在解释爱是这喜乐的根源之后，因为爱喜乐耶稣将以光荣天主圣父被承认，祂接着在话语中表达祂要求重新得回那光荣：\BibleQuote{因为这世界的首领来到，他在我内一无所用。}\BibleRef{若 14:30}这世界的首领在他内一无所用：因为既显为人的形状，祂居住在相似罪过的肉身中，却脱离了肉身的罪过，并在肉身中以罪罚定了罪。然后，以服从父的命令为唯一的动机，祂加上：\BibleQuote{但为使世界知道我爱父，并且父怎样命令我，我就照样行。起来，我们从这里走罢。}祂急于遵行父的命令，起身并急忙去完成祂身体受难的奥秘。但紧接着祂揭示了祂取肉身的奥秘。借着这取肉身，我们在祂内，如枝子在葡萄树上\BibleRef{若 15:1}；除非祂成了葡萄树，我们不能结出好果子。祂劝我们借着信德在祂所取的身体内住在祂内，既然圣言成了血肉，我们可以在祂肉身的本性内，如枝子在葡萄树内。祂将父威严的形体与所取肉身的卑微分开，称自己为葡萄树，所有枝子合一的根源，称父为细心的园丁，祂剪去无用和不结果的枝子丢在火里。在话语\BibleQuote{看见我的，就是看见了父}和\BibleQuote{我对你们所说的话，不是凭我自己说的，而是住在我内的父，祂做祂的工程}以及\BibleQuote{你们当信我，我在父内，父在我内}中，祂揭示了祂出生的真理和祂降生成人的奥秘。然后祂继续祂言论的线索，直到祂说到\BibleQuote{父比我大}；此后，为完成这些话的意义，祂急忙加上园丁、葡萄树和枝子的比喻，这比喻将我们的注意力引向祂对身体的卑微的服从。祂说，因为父比祂大，祂要往父那里去，并且爱应该喜乐，因为祂往父那里去，即从父那里得回祂的光荣：与祂一起，在祂内，得光荣不是以全新的尊荣，而是旧的，不是以陌生的尊荣，而是祂在起初就有的。那么，如果基督不带着光荣进入祂内，以住在天主的光荣中，你就可以贬低祂的本性；但如果祂所接受的光荣是祂天主性的证明，就应承认作为这证明的赐予者，父是更大的。\par

56. 为什么你将\OriginalTerm{降生成人}{Incarnation}扭曲为亵渎？为什么将救恩的\OriginalTerm{奥迹}{mystery}颠倒为毁灭的武器？显扬\Person{圣子}的\Person{圣父}是更大的：在\Person{圣父}内受显扬的\Person{圣子}并不更小。当祂在天主\Person{圣父}的光荣中时，怎能更小？\Person{圣父}又怎能不是更大的？因此\Person{圣父}更大，因为祂是\Person{圣父}；但\Person{圣子}，因为是\Person{圣子}，并不更小。藉\Person{圣子}的诞生，\Person{圣父}被确立为更大：祂由诞生而有的本性，不容许\Person{圣子}更小。\Person{圣父}更大，因为\Person{圣子}祈求祂将光荣赐予祂所取的\Person{人子}。\Person{圣子}并不更小，因为祂与\Person{圣父}一同收回祂的光荣。如此，诞生的奥迹与降生成人的\OriginalTerm{救恩安排}{dispensation}同时完成。\Person{圣父}，作为\Person{圣父}，作为显扬那现今为\Person{人子}的那位，是更大的：\Person{圣父}与\Person{圣子}是一体，因为\Person{圣子}，由\Person{圣父}所生，在取了地上的身体后，被带回\Person{圣父}的光荣中。\par

57. 因此，诞生并不使祂的\OriginalTerm{本性}{nature}低下，因为祂具有天主的形体，作为由天主所生者。虽然按其字面意义，\OriginalTerm{非受生者}{Unbegotten}与\OriginalTerm{受生者}{Begotten}似乎对立，然而受生者不能被排除于非受生者的本性之外，因为祂的\OriginalTerm{本体}{substance}无从来自其他。祂确实不分有非受生的至高威严；但祂从非受生的天主领受了\OriginalTerm{天主性}{divinity}的本性。如此，\OriginalTerm{信德}{faith}宣认\OriginalTerm{独生子天主}{Only-begotten God}的\OriginalTerm{永恒}{eternity}，尽管在祂身上，我们不能赋予受生或起始任何意义。祂的本性禁止我们说祂曾开始存在，因为祂的诞生超越时间的开始。但当我们宣认祂在万世之前已存在时，我们毫不犹豫地宣称祂在无始的永恒中受生，因为我们相信祂的诞生，虽然我们知道它从未有开始。\par

58. 为贬低祂的本性，\OriginalTerm{异端者}{heretics}抓住这样的话：\BibleQuote{圣父比我大}，或\BibleQuote{至于那日子和时辰，没有人知道，连天上的天使也不知道，子也不知道，惟独圣父知道}。这竟被用作对\OriginalTerm{独生子天主}{Only-begotten God}的责备：祂不知道那日子和时辰；也就是说，虽然祂是天主，由天主所生，但祂并不在天主本性的完美中，因为祂受制于\OriginalTerm{无知}{ignorance}的限制；也就是说，一种比祂更强的外力，仿佛胜过祂的\OriginalTerm{软弱}{weakness}，使祂被这\OriginalTerm{懦弱}{infirmity}所俘。的确，异端者看似有权声称这种\OriginalTerm{宣认}{confession}是无可避免的，他们癫狂地驱使我们走向如此\OriginalTerm{亵渎}{blasphemy}的解释。这些话是\Person{主}亲自说的，试问，还有什么比试图以我们的解释来曲解祂明确的断言更不圣洁的呢？\par

59. 但是，在我们探究这些话的意义与场合之前，让我们先诉诸常识的判断。祂，作为万有现时与未来的\OriginalTerm{根源}{Author}，难道会不知道万有吗？如果万有都是藉着\Person{基督}，并在基督内，而且如此藉着基督以致它们也在祂内，那么那既在祂内又藉着祂的，岂不也应在祂的\OriginalTerm{知识}{knowledge}中？那知识，凭藉一种不能无知的\OriginalTerm{本性}{nature}，惯常地领悟那既不在祂内也不藉着祂的事物。但那唯一从祂获得起源，并唯一在祂内拥有其现时状态与\OriginalTerm{未来}{future}发展的动力因的事物，岂能超出祂本性的认知？通过祂本性，万物得以成就，并在祂本性内包含万有，无论现在或将来。\Person{耶稣基督}知道\OriginalTerm{心思}{mind}的意念，如它现在被现时的动机所激动，也如它明日将被未来欲念的冲动所唤起。请听\Person{圣史}的见证：\BibleQuote{耶稣从起初就知道谁不信，谁要出卖祂}\BibleRef{若 6:64}。凭其德能，祂的本性能感知尚未诞生的未来，并预见尚潜伏在心思中的\OriginalTerm{情欲}{passions}的觉醒：你相信祂不知道那藉着祂并在祂内的事物吗？祂是万有属于他人的\Person{主}，难道祂不是自己所有的主吗？请记住关于祂所记载的：\BibleQuote{万有都是藉着祂、并在祂内受造的；祂在万有之先}\BibleRef{哥 1:16}；或者又说，\BibleQuote{因为圣父乐意使一切圆满居住在祂内，并藉着祂使万有与自己和好}，一切圆满都在祂内，万有都是藉着祂受造，并在祂内\OriginalTerm{和好}{reconciliation}，我们期待那和好的日子；那么，当时刻在祂手中，并由祂的奥迹所定，因为那是祂来临的日子，关于这日，\Person{宗徒}写道：\BibleQuote{当基督，你们的生命，显现时，你们也要与祂一同显现在光荣中}。没有人知道自己里面和自己所成就的事：基督要来，祂却不知道祂来临的日子吗？那是祂的日子，因为同一宗徒说：\BibleQuote{主的日子要像夜间的贼一样来到}\BibleRef{得前 5:2}；那么，我们能相信祂不知道吗？人的本性，尽其所能，预见他们决定要做的事：对预定目的的知识伴随着行动的意愿；那生为天主的祂，难道不知道在祂内和藉着祂的事物吗？时期是藉着祂，日子在祂手中，因为未来是藉着祂而构成的，祂来临的救恩安排在祂的权能中：难道祂的理解如此迟钝，以至于祂麻木的本性感觉不告诉祂自己所决定的事？祂岂像无理性的牲畜，没有理智或预见，甚至对自己的行为毫无意识，只被非理性欲念的冲动驱散，以其偶然与不确定的进程走向结局？\par

60. 但是，我们怎能相信光荣的主，因为祂能够不知道祂自己来临的日子，就具有不和谐和不完美的本性，受限于来临的必要，却不知道祂来临的日子呢？这会使得天主比无知的力量更软弱，这无知的力量从祂身上夺去了知识的特权。再者，如果我们不仅将软弱归咎于基督，还将缺陷归咎于天主圣父，说祂剥夺了祂所爱之子、独生子天主的预知，并恶意否认祂对将来终结的确定性：任凭祂知道祂苦难的日子和时辰，却向祂隐瞒祂德能的日子和祂在圣人们中光荣的时辰：在赐予祂预知死亡的同时，却夺去了祂对福乐的知识，那么我们岂不是加倍地加重了亵渎的理由？战兢的人心不敢如此揣测天主，或将人类反复无常的污点归咎于祂，以为圣父会拒绝给圣子任何东西，或圣子，祂是作为天主而生的，却拥有不完美的知识。\par

61. 但是，天主绝不可能是别的，只能是爱，只能是圣父：祂，那爱者，并不嫉妒；祂，那是圣父的，是完全而全然的圣父。这名号不容妥协：没有人能是部分父亲，部分不是。父亲在其整个位格上是父亲；他所有的一切都在孩子中，因为零碎的父亲身份是不可能的：这并非说父亲身份延伸到自我生发，而是说父亲在他所有的品质中完全地是父亲，对于从他而生的后裔。根据人身体的构造，人身体是由不同的元素组成，由各种部分构成，父亲必须是整个人的父亲，因为完美的生育将父亲中所有不同的元素和部分都传递给孩子。因此，父亲是他所有一切的父亲；生育从他的整个自我而出，并构成了孩子的全部。然而，天主没有身体，只有单纯的本质：没有部分，而是一个包罗万有的整体：没有被动之物，一切皆是活着的。因此，天主完全是生命，完全是合一，不由部分组合，在祂的单纯中是完美的，并且作为圣父，必须在祂自身所有的一切中，是祂所生之子的父亲，因为圣子完美的诞生使圣父在祂所有的一切中成为完美的父亲。所以，如果祂是圣子的真父亲，圣子就必须拥有圣父的一切属性。然而，如果圣子没有预知的品质，如果祂的诞生中缺少任何来自祂的根源的东西，这怎么可能呢？说天主有一项属性祂没有，几乎等于说祂没有任何属性。而什么是天主特有的，如果不是对未来的知识，一种拥抱不可见和尚未存在的世界的视野，并将尚未存在但将要存在的事物纳入其范围？\par

62. 此外，外邦人的导师保禄事先驳斥了那邪恶的谬误，即独生子天主部分无知。请听他的话：\BibleQuote{在爱中被教导，以达到理解的一切富藏，认识天主、即基督的奥迹，在祂内蕴藏着智慧和知识的一切宝藏。}\BibleRef{哥 2:2} 天主、即基督是这奥迹，智慧和知识的一切宝藏都隐藏在祂内。但是，一部分是一回事，整体是另一回事：部分不等于\Emph{全部}，也不能称\Emph{全部}为部分。如果圣子不知道那日子，知识的一切宝藏就不在祂内；但祂内藏有一切知识的宝藏，因此祂并非不知道那日子。但我们要记住，那些知识的宝藏是\Emph{隐藏的}藏在祂内，虽然隐藏，却不因此缺乏。如同在天主内，它们在祂内；如同在奥迹内，它们是隐藏的。但基督，天主的奥迹，在祂内隐藏着一切知识的宝藏，祂自己却不向我们的眼目和心智隐藏。既然祂自己就是奥迹，让我们看看，当祂不知道时，祂是否无知。如果别处祂宣告无知并不表示祂不知道，那么这里祂不知道时，称祂无知也是错误的。在祂内隐藏着一切知识的宝藏，因此祂的无知是一种\OriginalTerm{安排}{economy}，而非无知。这样我们就可以为祂的无知找出理由，而不必假设祂不知道。\par

63. 每当天主说祂不知道时，祂确实承认无知，但并非处于无知的缺陷之下。祂不知道，并非因为无知的软弱，而是因为说话的时候尚未到来，或天主的计划尚未行动。因此祂对\Person{亚巴郎}说：\BibleQuote{索多玛和哈摩辣的呼声实在很大，他们的罪极其深重。因此我要下去，看看他们是否完全照那呼声而行；若不然，我就知道。}\BibleRef{创 18:20} 在这里我们察觉到天主不知道那祂本来知道的事。祂知道他们的罪极其深重，但祂再次下去看看他们是否完全行了，并知道他们有没有行。于是我们看到，祂并非无知，虽然祂不知道，但当行动的时刻到来时，祂就知道了。因此，这知识不是从无知而来的改变，而是时机的满足。祂依然等待知道，但我们不能假设祂不知道：因此，祂不知道祂所知道的，并知道祂所不知道的，无非是言语和行为上的一种\OriginalTerm{安排}{economy}。\par

64. 我们不能因此怀疑，天主的认知依赖时机而非祂自身的任何改变：所谓时机，是指宣告认知的时机，而非获取认知的时机，正如我们从祂对\Person{亚巴郎}的话中所学到的：\Quote{不可在童子身上下手，也不要向他做什么，因为\Emph{现在我晓得}你敬畏天主，且为了我的缘故没有留下你的独生子。}\BibleRef{创 22:12} 天主现在知道，但那个\Emph{现在我晓得}是声称先前不知道：然而，直到现在天主不知道\Person{亚巴郎}的信德，这并不真实，因为经上记载：\Quote{\Person{亚巴郎}相信了天主，这就算为他的正义。}因此，这个\Emph{现在我晓得}标志着\Person{亚巴郎}得到这个见证的时刻，而非天主开始知道的时刻。\Person{亚巴郎}已通过献祭他的儿子证明了他对天主的爱，而天主在说话时就知道这事：但既然我们不能假定祂先前不知道，我们因而必须假定祂那时知道了，因为祂说了。\par

作为举例，我们从旧约中许多论及天主认知的段落中选出了这一处来思考，为的是表明，当天主不知道时，原因不在于祂的无知，而在于时机。\par

65. 我们在福音书中发现主耶稣知道许多事，却又不知道。例如，祂不认识那些以自己的大能作为和祂的名而夸耀的作恶者，因为祂对他们说：\Quote{那时我要发誓说：我从来不认识你们；你们这些作恶的人，离开我罢！}\BibleRef{玛 7:23} 祂甚至发誓说祂不认识他们，但祂知道他们是作恶者。祂不认识他们，不是因为祂不知道，而是因为他们的恶行使他们不配被祂认识，祂甚至以发誓的神圣性来确认祂的否认。凭借祂本性的德能，祂不可能无知；凭借祂旨意的奥秘，祂拒绝知道。再者，\OriginalTerm{无生的天主}{Unbegotten God}不认识那些愚拙的童贞女；当祂进入祂光荣来临的内室时，祂不认识那些粗心大意、没有预备好油的人。她们前来恳求，祂非但不认她们，反而喊道：\Quote{我实在告诉你们，我不认识你们。}\BibleRef{玛 25:12} 她们的到来和恳求迫使祂认出她们，但祂声称无知指的是祂的旨意，而非祂的本性：她们不配被那位无所不知者所认识。因此，为了不使我们把祂的无知归咎于软弱，祂立刻对宗徒们说：\Quote{所以，你们要警醒，因为你们不知道那日子，也不知道那时辰。}当祂吩咐他们警醒，因为他们不知道那日子或时辰时，祂指出祂不认识那些童贞女，是因为她们因沉睡和疏忽而没有油，因此不配进入祂的内室。\par

66. 主耶稣基督，\Emph{洞察人心肺腑}的\BibleRef{默 2:23}，祂的本性中没有软弱以至于祂不知道，因为我们察觉，甚至祂无知的事实也源自祂本性的全知。然而，如果有人将无知归咎于祂，让他们战栗吧，免得那洞察他们思念的祂对他们说：\Quote{\Emph{你们为什么心里思念恶事？}}\BibleRef{玛 9:4} 这位全知者，虽然并非不知思念和行为，有时却好像不知而询问，例如当祂问那女人是谁摸了祂的衣边，或问宗徒们他们为什么彼此争论，或问哀悼者\Person{拉匝禄}的坟墓在哪里：但祂的无知并非无知，只是言语上的。认为祂遥知\Person{拉匝禄}的死和埋葬，却不知其坟墓所在，这是不合情理的：认为祂能阅读心灵的思念，却不认识那女人的信德：认为祂无需询问任何事，却不知宗徒们的纷争。但是，祂知道一切，有时通过\OriginalTerm{救恩计划}{economy}的实践而声称无知，尽管祂并非无知。因此，在\Person{亚巴郎}的事例中，天主暂时隐藏了祂的知识；在愚拙的童贞女和作恶者的事例中，祂拒绝承认不配者；在人子的奥秘中，祂的询问仿佛不知，表达了祂的人性。祂使自己适应祂降生成人的现实，在每件事上都顺从了我们本性软弱的条件，但不是以祂的神性变得软弱的方式，而是天主降生为人，承担了人性的软弱，但并未因此将祂不变的本性降低为软弱的本性，因为那不变的本性正是祂奥秘地取了肉身的所在。祂原是天主，成为人，但既然是人，并没有停止是天主。因此，祂作为出生的人行事，并证明自己如此，尽管仍是天主圣言，祂常常使用人的语言（尽管天主，作为天主说话，也经常使用人的术语），并且不知道那还不宜宣告的事，或不配被祂承认的事。\par

67. 我们现在可以明白祂为什么说祂不知道那日子。如果我们相信祂真的无知，就与宗徒矛盾，宗徒说：\BibleQuote{\Emph{在祂内蕴藏着智慧和知识的一切宝藏}}。\BibleRef{哥 2:3} 有一种知识隐藏在祂内，而且因为它必须隐藏，有时就必须为这目的而被称为无知，因为一旦被宣告，就不再是秘密。因此，为了使那知识保持隐藏，祂声称自己不知道。但若祂不知道，是为了使知识保持隐藏，这种无知并非出于祂全知的本性，因为祂无知只是为了隐藏它。也不难看出为什么那日子的知识是隐藏的。祂劝勉我们要以不懈的信德时刻警醒，并剥夺我们确定知识的保障，使我们的心思被悬念的不确定性所绷紧，同时它们奔向并不断期待祂来临的日子，常在希望中警醒；并且，虽然我们知道那时刻必将来临，但其不确定性却能使我们谨慎警醒。因此主说：\BibleQuote{\Emph{为此，你们也应该准备，因为你们想不到的时刻，人子就来了。}}\BibleRef{玛 24:44} 又说：\BibleQuote{那仆人若被他的主人发现这样做，就有福了。} 所以，这种无知不是欺骗，而是鼓励恒心。被剥夺了拥有它并无益处的知识，并非损失，因为知识的保障可能会滋生对信德的疏忽，这信德如今是隐藏的，而期待的悬念使我们不断准备，如同家主因担心损失而时刻警醒，防备那选择睡觉时间作工的小偷的可怕来临。\par

68. 显然，天主的无知不是无知，而是一个奥秘：在祂的行动、言语和显示的安排中，祂不知道，同时也知道，或者知道，同时也不知道。但我们必须问，是否可能因圣子的软弱，祂不知道圣父所知道的？祂也许能读人心的思想，因为祂更强的本性可以与一个较弱的结合，在其一切活动中，并藉其力量的能力，可以说穿透那软弱的本性。但较弱的本性无力穿透较强的：轻的可被重的穿透，稀疏的可被致密的穿透，液体的可被固体的穿透，但重的不能被轻的穿透，致密的不能被稀疏的穿透，固体的不能被液体的穿透：强者不暴露于弱者，但弱者被强者所穿透。因此，异端者说，圣子不知道圣父的思想，因为祂自己软弱，不能接近那更强大的并进入祂内，或穿透祂。\par

69. 若有人胆敢不仅以他的口舌轻率地对独生天主\OriginalTerm{独生天主}{Unigenitus Deus}这样说，甚至在他心灵的邪恶中如此思想，就让他听听宗徒关于圣神所写的，他写给格林多人的话：\BibleQuote{但天主借着圣神将这一切启示给我们了，因为圣神洞察一切，连天主的深奥事理也洞悉。除了在人内的人灵外，有谁知道人的事呢？同样，除了天主的圣神外，也没有人知道天主的事。}\BibleRef{格前 2:10} 但让我们抛弃这些空虚的物质事物的比喻，以祂自己的德能，而非以世间的状况，来衡量天主所生的天主、圣神所发的圣神。让我们不以我们自己的感官，而以祂神圣的宣称来衡量祂。让我们相信祂所说的：\BibleQuote{看见了我，就是看见了父。}\BibleRef{若 14:9} 让我们不要忘记祂说：\BibleQuote{你们至少因我的工程而相信：父在我内，我在父内。} 又说：\BibleQuote{我与父原是一体。} 如果符合实在的名称在理智地使用时，能向我们传递任何真实的信息，那么藉着理解之眼在另一个人内被看见的那一位，在本性上就不异于那另一位；在种类上不异，因为祂住在父内，父也住在祂内；不分离，因为两者是一。请在它们本性的不可分性中感知它们的合一，并通过将这以一位视为另一位的镜子来领会那不可分本性的奥秘。但要记住，祂是镜子，并非如同由祂自身以外的一个本性的光辉所反映的像，而是作为一个活的本性，与父的活的本性不可区分，完全从父的全部本性中生出，因祂是独生子而拥有父在祂内，并因祂是天主而住在父内。\par

70. 异端者不能否认主曾用这些言词来表明祂诞生的奥迹，但他们试图将这些言词解释为旨意的和谐而逃避其含义。他们将天主圣父与天主圣子的合一视为并非在天主性上的合一，而仅仅是旨意的合一；仿佛神圣的教导在表达上如此贫乏，以致于主不能说：\Quote{我與父原是一心一意}；或者仿佛那些言词能与\Quote{我與父原是一體}意义相同；或者仿佛祂的意思是：\Quote{看见我旨意的，也看见了我父的旨意}；但因表达不熟练，便试图用\Quote{看见我的，也看见了父}这些话来表达那思想；或者仿佛神圣的词汇中不包含\Quote{我父的旨意在我内，我的旨意也在父内}这样的说法，但这思想却可以用\Quote{我在父内，父在我内}来表达。这一切都是令人作呕、不敬的胡言乱语；常识谴责这种愚蠢幻想的判断，即主不能说祂想说的，或者没有说祂所说的。诚然，我们发现祂用比喻和寓意说话，但用例证来加强言辞，或以暗示性的箴言来满足主题的尊严，或使语言适应时机的需要，则是另一回事。但我们所谈论的这段关于合一的经文，不允许我们在言词的平白意义之外寻求解释。如果父与子是一体，意思是他们在旨意上是一体，并且如果可分离的本性不能在旨意上合一，因为它们种类和本性的差异必然使它们陷入旨意和判断的差异，那么，他们若不在知识上合一，又怎能在一旨意上合一呢？无知与知识之间不可能有旨意的合一。全知与无知是相反的，相反的事物不能有相同的旨意。\par

71. 但或许有人认为，圣子宣认自己无知，证实祂说只有父知道。但除非祂明说只有父知道，否则对我们的理解来说将是极大的危险，因为我们可能认为祂自己不知道。因为，祂的无知是由于隐藏知识的计划，而不是由于能够无知的本性；既然祂说只有父知道，我们就不能相信祂不知道；因为，如我们上面所说，天主的认识不是发现祂所不知道的，而是宣告祂所知道的。只有父知道这一事实，并不证明圣子无知：祂说祂不知道，是为了使别人不知道；祂说只有父知道，是为了表明祂自己也知道。如果我们说，天主得知了亚巴郎的爱，是在祂不再隐藏祂的知识之时，那么，只有因祂没有向圣子隐藏这知识，才能说父知道那日子；因为天主不是通过突然的感知而学习的，而是借着时机宣告祂的知识。因此，如果圣子按照奥迹不知道那日子，是为了不启示它；另一方面，只有因祂启示了它，才能证明父知道那日子。\par

72. 愿我们远离在父与子身上想象身体变化的变迁，仿佛父有时对圣子说话，有时沉默。我们确实记得，有时从天上有声音为我们发出，为使父言语的能力为我们确认圣子的奥迹，正如主所说：\BibleQuote{这声音不是从天上为我而来的，而是为你们而来的。}\BibleRef{若 12:30} 但天主性可以免除人类机能所需的各种组合——舌头的运动、口形的调整、气息的促进和空气的振动。天主是单纯的存在；我们必须以虔诚理解祂，以敬畏宣认祂。祂应受朝拜，而非被感官追求，因为受造、软弱的本性无法以想象的猜测把握无限全能本性的奥迹。在天主内没有可变性，没有部分，如同合成的神性；在祂内，意志不应跟随无为，言语不应跟随沉默，工作不应跟随休息；祂也不应从某种其他精神状态转到意志，或在没有以声音打破沉默时说话，或在没有外出劳作时行动。祂不受自然律的支配，因为自然从祂领受了律法；祂行动时从不受软弱或变化之苦，因为祂的能力无限，正如主所说：\BibleQuote{父啊，一切于你都可能。}\BibleRef{谷 14:36} 祂能行超过人感官所能想象的事。主甚至不剥夺自己全能的属性，因为祂说：\BibleQuote{凡父所作的，子也照样作。}\BibleRef{若 5:19} 在没有软弱时，没有难事；因为只有软弱无力者才需要费力。困难的原因在于动力的软弱；无限能力强于无能的状态。\par

73. 我们确立这一点，是为了排除以下观念：天主在沉默之后才向圣子说话，或圣子在无知之后才开始知道。为达致我们的理解，必须使用适用于我们本性的术语：因此，除非藉着嘴上的言语，我们无法理解沟通；除非藉着知识，我们无法把握无知的反面。因此，圣子不知道那日子，是因为祂不启示它；祂说，只有圣父知道它，是因为圣父只向圣子启示它。但是，正如我们所说，基督并不意识到任何这样的自然障碍，即在祂能知道之前必须消除的无知，或圣父开始说话之前祂并不拥有的知识。祂作为独生子，藉着无可置疑的话语宣告祂与圣父在本性上的合一：\Quote{凡父所有的一切，都是我的。}这里并没有提到获得占有：拥有外在之物是一回事；自存自足是另一回事。前者是拥有天地和宇宙，后者是能藉着祂自己的属性来描述祂自己，这些属性是祂的，不是作为外在和从属之物，而是作为祂本身所赖以存在之物。因此，当祂说凡父所有的，都是祂的时，祂指的是天主性，而不是共同拥有赐予的恩赐。因为论到祂的话说圣神要从祂那里领受，祂说：\Quote{凡父所有的一切，都是我的；为此我说，祂要从我那里领受：}也就是说，圣神从祂那里领受，但也从圣父那里领受；如果祂从圣父那里领受，祂也从祂那里领受。圣神是天主之神，并不从受造物领受，而是教导我们祂领受所有这些恩赐，因为它们都是天主的。凡属于圣父的一切，也都是圣神的；但我们不可认为，凡祂从圣子领受的，祂就没有从圣父领受；因为凡父所有的，也同样属于圣子。\par

74. 因此，基督的本性不需要任何改变、询问或回答，以便从无知进到知识，或向那位一直沉默者提问，并等待领受祂的回答；反而，祂完美地住在与祂的奥秘合一中，从祂那里领受了全部的天主性存在，正如祂从祂那里获得了起源。并且，祂还领受了属于天主全部存在的一切，即祂的知识和祂的旨意。圣父所知道的，圣子并非通过问答来学习；圣父所意愿的，圣子并非因命令而意愿。既然凡父所有的，都是祂的，那么意愿和知道就是祂本性的属性，正如圣父意愿和知道一样。但是，为证明祂的诞生，祂常常解释祂位格的道理，例如当祂说：\Quote{我来不是为行我的旨意，而是为行那派遣我者的旨意。}祂行圣父的旨意，而非祂自己的，而祂所说的\Quote{那派遣我者的旨意}，指的是祂的圣父。但祂自己也同样意愿，这在不言而喻的话语中宣告了：\Quote{父啊！祢赐给我的人，我愿我在哪里，他们也同我在一起。}\BibleRef{若 17:24}圣父愿意我们与基督同在，按照宗徒的话，在祂内，祂在创世以前就拣选了我们\BibleRef{弗 1:4}；圣子也同样意愿，即我们与祂同在。因此，祂的旨意在本性上与圣父的旨意相同，尽管为表明诞生的事实，它被与圣父的旨意区分开来。\par

75. 因此，圣父所知道的一切，圣子都非无知，也不因唯独圣父知道，就推论圣子不知道。圣父与圣子住在本性的合一中，圣子的无知属于天主性沉默的计划，因为在他内蕴藏着一切智慧和知识的宝藏。主亲自证实了这一点，当祂回答宗徒们关于时期的提问时：\Quote{知道时期或季节，不是你们的事，那是父以自己的权柄所定的。}\BibleRef{宗 1:7}知识被拒绝给他们，不仅这样，连学习的焦虑也被禁止，因为知道这些时期不是他们的事。然而现在祂复活了，他们又再发问，尽管他们先前的提问曾得到答复，说连圣子也不知道。他们不可能字面理解圣子不知道，因为他们又向祂提问，好像祂知道一样。他们在祂无知的奥秘中看到了一个天主性沉默的计划，现在，在祂复活之后，他们重新提出这问题，以为说话的时候已经来到。圣子不再否认祂知道，却告诉他们，知道不是他们的事，因为圣父已将它置于祂自己的权柄之下。那么，既然宗徒们将圣子不知道那日子归因于天主性计划，而非软弱，我们岂能说圣子不知道那日子仅仅因为祂不是天主？请记住，天主圣父将那日子置于祂的权柄之内，以免被人知道，而圣子先前被问时，回答说自己不知道，但现在，不再否认祂的知识，回答说他们不该知道，因为圣父不是将时期置于祂自己的\Quote{知识}内，而是置于祂自己的\Quote{权柄}内。日子和时刻都包含在\Quote{时期}一词中：那么，那位将要复兴以色列王国者，祂自己难道不知道那复兴的日子和时刻吗？祂教导我们，在这圣父的专属特权中看到祂诞生的证据，然而祂并不否认祂知道；而当祂宣告这种知识不属于我们时，祂断言这属于圣父权柄的奥秘。\par

我们不可因此以为，因为祂说不知道那日子和那时刻，圣子就不知道。作为人，祂哭泣、睡眠和忧伤，但天主不会流泪、恐惧或睡眠。按照祂肉身的软弱，祂流泪、睡眠、饥饿、口渴、疲乏和恐惧，却无损祂独生子本性的真实性；同样，我们也必须将祂不知道那日子或时辰的话归于祂的人性。\par

\SourceRecord{Translated by E.W. Watson and L. Pullan.；Nicene and Post-Nicene Fathers, Second Series；Vol. 9.；Edited by Philip Schaff and Henry Wace.；Buffalo, NY: Christian Literature Publishing Co.,；1899.；<http://www.newadvent.org/fathers/330209.htm>.}

% structure: p0001 source=h1 target=section reason=page-title

\section{\WorkTitle{论天主圣三（卷十）}}

卷十。神学分歧并非诚实推理的结果，而是如亚略派那样，被先入为主的见解所扭曲的推理的结果；其原因是罪，其结果是伪善（§§ 1-3）。\Person{希拉利}落入了宗徒所预言的罪恶时代；真理被放逐，他本人也被放逐，然而他的苦难并不影响他在主内的喜乐（§ 4）。在前几卷中，他陈述了确切的真理，现在他给出了一个总结（§§ 5-8）。但进一步的异议被提出：既然天主是不受苦难的，基督在祂的受难中却遭受了恐惧和痛苦（§ 9）。但是，那教导他人不要惧怕死亡的那位，自己不可能惧怕死亡（§ 10）。祂出于自己的自由意志而死，知道在三天后祂的身体和灵魂将再次复活（§§ 11, 12）。祂也不惧怕身体的折磨，因为痛苦是居于我们身体内的软弱人类灵魂的一种情感，身体本身并不感受痛苦（§§ 13, 14）。而且，虽然童贞玛利亚圆满地完成了人类母亲的角色，但生祂的那一位是天主。基督在取了奴仆的形体时，仍然保持在天主的形体中，并且作为完全者而生，正如生祂的那一位是完全者，因为玛利亚不是祂人性的原因，而只是其工具（§§ 15, 16）。\Person{圣保禄}清楚地区分了第一人——出于地，和第二人——由圣神受孕，在祂内，一方面是肉躯，另一方面是从天降下的粮（§§ 17, 18）。因此，祂既是完全的人，也是完全的天主，并且没有继承亚当的肉躯或灵魂。祂的整个本性源于圣神，童贞玛利亚因圣神而受孕（§§ 19, 20）。再者（§ 21），亚略派论证说，圣言在耶稣内，正如圣神在先知们内一样，并指责天主教徒否认基督的真实人性。\Person{希拉利}回答说，正如基督是自己人性的身体产生之因，同样祂也是自己人性的灵魂的创造者，因为没有灵魂是传衍的。因此，祂的人性是完整的；祂取了奴仆的形体，但同时祂仍在天主的形体中，也就是说，既是天主也是人的那一位是同一基督，祂曾诞生、死亡并复活（§ 22）。在这一切中，祂承受了苦难，但没有痛苦，正如空气或水，若被一击穿透，却不受其影响。打击是真实的，苦难也是真实的；但这苦难不是施加在我们有限的人性上，而是施加在一个人性上，这人性可以在水面上行走，并穿过锁闭的门（§ 23）。若有人论证说祂曾哭泣、饥饿、口渴，\Person{希拉利}回答说，祂能够擦去眼泪并供应需要，因此祂并不受制于这些；尽管祂作为真实的人经历了这些，祂却不受其影响。这种痛苦是人的常态，祂经历它们是为了表明祂有一个真实的身体（§ 24）。因为祂确实有这样的身体，尽管（因为祂不是因罪而受孕）这个身体没有我们身体的缺陷；不是罪恶的肉躯，而只是罪恶肉躯的样式。因为祂是成了血肉的圣言，并且继续是真实的天主，如同祂先前一样（§§ 25, 26）。光荣的主在祂的受难中既没有遭受恐惧也没有遭受痛苦，正如祂在死亡边缘所行使的权能所示（§§ 27, 28）。祂在山园和十字架上的言语并非痛苦或恐惧的证据，因为它们可以被平静与希望的崇高表达所匹配（§§ 29-32）。因此，从受难的情境中无法得出恐惧、痛苦或软弱的证明。十字架也不是耻辱，因为它是祂从屈辱到荣耀的道路（§ 33），下降阴府也不是降级，因为祂始终在天上。十字架上的盗贼的信德与亚略派的信德是多么不同啊！（§ 34）。论点在 § 35 中总结。接下来讨论山园祈祷。基督并没有说祂因死亡而忧伤，而是说祂忧伤至死。这是为宗徒们的缘故而焦虑，免得他们的信德动摇；这种恐惧延伸至祂的死，而非超越，因为祂知道祂死后，祂的光荣将重振他们的信德。这就是天使安慰祂时所涉及的恐惧；祂自己则是无所畏惧的，因祂意识到自己的天主性（§§ 36-43）。祂没有痛苦和恐惧，因为正是罪恶的身体将这些情感传给灵魂。然而，即使人的身体有时也能超越它们，例如\Person{达尼尔}和其他信德的英雄们：更何况基督呢（§§ 44-46）。同样，我们必须理解祂担当我们的痛苦和我们的罪（§ 47），因为，正如\Person{圣保禄}所说，祂的受难本身就是一个凯旋（§ 48）。抱怨祂被圣父舍弃，也得到类似的解释（§ 49）。亚略派论证的目的，是以一个或另一个异端假设来取代基督是真天主与真人的真理，这一切都为教会所拒绝（§§ 50-52）。我们的理性必须承认其局限，并满足于相信那些表面上矛盾、却又无法理解的真理（§§ 53, 54）。基督为耶路撒冷哀哭，以及在拉匝禄墓前哭泣，同样无法解释，却又确定无疑（§§ 55, 56）。祂舍弃生命又取回生命，由两个不可分割地结合于一位格的本性来解释（§§ 57-62）。在简短总结（§ 63）之后，他回到两个本性的结合，这是世俗智慧的绊脚石（§ 64），并表明这是这些事实唯一合理的解释（§§ 65, 66）。正如\Person{圣保禄}所说，我们的信仰必须\Emph{按照圣经}；信德的必要性和赏报（§§ 67-70）。基督表面上的软弱是为了教导和救恩而接受的。\par

1. 显然，世人所说的一切无不受反对。意志不同意时，理智也不同意；在敌对判断的偏见下，它进行争辩，并否认它所反对的断言。尽管我们所说的每一句话若以真理的标准衡量都是无可辩驳的，但只要人们思想或感受不同，真理总是暴露在反对者的吹毛求疵之下，因为他们因错误或偏见的错觉而攻击他们误解或不喜欢的真理。因为已形成的决定会极其固执地固守；当意志不服从理智时，争论的热情就无法从其所走的道路上被驱离。对真理的探求让位于寻找我们希望相信的证明；欲望凌驾于真理之上。于是我们编造的理论建立在名称而非实物之上；真理的逻辑让位于偏见的逻辑：一种意志用来为其幻想辩护的逻辑，而不是通过理智对真理的理解来激发意志的逻辑。正是这些党派偏见的缺陷导致了对立理论之间的所有争论。随之而来的是真理自我主张与偏见自我辩护之间的顽强斗争：真理坚守其立场，偏见则抗拒。但是，如果欲望没有抢先于理智；如果对真理的理解促使我们去渴望真实的事物；我们就会让我们的学说支配我们的欲望，而不是试图把我们的欲望当作学说来树立；这样就不会有对真理的矛盾，因为每个人都会首先渴望真实的东西，而不是为自己所渴望的东西辩护其真理性。\par

2. 宗徒没有忘记这种意志的罪过，在写给弟茂德的信里，在多次嘱咐要为信德作证并宣讲圣言之后，他补充道：\BibleQuote{因为时候将到，人必不容忍健全的道理，反而耳朵发痒，随自己的私欲聚集许多教师，并掩耳不听真理，偏向无稽的传说}。\BibleRef{弟后 4:3} 因为当他们的不敬热情驱使他们超出容忍健全道理时，他们就会为自己的情欲聚集许多教师，也就是说，按照自己的欲望构造学说的计划，不愿受教，而是聚集那些会告诉他们所希望之事的人：以便他们搜罗并聚集起来的那一群教师，能以其喧嚣欲望的学说满足他们。如果这些狂人在其不虔敬的愚妄中不知他们以何种精神拒绝健全的、渴望腐败的学说，就让他们听听同一位宗徒致同一弟茂德的话：\BibleQuote{但圣神明白地说，在末后的时期，有些人要离弃信德，听从诱惑人的神和魔鬼的教理，这出于说谎者的虚伪}。\BibleRef{弟前 4:1} 发现自己所幻想的东西，而不是应该学习的东西，这算是什么学说进步呢？或者说，不渴望应该学习的东西，而是随从自己的欲望堆积学说，这算是什么学说中的虔敬呢？但这正是诱惑人的神所提供的。它们证实了假装虔敬的谎言，因为一种巧言令色的伪善总是继离弃信德之后而来：这样至少在言语上保留了良心已失去的敬意。连那种假装的虔敬，他们也以种种谎言使之成为不虔敬，用虚假学说的阴谋来违反信德的神圣：因为他们堆积学说以迎合自己的欲望，而不是根据福音的信德。他们以无法抑制的快感，任由发痒的耳朵因所偏爱的讲道的新奇而被搔痒；他们完全使自己远离对真理的听闻，完全投身于无稽的传说：以致他们既不能讲也不能理解真理，从而使他们的言论拥有了对他们而言像真理一样的外表。\par

3. 我们显然已经落在了宗徒所预言的邪恶时代；因为如今人们所寻求的教师，不是宣讲天主，而是宣讲受造物。人们对自己所渴望的事物，比健全信德所教导的更为热心。他们的发痒的耳朵已经如此强烈地促使他们去听他们所渴望的东西，以至于此刻在他们的众多博士中，只有那一种讲道占统治地位，它使独生子天主与天主圣父的权能和本性分离，并在我们的信德中使祂要么成为第二位的神，要么根本不是神；两者都是不虔敬的足以定罪的宣认，要么因设定不同等级的神性而宣认两位神；要么完全否认那位从天主藉由出生而获得本性者的神性。这样的学说取悦了那些耳朵远离真理的听闻而转向无稽传说的人，而我们这健全信德的听闻却不被容忍，并与它的宣讲者一道被强行流放。\par

4. 但是，尽管许多人可能随从自己的欲望聚集许多教师，并放逐健全的学说，然而真理的宣讲却永远不会从圣人们的团体中被流放。我们将在流放中通过这些著作讲话，而天主的圣言虽不能捆绑，却将不受阻碍地传播，警告我们这个宗徒所预言的时代。因为当人们表现出对真实信息的不断烦，并随从自己人性的欲望聚集许多教师时，我们就不再能对时代有所怀疑，而知道当健全学说的宣讲者被放逐时，真理也被放逐了。我们不抱怨时代；我们反倒欢喜，因为不义已在这个流放中显露出来，当它不能容忍真理时，就放逐健全学说的宣讲者，以便为自己随自己的欲望聚集教师。我们以流放为荣，在主内欢喜，因为宗徒的预言在我们身上得到了应验。\par

5. 在之前的书中，我们一直怀着（我深信是）真诚的信德宣认和纯正的真理，安排了整个答复的方式，使得（尽管我们有这样的局限，人的语言永远无法免于例外）没有人能在不公开承认无神的情况下反驳我们。因为我们如此彻底地阐明了他们狡猾地从福音中窃取并挪用为自己教训的那些经文的真正含义，以至于如果有人否认，他无法以无知为托词逃脱，而是被自己的话判定为无神。此外，我们按圣神的恩赐，在整个信仰证明过程中如此谨慎行事，以至于不可能对我们提出任何指控。因为他们的做法是用这样的声明塞满无知者的耳朵：当我们宣讲天主性的统一时，我们否认基督的诞生；而且他们说，通过经文\BibleQuote{我与父原是一体}\BibleRef{若 10:30}，我们承认天主是孤独的；因此，按照他们的说法，我们认为\OriginalTerm{无始的天主}{Unbegotten God}降入童贞玛利亚，并生为人，祂将开头的\Quote{我}一词归于祂血肉的经世，但又加上祂神性的证明\Quote{又与父}，作为祂自己（作为人）的父；此外，祂由两个位格（人性和神性）构成，论及自己时说\Quote{我们是一体}。\par

6. 但我们始终维护那超越时间的出生：我们教导，天主圣子是与天主圣父同一本性的天主，与无始者并非\OriginalTerm{同等}{co-equal}，因为祂自己并非无始，然而作为\OriginalTerm{独生者}{Only-begotten}，因受生而非不等；两者是一体，并非因给一个位格取双名，而是因真实的生与受生；在我们的信仰中，既没有两个不同种类的天主，天主也不是孤独的，因为祂是一，一如我们宣认独生天主的奥迹；而是圣子既在圣父的名中被指示，又存在於圣父内，圣父的名和本性都在祂内，而圣父以祂的名暗示并寓居於圣子内，因为一个儿子除非生自父亲，否则不能被言说或存在。此外，我们说祂是那活本性的\OriginalTerm{活的副本}{living copy}，神圣印记在神圣本性上的痕迹，在权能与种类上与天主如此无别，以至于祂的工作、言语和形象都无异于圣父；但是，由于肖象在本性上拥有其根源的本性，根源也通过祂本性的肖像工作、言说并显现。\par

7. 但在独生者这超越时间、不可言说的出生旁——它超越人类理解的感知——我们也教导了天主由童贞玛利亚腹中出生为人的奥迹，展示了按照\OriginalTerm{降生成人}{Incarnation}的计划，当祂空虚了天主的形体，取了奴仆的形体时，所取人性的软弱并未削弱天主性，而是神能赐予了人性，同时神性的德能并未在人的形体中丧失。因为当天主出生为人时，目的不是让天主性丧失，而是天主性保留，人生而成为天主。因此祂的名字是\OriginalTerm{厄玛奴耳}{Emmanuel}，即\BibleQuote{天主与我们同在}\BibleRef{玛 1:23}，好使天主不被降低到人的水平，而是人被提升到天主的水平。当祂求受光荣\BibleRef{若 17:5}时，这绝非光荣祂的天主性，而是祂所取的低贱本性：因为祂所求的光荣，是在世界受造以前，祂在天主面前所享有的光荣。\par

8. 我们正在答复一切，甚至他们最无知的言论，现在来到对未知时刻的讨论。现在，即使如他们所说，圣子不知道，这也不能成为攻击祂作为独生者的天主性的理由。照事物的本性，祂的出生并不能使祂的开端回溯，直到等同于那无始、无开端的存在；圣父保留作为祂的特权，为显示祂作为无始者的权威，来设定这个尚未确定的日期。我们也不能断定在祂的身位中，那个本性有任何缺陷，那个本性因出生之权包含了完美出生所能赋予的全部丰满。再者，在独生天主身上，对日子和时辰的无知不能被归咎于较低程度的神性。这是为了针对\OriginalTerm{撒伯流异端者}{Sabellian heretics}证明：圣父的权威是无生无始的，这\OriginalTerm{无始之权威的特权}{prerogative of unbegotten authority}并未赐予圣子。但是，如果我们所坚持的，当祂说祂不知道那日子时，祂保持沉默并非出于无知，而是依照天主的计划，那么所有不敬言论的借口都必须被消除，异端的亵渎教导被遏止，好使福音的真理能由那些看似掩盖它的话语本身得以阐明。\par

9. 因此，他们中的大多数人不肯承认祂拥有天主那不能受苦的本性，因为祂害怕祂的受难，并因屈服于苦难而显出软弱。他们断言，那害怕并感受痛苦者，不能享有超乎恐惧的能力之自信，也不能享有不知痛苦的精神之不朽性；然而，由于祂的本性低于天主圣父，祂在人性苦难前战栗恐惧，在肉体痛苦的暴力前呻吟。这些不敬的断言基于以下话语：\Quote{\Emph{我的灵魂忧闷得要死}} \BibleRef{玛 26:38}，以及\Quote{\Emph{父啊，若是可能，就让这杯离开我吧}}，还有\Quote{\Emph{我的天主，我的天主，祢为什么舍弃了我？}}，此外他们还加上\Quote{\Emph{父啊，我把我的灵魂交在祢手中}} \BibleRef{路 23:46}。他们将这些我们神圣信仰的话用于他们不敬的亵渎：即祂害怕，祂忧伤，甚至祈求那杯离开祂；祂感到疼痛，因为祂抱怨天主在祂苦难中离弃了祂；祂软弱，因为祂将灵魂交托给圣父。他们说，祂的疑虑和焦虑阻止我们将那相似于天主的性状归于祂——这性状本属于与天主同等的本性，因祂是作为圣父所生的独生子。祂藉着祈求挪去这杯、藉着被离弃的抱怨和交付祂的灵魂，宣告了祂自己的软弱和卑下。\par

10. 首先，在从这些经文本身证明祂并非因自己的缘故而受人性的恐惧或忧伤的软弱之前，让我们问：\Quote{我们在祂身上能找到什么可恐惧的，以致难以忍受的痛苦的恐惧抓住了祂？} 他们声称祂恐惧的对象，我想，是苦难和死亡。现在我问那些持此见解的人：我们能否合理假设，那从祂的宗徒身上驱走死亡恐惧的祂会害怕死亡？祂用这些话劝勉他们去殉道的荣耀：\Quote{\Emph{谁不背起自己的十字架跟随我，就不配是我的门徒}}，以及\Quote{\Emph{谁获得自己的性命，必要丧失性命；谁为我的缘故，丧失了自己的性命，必要获得性命}} \BibleRef{玛 10:38}。如果为祂而死就是生命，那么那以生命赏赐给为祂而死之人的祂，在死亡奥迹中能受什么痛苦呢？当祂劝勉门徒说\Quote{\Emph{你们不要害怕那些杀害身体的}} \BibleRef{玛 10:28}时，死亡能让祂害怕那加于身体的事吗？\par

11. 此外，死亡的痛苦对祂有什么恐怖呢？对祂而言，死亡是祂自由意志的行为。在人类中，死亡或是由于外在敌人对身体攻击而引起，如发烧、伤口、意外或跌倒；或是因为我们的身体本性被年龄征服而屈服于死亡。但是，独生子天主，祂有权舍弃自己的生命，也有权再取回它\BibleRef{若 10:18}，在喝了醋之后，祂见证了祂人性苦难的工程已完成，为在自己内完成死亡的奥迹，便低下头，交付了灵魂。如果我们的有死本性能够凭自己的意志呼出最后一口气，并在死亡中寻求安息；如果受打击的灵魂能够在身体尚未瓦解、精神尚未被肢体的破裂、穿刺和压碎所侵害（就好像在自己家中被强暴那样）时离去并逃逸；那么对死亡的恐惧可能会抓住生命的主；但这是如果祂的死不是出于自由意志。但祂若是自愿而死，并藉自己的意志交还了灵魂，死亡就没有恐怖；因为死亡在祂自己的权下。\par

12. 但或许出于人性的无知恐惧，祂害怕自己所拥有的死亡本身的能力；所以，尽管祂自愿而死，祂害怕是因为祂将死去。若有人这样想，让他们问：\Quote{死亡对哪一方是恐怖的，是对祂的灵还是对祂的身体？} 如果对身体恐怖，难道他们不知道圣者不应见到腐朽，三天之内祂要复活祂身体的殿宇吗？但如果死亡对祂的灵是恐怖的，那么基督应害怕地狱的深渊，而拉匝禄却在亚巴郎的怀中欢乐吗？认为祂会害怕死亡是愚蠢和荒谬的，因为祂能放下自己的灵魂，也能再取回它，祂为了完成人性的奥迹，将自愿死去。祂不能害怕死亡，因为祂死亡的能力和目的是只死片刻：恐惧与死亡的意愿和复活的能力是不相容的，因为这两者都夺走了死亡的恐怖。\par

13. 但是否是挂在十字架上的身体疼痛，或捆绑祂的粗糙绳索，或钉子钉入的残酷伤口，使祂惊慌呢？让我们看看那被钉、被绑、被刺的身体，耶稣这个人有什么样的身体，疼痛才能住在其中。\par

14. 我们身体的本性是这样的：当通过与有感觉的灵魂结合而被赋予生命和感觉时，它们变得比惰性、无感觉的物质更多。它们被触摸时有感觉，被刺痛时受苦，寒冷时发抖，温暖时感到愉快，饥饿时消瘦，食物时发胖。通过某种灵魂的灌注——灵魂支持和渗透身体——它们根据周围环境感到快乐或痛苦。当身体被刺或穿透时，是渗透身体的灵魂意识到并受苦。例如，肉体的伤口甚至能感觉到骨头，而当我们修剪从肉中突出的指甲时，手指没有感觉。如果因某种疾病肢体枯萎，它就失去了活肉的感觉：它可以被切割或烧灼，毫无痛感，因为灵魂不再与之混合。此外，当因某种严重需要必须切掉身体的一部分时，灵魂可以通过药物被催眠，这些药物克服痛苦，并使心灵对其感觉能力产生一种类似死亡的遗忘。然后，肢体可以无痛地被切除：肉体对所有感觉都死了，不理会刀的深刺，因为其中的灵魂睡着了。因此，正因为身体是通过与软弱的灵魂混合而活着，它才受到疼痛的软弱的支配。\par

15. 如果人耶稣基督祂肉身的生命以与我们肉身和灵魂相同的开始开始，如果祂作为天主，并非自己肉身和灵魂的直接创造者，当祂被造成人的模样和形状，作为人诞生时，那么我们可能认为祂感受到了我们肉身的痛苦；因为由于祂的开始，一个像我们一样的受孕，祂有一个由像我们一样的灵魂赋予生命的身体。但如果通过祂自己的行动，祂从童贞玛利亚那里取得了肉身，并且同样通过祂自己的行动将一个灵魂与如此受孕的身体结合，那么祂受苦的性质必定与祂身体和灵魂的性质相应。因为当祂空虚了自己，取了天主的形象，取了奴仆的形象，当天主之子也作为人子诞生时，没有失去祂自己和能力，天主圣言形成了完美的活人。因为天主之子如何作为人子诞生，祂如何取了奴仆的形象，同时仍保留天主的形象，除非（天主圣言本身就能从童贞玛利亚取得肉身，并为那肉身赋予灵魂，为了救赎我们的灵魂和肉身），那人基督耶稣是完美诞生的，通过取用童贞玛利亚所孕的身体，被造成奴仆的形象？因为童贞玛利亚独自由圣神受孕，她所孕的，虽然为了祂肉身的诞生，她从自己提供了妇女们总是贡献给播在他们体内的种子的那种元素，但耶稣基督并非由普通的人类受孕形成的。在祂的诞生中，其原因仅由圣神传递，祂的母亲扮演了与所有人类受孕中相同的角色：但由于祂的起源，祂从未停止是天主。\par

16. 主亲自在以下的话语中启示了祂摄取人性的这个深奥而美丽的奥秘：\BibleQuote{没有人升过天，除了那从天上降下来的人子，祂是在天上的。} \BibleRef{若 3:13} \Quote{从天上降下来}指的是祂源于圣神：因为虽然玛利亚贡献了祂在子宫成长和诞生中属于她性别的一切自然之物，但祂的身体并不源于她。\Quote{人子}指的是在童贞玛利亚内所孕的血肉的诞生；\Quote{祂是在天上的}暗示了祂永恒本性的能力：一个无限的本性，不能将自己限制在身体的界限内，而这身体本身的来源和基础就是祂。借着圣神的德能和天主圣言的能力，虽然祂住在奴仆的形象中，但祂始终以万有之主的身份临在，在天地之内和之外。所以，祂从天上降下，又是人子，却仍在天上：因为成为血肉的圣言并没有停止成为圣言。作为圣言，祂在天上；作为血肉，祂是人子。作为成为血肉的圣言，祂既是从天上来的，又是人子，且在天上，因为圣言的能力永恒地无身体而存在，仍然临在于祂所离开的天上：血肉的起源归因于祂，而非任何其他。因此，成为血肉的圣言，虽是血肉，却从未停止成为圣言。\par

17. 圣宗徒也完美地描述了基督身体这不可言喻的诞生奥秘，他说：\Emph{第一个人出于地，是属土的；第二个人出于天}。称祂为\Quote{人}，他表达了祂由童贞玛利亚所生，童贞玛利亚在行使她作为母亲的职责时，在人的受孕和诞生中履行了她性别的本分。当他说\Emph{第二个人出于天}时，他证明了祂源于圣神，圣神降临于童贞玛利亚。因此，祂既是人，又是从天而来，这个人是童贞玛利亚所生，由圣神受孕。宗徒如此说。\par

18. 主亲自再次启示祂诞生的这个奥秘，如此说：\BibleQuote{我是从天上降下的生活的食粮；谁若吃了我的食粮，将永远活着} \BibleRef{若 6:51}：祂称自己为食粮，因为祂是自己身体的根源。此外，为了不让人以为圣言为血肉放弃了祂自己的德能和本性，祂又说那是祂的食粮；因为祂是从天上降下的食粮，祂的身体不能被认为源于人类受孕，因为它被显示是来自天上的。祂关于祂食粮的言语是圣言取了身体的断言，因为祂接着说：\BibleQuote{你们若不吃人子的肉，不喝祂的血，在你们内便没有生命}。因此，既然那作为人子的存在也作为食粮从天降下，通过\Quote{从天降下的食粮}和\Quote{人子的肉和血}，必须理解为祂摄取血肉，这血肉由圣神受孕，并由童贞玛利亚所生。\par

19. 因此，带着这个身体的人耶稣基督，既是天主之子又是人子，祂空虚了自己，取了天主的形象，取了奴仆的形象。并非一个人子而另一个是天主之子；也非一个在天主的形象中，另一个在奴仆的形象中诞生为完人：因此，正如按天主为我们所定的本性，我们存在的创造者，人生而有身体和灵魂，同样地，耶稣基督凭借自己的能力，是有血肉和灵魂的天主和人，在祂内拥有完整而完美的人性，以及完整而完美的天主性。\par

20. 然而，许多人带着他们试图证明自己异端的技巧，惯常用这个错误欺骗无知者的耳朵：既然亚当的身体和灵魂都犯了罪，主也必须从童贞玛利亚取了亚当的灵魂和身体，并且她由圣神所孕的并非完人。如果这些人理解了降生成人的奥秘，他们也会同时理解人子也是天主之子的奥秘。仿佛从童贞玛利亚那里领受了那么多，祂也从她那里领受了灵魂；然而，虽然血肉总是由血肉所生，但每一个灵魂都是天主的直接作为。\par

21. 他们为了剥夺那与天主同在、在起初就是天主的圣言的\OriginalTerm{独生者天主}{Only-begotten God}的神性实体，将他仅仅视为天主声音的言词。他们认为圣子与天主圣父的关系，如同言语与说话者的关系。他们试图潜入一种立场，即不是那永远存在的圣言、居于\OriginalTerm{天主的形像}{form of God}内的天主，作为人的基督诞生，因此他的生命来自人的起源，而非来自一种属神受孕的奥迹；即他不是那由童贞女诞生而自取人性的\OriginalTerm{圣言天主}{God the Word}，而是居住在\Person{耶稣}内的天主圣言，如同先知们内居住的预言之神。他们指控我们说，基督身为人的诞生，其肉体和灵魂与我们的不同。但我们宣讲的是\OriginalTerm{圣言成了血肉}{Word made flesh}，基督空虚自己，舍弃天主的形像，而取了\OriginalTerm{奴仆的形像}{form of a servant}，按人的形态完美，依照我们的样式成为人而诞生：即祂是真实的\OriginalTerm{天主子}{Son of God}，也确实是真实的\OriginalTerm{人子}{Son of Man}，既不会因为生于天主而缺少人性，也不会因为由天主所生的人而缺少神性。\par

22. 但是，正如祂亲自从童贞女取得身体一样，祂也亲自从自己取得灵魂；尽管在普通的人出生中，灵魂从不来自父母。那么，如果童贞女仅仅从天主接受了她所怀的肉身，那么更肯定的是，那身体的灵魂只能来自天主。再者，如果同一个基督是\OriginalTerm{人子}{Son of Man}，同时又是\OriginalTerm{天主子}{Son of God}（因为整个的人子就是整个的天主子），那么，在\OriginalTerm{天主子}{Son of God}、\OriginalTerm{圣言成了血肉}{Word made flesh}之外，还宣讲另一个我不知道是谁的、被\OriginalTerm{圣言天主}{God the Word}所感召、如同先知的角色，是多么可笑；而我们的主\Person{耶稣基督}既是人子又是天主子。然而，因为他的灵魂忧闷至死，并且他有权将灵魂放下，也有权将它再取回，他们却想将灵魂从某个外来的源头得来，而非来自\OriginalTerm{圣神}{Holy Ghost}，他身体的受孕的创造者：因为\OriginalTerm{圣言天主}{God the Word}成为人，没有离开他本性的奥迹。他诞生也是为了一次不成为两个分离的存在，而是为表明：那在成为人之前是天主的，现在取了人性，就是天主和人。\Person{耶稣基督}、天主子，除了通过圣言成了血肉，即天主子虽在\OriginalTerm{天主的形像}{form of God}内，却取了\OriginalTerm{奴仆的形像}{form of a servant}，又怎能由\Person{玛利亚}所生呢？当那处于天主的形像内的取了奴仆的形像时，两个相反的东西被结合在一起。因此，他接受了奴仆的形像，与他停留在天主的形像内同样真实。使用同一个词\Emph{形像}来描述两个本性，迫使我们认识到他确实拥有两者。他是奴仆的\Emph{形像}，同时也是天主的\Emph{形像}。虽然他由于永恒的本性而是后者，并且前者是按照\OriginalTerm{天主的恩宠计划}{divine Plan of Grace}，但这词在两种情况中都有其真实含义，因为他两者都是：既真实地在天主的形像内，也真实地在人的形像内。正如取了奴仆的形像无异于生而为人，同样，在天主的形像内无异于就是天主：我们宣认他为同一位，不是因丧失神性，而是因摄取了人性：藉他的神性在天主的形像内，藉他的受孕\OriginalTerm{因圣神受孕}{conception by the Holy Ghost}而在人的形像内，成为人的样式。这就是为何在他作为\Person{耶稣基督}诞生之后，他的苦难、死亡和埋葬，他也复活了。我们不能在这些不同的奥迹中将他与他自己分开，以至于他不再是基督；因为基督，取了奴仆的形像的，正是那在天主的形像内的；死去的与诞生的同是一位；复活的与死去的同是一位；在天的与复活的同一位；最后，在天的与先前从天降下的同一位。\par

23. 因此，人\Person{耶稣基督}，\OriginalTerm{独生者天主}{Only-begotten God}，同时作为血肉和圣言，既是\OriginalTerm{人子}{Son of Man}又是\OriginalTerm{天主子}{Son of God}，在不停止成为自己（即天主）的情况下，取了相似于我们人性的真实人性。但是，当在这个人性中，他被打、受伤、被绳索捆绑、被高举时，他感受到苦难的冲击，却没有痛苦。正如一支箭穿过水，或穿透火焰，或击伤空气，它行使了它天性会做的事：它穿过、穿透、刺伤；但这一切对被击中的东西没有效果；因为在水上穿洞、在火焰中穿透、或击伤空气，是违背自然的，尽管箭的天性是穿孔、穿透和刺伤。同样，我们的主\Person{耶稣基督}遭受了打击、悬挂、钉十字架和死亡：但那攻击主的身体的苦难，虽然仍是苦难，却没有产生苦难的自然效果。它用尽所有暴力行使其惩罚的功能；但\OriginalTerm{基督的身体}{body of Christ}因其德能，承受了惩罚的暴力，却没有其知觉。诚然，主的身体本应像我们的本性一样能够感受疼痛，如果我们的身体拥有踏在水上、行走在波浪上而不被我们的脚步压沉或由于压力而分开的能力，如果我们能穿过坚固的物体，关闭的门对我们不构成障碍的话。但是，既然只有我们主的身体能借其灵魂的能力被支撑在水上，能在波浪上行走，并能穿过墙壁，我们怎能以人的身体类比那\OriginalTerm{因圣神受孕}{conceived of the Holy Ghost}的肉躯呢？那肉躯，即那\OriginalTerm{从天降下的食粮}{Bread from Heaven}，是从天降下的；那人性来自天主。他有一个身体可以受苦，并且他受苦了：但他没有一个能感受痛苦的本性。因为他的身体具有自己独特的本性；在山上它被改变成\OriginalTerm{天上的光荣}{heavenly glory}，它的触摸驱散了高烧，它的唾液赐予了新的视力。\par

24. 或许有人会说：\Quote{我们发现祂流泪、饥饿、口渴；难道我们不该认为祂也容易感受人性的一切其他情感吗？}但我们若不明白祂流泪、饥饿、口渴的奥迹，就请记住：那位哭泣的，也曾使人从死者中复活；祂没有为拉匝禄的死而哭泣，反而喜乐；那位口渴的，从祂自己内赐给了活水的江河。\BibleRef{若 7:38} 祂若能把水赐给口渴的人，就不能被口渴所焦灼。此外，那位饥饿的，能诅咒那棵树，因为它不为祂的饥饿结果实；但祂的本性若能被饥饿征服，又怎能以祂的一句话使那青翠的树枯干呢？而且，除了流泪、饥饿、口渴的奥迹之外，祂所取的血肉，即祂整个的人性，也曾暴露于我们的软弱之下：但即便如此，人性也未曾被那些软弱所带来的屈辱所苦。祂的哭泣不是为祂自己；祂的口渴不需要水来解；祂的饥饿不需要食物来止。经上从不说主在饥饿、口渴或忧愁时吃、喝或哭泣。祂顺应身体的习惯，以证明祂自己身体的真实性，并通过按照我们本性的方式行事，来满足人类身体的习惯。祂的吃喝并非出于祂自己的需要，而是迁就我们的习惯。\par

25. 因为基督确实有一个身体，但却是独特的，与祂的根源相称。祂并非通过人类受孕所伴随的种种情欲而进入存在：祂凭自己的大能行动进入了我们身体的形态。祂以仆人的形态承载了我们整体的人性，却脱离人类身体的罪过和缺陷：好使我们能在祂内，因为祂由童贞女所生，而我们的过错却不在祂内，因为祂是自己人性的根源，作为人而生，却非生于人类受孕的缺陷。祂诞生的这个奥迹，宗徒维护并证明，当他说：\BibleQuote{祂贬抑自己，取了仆人的形体，成了人的模样，形态上完全像人} \BibleRef{斐 2:7}：就是说，祂取了仆人的形体，就是作为人的形体而生；祂成了人的模样，形态上完全像人，祂身体的外貌和真实性证明了祂的人性，然而，虽然祂形态上完全像人，祂却不认识罪是什么。因为祂的受孕是按照我们本性的模样，而非具有我们的缺陷。因为恐怕\BibleQuote{祂取了仆人的形体}这话被理解为出于自然的诞生，宗徒加上说：\BibleQuote{成了人的模样，形态上完全像人}。这样，祂诞生的真实性就不致暗示我们软弱本性所伴随的缺陷，因为\BibleQuote{仆人的形体}暗示了祂诞生的真实性，而形态上完全像人则暗示了我们本性的模样。祂通过童贞女由自己作为人而生，形态上像我们有罪的败坏身体：正如宗徒在致罗马人书中所证明的：\BibleQuote{因为法律因肉性而软弱，无法成就的事，天主却做到了：祂派遣了自己的儿子，带着罪恶肉身的形状，为了罪而定了罪的罪}。\BibleRef{罗 8:3} 祂并非\Emph{形态上像人}，而是\Emph{形态上完全像人}；祂的肉不是罪的肉，而是罪的肉的形状。因此，肉的形态暗示了祂诞生的真实性，而罪的肉的形状则使祂脱离人性软弱的缺陷。所以，人耶稣基督作为人确实诞生了，基督在本性上没有罪：因为就祂的人性一面而言，祂被生，不能不成为人；就祂的神性一面而言，祂永不能停止成为基督。既然耶稣基督是人，祂便作为人服从了人的诞生：然而作为基督，祂脱离了我们败坏种族的软弱。\par

26. 宗徒们的信仰为我们理解这一奥迹作了准备，它证明耶稣基督在形态上完全像人，并且被派遣了带着罪的肉的形状。因为形态上完全像人，祂处于仆人的形体，却不在仆人本性的缺陷中；而处于罪的肉的形状中，圣言确实是肉，但却是罪的肉的形状，而非罪的肉本身。同样，耶稣基督作为人确实是具有人性的，但即便如此，祂也不能是别的，只能是基督，作为人通过祂身体的诞生而生，但在缺陷上却不是人，正如在根源上祂不是人一样。成了血肉的圣言不能不成为祂所成的肉；然而祂始终是圣言，即使祂成了肉。正如成了血肉的圣言不能放弃祂根源的本性，同样，由于祂本性的根源，祂不能不依然是圣言：但同时我们必须相信，圣言就是祂所成的那肉；然而始终带着保留：当祂住在我们中间时，肉不是圣言，而是住在肉中的圣言的肉。\par

虽然我们已经证明了这一点，但我们仍要看看，在祂所忍受的全部苦难范围内，我们能否在我们主身上随处发现身体疼痛的软弱。我们将暂时搁置对异端据以将恐惧归于我们主的那些段落的讨论；现在让我们转向事实本身：因为祂的话语不能表明恐惧，如果祂的行为显示信心。\par

27. 异端者，你以为荣耀的主害怕受苦吗？为什么当伯多禄因无知而犯此错误时，祂没有称他为\Quote{撒殚}和\Quote{绊脚石}\BibleRef{玛 16:22}？就这样，伯多禄，这位阻止受难奥迹的人，被温和的基督口中如此严厉的斥责巩固在信德中——此基督并非血肉所启示，而是天上的父所启示的。\par

你在追逐什么幻影般的希望，竟否认基督是天主，并把对受苦的恐惧归给祂？祂会害怕吗？那位主动走向武装的逮捕者队伍的主？祂的身体软弱吗？在祂走近时，追捕者们踉跄、行列散乱、仆倒在地，无法承受祂的威仪，而祂正把自己交给他们的锁链？什么软弱能禁锢祂的身体，祂的本性竟有如此能力？\par

28. 然而，祂或许是害怕伤口的疼痛。那么请问，对于那位仅凭触摸就使被砍掉的耳朵复原的主，钉子的刺入又有什么可怕的呢？你们这些断言主软弱的人，请解释一下，在祂的肉身软弱且受苦的时刻，这大能的行为是如何发生的。伯多禄拔出剑来砍去，大司祭的仆人站在那里，耳朵被削掉了。那耳朵的肉体是如何从赤裸的伤口上，因基督的触摸而复原的呢？在流血和剑劈的伤口之间，身体已如此残缺，那原本不存在的耳朵是从哪里长出来的？那先前没有的从哪里来？那缺失的从哪里复原？那创造了耳朵的手，会感到钉子的疼痛吗？祂阻止别人感受伤口的疼痛，难道祂自己会感受到吗？祂的触摸能使被割掉的肉体复原，难道祂会因害怕自己肉体的刺伤而忧伤吗？如果基督的身体有这种德能，我们岂敢在那本性中妄称软弱，这本性天然的力量能抵消所有人类天然的软弱？\par

29. 然而，也许由于他们误入歧途且邪恶的乖僻，他们从祂的灵魂忧闷至死这一点推断祂的软弱。\BibleRef{玛 26:38} 现在还不是责备你误解这段经文的时候，异端者。此刻我只问你：为什么你忘记当犹大出去出卖祂时，祂说：\BibleQuote{\Emph{现在人子受到光荣} \Emph{？}} \BibleRef{若 13:31} 如果受苦要光荣祂，那对受苦的恐惧怎能令祂忧伤呢？除非祂如此缺乏理性，以致害怕受苦，而受苦却要光荣祂？\par

30. 但也许人们会认为祂害怕到祈求那杯离去：\BibleQuote{\Emph{阿爸，父啊！一切为祢都可能：请将这杯从我这里撤去}} \BibleRef{谷 14:36} 从最狭隘的论证立场来说，难道你不能凭自己对以下这句话的阅读来驳斥这种迟钝的邪恶吗？\BibleQuote{\Emph{把你的剑收回鞘里：我父给我的杯，我岂能不喝吗？}} \BibleRef{若 18:11} 恐惧能使祂祈求那杯离去，而祂却因对天主计划的热情而急于完成它吗？说祂逃避祂所渴望的受苦是不一致的。你承认祂是甘愿受苦的：那么，承认你误解了这段经文，不是比以亵渎和鲁莽的愚蠢冲向断言祂祈求逃避受苦（尽管你承认祂甘愿受苦）更虔敬吗？\par

31. 然而，我想你也会用主的这句话武装自己进行你那不敬神的争辩：\Quote{\Emph{我的天主，我的天主，祢为什么舍弃了我？}} 也许你以为，在十字架的耻辱之后，祂圣父帮助的恩宠离开了祂，因此祂呼喊说自己在软弱中被遗弃。但如果你视基督的轻蔑、软弱和十字架为一种耻辱，你应当记得祂的话：\Quote{\Emph{我实在告诉你们：从今以后，你们要看见人子坐在大能者的右边，乘着天上的云彩而来。}}\par

32. 请问，在祂的受难中，你哪里能看到恐惧？哪里能看到软弱？或疼痛？或耻辱？不敬神的人说祂害怕吗？但祂亲口宣告祂甘愿受苦。他们坚持说祂软弱吗？祂彰显了祂的大能，当追捕祂的人惊慌失措，不敢面对祂时。他们争辩说祂感受到肉体伤口的疼痛吗？但祂在使被割伤的耳朵复原时显示了，尽管祂是肉身的，祂并未感受到肉体伤口的疼痛。触摸被割伤耳朵的那只手属于祂的身体：然而那只手从伤口中创造了一只耳朵：那么，那属于一个受软弱支配的身体的手，又怎能如此呢？\par

33. 但是，他们说，十字架对祂是一种耻辱；然而正是因着十字架，我们现在才能看见人子坐在大能者的右边，祂从童贞女的胎中作为人出生，如今带着天上的云彩威严地回来了。你的不虔敬使你盲目，看不到因与果的自然关系：不仅你充满其中的不敬神和错误的精神，从你的理解中隐藏了信德的奥迹，而且异端的迟钝把你拖到低于普通人智力的水平。因为按常理，凡我们所恐惧的，我们都会逃避；一个软弱的本性因其自身的脆弱而易陷于恐怖；凡感受疼痛的事物，其本性总是易于疼痛；凡耻辱之事总是一种贬低。那么，根据什么合理的原则，你竟然认为我们的主耶稣基督害怕祂所趋向之事；或吓住了勇敢的人，而祂自己却因软弱而战栗；或阻止了伤口的疼痛，却感受到了自己的疼痛；或因十字架的贬低而受辱，却因十字架而坐在天主之高天，并返回祂的王国？\par

34. 但或许你认为你的不虔敬仍有一个机会，可以从这句话：\BibleQuote{父啊！我把我的灵魂交托在祢手中}\BibleRef{路 23:46}，中看出祂害怕下到阴府，甚至害怕死亡的必然性。然而，当你读到这些话却无法理解时，难道保持沉默，或虔诚地祈求明白其意义，不是比以如此露骨的断言误入歧途、因自己的愚妄而疯狂到无法看清真理更好吗？当你听祂对十字架上的强盗所说的话：\BibleQuote{我实在告诉你，今天你就要同我一同在乐园里？}时，你怎能相信祂害怕深渊的深处、灼热的火焰或报应刑罚的坑穴呢？这样的本性拥有如此大的能力，不能被关闭在阴府的界限内，更不能受其恐惧的支配。当祂下降到\OriginalTerm{阴府}{Hades}时，祂从未离开\OriginalTerm{乐园}{Paradise}（正如祂作为人子在地上宣讲时，祂始终在天上），并向祂的殉道者应许那里的家，向他展示完全幸福的狂喜。身体的恐惧不能触及那确实下到阴府、却凭其本性的能力无所不在者。同样，地狱的深渊和死亡的恐怖也不能抓住那主宰世界、在属神能力的自由中无边无际、对乐园的狂喜充满信心的本性；因为那要下到阴府的主，也要住在乐园里。你若能，就从祂不可分割的本性中分出一部分能害怕惩罚的部分：把基督的一部分送到阴府去受苦，另一部分则必须留在乐园里为王；因为强盗说：\BibleQuote{当祢来临时，请记得我}。我猜想，正是当他听到钉子刺穿我们主的手时发出的呻吟，促使他作了这蒙福的信仰宣认：他从主那衰残受伤的身体认识了基督的王国！他恳求基督在祂来临时记得他；你却声称基督在十字架上垂死时感到害怕。主应许他：\BibleQuote{今天你就要同我一同在乐园里}；你却要使基督屈服于阴府和刑罚的恐惧。你的信仰有相反的期望。强盗在十字架上宣认基督在祂的王国里，并从十字架上获得了乐园的赏报；而你，将刑罚的痛苦和死亡的恐惧归给基督，将失去乐园和祂的王国。\par

35. 我们现在已经看到了基督言行中的能力。我们无可辩驳地证明，祂的身体并不分享自然身体的软弱，因为祂的能力能够驱除身体的软弱，以致当祂受苦时，苦难抓住了祂的身体，却没有给它带来痛苦的本性：这是因为，虽然我们身体的形像在主内，但祂因起源的缘故，并不在我们软弱和缺陷的身体内。祂是由圣神受孕，由童贞玛利亚所生，她尽了她性别的本分，却没有从人领受受孕的种子。她生出了一个身体，但那是因圣神受孕的；一个具有实在本质的身体，但其本性没有软弱。那身体确实是身体，因为它由童贞玛利亚所生；但它超越我们身体的软弱，因为它起始于属神的受孕。\par

36. 但即使现在我们已证明了宗徒的信仰，异端者仍以为可以用这句话来反驳：\Quote{我的心灵忧闷得要死}。他们说，这些话证明了本性的软弱的意识，使基督开始忧闷。现在，首先我诉诸普遍理智：我们所谓\Quote{忧闷至死？}是什么意思？它不能等同于\Quote{因死而忧闷}；因为哪里有因死而忧闷，死亡就是忧闷的原因。但\Quote{忧闷至死}则意味着死亡是忧闷的终点，而非原因。因此，如果祂忧闷\Emph{至}死，而非\Emph{因}死，我们必须追问，祂的忧闷从何而来？祂的忧闷不是某段时间，或人类无知不能确定的时期，而是一直持续到死。所以，祂的忧闷绝非由祂的死亡引起，反而被死亡消除了。\par

37. 为使我们明白祂忧伤的原因，让我们看看这忧伤的告白之前和之后的内容：因为在逾越节晚餐中，我们的主完全预示了祂苦难和我们信德的全部奥迹。在祂说了他们都要因祂而跌倒，但应许祂要在他们之前往加里肋亚去之后，伯多禄抗议说，即使其余的人都跌倒，他也要保持忠信，不跌倒。\BibleRef{玛 26:33} 但主以其天主性预知将要发生的事，回答说伯多禄要三次否认祂：好让我们从伯多禄知道其他人是如何跌倒的，因为他甚至因三次否认而陷入如此大的信德危险。之后，祂带了伯多禄、雅各伯和若望——前两位被拣选为祂的殉道者，若望则要坚强起来以宣讲福音——并说祂忧闷得要死。然后祂上前去祈祷说：\Quote{我的父，如果可能，就让这杯离开我；但不要照我，而要照祢的意愿。} 祂祈求这杯离开祂，但那时这杯确实已经在祂面前：因为新约的血为许多人的罪而倾流，那时已在实现。祂不是祈求这杯不临于祂，而是祈求它离开祂。然后祂祈求不要成就祂的意愿，并愿意祂所求的事不蒙应允。因为祂说：\Quote{但不要照我，而要照祢的意愿。} 藉着祂自发祈求杯的移去，表明祂与人类的焦虑相联，同时又使自己与祂与圣父不可分割地共有的旨意一致。此外，为表明祂不是为自己祈求，并且祂只寻求祂渴望和祈求之事的有条件成就，祂用这些话作为整个请求的前言：\Quote{我的父，如果可能。} 难道有任何事情对于圣父是可能性不确定的吗？但如果圣父没有不可能的事，我们就可以看到这个条件\Quote{如果可能}取决于什么：因为在这祈祷之后紧接着是这些话：\Quote{祂来到门徒那里，见他们睡着了，就对伯多禄说：你们竟不能同我醒寤一个时辰吗？醒寤祈祷吧，免得陷入诱惑：心神固然切愿，但肉体却软弱。} 这个忧伤和这个祈祷的原因还有什么可怀疑的吗？祂吩咐他们同祂一起醒寤祈祷，为的是使他们不陷入诱惑；\Quote{因为心神固然切愿，但肉体却软弱。} 他们是在忠实灵魂的恒心下许诺不跌倒，然而因肉体的软弱，他们将要跌倒。因此，祂忧伤并祈祷，不是为自己，而是为那些祂劝告要醒寤祈祷的人，免得苦难之杯成为他们的命运：免得祂祈求可能离开祂的那杯，却留在他们身上。\par

38. 祂祈求这杯若能如此就离开祂，其原因是：虽然在天主没有不可能的事，正如基督自己说：\Quote{父，一切为祢都是可能的} \BibleRef{谷 14:36}，但人无法承受对受苦的恐惧，只有通过试探才能证明信德。因此，祂作为人，为众人祈求那杯能过去，但作为出自天主的天主，祂的旨意与圣父有效的旨意一致。祂教导祂所说\Quote{如果可能}的意思，是在祂对伯多禄的话中：\Quote{看，撒殚寻求你们，要筛你们像筛麦子一样；但我已为你祈求，使你的信德不致丧失。} \BibleRef{路 22:31} 主的苦难之杯是给所有人的试探，祂为伯多禄祈求圣父，使他的信德不致丧失：好使当他因软弱否认时，至少不致失去痛悔的忧伤，因为悔改意味着信德得以幸存。\par

39. 主那时忧闷得要死，因为在死亡、地震、日色变暗、幔子裂开、坟墓开启和死人复活面前，门徒们的信德需要被坚固，这信德已被夜间被捕、鞭打、击打、唾面、茨冠、背十字架以及苦难的一切凌辱所动摇，但最甚者是被定罪于可咒骂的十字架。祂知道这一切在祂苦难之后都会结束，所以祂忧闷得要死。祂也知道，除非祂饮了这杯，这杯就不能过去，因为祂说：\Quote{我的父，这杯除非我喝，不能离开我：愿祢的旨意成就。} 也就是说，随着祂苦难的完成，对那杯的恐惧将要过去，这恐惧除非祂喝了杯就无法过去：那恐惧的终结只在祂苦难完成、恐惧消灭之后才来到，因为在祂死后，门徒软弱的绊脚石将被祂大能的荣耀除去。\par

40. 虽然祂用\BibleQuote{尔旨承行}这句话，将宗徒们交托于圣父旨意的决定，关于那杯的冒犯，即祂的苦难，但祂仍第二次、第三次重复了祂的祈祷。此后祂说：\BibleQuote{现在你们睡觉安歇吧。}\BibleRef{玛 26:45}那位曾责备他们睡觉的，如今却吩咐他们睡觉安歇，这并非没有某种隐秘理由的意识。路加被认为已将这命令的意义给了我们。在告诉我们撒殚如何像筛麦子一样筛宗徒，以及主如何被恳求使伯多禄的信德不至于消失之后，他补充说，主恳切祈祷，然后一位天使站在祂旁边安慰祂；当天使站在祂旁边时，祂祈祷更加恳切，以致汗水如血滴从祂身上流下。于是，天使被派遣来看顾宗徒们；当主因他得到安慰，以致不再为他们忧伤时，祂毫无悲伤恐惧地说：\BibleQuote{现在你们睡觉安歇吧。}玛窦和马尔谷对天使和魔鬼的请求保持沉默：但在祂灵魂的忧伤、对睡觉者的责备以及祈求这杯离开之后，紧随其后对睡觉者的命令必有某种好的理由；除非我们假设，那位即将离开他们、并从派给祂的天使那里得到安慰的，有意将他们弃于睡眠之中，不久将被捕并被监禁。\par

41. 我们确实不可忽视这一事实：在许多拉丁文和希腊文抄本中，都没有提到天使的来临或血汗。但当我们暂不作判断（这究竟是缺失之处的遗漏，还是出现之处的增添，因为抄本的不一致使问题不确定）时，不要让异端者以此自慰，认为这里面有祂软弱的证据，即祂需要天使的帮助和安慰。让他们记住，天使的造物主不需要受造物的支持。此外，祂的安慰必须以与祂忧伤相同的方式解释。祂为我们，即因我们而忧伤；祂也必须为我们，即因我们而得到安慰。如果祂因我们而忧伤，祂就因我们而得到安慰。祂安慰的对象与悲伤的对象相同。也不要有人胆敢将汗水归咎于软弱，因为流血汗违反自然。这不是软弱，因为祂的能力颠倒了自然法则。血汗绝不支持软弱的异端，反而确立了反对那虚构幻身异端的祂整个身体的真实性。既然祂的恐惧是因我们，祂的祈祷是为我们，我们不得不得出结论：这一切都是因我们而发生的，祂为我们恐惧，祂为我们祈祷。\par

42. 再者，福音彼此补充所缺：我们从一位学得一些，从另一位学得一些，如此类推，因为所有福音都是同一圣神的宣告。因此，若望——他在福音中特别揭示属灵原因的运作——保存了主为宗徒们的这篇祈祷，而其他人都略过了：即祂如何祈祷：\BibleQuote{圣父啊，求祢因祢的名保全他们……我在他们中间时，我因祢的名保全了他们；祢所赐给我的人，我保全了。}这祈祷不是为祂自己，而是为祂的宗徒；祂也不为自己忧伤，因为祂吩咐他们祈祷以免受诱惑；天使也不是被派到祂那里，因为祂若愿意，可以从天上调动一万二千天使\BibleRef{玛 26:53}；祂也不因死亡恐惧，尽管祂忧愁至死。再者，祂不是祈求这杯从祂身上过去，而是祈求它从祂身上消逝，尽管在它消逝之前祂必须喝它。但进一步说，\Quote{消逝}不仅指\Quote{离开该地方}，而且指\Quote{完全不复存在}；这在福音和书信的语言中有所显示：例如，\BibleQuote{天地要消逝，但我的话决不会消逝}；宗徒也说：\BibleQuote{看，旧的已成过去，都成了新的}\BibleRef{格后 5:17}。又说：\BibleQuote{这世界的局面将要消逝}\BibleRef{格前 7:31}。因此，祂向圣父祈求的那杯，除非被喝掉，否则不能消逝；而当祂祈祷时，祂是为那些祂与他们在时一直保全的人祈祷，如今祂将他们交给圣父去保全。既然祂即将完成死亡的奥迹，祂恳求圣父守护他们。派给祂的天使的出现（如果这解释正确的话）并非意义不明。耶稣表明祂的祈祷已蒙垂听，在祈祷结束时祂吩咐门徒睡觉安歇。这祈祷的效果以及促使祂下令\Quote{睡觉}的确信，被福音作者在苦难过程中注意到，当他说到宗徒们刚从追捕者手中逃脱时：\Quote{这应验了祂所说的话：\BibleQuote{祢所赐给我的人，我连一个也没有丧失}}\BibleRef{若 18:9}。祂亲自实现了祂祈祷的请求，他们都安全；但祂请求圣父现在因祂的名保全那些祂已保全的人。他们被保全了：伯多禄的信德没有消失：它曾畏缩，但悔改立即跟随。\par

43. 将若望福音中主的祈祷、路加福音中魔鬼的请求、玛窦和马尔谷福音中忧愁至死以及对睡觉的抗议、随后命令\BibleQuote{睡觉安歇}结合起来，所有困难都消失了。若望福音中祂将宗徒们托付给圣父的祈祷，解释了祂忧愁的原因以及祈求这杯离开的祈祷。主不是祈求苦难从祂自己身上除去；祂恳求圣父在祂即将来临的苦难中保全门徒。同样，圣路加中对抗撒殚的祈祷，解释了他以何种信心允许了祂刚刚禁止的睡眠。\par

44. 因此，在那超乎人性的本性中，没有人的焦虑和战栗的位置。它超越了尘世肉身的疾苦；那身体并非源于尘世元素，尽管祂作为\Person{人子}的起源是因圣神受孕的奥迹。至高者的能力将其能力赋予那由童贞女因圣神受孕而怀有的身体。这有生命的身体因与灵魂结合而获得自觉的存在，灵魂遍布全身，使之感受到来自外部的疼痛。因此，灵魂被自身属天的信德和望德的幸福热忱所警醒，超越其起源于尘世身体的开端，将身体提升至与自己在思想和精神上的结合，以致它不再感受到那它一直在忍受的苦难。那么，我们何必再多说主的身体、即那从天降下的\Person{人子}的本性呢？有时连尘世的身体也能变得对疼痛和恐惧的自然需要无动于衷。\par

45. 那犹太儿童在巴比伦的火窑中，面对投入其中的燃料所燃起的火焰，难道害怕了吗？那烈火的可怖岂能胜过他们的本性——尽管这本性与我们一样受孕\BibleRef{达 3:23}？当火焰包围他们时，他们感到疼痛了吗？也许你会说，他们没有感到疼痛，因为他们没有被烧着：火焰被剥夺了燃烧的本性。的确，身体害怕被烧、被火焚烧是自然的。但藉着信德的精神，他们尘世的身体（即按自然生育法则起源的身体）既不能被焚烧，也不能被恐吓。因此，在人身上是违反自然秩序的事，因对天主的信德而产生，在天主身上不能被视为自然的，而是圣神在其尘世起源开始时的活动。儿童在火中被捆绑，他们登上燃烧的柴堆时毫无恐惧；他们祈祷时感觉不到火焰；虽然身处火窑，却不能被焚烧。火和他们的身体都失去了各自的本性：火不再燃烧，身体不再被烧。然而在其他方面，火和身体都保留了本性：因为旁观者被烧死，惩罚的执行者自己受了罚。不敬神的异端者，你要基督忍受被钉子刺穿的痛苦，感受到伤口之苦，当钉子穿过祂的手时——请问，为什么那些儿童不怕火焰？为什么他们没有感到疼痛？他们身体中是什么样的本性胜过了火的本性？在信德的热情和真福殉道的光荣中，他们忘记了害怕那可怕的事；难道基督因为害怕十字架而忧愁吗？基督——即使祂是带着我们有罪的起源而受孕，在十字架上仍然是天主，祂将审判世界，并永远为王——难道祂能忘记如此赏报，而因不光彩的恐惧的焦虑而战栗吗？\par

46. 但以理，他的食物不过是先知微薄的份量\BibleRef{达 1:8}，并不害怕狮子坑。宗徒们为了基督的名，在受苦和死亡中欢喜。对保禄来说，他的祭献是正义的华冠。\BibleRef{弟后 4:6}殉道者们歌唱圣诗，将颈项伸给刽子手，吟咏圣咏攀登为他们堆积的燃烧的木柴。信德的意识消除了本性的软弱，转变了身体的感觉，使之感觉不到疼痛，因此身体因灵魂的坚定意志而坚强，除了其热忱的冲动之外什么也感觉不到。灵魂在渴求光荣时轻蔑的痛苦，身体感觉不到，只要灵魂振奋它。因此，在人身上，灵魂为光荣而燃烧的热忱会使他失去对痛苦的意识，忽略伤口，漠视死亡。但是耶稣基督、光荣的主，祂衣边能治愈人，祂的唾沫和言语能创造；那枯手的人因祂的命令伸手而痊愈，生来瞎眼的人不再感到天生的缺陷，被削掉的耳朵变得与另一只一样完好——难道我们敢认为祂被刺穿的身体处于那种痛苦和软弱中，而正是对祂的信德的精神拯救了光荣而真福的殉道者脱离那种痛苦和软弱吗？\par

47. 那么，独生天主在祂身上承受了我们所有软弱的一切攻击，但祂是以自身本性的能力承受了它们，正如祂以自身本性的能力而生，因为在诞生时祂并没有因诞生而失去全能的本性。虽然祂是在人的条件下诞生，但祂并非如此受孕；祂的诞生被人的环境所围绕，但祂的起源超越了这些。因此，祂在身体上承受了如同我们软弱身体的痛苦，却以祂自己身体的能力承受了我们身体的痛苦。先知对我们信仰的这一条作证说：\Quote{\Emph{祂承担了我们的罪，为我们忧伤；我们却以为祂受责罚，被天主击打、苦待了。祂为了我们的过犯而被刺穿，为了我们的罪而被击伤。}}因此，认为祂因受苦而感到疼痛，这是人判断的错误观点。祂承担了我们的罪恶，也就是说，祂取了我们有罪的身体，但祂自己无罪。祂被派遣在罪恶肉身的形状里，确实在祂的肉身中承担了罪恶，但却是\Emph{我们的}罪恶。同样，祂为我们感到忧伤，但不是以我们的感觉；祂被看作人形，有一个能感到疼痛的身体，但祂的本性不能感到疼痛；因为，虽然祂的形状是人，但祂的起源不是人的，而是因圣神的受孕而生。\par

由于上述原因，祂被视为\Quote{被打伤、击打和受苦}。祂取了奴仆的形体；而\Quote{由童贞女所生之人}向我们传达了一位其本性在受苦时感受到痛苦者的概念。但祂虽受伤，却是\Quote{因我们的过犯}。这创伤并非祂自己过犯的创伤；这苦难并非为祂自己的苦难。祂并非为自己而成为人，也非在自己的行动中犯了过犯。宗徒解释这神圣计划的原则时说：\Quote{我们求你们通过基督与天主和好。祂，那不知罪的，为我们成了罪}。为在肉身中定罪于罪，那不知罪的自己成了罪；也就是说，藉着肉身在肉身中定罪于罪，祂为我们成了肉身，却不知肉身；因此祂因我们的过犯而受伤。\par

48. 再者，宗徒对基督的畏惧痛苦一无所知。当他谈论受难的救恩计划时，他将之包含在基督天主性的奥迹中。\BibleQuote{赦免了我们的一切过犯，涂抹了那写在规条上反对我们、与我们为敌的债卷，把它撤去，钉在十字架上；脱去自己的肉身，将执政掌权者公开示众，仗着十字架在自身中凯旋得胜}\BibleRef{哥 2:13}请问，那能力——你认为——会屈服于钉子的创伤、在刺穿之痛下畏缩、转化为能感受痛苦的本性吗？然而，宗徒——他作为基督代言者说话\BibleRef{格后 13:3}——叙述藉着主我们得救的工程，将基督的死描述为\Quote{脱去自己的肉身，大胆地羞辱那些执政掌权者，并在自身中凯旋得胜}。如果祂的受难是本性的必然，而非你救恩的自由恩赐；如果十字架仅仅是伤口的受苦，而非将针对你的死亡判决固定在祂自己身上；如果祂的死是死亡施加的暴力，而非凭借天主的大能脱去肉身；最后，如果祂的死本身只是对掌权者的羞辱、大胆的行为、凯旋之外的任何东西——那么就把软弱归于祂，因为祂在其中屈服于本性和必然、暴力、恐惧和羞辱。但如果恰恰相反，在受难的奥迹中，正如我们所传讲的，那么，请问谁能如此无知而拒绝宗徒所教导的信仰、颠倒所有宗教情感、将本是自由意志的行为、奥迹、大能与勇敢的展示、凯旋扭曲为对本性软弱的可耻指控？这是何等的凯旋：当祂将自己献给那些寻求钉死祂的人，而他们无法忍受祂的临在；当祂被判处死刑，而祂不久将坐在大能的右边；当钉子穿透祂时，祂为迫害祂的人祈祷；当祂饮尽醋时完成了奥迹；当祂被列于罪犯之中，同时赐予乐园；当祂被高举在木头上时，大地震动；当祂挂在十字架上时，太阳和白天都逃遁；祂将自己的身体交出了，却将生命唤回他人的身体；被埋葬如尸体，却复活为天主；作为人为我们承受了一切软弱，作为天主在这一切中凯旋得胜。\par

49. 异端者还说，还有另一个严重而深远的软弱自白，尤其因为这是主亲口所说：\BibleQuote{我的天主，我的天主，祢为什么舍弃了我？}\BibleRef{玛 27:46}他们把这解释为一种痛苦的抱怨，表示祂被遗弃并交付于软弱。但这种不虔敬心灵的暴力解释是何等荒谬！与我们主的话语的整个主旨是多么抵触！祂赶紧走向那要荣耀祂的死亡，之后祂要坐在大能的右边；怀着所有这些蒙福的期望，祂会惧怕死亡，因此抱怨祂的天主将祂出卖给死亡的必然性吗？而死亡正是通往永恒福乐的门径！\par

50. 此外，他们异端的诡辩沿着自己不信的道路前进，甚至将天主圣言完全吸收进人的灵魂，从而否认耶稣基督、人子，与天主子是同一的。因此，要么天主圣言在赋予身体生命、履行灵魂功能时不再是祂自己，要么那出生的人根本不是基督，而是圣言住在他内，如同圣神住在先知们内一样。这些荒谬而乖谬的错误已变得大胆而不虔敬，以至于他们断言耶稣基督在由玛利亚出生之前并非基督。那出生的并非先存的存在，而是从那一刻才开始存在。\par

由此也产生这种错误：天主圣言，如同神圣大能的一部分，以不中断的延续扩展自身，住在那从玛利亚获得存在开端的人内，并赋予他神圣工作的能力：尽管那人是凭自己灵魂的本性生活而行动。\par

51. 藉着这种微妙而有害的学说，他们被引入错误，认为天主圣言变成了身体的灵魂，祂的本性藉着自我谦卑而改变自身，因此圣言不再是天主；或者认为人耶稣，因其本性的贫乏和远离天主，仅由祂自己人类灵魂的生命和活动所推动，而天主圣言，即祂发出的声音的能力，居住在其中。这样，就为各种不敬的臆测打开了道路；其总括是：要么天主圣言被并入灵魂中而不再是天主；要么基督在由玛利亚诞生之前并不存在，因为耶稣基督，一个仅有普通身体和灵魂的人，只从祂的人性诞生才开始存在，并藉着天主圣言伸展到祂身上的外在力量，被提升到那内在运行之能力的水平。然后，当圣言这次伸展之后被撤回时，祂呼喊说：\BibleQuote{我的天主，我的天主，祢为什么舍弃了我？}或者至少，当圣言的神性在祂内再次让位给一个人类灵魂时，祂——此前一直依靠祂圣父的帮助——现在与其分离，被弃于死亡，哀叹自己的孤独并责备那遗弃者。这样，在信仰中处处都产生致命的错误危险，无论是认为抱怨的呼喊表示天主圣言本性上的软弱，还是认为天主圣言并非先存，因为耶稣基督由玛利亚的诞生是祂存在的开始。\par

52. 在这些不敬且无根据的学说中，教会的信仰，受宗徒训诲的启发，承认基督的诞生，但非开始。它知道管理的计划，但不知分裂；它拒绝在耶稣基督内进行分离，以致耶稣是一，基督是另一；也不将人子与天主之子分开，以免人子不被视为也是天主之子。它不将天主之子吸收于人子之中；也不藉着三部分的信仰撕裂基督（祂的由头至尾织成的外衣未被分开），将耶稣基督分为圣言、身体和灵魂；另一方面，它也不将圣言吸收于身体和灵魂中。对祂而言，祂是完全的天主圣言，也是完全的基督之人。在我们的宣认奥迹中，我们只坚持这一点，即信仰基督不外乎耶稣，以及耶稣不外乎基督的教导。\par

53. 我并非不知道，神圣奥迹的伟大使我们软弱的理解力难以理解，以致语言几乎无法表达，理性无法定义，甚至思想无法把握。宗徒知道，地上的本性最难独自领悟天主的行动方式（因为这样我们的判断力将更敏锐地辨别，胜过天主的全能去实现），于是写信给他信仰上的真儿子（他自幼年就接受了圣经），说：\BibleQuote{我往马其顿去的时候，曾劝你仍住在厄弗所，好嘱咐那几个人不要传异教，也不要听从荒谬无凭的话语和无穷的家谱；这等事只生辩论，并不发明天主在信德上的建设。}\BibleRef{弟前 1:3}他吩咐他避免处理词藻的家谱和寓言，这些只会生出无穷的辩论。他说，天主的建设是在信德中；他将人的敬畏限制在对全能者的忠信敬拜中，并不容许我们的软弱在试图看到只令人眩目的东西时过度紧张。如果我们看太阳的光辉，视力就会紧张而减弱；有时当我们过于好奇地凝视光源时，眼睛会失去自然的能力，甚至视力会被摧毁。这样，由于试图看得太多，我们反而什么也看不见。那么，对于天主，正义的太阳，我们还能期望什么呢？那些过于聪明的人，难道不会得到愚昧作为报偿吗？难道迟钝而无脑的呆滞不会取代炽热之光的位置吗？较低的本性不能理解较高本性的原理；天上的思想方式也不能启示给人类的观念，因为在有限意识范围内的东西本身就是有限的。因此，天主的权力超过人类理智的能力。如果有限者竭力伸展到那么远，它就会变得比以前更软弱。它失去原有的确定性，非但看不见天上的事物，反而被它们弄瞎。没有理智能够完全理解天主；祂惩罚那些好奇者的顽固，剥夺他们的能力。如果我们想看太阳，就必须移开它的部分光芒，以便看见它；否则，因期望过多，我们反而达不到可能的程度。同样，我们只能希望在允许的范围内理解天主的计划。我们只能期望祂赐给我们领悟的东西；如果我们试图超越祂允许的限度，它就会完全被收回。在天主内，有我们\Emph{能}感知的东西：如果我们满足于可能，它是可见于所有人的。就像太阳，如果我们满足于看可见的东西，我们就能看到一些；但如果超越可能，我们就失去一切。天主的本性也是如此。有我们\Emph{能}理解的东西，如果我们满足于理解我们所能理解的；但若超出你的能力，你就会连达到力所能及的能力也失去。\par

54. 那另一无时间的诞生的奥迹，我暂且不谈；对此的论述需要比此处更广阔的空间。目前我只谈论降生成人。请问，你们这些窥探天上奥秘的人，基督由童贞女诞生及其本性的奥迹，你们将从何处解释祂是由童贞女受孕并诞生的？按照你们的争论，祂起源的物理原因是什么？祂如何在母腹中形成？祂的身体和人性从何而来？最后，那\Quote{从天降下却仍在天上的人子}\BibleRef{若 3:13} 是什么意思？按身体规律，同一物体不可能既停留又下降：一是向下运动的改变，另一是静止不动。婴孩啼哭，却在天上；孩童成长，却始终是不可测量的天主。凭人的何种理解，我们才能领悟祂升到祂先前所在之处，而祂仍在天上却下降？主说：\Quote{你们若看见人子升到他先前所在之处，又怎样呢？}人子升到祂先前所在之处，感官能领悟吗？人子从天降下，祂却在天上，理性能够应对吗？圣言成了血肉：言语能表达吗？圣言成为血肉，即天主成为人：人在天上，天主从天上而来。祂上升，祂曾下降；但祂下降，却又没有下降。祂始终是祂所是，却又并非始终是祂所是。我们回顾原因，却无法解释方式；我们感知方式，却无法理解原因。然而，即使我们如此理解基督耶稣，我们就会认识祂；如果我们试图进一步理解祂，我们就完全不会认识祂。\par

55. 再者，基督哭泣，祂的眼中因心灵的痛苦而充满泪水，这是何等言语和行动的巨大奥迹。\BibleRef{路 19:41} 祂灵魂中这缺陷从何而来，以致忧伤从祂身体中挤出泪水？何等痛苦的命运，何等无法忍受的痛苦，能感动那位从天降下的人子泪如泉涌？再者，祂里面是什么在哭泣？是天主圣言？还是祂的人性灵魂？因为尽管哭泣是身体的功能，身体只是仆人；眼泪就如同痛苦灵魂的汗水。再者，祂哭泣的原因是什么？祂是否欠\Place{耶路撒冷}眼泪的债？耶路撒冷，那个不虔敬的弑亲者，任何痛苦都无法赔偿她对宗徒和先知们的屠杀，以及对她的主本身的谋害？祂可能为降临人类的灾难和死亡而哭泣；但祂能为那注定灭亡的绝望种族的堕落而悲伤吗？请问，这哭泣的奥迹是什么？祂的灵魂因悲伤而哭泣；那派遣先知的，不就是这灵魂吗？那常常要将小鸡聚集在翅膀荫下的，不就是这灵魂吗？但天主圣言不能悲伤，圣神也不能哭泣；祂的灵魂在身体存在之前也不可能做任何事。然而，我们不能怀疑耶稣基督确实哭泣了。\par

56. 祂为拉匝禄所流的眼泪同样真实。\BibleRef{若 11:35} 这里第一个问题是，在拉匝禄的事上有什么可哭的呢？不是他的死亡，因为那病不是致死的，而是为光荣天主：因为主说：\Quote{这病不至于死，但为光荣天主，为叫天主子因此得荣耀。}那成为天主受光荣之原因的死亡，不能带来悲伤和眼泪。祂在拉匝禄死时不在场，也没有任何流泪的理由。祂清楚地说：\Quote{拉匝禄死了，我为你们欢喜，我不在那里，为使你们相信。}那么，祂的不在场有助于宗徒们的相信，并非祂悲伤的原因：因为凭借神圣的全知，祂从远方就宣布了病人的死亡。因此，我们找不到流泪的必要，然而祂却哭了。我再次问，我们必须将哭泣归于谁？归于天主，还是灵魂，或是身体？身体本身没有眼泪，除非在悲伤灵魂的命令下才流泪。天主更不能哭泣，因为祂要在拉匝禄身上得光荣。说祂的灵魂把拉匝禄从坟墓中召回也没有道理：一个与身体相连的灵魂，能凭其命令的能力，将另一个灵魂从已离开的死尸中叫回来吗？那将要受荣耀的会悲伤吗？那将要使死人复活的会哭泣吗？眼泪不属于那将要赐予生命的，悲伤不属于那将要接受荣耀的。然而，这哭泣和悲伤的，也正是生命的赐予者。\par

57. 如果我们对许多问题处理得简略，不是因为我们无话可说，或不知道已经说过的话；我们的目的是，通过避免过于繁琐的论证过程，使结果尽可能对读者有吸引力。我们知道我们主的言行，却又不认识它们：我们并非对它们无知，但它们无法被理解。事实是真实的，但背后的能力是奥迹。我们凭祂自己的话来证明这一点：\Quote{为此父爱我，因为我舍掉我的性命，为再取回它。没有人夺取我的性命，而是我甘愿舍掉它。我有权舍掉它，也有权再取回它。这是我由我父接受的命令。}祂甘愿舍掉自己的性命，但我问：\Emph{谁}舍掉了它？我们毫不犹豫地宣认：基督是天主圣言；但另一方面，我们知道人子由灵魂和身体组成：比较天使对若瑟说的话：\Quote{起来，带着孩子和他的母亲，往以色列地去；因为那些谋害孩子性命的人已经死了。}\BibleRef{玛 2:20} 这灵魂是谁的？是祂身体的，还是天主的？如果是祂身体的，身体有什么能力舍掉灵魂？因为身体只是因灵魂的活动才获得生命。再者，那离开灵魂就呆滞死亡的身体，怎么能从父领受命令？但另一方面，如果有人假定天主圣言撇下祂的灵魂，为能再取回它，他就必须证明天主圣言死了，即如同死尸一样没有生命和感觉，然后祂再取回祂的灵魂，由它再次获得生命。\par

58. 然而，但凡有理性的人，都不能将灵魂归于天主；尽管在许多地方经上记载，天主的灵魂憎恨安息日和月朔，也喜悦某些事物。但这只是按惯例的说法，应与论及天主有身体、有手、有眼、有指头、有臂、有心时一样理解。主说：\BibleQuote{神魂没有骨肉}\BibleRef{路 24:39}。那\BibleQuote{自有永有、永不改变者}\BibleRef{拉 3:6}，不可能有可触摸之身体的肢体和部分。祂是单纯而蒙福的本性，是单一、完全、包容万有的整体。因此，天主并非像身体那样，靠内在灵魂的行动而得生命，而是祂自己就是自己的生命。\par

59. 那么，祂又如何放下灵魂，或再取回它呢？祂所接受的这条命令是什么意思？天主不能放下灵魂，即死亡，也不能再取回它，即复活。但身体也没有接受再取回它的命令；它不能靠自己做到这一点，因为祂论及自己身体的圣殿时说：\BibleQuote{你们拆毁这座圣殿，三天之内我要把它重建起来}\BibleRef{若 2:19}。所以，是天主重建祂身体的圣殿。是谁放下灵魂以便再取回？身体不能靠自己再取回灵魂：它是由天主复活的。那复活起来的必须曾死过，而活着的不会放下灵魂。那么，天主既没有死也没有被埋葬；然而祂说：\BibleQuote{她把这香液倒在我身上，是为安葬我而做的}\BibleRef{玛 26:12}。既然倒在祂身体上，是为祂的安葬而做的；但\Quote{祂的}与\Quote{祂}并不相同。当我们说\Quote{这是为祂的安葬而做的}和说\Quote{祂的身体被傅油}时，代词用法完全不同；\Quote{祂的身体被埋葬}与\Quote{祂被埋葬}意思也不相同。\par

60. 要领会这神圣的奥迹，我们必须在祂身上看见天主而不忽视人，看见人而不忽视天主。我们不可将耶稣基督分割，因为\OriginalTerm{圣言}{Verbum}成了血肉；然而我们不可称祂被埋葬，尽管我们知道祂复活了自己；不可怀疑祂的复活，尽管我们不敢否认祂曾被埋葬。耶稣基督被埋葬了，因为祂死了；祂死了，甚至在死亡那一刻呼喊：\BibleQuote{我的天主，我的天主，祢为什么舍弃了我？}然而说这话的那一位也说过：\BibleQuote{我实在告诉你，今天你就要同我一起在乐园里}\BibleRef{路 23:43}，那应许给强盗乐园的那一位大声呼喊：\BibleQuote{父啊！我把我的灵魂交托在祢手中；说了这话，便断了气。}\par

61. 你们这些把基督分割为圣言、灵魂和身体的人，或者把整个基督，甚至天主圣言，贬低为我们种族的一个成员，请为我们阐述这个在肉身中显现的\BibleQuote{伟大的虔敬奥迹}\BibleRef{弟前 3:16}。基督交付的是哪一种神魂？是谁将祂的灵魂交托在圣父手中？是谁当天要与那强盗同在乐园里？是谁抱怨被天主舍弃？被舍弃的呼喊表明垂死者软弱；乐园的应许表明永生天主的至高权能。将灵魂交托表示信赖；交付灵魂表示祂因死亡而离去。那么我问，是谁死了？当然是那交付灵魂的。但谁交付了灵魂？当然是那将灵魂交托给祂圣父的那一位。如果交付灵魂的那一位与将灵魂交托的那一位相同，且祂死了，那么是身体交托了它的灵魂，还是天主交托了身体的灵魂？我说\Quote{灵魂}，因为毫无疑问它常与\Quote{神魂}同义，就如仅从这里的措辞可看出的：耶稣在临终时交付了祂的\Quote{神魂}。因此，如果你确信身体交托了灵魂，可朽坏的交托了永生的，可朽的交付了不朽的，那将要被复活的交付了那保持不变的那一位，那么，既然将灵魂交托给圣父的那一位当天也要与强盗同在乐园里，我愿知道，当坟墓接纳祂时，祂是否居住在天上？或者，当祂呼号天主舍弃了祂时，祂是否居住在天上？\par

62. 是同一的、唯一的\Person{耶稣基督}，圣言成了血肉，祂在所有这些话语中表达自己。当祂说自己被抛弃至死时，祂是人的样子；然而，作为人，祂仍在乐园中作为天主统治；并且，尽管在乐园中统治，作为天主子，祂将灵魂交托给圣父；作为人子，祂将交托给圣父的灵魂交付于死亡。那么，我们为什么把奥迹视为耻辱呢？我们看见祂抱怨被留下等死，因为祂是人；我们看见祂死时宣告自己在乐园中统治，因为祂是天主。为什么我们要强调祂为让我们理解祂的死亡而说的话，却隐瞒祂为证明祂的不朽而宣告的话，以支持我们的不敬？那抱怨被舍弃的话和那宣告统治的话同属祂；我们是凭什么异端逻辑分割我们的信仰，否认死去的祂同时是统治的呢？当祂交托灵魂和交付灵魂时，祂岂不是同样为这两者作了见证？但如果那交托灵魂和交付灵魂的是同一位，如果祂在统治时死去、在死时统治，那么天主子和人子的奥迹意味着祂是同一的，祂在死去时统治，在统治时死去。\par

63. 那么，所有不敬神的不信者，请走开吧，你们无法承受这神圣奥迹：你们不知道基督哭泣不是为自己，而是为我们，为证明祂所取的人性实在，祂顺从了人性共通的情感；你们不明白基督死亡不是为自己，而是为我们的生命，为通过不死天主的死亡更新人的生命；你们不能协调被舍弃者的抱怨与统治者的信赖；你们要教导我们，因为祂作为天主统治、作为人抱怨自己死去，所以我们这里有一个死去的人和一个统治的天主。因为死去的不是别的，正是统治的那一位，交托灵魂的不是别的，正是交付灵魂的那一位；被埋葬的那一位复活了；无论是上升还是下降，祂完全是同一的。\par

64. 请聆听宗徒的教导，并在其中看到一个不是由血肉的理解，而是由圣神的恩赐所教导的信德。\Quote{希腊人寻求智慧，他说，犹太人要求标记；但我们宣讲被钉十字架的基督，对犹太人是绊脚石，对外邦人是愚妄；但是对那些蒙召的，不论犹太人或希腊人，基督耶稣是天主的德能，天主。}\BibleRef{格前 1:23} 这里的基督被分裂了吗？被钉的耶稣是一位，而基督——天主的德能和智慧——又是另一位？这对犹太人是绊脚石，对外邦人是愚妄；但对我们，基督耶稣是天主的德能和天主的智慧：然而这智慧不为世界所知，也不为世俗哲学所理解。请听同一位蒙福宗徒宣告它未被理解时所说的话：\Quote{我们讲论的，是天主在奥秘中所隐藏的智慧，这智慧天主在万世之前就预定为使我们得光荣；今世的掌权者没有一个认识它，因为他们若认识它，就不会把光荣之主钉在十字架上了。}\BibleRef{格前 2:7} 难道宗徒不知道天主的这智慧隐藏在奥秘中，不能为今世的掌权者所知吗？他是否将基督分割为一位光荣之主和一位被钉的耶稣？不，他反而用这些话驳斥这最愚蠢和不敬的思想：\Quote{因为我曾决定，在你们中间不知道别的，只知道耶稣基督，并他被钉十字架。}\par

65. 宗徒不知道别的，他也决定不知道别的：我们这些智力更弱、信德更弱的人，分裂、分割并双重耶稣基督，自命为未知事物的审判者，亵渎隐藏的奥秘。对我们来说，被钉十字架的基督是一位，作为天主智慧的基督又是另一位；被埋葬的基督不同于从天降下的基督；人子不同时也是天主子。我们教导我们所不了解的，我们试图反驳我们所不能掌握的。我们人类改进天主的启示，不满足于同宗徒一起说：\Quote{谁能控告天主的选民呢？天主是成义者，谁又能定罪呢？是基督耶稣死了，甚至复活了，如今在天主的右边，也为我们转求。}\BibleRef{罗 8:33} 难道为我们转求的，不同于在天主右边的那一位吗？在天主右边的那一位，难道不是同一位复活者吗？复活者与死者不同吗？死者与定罪者不同吗？最后，定罪者岂不也是使我们成义的天主吗？你若是能够，就把基督——我们的控告者——与天主——我们的辩护者——区别开来；把死去的基督与定罪的基督区别开来；把坐在天主右边为我们祈祷的基督与死去的基督区别开来。所以，无论是死了、埋葬了、下降阴府或升上天堂，一切都是同一位基督：正如宗徒所说：\Quote{那么\InnerQuote{他升上去了}是什么意思呢？岂不是说他也曾下降到地下吗？那下降的，也正是那上升远超越诸天的，为要充满万有。}\BibleRef{弗 4:9} 我们还要将喋喋不休的无知和亵渎推到多远，自称要解释隐藏在神秘中的东西呢？\Emph{那下降的，也正是那上升的。}我们还能怀疑基督耶稣这人从死者中复活了，升到诸天之上，并坐在天主的右边吗？我们不能说他的身体（那躺在坟墓中的）下降到了阴府。如果那下降的与那上升的是同一的；如果他的身体并没有下到阴府，却真正从死者中复活了，升上了天堂，那么除了相信这隐藏的奥秘（它向世界和今世的掌权者隐藏），并承认无论上升或下降，他只是一，为我们而言是一位耶稣基督，天主子与人子，天主圣言与成了肉身的人，他受难、死亡、被埋葬、复活、被接升天，并坐在天主的右边：他按照天主的计划和本性，在一个单独的自体内，既具有天主的形体，又具有仆人的形体，人性和神性不可分离或分割，还有什么呢？\par

66. 所以宗徒将我们无知和随意的想法塑造成符合真理，论到这信德的奥秘说：\Quote{他固然由于软弱而被钉十字架，却由于天主的德能而活着。}\BibleRef{格后 13:4} 他宣讲人子和天主子，按照神圣计划是人，按照永恒本性是天主，他说，那因软弱而被钉十字架的，就是那因天主的德能而活着的。他的软弱来自于\Emph{仆人的形体}，他的本性则因\Emph{天主的形体}而存留。他取了仆人的形体，虽然他原具有天主的形体：因此，对于他既受苦又活着的奥秘，无可怀疑。在他身上同时存在受苦的软弱与赐予生命的天主的德能：所以那受苦又活着的，不能多于一位，也不能不是他自己。\par

67. 独生天主确实受尽了人所能受的一切苦；但让我们用宗徒的话和信德来表达。他说：\Quote{因为我首先传授给你们的就是我所领受的，就是基督照圣经记载的为我们的罪死了，被埋葬了，第三天照圣经记载的复活了。}\BibleRef{格前 15:3}这不是他个人的无根据的陈述，可能引致错误，而是警告我们，要宣认基督确实死了并复活了，而非名义上的，因为事实得到圣经权威的充分保证；而且我们必须按照圣经所宣告的确切含义来理解祂的死。他顾念那些软弱而敏感的信徒的困惑和顾虑，在宣告死亡和复活时，加上了这些庄重的结束语：\Emph{照圣经记载}。他不愿我们变得更软弱，被各种虚妄学说的风吹动，或被空洞的诡辩和虚假的疑虑所困扰；他召唤信德在遭遇沉船之前回到虔敬的港湾，相信并宣认耶稣基督——人子和天主子——的死亡和复活，\Emph{照圣经记载}，这是对对抗者攻击的敬畏护盾，要照圣经所记载的来理解耶稣基督的死亡和复活。信德中没有危险：对天主隐藏奥秘的敬畏宣认永远是安全的。基督由童贞玛利亚出生，却由圣神怀孕，\Emph{照圣经记载}。基督哭泣，但\Emph{照圣经记载}：使他哭泣的原因也是喜乐的缘由。基督饥饿，但\Emph{照圣经记载}，当祂没有食物时，祂用祂作为天主的权能对付那棵不结果子的树。基督受苦，但\Emph{照圣经记载}，祂即将坐在大能的右边。祂抱怨被抛弃至死，但\Emph{照圣经记载}，同一时刻，祂在乐园的国中接待了那宣认祂的盗贼。祂死了，但\Emph{照圣经记载}，祂复活了，并坐在天主的右边。相信这奥秘就有生命：这宣认抵抗一切攻击。\par

68. 宗徒小心不留下任何怀疑的余地：我们不能说：\Quote{基督出生、受苦、死了并被埋葬了，又复活了：但是怎样，藉着什么能力，藉着祂自己的什么部分分开？谁哭了？谁喜乐了？谁抱怨了？谁降下了？谁升上了？}祂将信德的功劳完全建立在毫无疑问的敬畏宣认上。他说：\Quote{\Emph{义德}本于信德却说：你心里不要说：谁升到天上去？那就是要带基督下来；谁下到深渊里去？那就是要带基督从死人中上来？但经文怎么说？你的话离你很近，就在你口里，在你心里；这就是我们所传信德的话：因为你若口里宣认耶稣为主，心里相信天主使他从死人中复活，你必得救。}\BibleRef{罗 10:6}信德成就义人：正如经上记载：\BibleQuote{亚巴郎相信了天主，这就算为他的义德。}亚巴郎在得到应许——外邦人的产业和如同海沙星辰一样无数的永久后裔时，是否曾质疑天主的话？对于虔敬的信德，它完全信赖天主全能，人类软弱的限制不是障碍。轻视一切自身软弱属地的事，它相信天主的应许，即使超过了人本性的可能。它知道支配人的法律不是天主能力的阻碍，祂在履行上慷慨如同在应许上恩惠。没有什么比信德更正义。因为正如在人类行为中，我们赞赏公平和自制，同样在天主身上，对于人来说，有什么比将全能归于祂更正义呢？祂是谁？人感知到祂的权能没有界限。\par

69. 宗徒于是寻找我们内那本于信德的义德，就斩断了不信的怀疑和无神的不信的根。他禁止我们将焦虑思想的挂虑接入心中，并指向先知话的权威：\Emph{你心里不要说：谁升到天上去？}然后他补充先知话的思想：\Emph{那就是要带基督下来。}人理智的知觉不能达到对天主的认识；但虔敬的信德也不能怀疑天主的作为。基督不需要人的帮助，好使谁升到天上将祂从蒙福的家乡带到祂地上的身体。没有外力驱使祂降到地上。我们必须相信祂来了，正如祂确实来了：宣认耶稣基督不是被带下来，而是降下来，这是真宗教。祂来临的时间和方式的奥秘，只属于祂自己。我们不能因为祂在不久前来了，就以为祂必定是被带下来的，也不能认为祂在时间中的来临依赖于另一个带祂下来的人。\par

也不论另一个方向，宗徒没有给不信留余地。他立即引用先知的话：\Emph{或谁下到深渊里？}并随即加上解释：\Emph{那就是要带基督从死人中上来。}祂能自由地返回天上，祂也能自由地降到地上。于是所有犹豫和怀疑都消除了。信德揭示全能者所计划的事：历史叙述效果，全能的天主是原因。\par

70. 但人们要求我们有不摇动的确信。宗徒阐述圣经的整个奥秘，继续说道：\BibleQuote{你的话离你很近，就在你口里，就在你心里}\BibleRef{申 30:14}。我们宣认的话不可迟缓或故意含糊：心与唇之间不可有间隔，免得本当是真实敬畏的宣认变成不信的托辞。话必须离我们近，在我们内；心与唇之间没有迟延；一种既用言语也用确信的信德。心与唇必须和谐，并在思想和言语中揭示一种不摇摆的宗教。此处，如前所述，宗徒又加上了先知话的解释：\BibleQuote{这就是我们所传的信德之言；因为如果你口里宣认耶稣为主，心里相信天主使祂从死者中复活了，你必会得救}\BibleRef{罗 10:8-9}。虔敬在于拒绝疑惑，正义在于相信，救恩在于宣认。不要玩弄模棱两可的话，不要被激起虚妄的空谈，不要以任何方式辩论天主的权能，或限制祂的大能，不要再三探究不可探测的奥秘的原因：反而宣认耶稣是主，并相信天主使祂从死者中复活了；这就是救恩。贬低基督的本性和身份是何等愚蠢，因为唯独认识祂是主，才是救恩。再者，为祂的复活而争辩是何等人性虚妄的错误，因为相信天主使祂复活，于永生已经足够。所以，信德在于纯朴，正义在于信德，真正的虔诚在于宣认。因为天主并非通过艰苦的探询召叫我们走向幸福生活。祂不用各种修辞技巧试探我们。通往永生的道路是平坦容易的：相信耶稣被天主从死者中复活，并宣认祂是主。因此，没有人可以把因我们的无知而说出的这番话扭曲为不虔敬的借口。我们必须证明耶稣基督死了，好使我们能在祂内生活。\par

71. 那么，如果祂说：\BibleQuote{我的天主，我的天主，祢为什么舍弃了我？}\BibleRef{谷 15:34}，以及\BibleQuote{父啊！我把我的灵魂交在祢手里}\BibleRef{路 23:46}，为使我们确信祂确实死了，这岂不是在祂对我们信德的关怀中，驱散我们的疑惑，而非承认祂的软弱？当祂要复活拉匝禄时，祂向父祈祷：但祂有何祈祷的需要呢？祂说：\BibleQuote{父啊！我感谢祢，因为祢俯听了我。我本来知道祢常常俯听我，但我是为了四周站着的群众说的，好叫他们相信是祢派遣了我}\BibleRef{若 11:41}\Emph{?} 祂是为我们祈祷，好使我们认识祂是圣子；祈祷的话对祂毫无用处，但祂是为增进我们的信德而说。祂不需要帮助，我们需要教导。祂又祈祷受光荣，随即从天上传来天主圣父光荣祂的声音：但当他们对那声音感到惊奇时，祂说：\Quote{这声音不是为我而来，而是为你们来的。}圣父为我们而被恳求，为我们说话：愿这一切引导我们相信并宣认！荣耀者的回应不是赐予对光荣的祈祷，而是赐予旁观者的无知：那么，祂在苦难中最大喜乐时，我们岂不应当将受苦的哀叹视为建立我们信德之举？基督为迫害祂的人祈祷，因为他们不知道自己在做什么。祂从十字架上应许了乐园，因为祂是天主君王。祂在十字架上喜乐，因为祂喝醋时一切都完成了，因为祂在死前已满足了全部预言。祂为我们而生，为我们受苦，为我们而死，为我们而复活。为我们的得救，唯独宣认从死者中复活的天主圣子，这是必要的：那么我们为何要在这不虔敬的不信状态中死去呢？如果基督——永远确知祂的天主性——向我们显明了祂的死，祂自己漠视死亡，却藉死亡来确保祂所取的是真实的人性：那么我们为何要用同一圣子为我们成为人子而死的宣认，作为否认祂天主性的主要武器呢？\par

\SourceRecord{Translated by E.W. Watson and L. Pullan.；Nicene and Post-Nicene Fathers, Second Series；Vol. 9.；Edited by Philip Schaff and Henry Wace.；Buffalo, NY: Christian Literature Publishing Co.,；1899.；<http://www.newadvent.org/fathers/330210.htm>.}

% structure: p0001 source=h1 target=section reason=page-title

\section{论天主圣三（卷十一）}

卷十一。信德是唯一的，正如天主是唯一的；但异端者的信德却有许多（§§1,2）。希拉流现已证明关于基督的真理，使之无可否认；即使在当代，奇迹也为之作证（§3）。亚略人宣讲另一位受造的基督；他们使基督成为受造物，从而宣告另一位神，不是父而是造物主（§4）。圣子作为肖像，与圣父同性同体；若祂较低，则祂不是肖像（§5）。但亚略人藉着祂屈尊就我们境遇的论证来解释合一（§6），并且甚至在复活后，他们还主张祂承认自己的不平等。他们据格林多前书第十五章二十四至二十八节如此辩论\BibleRef{格前 15:24-28}，本书其余部分致力于这段经文（§§7,8）。但我们必须承认真理的奥秘，接受它的两面，两者都清楚启示，尽管我们不能调和它们（§9）。他们只看一面；希拉流在回应中再次证明基督既从天主而生，祂自己也是天主（§§10-12）。但在祂降生成人时，祂开始以那永恒为祂父的天主为主（§13），当祂说祂要上升到\Emph{祂的天主}时，祂说话就如祂称我们为弟兄时一样（§§14,15）。因此，天主是基督之父有两种意义；那为圣子基督之父的，是为仆人基督的主（§§16,17）。正是作为仆人，圣咏作者说：\Emph{你的天主以油傅了你}，这些话若对作为圣子的祂说就没有意义（§§18,19）。正是藉着这较低的本性，祂成为我们的弟兄，天主成为我们的父，祂成为中保（§20）。但有人辩称，祂最终的臣服和将王国交给父，是不平等的证明。这段经文必须整体理解（§§21,22）。有些真理是人难以把握的，若我们误解了，不应羞于认错（§§23,24）。在这段经文中，亚略人通过改变预言的顺序来支持他们的论点（§§25-27）。\Emph{终结}意为最终恒久的状态，而非结束（§28），并且虽然祂交出王国，祂并不停止为王（§29）。接下来讨论祂对父的臣服和万有对祂的臣服；一方面这是比喻性语言，另一方面它证明父与子的合一。圣子的臣服意味着祂分享父的光荣（§§30-36）。显圣容显示基督身体的光荣，这光荣信众也将分享（§§37,38）。义人是祂的王国，祂作为人将把这王国交给父，因为\Emph{因着人也有了死人的复活}（§39）。最后天主将成为万有之中的万有，在基督内的人性不被抛弃，而是受光荣并被接纳入天主性（§40）。基督与圣保禄都已预言了这事（§§41,42）。亚略人对这真理的曲解纯粹是愚昧（§43）。任何合理的解释必须假定天主的威严不能被增加，正如它不能被衡量（§§44,45），而我们的理智是有限的，与天主的无限相对照。天主在成为\Emph{万有之中的万有}时不能变得比先前更大。父和子，之后和之前一样，各自必须如祂所是（§§46-48）。一切都是为我们做的，好使我们得光荣，被塑造成为祂——父的肖像——的相似（§49）。\par

1. 宗徒在写给厄弗所人的书信中，以多方面回顾福音的圆满完美奥秘，在认识天主的其他教导中混杂了以下内容：\BibleQuote{你们也蒙召于同一望德中; 一主、一信德、一圣洗，一位天主，众人的父，且超越一切，并在我们众人之内。} \BibleRef{弗 4:4} 祂没有将我们留在模糊而误导的无限教义之路上，也没有将我们遗弃于想象的变化无常之中，而是藉着设立约束的屏障，限制了理智和欲望的无节制放纵。祂不给我们机会智慧超过他所宣讲的，而是以精确严谨的语言界定了永恒固定的信德，使信仰不稳定的借口无处可寻。祂宣告一个信德，如同祂宣讲一个主，祂宣布一个圣洗，如同祂宣告一个主的一个信德，为要如同有一个主的一个信德，同样有一个信德在唯一的主内的唯一圣洗。既然圣洗与信德的整个奥秘不仅在于一个主，也在于一个天主，祂藉着宣认一个天主完成我们望德的圆满。唯一的圣洗和唯一的信德属于唯一的天主，如同属于唯一的主。主与天主各自为一，不是藉着位格的合一，而是藉着特性的区别：因为，一方面，各自为一的特性——父在其父职中，子在其子位中；另一方面，个性的那一特性——各自所拥有的——构成各自与另一位合一的奥秘。因此，唯一的主基督不能夺去天主父的主权，唯一的天主父也不能否认唯一的主基督的天主性。如果因为天主是唯一的，基督不是本性为天主，那么我们也不能承认唯一的天主是主，因为有一位主基督：也就是说，假设由他们的\Quote{唯一}所意指的不是奥秘，而是排他性的合一。所以，有一个圣洗和一个信德属于一个主，如同属于一个天主。\par

2. 然而，若不能坚定而真诚地宣认唯一的主和唯一的天主父，它如何还能是唯一的信德？而并非唯一信德的信德，如何能宣认唯一的主和唯一的天主父？再者，当宣讲者彼此如此分歧时，信德如何能是唯一的？一人前来教导说，主耶稣基督在我们本性的软弱中，当钉子刺穿祂的手时，祂痛苦呻吟，祂失去了自己能力与本性的德能，并因面临的死亡而颤抖退缩。另一人甚至否认生成的基要教义，称祂为一个受造物。另一人将称祂为天主，却并不认为祂是天主，理由是人们可以谈论多位天主，而祂，我们所承认为天主的，必须意识到分享天主性。同样，主基督如何能是唯一的，当有人说祂作为天主没有感觉，另一些人使祂软弱胆怯；对有些人祂是天主之名，对另一些人祂是天主之性；对有些人祂是由生成而来的子，对另一些人祂是由称谓而来的子？若果真如此，天主父如何在信德中是唯一的，因为对有些人祂是因权威而为父，对另一些人祂是因生成为父——即天主是万有之父的意义？\par

然而，谁能否认，凡不是唯一的信德，就根本不是信德？因为在唯一的信德中，主基督是唯一的，天主父也是唯一的。但唯一的主耶稣基督，除非祂是子，除非祂是天主，除非祂是不变的，除非祂的子性与天主性永远存在于祂内，否则祂在宣认的真理中并非唯一，而只是名义上唯一。凡宣讲基督不同于祂所是的——即不同于子和天主——就是宣讲另一位基督。他也并不在唯一洗礼的唯一信德中，因为在宗徒的教导中，唯一的信德正是那唯一洗礼的信德，在其中唯一的主就是基督，天主子，祂也是天主。\par

3. 然而不能否认基督就是基督。祂不可能不为人类所认识。先知书已为祂盖上了印：日渐圆满的时期为祂作证；借着德能的工作，宗徒和殉道者的坟墓宣告祂；祂圣名的德能启示祂；不洁的神灵承认祂，恶鬼在痛苦中嚎叫，呼喊着祂的名。在一切之中，我们看到了祂权能的安排。但我们的信德必须宣讲祂所是的，即唯一的主，不在名义上，而在宣认中，在唯一洗礼的唯一信德中：因为我们对唯一主基督的信德，取决于我们对唯一天主父的宣认。\par

4. 但这些宣讲一位新基督的教师，否认祂的一切所有，他们宣讲另一位主基督，以及另一位天主父。一位不是生者，而是创造者；另一位不是受生，而是受造。因此基督并不是真正的天主，因为祂并非由出生而为天主；信德无法在天主内认出父，因为没有生成为构成父。他们确实尊崇天主父，正如祂的权利与应得的，当他们将一种不可接近、不可见、不可侵犯、不可言说、无限的，赋有全知和全能，充满爱，运行并渗透万有，内在又超越，在一切有知觉的存在中有知觉的本性归于祂时。但当他们进而将单独善、单独全能、单独不朽的独特光荣归于祂时，谁能不感到这虔诚的赞美意在将主耶稣基督排除于真福之外——这真福因\Quote{单独}这保留而局限于天主的光荣？难道这没有使基督留在罪恶、软弱和死亡中，而父却独享完美吗？难道这没有否认基督从天主父而来的自然起源，因为害怕祂会因一次出生——这出生赋予受生者与生者相同本性的德能，而那是唯独天主父自然享有的真福——而承受了什么吗？\par

5. 他们在福音和宗徒的教导上无知，却颂扬天主父的光荣，然而并非出于虔诚信徒的真诚，而是借着不虔敬的诡计，从中攫取他们邪恶异端的论据。他们说，没有任何事物可与祂的本性相比；因此独生的天主被排除在比较之外，因为祂拥有较低和较软弱的本性。他们这样说天主，即活天主的活肖像，祂真福本性的完美形象，祂未生实体的独生子；除非祂拥有父真福的完美光荣，并在精确中再现祂全部本性的相似，否则祂就不是天主的真实肖像。但如果独生的天主是未生天主的肖像，那完美和至高本性的真理就居于祂内，并使得祂成为那真正天主的肖像。父是全能吗？软弱的子就不是全能的肖像。父是善的吗？这位天主性较低的子，并没有在祂有罪的本性中反映善的肖像。父是无形的吗？这位灵性局限于身体界限的子，就不在无形者的形象中。父是不可言说的吗？这位语言可定义、口舌可描述本性的子，就不是那不可言说者的肖像。父是真正的天主吗？子只拥有虚构的天主性，虚假不能是真实的肖像。然而宗徒并没有将肖像的一部分或形象的一部分归于基督，而是毫无保留地称祂为不可见的天主的肖像和天主的形象。他如何能更明确地宣告天主子是天主性呢？岂不是说基督甚至就其不可见性而言是不可见天主的肖像：因为如果基督的实体是可看见的，祂如何能是那不可见本性的肖像？\par

6. 但正如我们在前几卷所指出的，他们利用所取人性的安排作为贬低祂天主性的借口，并把救恩的奥迹扭曲为亵渎的机会。如果他们坚守宗徒的信德，就不会忘记那具有天主位格者取了仆人的形体，也不会利用仆人的形体来羞辱天主的位格（因为天主的位格包含天主的圆满性），反而会合理而虔敬地注意场合和奥迹的区别，既不羞辱天主性，也不因基督的降生成人而误入歧途。然而现在，我们深信已竭尽所能证明了这一切，并指出了在所取身体诞生之下的天主性本体的德能，再无怀疑的余地。祂既是人又是独生天主，以天主的德能行一切事，并藉着真实的人性在天主的德能中成就一切。祂作为天主所生，拥有全能的天主性本体；作为童贞女所生，拥有完美而完整的人性。祂虽有真实的身体，却存于天主性本体中；祂虽存于天主性本体中，却停留在真实的身体内。\par

7. 在我们的答复中，我们跟随祂直到祂光荣的死亡，逐一审视了他们不敬神学的主张，并根据福音和宗徒的教导驳斥了它们。但即使在祂光荣的复活之后，仍有某些事情他们竟敢解释为较低本性的软弱证据，对此我们现在必须答复。让我们再次采用我们惯用的方法，从话语本身中引出其真实含义，以便我们正好在他们认为能推翻真理的地方发现真理。因为主是凭着简单的言语来教导我们信仰，祂的话语不需要从外在和无关的言论中获得支持或注释。\par

8. 在其他的罪过中，异端者常引用主的这句话作为论据：\BibleQuote{\Emph{我升到我的父和你们的父那里去，升到我的天主和你们的天主那里去}}\BibleRef{若 20:17}。祂的父也是他们的父，祂的天主也成了他们的天主；因此祂并不具有天主的性体，因为祂如自己一样把天主称为他人的父，而当祂与他人分享使其成为子与天主的本性与起源时，祂那独一无二的儿子身份就停止了。但他们还要加上宗徒的话：\BibleQuote{\Emph{但当祂说\InnerQuote{\Emph{万有都屈服了}}时，显然是把那使万有屈服于祂的除外。万有都屈服于祂以后，那时，子也要屈服于那使万有屈服于祂的父，好叫天主成为万物中的万有}}\BibleRef{格前 15:27}。由于他们认为这种屈服是软弱的证据，他们便可以剥夺祂父性本性的德能，因为祂本性的软弱使祂屈服于更强本性的统治。在此之后，让他们采取他们最强的立场和不可攻破的防御，即天主圣子降生的真理要被摧毁：即，如果祂是屈服的，祂就不是天主；如果祂的天主和父也是我们的，祂就与受造物分享一切，因此祂自己也是一个受造物：由天主所创造而非所生，因为受造物从无中获得其本体，但所生者拥有创造者的本性。\par

谎言总是可耻的，因为说谎者松开羞耻的缰绳，敢于反驳真理，或者有时他躲在某种借口的帷幔之后，好显得他以谦逊来捍卫那意图无耻之事。但在现在的情况下，当他们亵渎地利用圣经来贬低我们主的尊严时，没有脸红或虚假借口的余地；因为有时连对无知的宽恕也被拒绝，而故意的曲解则在其赤裸的亵渎中被揭露。让我们暂时推迟对这一福音段落的解释，先问他们是否忘记了宗徒的宣讲，他说：\Quote{毫无疑问，虔敬的奥秘是伟大的，祂在肉身内显现，在圣神内成义，被天使看见，在万邦中宣讲，在世上被相信，在光荣中被接升。} \BibleRef{弟前 3:16} 谁是如此迟钝，不能理解虔敬的奥秘不过是主所取了的血肉之身的\OriginalTerm{救恩计划}{Dispensation}？那么从一开始，不同意这一宣认的人就不在天主的信仰中。因为宗徒毫不怀疑，所有人都必须承认，我们救恩的隐藏秘密不是天主的羞辱，而是伟大虔敬的奥秘，并且是一个不再隐藏在我们眼前的奥秘，而是在肉身中显现的；不再因肉身的本性而软弱，而是在圣神内成义的。因此，借着圣神的成义，肉身软弱的观念从我们的信仰中被移除；借着肉身的显现，那秘密的被揭示，而在那秘密的未知原因中，包含了唯一的宣认，即伟大虔敬的奥秘的宣认。这就是宗徒按适当顺序陈述的整个信仰体系。从虔敬产生奥秘，从奥秘产生在肉身中的显现，从在肉身中的显现产生在圣神内的成义：因为那在肉身中显现的虔敬的奥秘，要真正成为一个奥秘，是通过圣神的成义在肉身中显现的。再者，我们不可忘记，在圣神内的这种成义是怎样的，这种在肉身中的显现：因为那在肉身中显现、在圣神内成义、被天使看见、在万邦中宣讲、在世上被相信的奥秘，同样的奥秘是在光荣中被接升的。因此，它完全是伟大虔敬的奥秘，当它在肉身中显现，当它在圣神内成义，当它被天使看见，当它在万邦中宣讲，当它在世上被相信，并且当它在光荣中被接升。宣讲紧随看见，相信紧随宣讲，而所有这一切的圆满是在光荣中被接升：因为被接升到光荣中是伟大虔敬的奥秘，并且借着对救恩计划的信心，我们被准备好被接升，并符合主的荣耀。因此，肉身的被取也是伟大虔敬的奥秘，因为借着肉身的被取，奥秘在肉身中被显现。但我们必须相信，在肉身中的显现也正是这伟大虔敬的奥秘，因为祂在肉身中的显现是祂在圣神内的成义，以及祂被接升到光荣中。那么，我们的信仰还留下什么空间让任何人认为，这虔敬的救恩计划的秘密是神性的削弱呢，因为借着光荣的被接升，伟大虔敬的奥秘要被宣认？那曾是\Quote{软弱}的，现在成了\Quote{奥秘}；那曾是\Quote{必然}的，成了\Quote{虔敬}。现在，让我们转向福音作者话语的含义，以免我们救恩和荣耀的秘密转化为亵渎的藉口。\par

你赋予主的话以不可抗拒的权威，异端者，祂说：\Quote{我升到我的父和你们的父那里，到我的天主和你们的天主那里。} \BibleRef{若 20:17} 你说，同样的父是祂的父和我们的父，同样的天主是祂的天主和我们的天主。因此，祂分享了我们的软弱，因为在拥有同样父的时候，我们作为儿子并不逊色，而在服务同样天主的时候，我们作为仆人平等。既然我们是被造起源和仆人的本性，却与祂共有父和天主，那么祂就与我们的本性共有受造物和仆人的身份。这就是这愚昧而亵渎的教导。它还引出先知的话：\Quote{你的天主以喜油傅你，胜过你的同伴}，以证明基督没有分享那属于天主的荣耀本性，因为傅油的天主比祂更被尊崇为祂的天主。\par

我们不认识天主基督，除非我们认识受生的天主。但生而为天主，就是属于天主的本性，因为\OriginalTerm{受生者}{Begotten}这一名称确实标示祂起源的方式，但并不使祂与\OriginalTerm{生者}{Begetter}种类不同。如果是这样，受生者确实要对祂的根源存有感激，但并不失去那根源的本性，因为天主的诞生只能来自一个起源，并且只有一个本性。如果它的起源不是来自天主，那就不是诞生；如果它除了诞生以外是别的什么，那么基督就不是天主。但祂是出自天主的天主，因此天主圣父对天主圣子而言，是祂诞生的天主和祂本性的父，因为天主的诞生是出自天主，并且是在天主的特定本性之内。\par

12. 请注意，在整个言论中，主如何谨慎地调和祂对债务的虔诚承认，使得既不以降生之承认来贬损祂的天主性，也不以祂的恭敬服从来侵犯祂的主权本性。祂作为受生者并未省略应归于祂的敬意，祂将自身的存有归功于祂的本源，但祂也以自信的态度表明了对那本性的意识，这本性因祂作为天主所受生的起源而属于祂。例如，祂说：\BibleQuote{看见了我，就是看见了父}\BibleRef{若 14:9}，以及\BibleQuote{我所说的话，不是由我自己说的}。祂不是由自己说的：所以祂从祂的本源领受了祂所说的话。但如果有人看见了祂，他们也看见了父：他们通过这证据意识到，天主在祂内，一个与天主同类本性是从天主受生而存有作为天主的。又如祂说：\BibleQuote{父所赐给我的，比一切都大}，以及\BibleQuote{我与父原是一体}。说父赐给，是承认祂的起源：但祂与父的一体性是源自那起源的本性特征。再如：\BibleQuote{父将一切审判交给了子，为叫众人尊敬子如同尊敬父}。祂承认审判被赐给祂，因此祂没有遮掩祂的受生：但祂要求与父同等的尊敬，因此祂没有放弃祂的本性。又如：\BibleQuote{我在父内，父也在我内}，以及\BibleQuote{父比我大}。一个在另一个内：要认出天主的、由天主所生的天主性：父比祂大：要觉察祂对父权柄的承认。同样，祂说：\BibleQuote{子不能由自己做什么，除非看见父做什么：因为父所做的事，子也照样做}。祂不能由自己做什么：即按祂的受生，父推动祂的行动：但父所做的事，子也照样做；即祂存有丝毫不小于天主，因父全能的本性居于祂内，祂能做天主父所做的一切。一切言论都与祂与父在圣神内的合一以及那本性（祂因受生而拥有）的特征相符。那带来祂存在的受生，使祂成为天主，祂的存在揭示了那天主本性的意识。天主圣子承认天主祂的父，因为祂是由父所生；但也因为祂受生，祂继承了天主的全部本性。\par

13. 所以，伟大而虔诚奥迹的计划使祂——祂已经是天主圣子的父——也在祂所取受造的形像中成为祂的主，因为那在天主的形像中的祂，也以奴仆的形像出现。然而祂并非奴仆，因为按圣神而言，祂是天主，天主之子。每人也会同意，没有主就没有奴仆。天主在独生天主的受生中确实是父，但只有在另一位是奴仆的情况下，我们才能称祂为主和父。子起初并非按本性为奴仆，但后来开始按本性成为祂原来不是的某物。因此，父为子的主与子为奴仆的依据相同。因祂本性的计划，子有了一个主，当祂藉取人性而使自己成为奴仆时。\par

14. 因此，在奴仆的形像中，那原先在天主形像中的耶稣基督，作为人说道：\BibleQuote{我升到我的父和你们的父那里，升到我的天主和你们的天主那里}。祂是作为奴仆向奴仆说话：我们怎能将这话从作为奴仆的基督剥离，转而归于那毫无奴仆成分的本性呢？因为那存留在天主形像中的祂，取了奴仆的形像，这形像是祂与奴仆共融作为奴仆的必要条件。正是在这个意义上，天主是祂的父也是人的父，祂的天主也是奴仆的天主。耶稣基督是作为人、在奴仆的形像中对人和奴仆说话；那么，在祂的人性方面，父是祂的父如同是我们的父，在祂奴仆的本性中，天主是祂的天主如同众人的天主，这观念有何困难呢？\par

15. 于是，祂以这些开场白说：\BibleQuote{你往我弟兄那里去，对他们说：我升到我的父和你们的父那里，升到我的天主和你们的天主那里}。我问：这些弟兄应理解为就天主的形像而言，还是就奴仆的形像而言？我们的血肉与祂内住的天主性丰满有亲属关系，以致我们在祂的天主性方面算为祂的弟兄吗？不，因为预言之神清楚认识到我们在哪方面是独生天主的弟兄。正是作为\BibleQuote{虫，不是人}，祂说：\BibleQuote{我要向我的弟兄宣扬你的名}。祂作为虫，生于不经过正常受孕的过程，或从地底深处以活着状态来到世界，这里祂说话是表明祂取了血肉，并从阴府中活生生地提起了它。整篇圣咏，祂藉预言之神预言祂受难的奥迹：因此，就祂受苦的计划而言，祂有弟兄。宗徒也认识到这弟兄关系的奥迹，因为他称祂不仅为从死者中首生者\BibleRef{哥 1:18}，也为许多弟兄中的首生者\BibleRef{罗 8:29}。基督是许多弟兄中的首生者，与祂是从死者中首生者意义相同：正如死亡的奥迹关乎祂的身体，弟兄关系的奥迹也关乎祂的血肉。因此，天主按祂的血肉有弟兄，因为\BibleQuote{圣言成了血肉，居住在我们中间}\BibleRef{若 1:14}：但独生子，作为独生者是独一无二的，没有弟兄。\par

16. 然而，祂取了血肉，便完全获得了我们的本性，成了我们所是的一切，却没有失去祂先前所是的。先前因祂天上的根源，如今因祂地上的构成，天主都是祂的父。因祂地上的构成，天主是祂的父，因为万物都来自天主父，且天主是万物的父，因为万物都出于祂并在于祂。但对于独生子天主，天主是父，不仅因为圣言成了血肉；祂的父职也延伸到那作为天主圣言、在起初就与天主同在的那一位。因此，当圣言成了血肉时，天主既是因天主圣言的出生，也是因祂肉身的构成而成为祂的父：因为天主是一切血肉的父，尽管不像祂是天主圣言的父那样。但天主圣言虽未停止是天主，却的确成了血肉；当祂如此居住时，祂仍然真是圣言，正如圣言成了血肉时，祂仍然真是天主也是人。因为\Quote{居住}只能指那寓居于某物中的那一位；而\Quote{成了血肉}指那被生的那一位。祂住在我们中间；就是说，祂取了我们的血肉。圣言成了血肉，住在我们中间；就是说，祂在我们身体的事实中是天主。如果基督耶稣，那按血肉的人，剥夺了天主圣言的神性，或没有按虔敬的奥迹也是天主圣言，那么，天主是祂的父，也是我们的父，祂的天主和我们的天主，便将祂的本性降低到我们的水平。但如果天主圣言在成为人基督耶稣时，并没有停止是天主圣言，那么天主就同时是祂的父和我们的父，祂的天主和我们的天主，只是就那本性而言，圣言由此成为我们的弟兄，而给祂弟兄们的讯息：\BibleQuote{我升到我的父和你们的父那里，我的天主和你们的天主那里}——这并非独生子天主圣言的讯息，而是成了血肉的圣言的讯息。\par

17. 宗徒在此用审慎措辞说话，其明确性不会给不虔敬者任何把柄。我们看到福音作者让主在讯息的序言中使用\Quote{弟兄}一词，以此表明整个讯息是向祂的弟兄们说的，涉及祂与那本性——这本性使祂成为他们的弟兄——的共融。因此，他显明此处所宣告的虔敬的奥迹并非祂神性的贬抑。与祂的共融——藉此天主成为我们的父和祂的父，我们的天主和祂的天主——存在于肉身的\OriginalTerm{经济}{Dispensation}中：我们被算为祂的弟兄，因为祂生于肉身。无人否认天主父也是我们主耶稣基督的天主，但这虔敬的宣认并不提供不敬的借口。天主是祂的天主，但并非拥有与祂不同的神性秩序。祂受生为出自父的天主，又藉着\OriginalTerm{经济}{Dispensation}生为仆人：因此，天主是祂的父，因为祂是出自天主的天主；而天主是祂的天主，因为祂是童贞女的肉身。宗徒用一句简短而决定性的句子确认了这一切：\BibleQuote{在祈祷中时常提到你们，求我们主耶稣基督的天主，荣耀的父，赐给你们智慧和启示的神恩。}\BibleRef{弗 1:16} 当他称祂为耶稣基督时，他提起祂的天主；当他涉及基督的荣耀时，他称天主为祂的父。对于拥有荣耀的基督，天主是父；对于作为耶稣的基督，天主是天主。因为天使在谈论应生于玛利亚的基督主时，称祂为\Quote{耶稣:}；但在先知们看来，基督主是\Quote{神}。宗徒在此处的言词因拉丁文而对许多人显得有些晦暗，因为拉丁文没有冠词，而希腊文优美且合乎逻辑的用法却有冠词。希腊文是\OriginalTerm{我们主耶稣基督的天主，荣耀的父}{\GreekText{ὁ} \GreekText{Θεὸς} \GreekText{τοῦ} \GreekText{Κυρίου} \GreekText{ἡμῶν} \GreekText{᾿Ιησοῦ} \GreekText{Χριστοῦ}, \GreekText{ὁ} \GreekText{πατὴρ} \GreekText{τῆς} \GreekText{δόξης}}，我们若允许使用冠词，可译为拉丁文\Quote{Ille Deus illius Domini nostri Jesu Christi, ille pater illius claritatis}（我们主耶稣基督的天主，荣耀的父）。在这种形式下，\Quote{\Emph{那位}耶稣基督的\Emph{那}位天主}和\Quote{\Emph{那}荣耀的\Emph{那位}父}，这句话表达了我们所能理解的祂本性的某些真理。涉及基督的荣耀时，天主是祂的父；涉及基督是耶稣时，父是祂的天主。在使祂成为仆人的\OriginalTerm{经济}{Dispensation}中，祂以那一位为天主，而在使祂成为天主的荣耀中，祂以那一位为父。\par

18. 对于圣神来说，时间和时代的流逝并无区别。基督是同一基督，无论在身体里，还是藉圣神住在先知们里面。祂藉着圣祖达味的口说：\BibleQuote{你的天主，啊天主，用喜乐油傅油了你，超过你的同伴}，这所指的不亚于祂取肉身的\OriginalTerm{经济}{Dispensation}这一奥迹。那现在向祂弟兄们传送讯息——他们的父是祂的父，他们的天主是祂的天主——的那一位，那时就宣告自己被祂的天主傅油，超过祂的同伴。无人是独生子基督、天主圣言的同伴：但我们知道，藉着那使祂成为血肉的取用，我们是祂的同伴。那次傅油并没有高举那受生的、有福而不朽的、居于天主本性中的那一位，而是建立了祂身体的奥迹，并圣化了祂所取了的人性。宗徒伯多禄为此作证：\BibleQuote{在城中他们果然聚集反对你的圣子耶稣，你所傅油的那一位}\BibleRef{宗 4:27}；在另一场合：\BibleQuote{你们知道，那从加里肋亚开始，遍及全犹太的传言，就是若翰宣讲洗礼以后，纳匝肋人耶稣——天主怎样用圣神和德能傅油了祂}。因此，耶稣被傅油，为成就肉身重生的奥迹。当我们听到父在祂从约旦河上来时所说的话：\BibleQuote{你是我的儿子，我今天生了你}，我们就不怀疑祂如何这样被天主的圣神和德能所傅油。如此，祂肉身的圣化得到了见证，在这见证中，我们必须认出祂被圣神之德能所傅油。\par

但圣言是天主，且在起初就与天主同在，因此，如果傅油是指那本性——关于它，除了它在起初外，我们一无所知——那么它既不能被关联也不能被解释。事实上，那即是天主的一位，不需要用天主的圣神和德能傅自己，因为祂本身就是天主的圣神和德能。所以，祂作为天主，被祂的天主傅油，\Emph{胜过祂的同伴}。而且，虽然在肉身的安排之前，按照法律有许多基督（即受傅者），但那傅油超过祂同伴的基督，是在他们\Emph{之后}来的，因为祂被高举于祂受傅的同伴之上。相应地，预言的言语揭示了傅油是在时期内发生，且相对来说在较晚的时期内发生。\BibleQuote{因为你喜爱正义，憎恨不义，因此天主，你的天主，用喜乐油傅你，胜过你的同伴}。现在，一个随后发生的事实，不可能被定在它们之前。应得某种酬报，预先假设那能应得此酬报者的存在，因为赚得的功绩暗示着曾经有一个能获取它的人。因此，如果我们把独生天主之诞生归于这傅油——这傅油是因喜爱正义和憎恨不义而得的酬报——那么我们必将视祂不是生而为，而是因傅油被提拔为独生天主。然而，这将意味着祂是通过逐渐的进步和晋升而达至完善的元神性，并且祂不是生而为天主，而是后来因祂的功绩被傅油为天主。这样，我们将使基督作为天主本身受到条件限制，而祂是万有条件的最终原因；那么宗徒的话\BibleQuote{万物是藉着祂，并在祂内受造；祂在万有之先，万有都靠祂而存在}\BibleRef{哥 1:16}\Emph{？}主耶稣基督并非因任何事物或藉任何事物而被神化，而是生而为天主：原初就是天主，并非因祂诞生后的任何原因被提升到神性，而是作为圣子；且因由祂所生而与天主同体。然而祂的傅油，虽是某个原因的结果，但并不增强祂内那不能变得更完善的东西。它涉及祂身上那将要通过奥迹的完善而得以完善的部份：也就是说，我们的人性在基督内藉着傅油而被圣化。如果先知在此也教导我们仆人的安排——为此，基督被祂的天主傅油胜过祂的同伴，是因为祂喜爱正义，憎恨不义——那么先知的这些话必定是指基督内那本性，藉着那本性，祂因取了肉身而有同伴。当我们注意到预言之神如何谨慎地选择祂的言语时，我们还能怀疑吗？天主被祂的天主傅油；也就是说，按祂自己的本性，祂是天主；但在傅油的安排中，天主是祂的天主。天主被傅油：但请告诉我，那在起初就是天主的圣言被傅油了吗？显然不是，因为傅油是在祂神圣的诞生之后。所以，被傅油的不是那受生的圣言，那在起初就与天主同在的天主，而是天主内那藉着安排而来、晚于祂神性的本性；当祂的天主傅油祂时，是傅油祂内那仆人的整个本性，是祂在祂肉身的奥迹中所取了的。\par

因此，愿无人以他不敬的解释玷污那在肉身中显现的伟大虔敬奥迹，或自认为与独生子在祂的神性本质方面同等。愿祂成为我们的兄弟和同伴，因为圣言成了血肉居住在我们中间；因为耶稣基督作为人，是天主与人之间的中保。愿祂按照仆人的方式，与我们有一个共同的父亲和共同的天主；并作为傅油超过祂同伴的，愿祂与祂受傅的同伴具有同一本性，尽管祂的傅油是具有特殊特权的。在中保的奥迹中，愿祂同时是真人又是真天主：祂自己是从天主出的天主，但在那使我们成为祂兄弟的共融中，祂与我们有一个共同的父亲和天主。\par

但也许这种屈服，这种交出王国，以及最后这种终结，标志着祂本性的瓦解，或祂权能的丧失，或祂神性的衰弱。许多人这样争论：基督被包括在万物对天主的共同屈服中，并且因屈服的状态而失去祂的神性：祂交出祂的王国，因此祂不再是君王：那降临于祂的终结，其后果是祂权能的丧失。\par

在此，如果我们回顾宗徒在此主题上教导的全面含义，并非不恰当。那么，让我们逐句解读并阐明，以便通过完整理解来掌握整个奥迹。宗徒的话是：\BibleQuote{因为死亡既因一人而来，死者的复活也因一人而来；就如整个的人都在亚当内死了，照样，整个的人也都在基督内复活。但各人要按着次序：首先是作为初果的基督，然后是在基督再来时属于基督的人。再后末期来临，那时基督将消灭一切率领者、一切掌权者和大能者，并把王权交予天主父。因为基督必须为王，直到把一切仇敌屈伏在祂的脚下。最后被毁灭的仇敌便是死亡。因为圣经说：\InnerQuote{天主使万物屈伏在祂脚下。}既然万物都屈伏了，那么，使万物屈伏于祂的那位，明显地不在那屈伏的万物之内。万物都屈伏于祂以后，子自己也要屈伏于那使万物屈伏于祂的父，好叫天主成为万物中的万有。}\BibleRef{格前 15:21}\par

那位\BibleQuote{不是从人来的，也不是借着人，而是借着耶稣基督}被选为外邦人导师的宗徒\BibleRef{迦 1:1}，以他所能调遣的最明确言语，阐述了天上奥迹的奥秘。他曾被提到第三层天上，听见了不可言说的话语，并将人性所能接纳的一切启示于人类理智的感知。但他并未忘记，有些事物在听闻的瞬间无法理解。人的软弱需要时间来审视那不加分辨灌入耳中的话语，使其在心灵真实而完美的法庭上得到审断。理解比听闻更慢地跟随言语，因为是耳朵听闻，而理智理解，尽管是天主将内在含义启示给寻求祂的人。我们从写给弟茂德——那位因祖母和母亲光荣的信德从幼年就受教于圣经的门徒——的许多劝勉中读到这句话：\BibleQuote{\Emph{你要明白我所说的话，因为主必在一切事上赐给你领悟。}} \BibleRef{弟后 2:7}对理解的劝勉出于理解的困难。但天主的领悟恩赐是信德的赏报，因为通过信德，感官的软弱因启示的恩赐而得补偿。那位被宗徒称为\Quote{属天主的人}、保罗在信德中的真子的弟茂德\BibleRef{弟前 5:11}，被劝勉要理解，因为主必在一切事上赐给他领悟。因此，既然知道主必在一切事上赐给我们领悟，我们当记念宗徒也劝勉我们理解。\par

倘若因人性常有的错误，我们执著于自己的某种先入之见，请不要拒绝借着启示恩赐在知识上的进步。如果我们至今只用自己的判断，也不要羞于为更好而改变决定。为明智而谨慎地引导这一进步，同一蒙福的宗徒写给斐理伯人：\BibleQuote{\Emph{为此，我们凡是成全的人，都应怀有这种心情；若你们在什么事上另有看法，天主也会将这事启示给你们。不过，我们已经达到什么程度，就应照这程度前进。}} \BibleRef{斐 3:15}理性不能用先入之见预知天主的启示。因为宗徒在此已向我们指明，那些拥有完全智慧的人之智慧何在；对于另有看法的人，他等待天主的启示，好使他们获得完全的智慧。因此，若有人对这隐藏知识的深奥奥迹另有构想，而我们提供给他们的在某个方面更正确或更可取，愿他们不羞于按照宗徒的劝诫，借天主的启示接受完全的智慧；若他们憎恶住在虚妄中，就不要偏爱无知胜过真理。若对他们——那些另有智慧的人——天主也启示了这事，宗徒就劝勉他们急忙走上已开始的道路，抛却先前无知的概念，借着他们热切进入的道路获得完全理解的启示。因此，让我们继续走在我们已经急忙踏上的道路；或者，若我们游荡脚步的错误延误了热切的匆忙，仍要通过天主的启示重新向我们的渴望目标出发，不要将脚从道路上转开。我们已急忙奔向基督耶稣——光荣的主、永恒万世的君王，在祂内万有在天上和地上得以恢复，万有靠祂而存立，在祂内并同祂我们将永远存留。只要我们走在这条路上，我们就拥有完全智慧；若我们有别的智慧，天主会启示我们何为完全智慧。因此，让我们在宗徒信德的光照下，审视眼前话语的奥秘；让我们的处理——如同一直以来的——是根据宗徒宣认的真理，对每一种他们不敬地强加于其话语的解释加以驳斥。\par

此处争议三个主张，按宗徒陈述的顺序，首先是终向，其次是交出，最后是服从。其目的在证明基督在终向时不再存在，祂交出王国时失去王国，祂服从天主时剥夺了祂的神性。\par

首先请注意，这并非宗徒教导的顺序，因为在他的顺序中，交出王国在先，其次是服从，最后是终向。但每一个原因本身是其特定原因的结果，因此在因果链中，每个原因本身产生一个结果，不可避免地有其潜在的前因。因此，终向将会来到，但那是当祂将王国交给天主时。祂将交出王国，但那是当祂废除一切权柄和能力时。祂将废除一切权柄和能力，因为祂必须为王。祂将为王，直到祂把一切仇敌放在祂脚下。祂将把一切仇敌放在祂脚下，因为天主已将万有置于祂脚下。天主如此置于祂脚下，以致死亡成为祂所征服的最后仇敌。然后，当万有都服从于天主——除了那将万有置于祂之下的那一位——祂自己也要服从于那使自己服从万有的那一位。但服从的原因无非是使天主成为万物中的万有；因此，终向就是天主是万物中的万有。\par

在进一步讨论之前，我们现在必须探究：终向是否是一种消解，交出是否是一种丧失，服从是否是对基督的一种削弱。如果我们发现这些是相反的事，不能作为原因与结果相连，我们就能按这些话被说出的真实意义来理解它们。\par

28. \BibleQuote{基督是法律的終結}\BibleRef{羅 10:4}；但是，請告訴我，祂是來摧毁法律，還是來成全法律？如果基督，法律的終結，並不摧毁法律，反而成全法律（正如祂說：\BibleQuote{我來不是爲摧毁法律，而是爲成全法律}\BibleRef{瑪 5:17}），那麼，法律的終結豈非不是它的廢除，而恰恰相反，乃是它的最終圓滿？萬物都向着一個終結前進，但這個終結是在圓滿中的安息狀態，是它們前進的目標，而非它們的消滅。此外，萬物都是爲了終結而存在，但終結本身並非達到任何超越者的手段：它是一個最終的、涵蓋一切的整體，安息於自身之内。因爲它自足自存，不爲其他時刻或對象而運作，所以終結總是我们希望所指向的。因此，主勸勉我們以忍耐和敬畏的信心等待，直到終結來臨：\BibleQuote{那堅持到底的人是有福的}。等待我們的並非有福的消解，也不是虛無的果實和滅亡作爲信仰的指定賞報：而是終結是最終獲得所應許的真福，那些堅持到底，直到完美幸福的目標達成的人是有福的，那時忠信之望已無進一步的期望。他們的終結是在那朝向的狀態中，以不間斷的安息長存。同樣，作爲警告，宗徒告訴我們惡人的終結：\BibleQuote{他們的終結是滅亡……但我們的期望在天上}。那麼，假設我們把終結解釋為消解，我們就不得不承認，既然無論對有福者還是惡人都有終結，結果就使敬神者與不敬神者平等，因爲兩者指定的終結都是共同的虛無。如果我們和惡人一樣，終結都是存在的停止，那麼我們在天上的期望又是什麼？但是，即使爲聖人保留著期望，而惡人等待他們應得的終結，我們也不能把那個終結理解為最終的消解。對惡人來說，如果因消解而失去了承受報復痛苦的能力，那還有什麼懲罰可言？因此，終結是等待我們的一個最終且不可逆轉的狀態，爲有福者保留，爲惡人預備。\par

29. 因此我們不再懷疑，終結所指的是最終和末後的狀態，而非消解。我們在講解這段經文時還會對此有更多論述，但現在這已足夠説明我們的意思。所以，讓我們現在轉向王國的交付，看看它是否意味着統治權的放棄，是否聖子通過交付而喪失祂交付給聖父的。如果這是惡人在他們不理智的迷惑中所爭論的，他們必須承認，聖父在交付萬物給聖子時也喪失了一切，如果交付意味着放棄所交付之物的話。因爲主說：\BibleQuote{我父將一切都交付給了我}\BibleRef{路 10:22}，又說：\BibleQuote{天上地下的一切權柄都交給了祂}\BibleRef{瑪 28:18}。因此，如果交付是放棄所有權，那麼聖父就不再擁有祂所交付的。但如果聖父並未因交付而停止擁有祂所交付的，那麼聖子也不會放棄祂所交付的。因此，如果祂沒有因交付而失去祂所交付的，我們就必須認識到，只有降生奧蹟才能説明聖父如何仍然擁有祂所交付的，而聖子也沒有喪失祂所給予的。\par

30. 至於服從，還有其他事實有助於我們的信仰，防止我們因此事侮辱基督，但首先這段經文本身就包含了自己的辯護。然而，我首先要訴諸常理：服從是否仍被理解為奴役對主權、軟弱對權能、卑賤對尊榮的隸屬，這些品質互相對立？聖子是否因本性不同而以這種方式服從於聖父？如果我們真要這樣想，我們會在宗徒的話中找到防止這類想像錯誤的預防。他說：\BibleQuote{當萬物都服從祂時，那時，子也要服從那使萬物服從自己的父}；他用\Quote{那時}表示時間性的降生奧蹟。如果我們對服從作任何其他解釋，那麼基督雖然將要服從，但現在卻不服從，這樣我們就使祂成為一個傲慢不敬的叛逆，時間的必要性將如同折斷並制服祂褻瀆和過度的驕傲，使祂歸於遲來的服從。但是祂自己說：\BibleQuote{我不是來行我的意願，而是行那派遣我者的意願}\BibleRef{若 6:38}；又說：\BibleQuote{因此父愛我，因為我一切都按祂所喜悅的去做}；以及：\BibleQuote{父啊，願你的旨意成就}。或者聽宗徒說：\BibleQuote{祂貶抑自己，成了聽命的，以至於死}\BibleRef{斐 2:8}。雖然祂貶抑自己，祂的本性卻不知屈辱；雖然祂聽命，卻是自願的聽命，因為祂藉貶抑自己而成爲聽命的。獨生子天主貶抑自己，聽從祂的父，以至於死在十字架上：但是，是作爲人還是作爲神，在萬物都服從祂時，祂要服從父？確實，這種服從並非新聽命的標誌，而是奧蹟的降生，因為忠誠是永恆的，服從是時間内的事件。因此，服從在意義上只是奧蹟的顯示。\par

31. 此一意义必须从我们信德的同一望德来理解。我们不可不知，主耶稣基督已由死者中复活，并坐在天主右边，因为我们也有宗徒的见证：\Quote{按祂强力的德能，即祂在基督身上所运行的，当祂使祂从死者中复活，并使祂坐在自己的右边，超越一切率领者、掌权者、异能者、宰制者，以及一切现世及来世可称的名号，又将万有置于祂脚下。}宗徒的话语，符合天主的德能，将未来之事当作已经过去来述说：因为那将由时间圆满所实现的事，已在基督内存在，在祂内有一切圆满，而\Quote{未来}仅指救恩工程的时间秩序，并非新的发展。因此，天主已将万有置于祂脚下，尽管它们尚待被屈服。藉着已被视为过去的屈服，表达基督不可改变的德能；藉着作为未来的屈服，预示它们因时间圆满而在诸世代终结时的完成。\par

32. 消灭一切反对祂的权势的意义并不模糊。空中的元首、属灵邪恶的权势，将被交付于永久的毁灭，正如基督所说：\Quote{你们这些被咒骂的，离开我，到那为魔鬼及其使者所预备的永火里去吧！}见：\BibleRef{玛 25:41} 消灭与屈服不同。消灭敌人的权势，就是永远扫除其特权权势，使其权势的消灭终结其王国的统治。主为此作证时说：\Quote{我的国不属于这世界}见：\BibleRef{若 18:36}：正如祂先前曾见证，那王国的统治者是世界的元首，其权势将因祂王国统治的消灭而被摧毁。另一方面，屈服意味着服从与效忠，是顺服与可变的明证。\par

33. 因此，当他们的权威被消灭后，祂的仇敌将被屈服：并且如此屈服，以致祂要将他们屈服于自己。此外，祂要将他们如此屈服于自己，以致天主将把他们屈服于祂。宗徒，你以为，难道不明白福音中这些话的力量吗？\Quote{除非父吸引他，没有人能到我这里来}与另一些话并列：\Quote{除非藉着我，没有人能到父那里去}——正如在这封书信中，基督将仇敌屈服于自己，然而天主将他们屈服于祂，而祂在整个屈服工程中见证，天主在祂内运作。除非藉着祂，没有人能接近父，但除非父吸引我们，也没有人能接近祂。我们既理解祂是天主子，便在祂内认出父的真正本性。因此，当我们学习认识子时，天主父召唤我们；当我们相信子时，天主父接纳我们；因为我们对父的认知和知识在于子，祂在自己内向我们显示天主父，祂若我们是虔诚的，便以父的爱吸引我们与祂的子建立相互的联系。因此，父吸引我们，因为首要条件是祂被承认为父；但没有人能藉着子到父那里，除非对子的信德在我们内运作，因为我们无法以敬礼接近父，除非我们先朝拜子，而如果我们认识子，父便吸引我们获得永生并接纳我们。但每一结果都是子的工程，因为藉着父的宣讲——子所宣讲的父——父将我们带到子面前，子又带领我们到父面前。陈述这一奥迹，是为了更完美地理解当前段落，以显示通过子，父吸引并接纳我们；好使我们理解两个方面：子将万有屈服于自己，以及父将万有屈服于祂。藉着生育，天主的本性在子内永存，并做祂自己所做的事。祂所做的事，天主也做；但天主在祂内所做的事，祂自己做：意思是，在祂亲自行动之处，我们必须相信天主子在行动；而在天主行动之处，我们必须看出父的本性在祂作为子内固有的属性。\par

34. 当率领者和掌权者被消灭后，祂的仇敌将被屈服于祂的脚下。同一宗徒告诉我们这些仇敌是谁：\Quote{就福音而论，他们为你们的缘故是仇敌；就拣选而论，他们为祖宗的缘故是蒙爱的。}见：\BibleRef{罗 11:28} 我们记得他们是基督十字架的仇敌；也记得因他们为祖宗的缘故是蒙爱的，他们被保留来服从，正如宗徒所说：\Quote{弟兄们，我不愿你们不知道这奥迹，免得你们自作聪明：以色列人有一部分是顽梗的，直到外邦人全数进入，这样全以色列都要得救，正如经上所载：\InnerQuote{必有一位拯救者从熙雍出来，将不虔敬从雅各布伯那里除去；这就是我与他们所立的盟约，当我除去他们的罪过时。}}因此，祂的仇敌将被屈服于祂的脚下。\par

35. 但我们不可忘记隶属之后的事，即\BibleQuote{最后被毁灭的仇敌便是死亡}。\BibleRef{格前 15:26} 这战胜死亡无非就是由死者中复活；因为当死亡的腐朽被止住时，那得复活的、如今已是属天的本性便成为永恒的，正如经上所载：\BibleQuote{这必朽坏的必须穿上不朽坏的，这必死的必须穿上不死的。但当这必死的穿上了不死的时候，经上所记的话就应验了：\InnerQuote{死亡在争战中被吞灭了。死亡啊，你的刺在哪里？死亡啊，你的争战在哪里？}} 在祂仇敌的隶属中，死亡被战胜；死亡被战胜后，不朽的生命随之而来。宗徒也告诉我们这隶属所获得的特殊酬报，这隶服从信仰的隶属得以成全：\BibleQuote{祂要运用祂的大能，使万物屈伏于自己，并以此大能，将我们卑微的身体改造，相似祂光荣的身体}。\BibleRef{斐 3:21} 因此还有另一种隶属，即从一种本性转变为另一种本性；我们的本性就其当前形态而言停止存在，并隶属于祂，我们的本性进入了祂的形态。但\Quote{停止}并不意味着存在的终结，而是提升到更高的层面。因此，我们的本性因融入所接受的另一种本性的肖像，并通过新形态的加冕而成为隶属。\par

36. 因此，为使对这奥迹的解释完整，宗徒在说了死亡是最后被征服的仇敌之后，补充说：\BibleQuote{但当祂说\InnerQuote{万物都屈伏于祂，}显然那使万物屈伏于祂的除外，那时子自己也要屈伏于那使万物屈伏于祂的，好叫天主成为万物之中的万有}。\BibleRef{格前 15:27} 这奥迹的第一步是万物都屈伏于祂；然后祂屈伏于那使万有屈伏于自己的那位。正如我们因祂身体统治的光荣而成为隶属，同样，祂也在祂身体的光荣中为王，藉这同一奥迹转而屈伏于那使万有屈伏于自己的那位。我们隶属于祂身体的光荣，为能分享祂在身体中统治的光辉，因为我们必将相似祂的身体。\par

37. 福音书亦未对祂现今统治身体的光荣保持沉默。经上记载主说：\BibleQuote{我实在告诉你们，站在这里的，有人在没有尝到死味以前，必要看见人子来到祂的国里。过了六天，耶稣带着伯多禄、雅各伯和他的兄弟若望，单独带领他们上了一座高山。耶稣在他们面前变了容貌，祂的面貌发光有如太阳，祂的衣服洁白如雪}。这样，宗徒们被显示了基督进入祂国度的身体光荣；因为在祂光荣显圣容的形态中，主以祂统治身体的辉煌显明了出来。\par

38. 祂也应许宗徒们分享祂的光荣。\BibleQuote{世界穷尽时也将如此。人子要派遣祂的天使，把祂国内一切使人跌倒的事和作恶的人收集起来，扔到火窑里；在那里要有哀号和切齿。那时，义人在他们父的国里，要发光如同太阳。有耳的，听罢！} 难道他们的肉体耳朵对听这些话是闭着的，以致主需要提醒他们听吗？然而，主暗示着对奥迹的认识，命令他们听从信仰的教导。在世界穷尽时，一切使人跌倒的事都要从祂的国中被除去。我们于是看见主以祂身体的光荣为王，直到那些使人跌倒的事被除去。随之，我们看到自己因此相似祂身体的光荣，在父的国里发光如同太阳，就是祂在山上显圣容时向宗徒们显示祂国度形态的那种光辉。\par

39. 祂将国度交给天主父，并不是说祂因交出而放弃权能，而是说我们既相似祂身体的光荣，便构成了天主的国度。经上不是说\Quote{祂把自己的国度交出}，而是说\BibleQuote{祂把国度交出} \BibleRef{格前 15:24}，即把藉祂身体光荣化而成为国度的我们交给天主。祂将我们交付于国度，正如福音所说：\BibleQuote{我父所祝福的，你们来吧！承受自创世以来为你们预备了的国度吧！} \BibleRef{玛 25:34} 义人将在他们父的国里发光如同太阳，而子将把那些祂召叫进入祂国度的人作为祂的国度交给父，祂也向他们应许了这奥迹的真福：\BibleQuote{心里洁净的人是有福的，因为他们要看见天主}。\BibleRef{玛 5:8} 当祂为王时，祂将除去一切使人跌倒的事，然后义人要在父的国里发光如同太阳。之后祂将国度交给父，那些祂作为国度交给父的人将看见天主。祂亲自向宗徒们见证这是怎样的国度：\BibleQuote{天主的国就在你们中间}。\BibleRef{路 17:21} 因此，祂是以君王的身份交出国度；若有人问交出国度的是谁，愿他听：\BibleQuote{基督从死者中复活了，成为死者的初果；因为死亡既因一人而来，死者的复活也因一人而来}。\BibleRef{格前 15:20} 所有关于当前论点的言论都涉及身体的奥迹，因为基督是死者的初果。我们也从宗徒的话中领会基督藉什么奥迹从死者中复活：\BibleQuote{你要记着：耶稣基督从死者中复活了，祂是达味的后裔}。\BibleRef{弟后 2:8} 这里他教导说，死亡和复活只归因于基督成为血肉的安排。\par

40. 在祂的身体内——虽然这同样的身体现已变得光荣——祂为王，直到权势被废除，死亡被征服，祂的仇敌被制伏。宗徒谨慎地保存了这区别：权势和异能是\Emph{被废除的}，仇敌是\Emph{被制伏的}。\BibleRef{格前 15:24}然后，当他们被制伏时，祂，即主，将制伏万物于己之圣父所制伏，使天主成为万物中的万有：圣父神性的本性将其自身强加于我们所蒙取的身体的本性上。天主就是这样成为万物中的万有：按照降生奥迹，祂藉其天主性和人性成为人与天主之间的中保，并因此藉降生奥迹获取肉身的本性，并藉此制伏将在万物中获得天主的本性，以致祂并非部分地，而是完全且彻底地成为天主。于是，制伏的终结正是使天主成为万物中的万有，使祂内再无丝毫现世身体的本性存留。虽然在此之前，这两者已在祂内结合，如今祂必须只成为天主；然而，并非藉弃绝身体，而是藉制伏使之转化；并非藉分解而失去身体，而是藉光荣使之变形：将人性加于其神性，而非以人性剥除其神性。祂被制伏，并非为停止存在，而是为要使天主成为万物中的万有，在制伏的奥迹中，祂要继续成为祂不再是的，而非藉分解被剥夺祂自己，即失去其存有。\par

41. 对于这信仰，我们在宗徒的权威中有充足而神圣的保证。藉着降生奥迹，并在此时期内，主耶稣基督，即长眠者之初果，将被制伏，使天主成为万物中的万有；此制伏并非其神性的贬抑，而是其蒙取本性的提升，因为祂原是神而人者，现在彻底成为天主。但有人可能认为，当我们说祂在身体为王时同时已经在身体内受光荣，并且日后将被制伏以使天主成为万物中的万有时，我们的信仰在福音书及书信中都找不到支持。因此，我们不仅要藉宗徒的话，也要藉我们主的口，来提出我们信德的见证。我们将显示，基督首先亲口所说的话，后来藉保禄的口说出。\par

42. 祂岂非以这些话的明确意义，向祂的宗徒启示了这光荣的降生奥迹吗？\Quote{现在人子受光荣，天主也在祂内受光荣。如果天主在祂内受了光荣，天主在祂自己内也光荣了祂，并且立刻光荣了祂。}在\Quote{现在人子受光荣，天主也在祂内受光荣}这些言词中，我们首先有人子的光荣，然后有天主在人子内的光荣。所以，首先是身体的光荣被预示，这光荣是身体从它与神性本性的结合中借来的：接着是提升到更圆满的光荣，这光荣来自对身体光荣的添加。\Quote{如果天主在祂内受了光荣，天主在祂自己内也光荣了祂，并且立刻光荣了祂。}天主在祂自己内光荣了祂，因为祂已经在祂内受了光荣。\Quote{天主在祂内受了光荣：}这指的是身体的光荣，因为藉这光荣，天主的光荣在一个人的身体中被表达出来，在人子的光荣中可见到神性的光荣。\Quote{天主在祂内受了光荣，因此天主在祂自己内光荣了祂：}也就是说，藉着祂被提升至神性，其光荣在祂内增加了，天主在祂自己内光荣了祂。在此之前，祂已经在源自神性光荣的光荣中为王；但从今以后，祂自己要进入神性的光荣。\Quote{天主在祂自己内光荣了祂：}即，在那使天主之所以为天主的本性内。\Quote{使天主成为万物中的万有：}使祂的整个存有，脱离祂为人的降生奥迹，永远转化为神性。这时间对我们也不是隐藏的：\Quote{天主在祂自己内光荣了祂，并且立刻光荣了祂。}当犹达斯起来出卖祂的那一刻，祂将祂在苦难后通过复活所将获得的光荣表明为现在，但把天主将与祂一同光荣祂的光荣归之于将来。天主的光荣在祂内藉着复活的大能显现，但祂自己，出于制伏的降生奥迹，将永远被带入天主的光荣中，即进入那万物中的万有的天主内。\par

43. 异端者们竟认为天主无法达到人所希望达到的目标，这岂不是愚蠢至极的荒谬吗？这无异于暗示祂无力在自身中成就那祂有能力在我们身上成就的事。说天主被祂本性的某种必然性所束缚，只能顾及我们的幸福，却不能将同样的祝福赐予祂自己，这既不合理性也不合常识。天主确实不需要更多的幸福，因为祂的本性和能力在其永恒圆满中坚立不变。然而，在那伟大敬虔的奥迹的时期中，祂虽是神，却成了人，祂并非无力使自己再次完全成为神，因为祂无疑也要将我们这些尚不是的，变为那样的存在。人生命与死亡的最终结果就是复活；我们争战的确切奖赏是不朽与不坏，而非永恒惩罚的无休止持续，而是永恒荣耀的不断享受与幸福。这些属地的身体将被提升到更高本性的样式，并符合主身体的荣耀。但那位取了奴仆形体的神又如何呢？虽然祂仍在奴仆形体中，已在身体内得荣耀，难道祂不会也被模合成神的样式吗？祂将祂荣耀身体的样式赐给我们，却无法为祂自己的身体做更多的事，超过祂为我们共同所做的？异端者们大多将\Emph{\BibleQuote{那时，子也要自己屈服于那使万有屈服于自己的父，好叫天主成为万物之中的万有}}这句经文解释为圣子要屈服于天主圣父，以便通过圣子的屈服，天主圣父成为万物之中的万有。但难道天主还缺少某种完美，要通过圣子的屈服才能获得吗？他们能相信，因为说在时期圆满来临时祂会成为万物之中的万有，所以天主尚未拥有那最终的、有福的神性吗？\par

44. 对我来说，我认为除非藉着虔敬，否则无法认识天主，因此回答这类异议似乎与支持它们同样不圣洁。何等僭妄，竟以为言语能充分描述祂的本性，而思想常常深于言语，祂的本性甚至超乎思想的构思！谈论天主是否有所欠缺，祂自己是圆满的，还是需要比圆满更圆满，这是何等亵渎！如果天主，祂是自己永恒神性的根源，竟然能够进步，以至于今天比昨天更大，那么祂永远无法达到一无所缺的时刻，因为还能进步的本性在其进程中总是留下一些未踏足之地：如果它受进步律的支配，那么尽管不断进步，它必然总是能够进一步进步。但对于那居于完全圆满中、永远存在的那一位，祂没有剩余的圆满可使祂更圆满，因为完全的圆满不能接受更多的圆满。这就是虔敬默观天主的思维态度，即祂一无所缺，祂是圆满的。\par

45. 但是宗徒并没有忽略我们应当以怎样的宣认方式来为天主作证。\Emph{\BibleQuote{哦，天主的智慧和知识多么深奥！祂的判决何其难测，祂的道路何其难寻！谁曾知道主的心意？谁是祂的谋士？谁先给了祂，使祂偿还呢？因为万有都出于祂、藉着祂、归于祂。愿光荣归于祂至于无穷。}}\BibleRef{罗 11:33} 没有属地的心智能定义天主，没有理解能以其洞察力测透祂智慧的深度。祂的判决回避了受造物的探查：祂知识的无迹路径挫败了所有追寻者的热忱。祂的道路淹没在不可理解性的深处：在属神之事中，没有一样能被彻底测透或追踪到底。没有人曾受教认识祂的心意，除了祂自己之外，没有人曾被允许分享祂的谋划。但这一切只适用于我们人，而不适用于祂——万有都是藉着祂而有的，\Emph{伟大谋士的天使}，祂说过：\Emph{\BibleQuote{除了父，没有人认识子；除了子和子所愿意启示的人，也没有人认识父。}}\BibleRef{玛 11:27} 正是为了约束我们软弱的理智，当它试图以描述和定义来探究神圣本性的深度时，我们必须回应宗徒的惊叹语，以免我们以鲁莽的猜测，试图从天主那里夺取祂未曾乐意启示给我们的事物。\par

46. 自然哲学中有一条公认的公理：除非某物服从于感官的观察，否则它不属于感官的范围，例如放在眼前的对象，或发生在人类感官和智力产生之后的事件。前者我们能看到和触摸，因此心灵有能力对其作出判断，因为它可以通过触觉和视觉来检查。后者，即时间中的事件，产生或构成于人类起源之后，它属于知觉感官可以要求作出判断的范围，因为它并非时间上先于我们的知觉和理性。因为我们的视力不能看到不可见之物，它只能分辨可见之物；我们的理性不能投射到它尚未存在的时刻，因为它只能判断它时间上先于的事物。而即使在这些范围内，与本性相连的脆弱也使它对因果序列缺乏绝对确定的认知。那么，它又如何能回溯到它起源之前的时间，并以它的知觉理解那些在永恒领域内先于它存在的事物呢？\par

47. 宗徒於是認識到，除非在感官能力之後，沒有任何事物能成為我們的知識對象。因此，當他斷言天主智慧的深奧、祂不可揣測的審判的無限、祂莫測之道的奧秘、祂不可理解的心思的奧秘、祂未傳達的旨意的不可知性時，他繼續說：\Emph{因為誰先給了祂，以致祂必須償還呢？因為萬物出於祂、藉著祂、並在祂內。}永恆的天主既不受限制，人類的理性和理智在祂存在之前也未行使功能。因此，祂的全部存有是一個我們既不能考察也不能滲透的深淵。我們說祂的\Quote{全部}存有，不是為了定義其有限，而是為了理解其無限的無限性：因為祂不是從任何人領受祂的存有，沒有預先的施予者可以要求祂回報所賜的禮物：\Emph{因為萬物出於祂、藉著祂、並在祂內。}祂不缺少出於祂、藉著祂、並在祂內的東西。萬有的源頭和製造者，包含萬有、超越萬有的，不需要祂內在的東西，創造者不需要受造物，擁有者不需要所擁有的。沒有事物先於祂，沒有事物來自祂以外的任何源頭，沒有事物超越祂。天主還缺少什麼圓滿的元素，需要時間來供應，使祂成為萬有中的萬有？祂從哪裡可以領受呢？如果祂之外一無所有，而祂之外一無所有時，祂永恆地是祂自己嗎？如果祂永恆地是祂自己，且祂之外一無所有，那麼祂將以何種增長而成為圓滿？以何種添加而成為與祂不同的？祂豈不是說：\BibleQuote{我是不變的？}\BibleRef{拉 3:6}在祂內還有什麼改變的可能性？有什麼進步的空間？什麼是先於永恆的？什麼比天主更神聖？因此，子的服從並不會使天主成為萬有中的萬有，也不會使任何原因使祂完美，因為萬有出於祂、藉著祂、並在祂內。祂仍然是祂所是的天主，祂不再需要什麼，因為祂所是的，祂是永恆地出於自己並為了自己。\par

48. 但獨生者天主也不必改變。祂是天主，這是圓滿和完美神性的名號。因為正如我們之前所說，重複的榮耀和服從的原因，是為使天主成為萬有中的萬有：但這是奧蹟，而非必然性，即天主將成為萬有中的萬有。基督在取了僕人的形體時，仍停留在天主的形體內，沒有遭受改變，而是空虛了自己；隱藏祂自己，並且雖被空虛，仍是自己的主人。祂甚至約束自己成了一個人的形體和樣式，以免所取的謙卑的軟弱不能承受祂本性不可測量的能力。祂無限的大能自限，直到能完成順服的職責，甚至忍受與之結合的身體。但既然祂在空虛自己時仍然是自足的，祂的權威沒有受到削減，因為在空虛的貶抑中，祂在自己內行使了那被空虛的權威的能力。\par

49. 因此，為了我們——所取的人性——的進步，天主將成為萬有中的萬有。那位雖然具有天主的形體，卻以僕人的形體被發現的，如今再次要在天主父的光榮中被承認：這就是說，無疑地，祂居住在天主的形體內，在祂的光榮中祂將被承認。因此，這整個只是救恩計劃，而非本性的改變；因為祂仍然留在祂所永遠在的那一位內。但有一個新的本性介入，它隨著祂的人性誕生而在祂內開始，因此祂所獲得的一切都是為那以前不是天主的本性，因為在救恩計劃的奧蹟之後，天主成了萬有中的萬有。所以，是我們受益，我們被提昇，因為我們將被同化於天主身體的光榮。此外，獨生者天主，儘管有祂的人性誕生，仍是天主，祂是萬有中的萬有。身體的服從——藉此祂內一切屬血肉的被吞沒於屬靈的本性——將使祂成為天主和萬有中的萬有，因為祂既是人也是天主；而祂的人性——向這目標前進——也是我們的。我們將被提昇到與那位為我們而成為人的祂相稱的光榮，被更新而認識天主，並重新被造在創造者的肖像內，正如宗徒所說：\BibleQuote{脫去了舊人及其行為，穿上了新人，這新人是在認識天主上，照創造者的肖像而更新的。}\BibleRef{哥 3:9}這樣，人成了天主的完美肖像。因為，被同化於天主身體的光榮，他高升到創造者的肖像，按給第一人指定的模式。他離開罪和舊人，被造成新人以認識天主，並抵達其構造的完美，因為藉著認識他的天主，他成了天主的完美肖像。藉著虔敬他被提昇到不朽，藉著不朽他將永遠作為他創造者的肖像而生活。\par

\SourceRecord{Translated by E.W. Watson and L. Pullan.；Nicene and Post-Nicene Fathers, Second Series；Vol. 9.；Edited by Philip Schaff and Henry Wace.；Buffalo, NY: Christian Literature Publishing Co.,；1899.；<http://www.newadvent.org/fathers/330211.htm>.}

% structure: p0001 source=h1 target=section reason=page-title

\section{\WorkTitle{论天主圣三（卷十二）}}

卷十二：希拉里对亚略派大文本《主在祂道路的起始就创造了我》作出最后解释；这些词语不可按字面理解。基督并非受造，而是造物主（§§1-5）。如果祂是受造物，则圣父也是受造物，因为祂们在性体与尊荣上是一体（§§6,7）。类似段落《我从胎中生了祢》是比喻性的；别处也说到天主的臂与眼。意思是圣子是出自天主的天主（§§8-10）。基督也不是被造的；祂是圣父的圣子，而非作品（§§11,12）。而且祂的儿子身分是直接的，不像我们那样是间接的，也不像以色列是首生子。后一种儿子身分有明确的受生开始，且出自无中（§§13-16）。亚略派的论证未能证明基督的儿子身分具有这些特征（§§17,18）。真理不是靠自信的论辩，而是靠信德达至（§19），但我们仅仅躲避他们的推理还不够；还必须推翻它们（§20）。圣子是永恒而生的，因为是永恒圣父的圣子（§21）。儿子身分包含开始这一反对意见在祂身上不成立（§§22,23）。圣子拥有圣父所有的一切；因此祂拥有永恒和无条件的存在（§24）。祂出自永恒者，因此祂自己是永恒的；出自永恒者，因此不是出自无中。理性无法把握，因而无法驳斥这一点。我们不得断言在祂出生之前有一段时间，有一段时间祂不存在（§§25-27）。我们不得以人受生的类比来论证真理不可能（§28），也不得因为祂永恒存在而说圣子没有被生（§§29-32）。再者，亚略派否认天主永恒的父性；他们说祂常存在，但不常是父。这违背了圣经（§§33,34）。他们论证说，智慧被称为天主受造物的首生者；但在此意义上，受造是受生的同义词，而智慧先于受造（§§35-38）。智慧与天主同永恒（§39），并分享祂永恒的创造旨意（§§40,41）。我们也不可相信基督仅仅为执行创造之工而被生，作为天主的仆役，因为智慧既参与设计也参与执行（§§42,43）。再者，智慧被称为受造的，以表示祂对受造之物的掌管（§§44）。作为天主道路之起始的创造是与永恒受生分开的事件。它意味着基督作为生命之路，在旧约下取了受造的人的外表，在新约下取了实质，以引导我们进入此路。两种意义不可混淆（§§45-49）。然而，没有异端意图的单纯措辞不精确并非不可宽恕（§50）。在最后重申（§51）对基督作为出自天主的天主、永恒圣子的信仰后，希拉里呼求全能圣父，宣告他的信经、对人软弱和信德需要的意识（§§52,53）。圣子是天主的独生子，第二位因为祂是子（§54）。圣神由圣父发出，并由圣子派遣。祂也不是受造物，而是与认识祂奥迹的天主同一性体，并像祂所是的圣神一样不可言喻（§55）。最后，希拉里祈祷，正如他受洗一样，愿他持守三位一体天主的信仰。\par

1. 终于，在圣神加快我们行程之际，我们正临近坚固信德那安全平静的港湾。我们处于长期被海风颠簸之人的境地；他们常常遇到这样的情况：巨大的涌浪暂时将他们困在港口附近的海岸周围，最后那浩瀚可怖的波涛本身却将他们推进一个可靠、熟悉的锚地。我希望，在我们与异端风暴搏斗的这第十二卷书中，这也会发生在我们身上：当我们在这卷书中将可靠的船冒险地驶入这严重不敬的波涛时，正是这波涛本身可将我们带入我们所渴望的安息之所。因为当所有人都被教义的不定之风驱散时，这里有恐慌，那里有危险，然后又常常有船难，因为基于先知权威，人们坚持天主独生子是一个受造物——这样，祂所有的不是出生而是受造，因为曾有论及智慧的话说：\Quote{主在祂道路的起始就创造了我。}\BibleRef{箴 8:22}这是他们掀起的风暴中最大的浪，是旋转风暴的大浪：但当我们面对它，它没有损坏我们的船而破碎时，它将推动我们甚至进入我们所渴望的海岸那全然安全的港湾。\par

2. 然而，我们不像水手一样依靠不确定或空洞的希望：他们按自己的愿望而非确定的知识调整航向，时而变化无常、飘忽不定的风会抛弃他们或驱离航线。但我们身边有那永不失落的信德之神，借着独生天主的恩赐常与我们同在，引我们沿着坚定的航线进入平静的水域。因为我们认出主基督不是受造物，的确祂也不是；也不是什么受造的东西，因为祂本身就是一切受造之物的主；但我们知道祂是天主，是圣父天主的真实受生。我们众人，正如祂的良善所愿，被命名为天主的儿子并被收养：但祂对圣父天父是那唯一、真实的儿子，是真实而完美的受生，这仅存于圣父与圣子的认知中。但这仅此一点，就是我们宗教的全部：宣认祂是儿子，不是被收养的，而是受生的，不是被拣选的，而是受生的。因为我们不说祂是受造的，也不说祂是未受生的；因为我们既不将造物主与祂的受造物相比，也不虚假地说没有受生就出生了。祂不是自有的，祂是借着受生而存在的；祂也不是未受生的，因为祂是儿子；也不能说，祂是儿子，却以受生以外的方式存在，因为祂是儿子。\par

3. 再者，无人怀疑，不敬之言总是与信德之言相对抗、相抵触；先前已被判为不虔之思的，现在不能再虔信地持守；例如，这些宗徒信仰的新修正者所主张的，福音作者之灵与先知之灵之间存在差异与分歧；他们声称先知预言了一回事，而福音作者宣讲了另一回事，因为\Person{撒罗满}召唤我们朝拜受造物，而\Person{保禄}谴责那些事奉受造物的人。按照这种亵渎的理论，这两处经文显然互不协调：由法律所训练、由天主所分别、并藉\Person{基督}在他内发言而说话的\Person{保禄}，要么不知晓那预言，要么知晓却与之矛盾；因此，当他称\Person{基督}为创造主时，并不知祂是受造物；并禁止崇拜受造物，警告我们唯独事奉创造主，说：\BibleQuote{他们将天主的真理变为虚妄，去事奉受造物，而越过创造主——祂是永受赞颂的。}\BibleRef{罗 1:23}\par

4. 那在\Person{保禄}内说话的天主\Person{基督}，难道不能驳斥这种虚假的不敬吗？难道不能谴责这种歪曲真理的谎言吗？因为万物是藉主\Person{基督}受造的；因此，祂本有的名号就是创造主。难道创造之工的现实与名号不属于祂吗？\Person{默基瑟德}是我们的证人，如此宣告天主是天地创造主：\BibleQuote{愿至高者天主祝福\Person{亚巴郎}，祂是创造天地的主。}\BibleRef{创 14:19} \Person{欧瑟亚}先知也是证人，说：\BibleQuote{我是上主你的天主，祂铺设诸天，创造大地，祂的双手造了天上的万军。}\Person{伯多禄}也是证人，如此写道：\BibleQuote{将你们的灵魂交托给忠信的创造主。}\BibleRef{伯前 4:19}为什么我们将作品的名号归于那作品的创造者？为什么我们将同样的名号给予天主和我们的同类？祂是我们的创造主，祂是所有天军的创造主。\par

5. 既然按宗徒与福音作者的信仰，这些陈述在意义上都指向那藉祂万物受造的圣子，那么祂如何能与祂亲手所造的工本身相等、并与其他一切事物属于同一性质类别呢？首先，我们人类的理智拒绝这种说法：创造主是受造物；因为受造之物是因创造主而存在。但如果祂是受造物，那么祂就必是腐败的、处于等待的悬疑中、并受奴役。因为同一蒙福的宗徒\Person{保禄}说：\BibleQuote{受造物怀着热切的期望，等待天主子民的显扬。因为受造物屈服于虚幻，并非出于自愿，而是由于那使它屈服的那一位，但怀着希望，因为受造物本身也要从腐败的奴役中得释放，进入天主子女光荣的自由。}\BibleRef{罗 8:19} 因此，如果\Person{基督}是受造物，那么祂必是在不确定性中，始终以烦人的期待希望着，并且是祂的长期期待，而非我们的，正在等待；同时，祂在等待时屈服于虚幻，并因一种必要性的屈服、并非自愿地屈服。此外，既然祂并非自愿屈服，祂也一定是个奴仆；再者，既然祂是奴仆，祂也必然居住于腐败的本性中。因为\Person{保禄}教导，这一切都属于受造物，当它通过长期的期待从这些中获得释放时，它将闪耀着属于人的荣耀。但关于天主，这是何等轻率不敬的论断，竟然想象祂因受造物所忍受的侮辱，而遭受这样的嘲弄，即祂应当希望、事奉、被迫、被承认，并以后被释放进入我们而非祂的状况；而事实上，我们的一点点进步都是出自祂的恩赐。\par

6. 然而，我们的不敬，凭借这种禁言的放肆，以更深的不信日益增长；声称既然圣子是受造物，那么就必须坚持圣父也不异于受造物。因为\Person{基督}，祂虽然在天主的形体内，却取了奴仆的形体；如果那在天主形体内的是受造物，则天主绝不能与受造物分开，因为在天主形体内有一个受造物。但处于天主的形体只能被理解为：存在于天主的本性内；由此，天主也是受造物，因为有一个具有祂本性的受造物。但那在天主形体内的，并不以自己与天主同等为强夺的，因为从天主同等（即天主的形体）降到了奴仆的形体。但祂除非作为天主，空虚自己，脱离天主的形体，就不能从天主降下成为人。然而当祂空虚自己时，祂并未被抹去以致不存在；因为那时祂便会变成与先前不同的种类。因为那位在自己内空虚自己的，并未停止成为祂自己；因为祂大能的权能甚至存于空虚自己的权能中；而过渡到奴仆的形体并不意味丧失天主的本性，因为脱下天主的形体本身正是天主大能的一个大能作为。\par

7. 但以这种方式处于天主的形象中，无非就是与天主平等：因此，对于处于天主形象中的主耶稣基督，应给予同等的尊荣，正如他亲自所说：\BibleQuote{叫众人尊敬子如同尊敬父；不尊敬子的，就是不尊敬派遣他的父。}\BibleRef{若 5:23} 事物之间的差异总是意味着不同程度的尊荣。同样的事物值得同样的敬意；否则，最高的尊荣将不配地加于低等事物，或侮辱高等事物而使低等事物与之同尊。但如果子被视为受造物而非所生者，并被施予与父同等的敬意，那么我们就未给予父任何特别的尊荣，因为我们只以对待受造物的敬意对待父。然而，由于子与天主父平等，因为他是从天主所生的天主，因此他在尊荣上也与父平等，因为他是子，而非受造物。\par

8. 关于他，父又有这一显著的言论：\BibleQuote{在晨星之前，我从胎中生了祢。}这里，正如我们经常所说，为了迁就我们软弱的理解，并无贬低天主之意；正如他说\Quote{从胎中}生了祢，就以为他由内外部分构成，这些部分组成他的肢体，并且他的存在如同地上的身体一样，受时间内的相同原因支配；然而，他的存在不受任何自然必然性约束，他是自由、完美、永恒且是万有的主，为了解释其独生子的真正出生特性，他指向自己不变本性的能力。因为尽管神由神所生（要记住，这符合神的真正特性，由于这特性，他本身也是神），但其所生的唯一原因在于那些完美不变的原因之内。尽管它来自完美不变的原因，但必须按照该原因的真正特性从该原因而生。人类出生的必然过程受作用于子宫的原因支配。但既然天主并非由部分组成，而是不变的，因为他是神（天主是神），他内在不受任何自然必然性约束。但由于他向我们讲述神由神所生，他便通过我们中间起作用的原因的示例来教导我们的理解：并非为了给出出生方式的例子，而是为了宣告生育的事实；并非要示例证明他受必然性约束，而是为了启迪我们的心智。因此，如果天主的独生子是受造物，那么用人类出生的常见事实来启示他是神圣所生的启示，还有什么意义呢？\par

9. 因为天主常常藉着我们身体的这些肢体，为我们说明他自己运作的方式，通过使用通常理解的术语来启迪我们的智力：例如他说：\BibleQuote{谁的手创造了天上的万军}；又或：\BibleQuote{上主的眼目看顾义人}；或：\BibleQuote{我找到了叶瑟的儿子达味，一个合我心意的人}。这里，\Quote{心}表示渴望，达味因其品格的正直而蒙天主悦纳；对宇宙万物的认识，即无物超出天主的认知，用\Quote{眼目}一词表达；他的创造活动，即无物不由天主而来，用\Quote{手}这一名称来理解。因此，既然天主愿意、预见并做成一切，甚至在运用表示身体行动的术语时，也必须理解为无需身体的协助；那么，在他说\Quote{从胎中生了}这句话时，所提出的观念并非通过身体行为产生的人类的起源，而是必须理解为神性的出生，因为在其祂提及肢体的情形中，这样做是为了向我们代表天主内其祂的运作能力。\par

10. 因此，既然\Quote{心}代表渴望，\Quote{眼}代表视觉，\Quote{手}代表完成的工作——而天主虽非由任何部分组成，却愿意、预见并行动，这些运作本身就用\Quote{心}、\Quote{眼}、\Quote{手}这些词来表达——\Quote{从胎中生了}这句话的意思难道不是对出生真实性的断言吗？并非他真从子宫生了子，正如他并非用手行动、用眼观看、用心渴望一样。但由于使用这些术语，表明他确实行动、看见并愿意一切，因此从\Quote{胎}这个字，表明他确实从自己生了所生者；并非他使用了子宫，而是他想表达真实性。同样，他并非通过身体官能愿意、看见或行动，而是使用这些肢体的名称，以便我们通过有形力量所履行的服务，理解无形力量的能力。\par

11. 人类社会的构成不容许，主教导的话也不许可，门徒高于师傅，或奴隶管辖主人；因为在所有这些对立地位中，顺服知识是无知的适宜状态，无条件服从是奴役的指定命运。既然众人普遍认为如此，那么是谁的鲁莽会诱使我们说或想天主是受造物，或子是被造的？因为我们从未发现我们的主和师傅这样对祂的仆人和门徒谈论自己，或教导祂的出生是一种受造或被造。此外，父也从未作证祂是别的，只是子，子也从未承认天主是别的，而是祂自己真正的父，确凿地肯定祂是生的，而非受造或被造，正如祂说：\BibleQuote{凡爱父的，也爱那由祂生的子。}\BibleRef{若一 5:1}\par

12. 另一方面，祂在创造中的工程是制作的行动，而非因生育而出生。因为天并非儿子，地也非儿子，世界也非生育；因为关于这些，经上说：\BibleQuote{万有是藉着祂而造成的}\BibleRef{若 1:3}；并藉先知说：\BibleQuote{诸天是祢手的化工}；又藉同一位先知说：\BibleQuote{不要忽略祢手的化工}。画家的画是画家的儿子吗？铁匠的剑是铁匠的儿子吗？建筑师的房屋是建筑师的儿子吗？这些都是他们制作的工程；但唯独那由父所生的，才是圣父的儿子。\par

13. 我们确实是天主的子女，但成为子女是因为圣子使我们如此。因为我们曾是义怒之子，但藉着义子的圣神，被立为天主的子女，并因恩宠而非因生身权利获得了这名分。既然一切受造之物，在受造之前并不存在，同样，我们虽非儿子，却成了现在的样子。因为从前我们不是儿子，但获得这名号之后，我们便是了。此外，我们并非受生，而是受造；并非被生，而是被赎。因为天主为自己赎买了一子民，并藉此行动生了他们。但我们从未得知天主在严格的意义上生了儿子。因为祂并没有说：\Quote{我生了并抚养了\Emph{我的}众子}，而只是说：\Quote{我生了并抚养了儿子}。\par

14. 然而，或许因为祂说：\BibleQuote{我的长子以色列}\BibleRef{出 4:22}，有人会解释祂说\BibleQuote{我的长子}这件事，是为了剥夺圣子受生的特性；好像因为天主也将\BibleQuote{我的}这个称谓用于以色列，那些被立为义子的便错被当作实际受生，因此用在这位身上的话——\BibleQuote{这是我的爱子}\BibleRef{玛 17:5}——就不单单适用于天主的受生，因为（据称）\BibleQuote{我的}这个称谓是与那些显然不是生于义子的人共享的。但是，他们虽然被说成是受生的，实际上并非受生，这一点甚至从那段话中可以看出，经文说：\Quote{有一民族将要诞生，是主所造的}。\par

15. 因此，以色列子民是这样受生的：它被造成；我们并不把受生这件事视为与受造矛盾。因为它是藉着义子名分，而非藉着生育；这并非它的实际本性，而只是它的称号。因为尽管经文写着\BibleQuote{我的长子}，但\BibleQuote{我的爱子}与\BibleQuote{我的长子}之间，却有着巨大而广泛的差异。因为哪里有生育，那里我们便见到\BibleQuote{我的爱子}；但哪里若从万民中拣选，并藉着意志的行为而收养，那里便是\BibleQuote{我的长子}。在此，子民属于天主，在于它作为长子的特性；而在前一种情形中，祂属于天主，则在于祂作为子的特性。再者，在生育的情况下，父亲的拥有权在先，然后才是他的爱；在收养的情况下，首要的事实是子被立为长子，然后才拥有权。因此，对于从地上万民中被收养为子的以色列，长子的特性正当地属于它；但对于那独一的天主受生者，子的特性正当地属于祂。因此，那里若仅归于子之名而非实际，便没有真正而完全的生育：因为那不争的事实是，那进入子民状态的子民也是被造的。但由于它不会成为它现在所是的样子，而且它的受生只不过是它被造的一个名号，它就没有真正的受生，因为它在受生之前曾是别的东西。为此缘故，它在受生之前并不存在，也就是说，在被造之前并不存在，因为那从万民中成为子的，在成为子之前是一个民族：因此它不是真正的子，因为它不总是子。但天主的独生子，既非在任何时候不是子，亦非在成为子之前有任何别的身份，祂自身亦非子以外的任何东西。所以，那永远是子的，使我们不可能想像祂曾有不存在的时候。\par

16. 事实上，人的受生涉及先前的非存在，因为，第一个理由，所有人都是从先前不存在的人所生。因为虽然每个受生者都来自一个曾存在的人，但那个生他的父母，在他出生之前并不存在。其次，第二个理由，受生者在不存在之后才被生，因为时间在他出生之前就已存在。如果他今天出生，在昨天的时间里，他并不存在；他从非存在的状态进入存在的状态；我们的理性迫使我们认为今天出生的事物昨天并不存在。因此，他的出生（藉此他存在）发生在非存在状态之后；因为今天必然意味着昨天的先前存在，所以可以说，曾有一个时候他并不存在。这些事实适用于人类一切事物的起源：万物都有一个开始，在此之前它们并不存在：首先，如我们所解释的，就时间而言，然后就原因而言。就时间而言，无疑现在开始存在的事物先前并不存在；就原因而言也是如此，因为它们的存有并非来自自身内部的原因。因为请思考一切开始的原因，将你的理解引向它们的前因：你会发现没有任何事物是自我起因的，因为没有任何事物是由于父母自由的行为而出生，而是所有事物都因天主的德能被造成它们所是的样子。因此，藉着实际的遗传，每一类事物的自然属性就是：它曾经不存在，然后开始存在，在时间开始之后开始，并在时间内存续。而所有存在的事物都有一个晚于时间开始的开端，它们的原因也依次曾不存在，从曾经不存在的事物出生。甚至人类始祖亚当，也是由泥土形成的，这泥土是从无中造成的，而且在时间之后，也就是说，在天、地、日、月、星辰之后，他的出生没有最初的开端，而他曾不存在时开始存在。\par

17. 但对于天主独生子，祂之前没有时间在先，因此不可能有不存在的时期，因为这样\Quote{某个时期}就变成了在祂之先；再者，断言祂不存在就涉及时间概念：如此时间将不是在祂之后开始，而是祂自身将在时间之后开始，并且，既然祂在受生之前不存在，那么祂不存在的那个时期本身就会先于祂。此外，那从真正\Quote{是}的那一位而生者，不能被理解为是从不存在者而生的：因为真正\Quote{是}的那一位是祂存在的根源，而祂的诞生不可能起源于不存在者。因此，既然在祂身上，无论是关于时间（祂从未不存在），还是关于圣父（即祂存在的根源）说祂是从无中产生的，都是不真实的，那么祂就没有留下任何余地，使人可以认为祂是从无中生出，或认为祂在受生之前不存在。\par

18. 我并非不知道，大多数人因不敬虔而心智迟钝，不接受天主的奥秘，或因敌对之灵的强烈影响，准备在虔诚的伪装下表现出对天主的疯狂贬低热情，他们惯常在头脑简单的人耳中说出奇怪的主张。他们断言，由于我们说圣子永远存在，并且祂从未不是祂永远所是的，因此我们就是在宣称祂是无生的，因为祂永远存在；因为根据人类理性的运作，永远存在的东西不可能被生出：因为（他们争辩说）一件事物被生的原因是，某种不存在的东西可以存在，而某种不存在之物的存在，根据常识的判断，无非就是被生。他们还会加上一些巧妙悦耳的论点——\Quote{如果祂被生，祂就开始存在；在祂开始存在的时候，祂不存在；而当祂不存在的时候，就不能说祂在。}通过这些证明，他们主张合理虔诚的说法是：\Quote{祂在出生之前不存在：因为为了祂得以存在，被生的是不存在的一位，而不是存在的一位。存在的一位不需要出生，虽然不存在的一位被生，是为了祂得以存在。}\par

19. 现在，首先，那些自称虔诚认识神圣事物的人，在福音和宗徒所传讲的真理指明道路的事情上，本应放下狡诈哲学的复杂问题，而应追随那安息于天主的信德：因为诡辩式问题的狡辩容易解除软弱理解的信仰防护，因为狡猾的断言引诱无防备的辩护者，他试图通过调查事实来支持自己的论点，最终通过他自己的调查剥夺了他的确定性；以致回答者不再在意识中保留他通过承认而放弃的真理。因为什么答案能如此符合提问者的目的，如同当我们被问\Quote{一个东西在它被生之前存在吗？}时我们承认说，那被生的先前并不存在？因为无论是就自然还是就必然理性而言，一个已经存在的东西被生是矛盾的：因为一个东西必须被生才能存在，而不是因为它已经存在。但当我们做出这个让步时（因为它是正确的），我们就失去了我们信德的确定性，并且被诱捕而落入他们不敬虔和非基督教的计谋中。\par

20. 但蒙福的宗徒保禄提防这一点，正如我们常显示的，警告我们要谨慎，说：\BibleQuote{你们要谨慎，恐怕有人用哲学和虚空的欺骗，照人的传统，照世上的元素，而不照基督掳掠你们；因为在基督内居住着整个天主性圆满的实体。} \BibleRef{哥 2:8} 因此我们必须提防哲学，而基于人传统的方法，我们与其躲避不如驳斥。我们所做的任何让步，不是表明我们被驳倒了，而是表明我们被混淆了，因为那些宣称基督是天主的德能和天主的智慧的人，理应不逃避人的学说，反而要推翻它们；并且我们必须管束和教导那些头脑简单的人，以免被这些教师们掳掠。因为既然天主能做一切事，且凭其智慧能做一切明智的事，因为祂的旨意并非没有能力武装，祂的能力也并非没有旨意引导，那么，那些向世界宣讲基督的人，就应当以那明智的全能所传授的知识，面对世界不虔敬和有缺陷的学说，正如蒙福宗徒所说的：\Emph{因为我们作战的武器，不是属于血肉的，而是在天主前有能力摧毁堡垒的，并且摧毁一切反对天主知识的推理和一切高耸的事。} 宗徒并没有留给我们一种赤裸且缺乏理性的信德；因为虽然赤裸的信德在救恩上可能极为有力，然而，除非它被教导所训练，它在敌人中间固然有安全的退路可寻，却无法维持一个安全而坚固的抵抗阵地。它的位置就像一座营寨在溃败后为弱旅所提供的，不像那些有营寨可守之人无畏的勇气。因此，我们必须击倒那些反对天主的傲慢论点，摧毁虚假推理的堡垒，粉碎那些以不敬虔自高自大的狡猾心智，用并非属血肉而是属神的武器，不是用世俗的学问而是用天上的智慧；以便如同神圣事物与人类事物不同，天上哲学胜过地上竞争一样多。\par

21. 因此，让不信放弃其努力；不要因为它不理解，就以为我们否认一个我们实际上正确理解并相信的真理。因为当我们用这么多的话语宣告祂被生时，我们并没有断言祂曾经未被生。因为未被生与被生不是同一回事：后者表示源自另一个的起源，前者表示不源自任何存在。而永远存在，作为永恒者，没有任何存在的源头是一回事，而与圣父同永恒，以祂为存在的源头是另一回事。因为凡圣父是存在的源头之处，那里也有生；此外，凡存在的源头是永恒的，生也是永恒的：因为生来自存在的源头，来自永恒存在源头的生必然是永恒的。凡永远存在的，也是永恒的。但并非一切永恒的也是未被生的；因为从永恒所生的永远具有被生的性质；但未被生的既是永恒的，也是\OriginalTerm{非受生}{ingenerate}的。但如果那从永恒者所生的不是永恒生的，那么圣父也不可能是永恒的存在源头。因此，如果从永恒圣父所生的缺少任何永恒性，显然祂存在的创造者也缺少同等的永恒性；因为无限属于生育者，也无限属于被生者。因为理性和理智不允许在天主圣子的生与天主圣父的生育之间有任何间隔；因为生育在于生，而生在于生育。因此，这两个事件完全吻合；除非两者都发生，否则任何一方都不会发生。因此，那因这两个事件而存在的事物，除非两者都是永恒的，否则不能是永恒的；因为两个关联物之一，离开另一个没有任何实在，因为一个不可能没有另一个而存在。\par

22. 但有某人，不能接受这一神圣的奥秘，会说：\Quote{凡被生的，曾经不存在；因为它被生是为了得以存在。}\par

23. 但难道有人怀疑所有曾被生的人，都曾有一段时间不存在吗？然而，从一个曾经不存在的人生出来是一回事，从一位永远存在者生出来是另一回事。因为每一种婴幼儿状态，由于先前不存在，都是从某个时间点开始的。而这又长大为童年，后来更推动青年成为父亲。但人并不总是父亲，因为他通过童年进入青年，通过最初的婴幼儿进入童年。因此，那不总是父亲的人，也不总是生育：但哪里圣父是永恒的，圣子也是永恒的。所以，如果你坚持——无论通过论证或直觉——认为天主，在认识祂的奥秘中我们认识到祂的一个属性是圣父，祂并不总是被生圣子的圣父，那么你也就作为理解和知识的结果而认为那被生的圣子并不总是存在。但如果\OriginalTerm{父性}{fatherhood}是与圣父\OriginalTerm{同永恒}{co-eternal}的，那么\OriginalTerm{子性}{sonship}也必然是与圣子同永恒的。那么，我们如何能用语言或理解来主张：那其属性是祂永远是他所被生的那一位，在祂被生之前不存在？\par

24. 因此，天主的\OriginalTerm{独生者}{Only-begotten}，在自己内包含无形天主的样式和肖像，在一切属于天主圣父的属性上，由于祂内真实天主性的圆满而与祂平等。因为，正如我们在前几卷中所指出的，在能力和尊崇方面，祂与圣父一样强大、一样值得尊敬：同样地，既然圣父永远是圣父，祂作为圣子，也拥有永远为圣子的同样属性。因为根据对梅瑟所说的\Emph{那自有者派遣了我到你们这里来}，我们获得明确的概念：\OriginalTerm{绝对存在}{absolute being}属于天主；因为那\Emph{是}的，不能被思想或说成是非存在。因为存在与不存在是相反的，且这些互相排斥的描述不能同时对一个对象为真：因为一个存在时，另一个必定不存在。因此，凡某物\Emph{是}的地方，思想和语言都不会承认其不存在。当我们的思想向后追溯，并不断被带回到更远的地方以理解那\Emph{是者}的本性时，关于祂的唯一事实，即祂是，永远先于我们的思想；因为那无限存在于天主的性质，总是从我们思想的后顾之眼中退去，尽管我们的思想达到无限远。结果是，我们思想的向后努力永远无法把握任何先于天主绝对存在的属性；因为没有什么呈现出来，使我们能够理解天主的本性，即使我们继续寻求到永恒，除了总是天主永远存在这个事实之外。那么，梅瑟关于天主所宣告的，我们人类智力不能进一步解释的那属性，正是福音书所证实的属于天主独生子的属性：因为\BibleQuote{在起初已有圣言}，因为\BibleQuote{圣言与天主同在}，因为\BibleQuote{祂是真光}，因为\BibleQuote{天主独生者是在圣父怀里}，因为\BibleQuote{耶稣基督是超越万物的天主}\BibleRef{罗 9:5}。\par

25. 因此，祂\Emph{曾是}，祂\Emph{是}，因为祂出自永远是其之所是的那一位。但出自祂——也就是说，出自圣父——便是出生。此外，永远出自永远存在者，便是永恒；但这永恒不是来自祂自己，而是来自永恒者。从永恒者而来的，只能是永恒的；因为如果后裔不是永恒的，那么作为受生源头的圣父也不是永恒的。既然祂的存在特性是祂的圣父永远存在，而祂永远是圣子，并且永恒性体现在\Quote{\OriginalTerm{自有者}{Qui est}}这个名号中，因此，既然祂拥有绝对的存在，祂也拥有永恒的存在。再者，无人怀疑生育意味着出生，而出生的指一个从那时开始存在者，而非一个不持续存在者。此外，毫无疑问，无人能出生已经存在者。因为对于自己持续永恒者，出生原因不能加给祂。但\OriginalTerm{独生子天主}{Deus Unigenitus}，即\OriginalTerm{天主的智慧}{Sapientia Dei}、\OriginalTerm{天主的德能}{Virtus Dei}和\OriginalTerm{天主的圣言}{Verbum Dei}，既然祂被生，就见证圣父是祂存在的源头。既然祂从永远存在者所生，祂不是从无而生。既然祂在万世之前被生，祂的出生必然先于一切思想。没有余地给文字游戏：\Quote{祂在出生之前并不存在。}因为如果祂在我们的思想范围内，即在出生之前祂并不存在，那么我们的思想和时间都先于祂的出生；因为一切曾不存在的事物，由于\Quote{曾不存在}这一断言的含义，都在思想和时间的范围之内，它在时间中划分出一个它不存在的时期。但祂出自永恒者，却一直存在；祂不是\OriginalTerm{非受生的}{ingenerate}，却从未不存在；因为一直存在超越时间，而被生即是出生。\par

26. 因此，我们宣认\OriginalTerm{独生子天主}{Deus Unigenitus}被生，但在万世之前被生；因为我们必须在宗徒和先知明确宣讲所指定的范围内作出我们的宣认；但同时，人的思想无法把握任何超越时间的出生的可理解概念，因为这与尘世受造物的本性不符：它们中的任何一个都不应在一切时间之前被生。但当我们作此断言时，我们如何能将其与同一条道理中矛盾的陈述——祂在出生前并不存在——相协调呢？因为按照宗徒，祂在万世之前就是\OriginalTerm{独生子天主}{Deus Unigenitus}。因此，如果相信祂在万世之前被生不仅是人类理智的合理结论，而且是深思熟虑的信德的宣认，那么，既然出生暗示着存在的源头，而超越一切时间的是永恒的，凡在万世之前被生的都超越尘世的感知，我们若按肉性的意识坚持祂在出生前并不存在，就是借着不虔敬的私意抬高人的理性观念，因为祂被生为永恒者，超越人的感知或肉性的理解。再者，凡超越时间的都是永恒的。\par

27. 因为我们可以在想象或知识中把握一切时间，因为我们知道，今天是现在的东西昨天并不存在，因为昨天存在的东西现在不存在；另一方面，现在存在的东西只是现在，昨天也不存在。借着想象，我们可以回溯过去，以至于毫不怀疑在某座城市建立之前，存在着该城尚未建立的时间。因此，既然一切时间都是知识或想象的范围，我们凭人类理性的感知来判断它；所以我们被认为合理地断言任何事物：\Quote{它在出生之前并不存在，}因为先行的时间先于每一事物的起源。但另一方面，在天主的事情上——也就是说，关于天主的出生——没有任何事物不在万世之前：所以对祂使用\Quote{在祂出生之前}这一说法是不合逻辑的，也不能设想那在万世之前拥有永恒应许者，仅仅（用真福宗徒的话来说）是在永生的希望中——即那不说谎的天主在万世之前所应许的——或者说祂曾一度不存在。因为理性拒绝这样的观念：祂在万世之前就已存在，却是在某事物之后才开始存在，而我们必须宣认祂在万世之前就已存在。\par

28. 我们可以承认，任何事物在万世之前被生不是人的本性之道，也不是我们能够理解的事；但在此我们相信天主关于祂自己的宣示。那么，我们当代的不信者怎能按照人类理智的概念断言，那在出生前不存在的事物，而宗徒信仰告诉我们，它以一种人类理智无法理解的方式总是被生，或者说在万世之前就已存在？因为凡在时间之前被生的，总是被生；凡在万世之前存在的，总是存在。但总是被生者，绝不可能在任何时候不存在；因为在某个时间不存在直接与存在的永恒性相矛盾。再者，总是存在排除了并非总是存在的观念。而且，由于\Quote{祂总是被生}这一前提排除了\Quote{并非总是存在}的观念，我们无法设想祂在出生前不存在的假设。因为很明显，那在万世之前被生的，总是被生，尽管我们无法形成任何在一切时间之前被生的事物的积极概念。因为如果我们必须宣认（显然必要）祂在一切受造物——无论无形或有形——之前，并在一切世代和万世之前，并在一切感知之前被生，而祂因如此被生而总是存在——那么无论如何思想都无法设想祂在出生前不存在；因为那在万世之前被生的，先于一切思想，我们绝不能认为祂曾不存在，因为我们必须宣认祂总是存在。\par

29. 但我们的对手狡猾地抢先提出这个吹毛求疵的异议：\Quote{如果，}他反驳说，\Quote{无法想象祂在出生前不曾存在，那么必然可以想象一位已经存在者被生。}\par

30. 我要反问这位反对者：他是否记得我曾称祂为受生者以外的东西？我是否没有说过，在祂身上，永远存在与受生具有相同的含义？因为对于一个已存在者而言，受生并非真正的出生，而是通过出生实现的自我改变，而受生者的永远存在意味着在祂的出生中，祂先于任何时间概念，并且心灵没有余地设想祂在某个时间未曾受生。因此，在万世以先的永远受生，并不等同于在出生之前存在。但永远在万世以先受生，排除了在出生之前不存在的可能性。\par

31. 同样，这一事实也排除了说祂在出生之前就已存在的可能性；因为超越知觉的祂在每个方面都是超越的。如果虽然永远存在，但受生的概念超越思想，那么说祂在出生之前不存在的概念也同样不可能成为思想的对象。因此，既然我们必须承认永远受生对于我们来说无非是出生这一事实，那么祂在出生之前是否存在这一问题便无法在我们的思想条件下确定；因为祂在万世以先受生这一事实本身永远逃避我们思想的把握。因此，祂受生且永远存在；除了祂是受生者之外，祂不容许其他任何关于祂的理解或言说。因为既然祂先于思想所在的时间本身（因为永远的时间先于思想），祂便禁止思想对祂在出生之前是否存在作出判断；因为出生前的存在与受生的概念不相容，而先前的非存在则涉及时间的概念。因此，永远时间的无限性摧毁了任何涉及时间概念的解释——即祂不存在的概念；而祂的受生同样禁止任何与之不符的概念——即祂在出生之前就已存在的概念。因为如果祂的存在或非存在可以在我们的思想条件下确定，那么受生本身必在时间之后；因为不永远存在者必然在某个时间点之后才开始存在。\par

32. 因此，信德、论证和思想所达成的结论是：主耶稣既受生又永远存在；因为如果心灵追溯过去以寻求关于圣子的知识，那么这一事实——祂受生且永远存在——将不断出现在探索者的知觉中。因此，正如无受生而存在是圣父天主的属性，同样，通过受生而永远存在也必然是圣子的属性。但受生只能宣告有圣父，而圣父的名号也只能宣告有受生。因为无论是那些名号还是事情的本质，都不容许任何中间状态。要么祂并非永远是圣父，除非永远也有圣子；要么如果祂永远是圣父，那么永远也有圣子；因为凡是否认圣子永远为子的一段时间，正是圣父未能永远为父的那段时间；因此，虽然祂永远是天主，但在祂是天主的同一永远期间，祂却不能也成为圣父。\par

33. 然而，不虔敬的声明甚至走得更远，不仅将时间中的受生归于圣子，还将时间中的生成归于圣父；因为生成过程和受生发生在同一时间段内。\par

34. 但是，异端者，你认为承认天主确实永远存在，却并非永远是圣父，这是虔诚和敬虔的吗？因为如果你这样认为算是虔诚，那么你就必须谴责保禄为不虔敬，因为他说圣子在万世以先就已存在\BibleRef{铎 1:2}；你还必须控告智慧本身，因为它自己见证说它在万世之前已被奠定；因为当圣父预备诸天时，它就在圣父身边。但为了让你能指定天主开始为父的起点，首先确定时间必须开始的那个起点。因为如果时间有起点，那么宗徒宣称它们是永远的就是说谎。因为你们所有人都习惯于从太阳和月亮的创造来推算时间，正如论到它们所记载的：\BibleQuote{让它们作为记号、时节、日子和年份。}\BibleRef{创 1:14}然而，那先于诸天——按你们的观点甚至在时间之前——的祂，也先于诸世代。祂不仅先于诸世代，而且先于诸世代之前的世世代代。你为何用可朽坏、属地和狭隘的事物来限制神圣和无限的事物？关于基督，保禄只知道时间的永远。智慧不说它在某物之后，而是说它在万物之前。按照你的判断，时间是由太阳和月亮确立的；但达味却显示基督先于太阳，说：\BibleQuote{祂的名先于太阳。}并且，免得你以为天主的事始于这个宇宙的形成，他又说：\BibleQuote{在月亮之前，直到万代。}这些配得先知感召的伟大人物藐视时间：每一个通道都被关闭，人的知觉无法穿透那超越永远时间的出生。然而，让虔诚想象的信心接受这一点作为其思虑的界限，记住主耶稣基督，独生子天主，是以一种被承认为完美出生的方式受生，并在对其神性的敬拜中，不忘祂是永远的。\par

35. 但我们被指控说谎，连同我们也被攻击的是宗徒所宣讲的道理，因为它虽然承认受生，却坚持受生的永远性；结果是，虽然受生见证了存在的根源，但永远性的断言在神圣受生的奥迹中超越了人类思想的界限。因为有人针对我们提出智慧关于它自己的声明，即当它教导说它受造时，用了这些话：\BibleQuote{主在祂道路的起点就创造了我。}\par

36. 啊，不幸的异端者！你把赐给教会对抗会堂的武器转向反对教会的宣讲之信，曲解那拯救性教义的确定意义，以对抗众人的共同救恩。因为你以这些话坚持说基督是受造物，而不是用永生活智慧的话语去堵住那否认基督在永恒诸世纪之前就是天主、否认祂的能力在人子的一切行动和教导中活跃的犹太人！因为智慧在此段经文中声称，祂为了天主道路的开始和祂从太初的作为而被创造，免得有人以为祂在\Person{玛利亚}之前并不存在；然而并没有用\Quote{受造}这个词来表示祂的诞生是一种创造，因为祂是为了天主道路的开始和祂的作为而被创造的。不，更确切地说，免得有人以为这道路的开始——它确实是人类认识神圣事物的起点——意味着要将无限的诞生置于时间条件之下，智慧宣称自己在诸世纪之前已被立定。因为，既然为道路的开始和天主的作为而被创造是一回事，在诸世纪之前被立定是另一回事，那么立定就被理解为先于创造；而祂在天主的作为之前、在诸世纪之前被立定这一事实，正是为了指明创造的奥迹；因为立定是在诸世纪之前，而创造是为了道路的开始和天主的作为，是在诸世纪开始之后。\par

37. 但现在，免得\Quote{受造}和\Quote{立定}这两个词成为信仰神性诞生的障碍，接着就有这些话：\BibleQuote{在祂造地以前，在祂坚定群山以前，在一切山丘以前，祂生了我}。因此，在诸世纪之前被立定的那一位，在地之前被生；不仅在地之前，也在群山和山丘之前。确实，在这些表达中，因为智慧谈到自己，所以含义多于所言。因为所有用来表达无限概念的对象，都必须是那种在时间上不后于任何单一事物或任何种类事物的东西。但在时间中存在的事物，不可能适当地指示永恒；因为，从它们后于其他事物这一事实来看，它们本身在时间中有自己的开端，因此无法引发无限作为开端的想法。因为，如果天使的起源被发现先于地的创造，那么天主在造地之前生了主基督有什么可惊奇的呢？或者，为什么那被说成在地之前被生的，也被宣示在群山之前出生，不仅在山之前，也在山丘之前；山丘在群山之后被提及，似乎是一种后思，而理性要求必须先有世界，然后才有群山？基于这些理由，我们不能认为这些话仅仅是为了让人理解祂存在于山丘、群山和地之前——祂以自己的无限永恒超越了那些本身先于地、山和山丘的事物。\par

38. 但这神圣的论述并未让我们的理解失去光照，因为它用接下来的话解释了这句话的理由：——\BibleQuote{天主造了区域，既包括无人居住的部分，也包括天下有人居住的高处。当祂预备穹苍时，我与祂同在；当祂安置自己的宝座时。当祂在云霄之上使云层在高层大气中变得巨大时，当祂稳妥地安置穹苍之下的泉源时，当祂坚定大地的根基时，我就在祂身边，将一切结合在一起}。这里有什么时间时期呢？或者，人类智力的概念能超越天主独生子的无限诞生多远？凭借那些我们可以在头脑中构想其受造之物，是无法理解那先于这一切者的诞生的；因此，我们不能坚持说祂确实在时间上居先，但却不是无限的，因为赐给祂的唯一特权是在时间之物之前的出生。因为，如果那些受造之物就其构成而言受时间条件限制，那么祂虽然先于它们全体，也同样受时间条件限制，因为它们在时间内的被造就规定了祂诞生的时间，即祂在那之前出生；因为那先于时间之物的事物与时间的关系与它们相同。\par

39. 但是天主的声音，在我们真智慧中的教训，所说的是完美的，表达的是绝对真理，当它教导说，它自身不仅先于时间的事物，甚至先于无限的事物。因为当诸天被准备时，它已与天主同在。诸天的准备是天主在时间内的行动吗？以至于思想的冲动忽然惊动祂的心智，好像先前迟钝、惰性，且按照人的方式，祂为建造诸天寻求材料和工具？不，先知对天主工作的观念大不相同，他说：\BibleQuote{藉着主的话诸天得以建立，并藉着祂口中的气诸天的一切力量}。然而诸天需要天主的命令才能被建立；因为它们在这个坚固不可动摇的构造中所显示的安排和卓越，并非来自某种物质的混合与掺和，而是来自天主口中的气。那么，这有什么意义呢？天主的智慧在祂准备诸天时与祂同在？因为诸天的创造既不是物质的准备，也不符合天主在祂工作的初步思想中徘徊的性质。因为一切存在于受造物中的事物，总是与天主同在：虽然这些事物就它们的受造而言有一个开始，但就天主的认识和能力而言，它们没有开始。这里先知是我们的见证，说：\BibleQuote{天主啊，祢创造了所有将要存在的事物}。因为虽然未来的事物，就它们将要受造而言，尚待制造，但对于天主来说，在祂那里没有新事或突发之事，它们已经被制造；因为有一个时间秩序用于它们的受造，并且在天主能力的预知工作中，它们已经被制造。因此，智慧在这里教导说，她在万世之前出生，教导说，她不仅先于已经受造的事物，而且与永恒的事物同为永恒，即与诸天的准备和天主居所的区分同永恒。因为这居所并不是在它实际制造时被区分的，因为区分和塑造居所是不同的。诸天也不是在它（理想地）被准备时形成的，因为智慧与天主同在，无论祂是准备诸天还是区分诸天。之后，她在创造诸天的天主旁边塑造诸天：她藉着在祂准备时的同在证明她的永恒；当她在塑造的天主旁边塑造时，她揭示了她的职能。因此，在前面的经文中，它说她在大地、高山和丘陵之前就已出生，因为它要教导说，她在诸天的准备时就已临在；为要表明，即使当诸天被准备时，这项工作在天主的计划中已经完成，因为对祂没有新事。\par

40. 因为创造的准备是永久而永恒的：这个宇宙的结构实际上并非由孤立的思想行为造成的，好像先想到了诸天，然后天主的心里才出现关于大地的思想和计划；祂逐一思想每一部分，以致先铺展大地为平原，然后通过更好的计划使之升起为高山，再又变化出丘陵，第四步甚至使高处也可居住；如此，诸天被准备，天主的居所被区分，上空巨大的云层容纳了被风收集的蒸发物；然后，可靠的泉水开始在诸天之下奔流，最后，大地以坚固的根基得以稳固。因为智慧宣告她先于这一切。但由于诸天之下的一切都是通过天主造的，基督在塑造诸天时临在，甚至先于被准备的诸天的永恒，这一事实不容许我们思想天主时认为祂对细节有不相连的思绪，因为这一切的准备与天主同永恒。因为，虽然如梅瑟所教导，每个创造行为都有其恰当的顺序——使穹苍坚固、露出旱地、聚集海水、排列星宿、水和大地生出活物——然而诸天和大地及其他元素的创造在天主的工作中并无丝毫间隔，因为它们的准备在天主的计划中已在同等的永恒中完成。\par

41. 这样，虽然基督以这些无限而永恒的旨意与天主同在，但祂赐给我们的无非是对祂诞生这一事实的认识；为的是，正如对诞生的理解是导向对天主的信德的方法，同样，对祂诞生永恒性的认识也能维持虔诚；因为理性或经验都不允许我们称任何不是永恒圣子出自永恒圣父。\par

42. 但也许\Quote{创造}一词及其对祂的使用困扰我们。若万世之前的诞生和为天主道路的开端及祂的作为而受造未被肯定于祂，那么\Quote{创造}一词确实会困扰我们。因为诞生不能被理解为表示创造，因为诞生先于成因，而创造通过成因发生。因为在诸天准备之前、在世代开始之前，祂就被奠定了，祂为天主道路的开端及祂的作为而受造。可能为天主道路的开端及祂的作为而受造与万有之先诞生是同一回事吗？不：这些观念之一涉及行动中的时间，而另一个具有与时间无关的含义。\par

43. 或许你希望将\Quote{祂为工作而被造}这一说法理解为：祂是因工作而被造的；换句话说，基督是为了执行工作而被造的。这样一来，祂就存在为宇宙的仆人和建造者，而非生为\OriginalTerm{荣耀之主}{Lord of Glory}；祂被造是为了服侍万代的形成，而非永远是钟爱之子、万代之王。但是，尽管基督徒的一般理解反驳你这不敬的想法，承认：为天主道路的开始并为祂的工作而被造，与在万世之前受生，是两回事；然而，同样的这段经文也挫败了你错误断言主基督因宇宙的形成而被造的企图，因为它表明天主圣父是宇宙的创造者和形成者，并且令人信服地表明，因为基督亲自参与，在形成万有的那一位旁边一同塑造。但是，尽管全部圣经都旨在谈论主耶稣基督为宇宙的创造者，\OriginalTerm{智慧}{Wisdom}为杜绝一切不敬的借口，在这里宣告：虽然天主圣父是宇宙的构造者，然而智慧本身在祂构造时并未离开祂，因为当祂预先准备时，智慧已与祂同在；当圣父形成宇宙时，智慧也在形成的那一位旁边一同塑造，并在他预备时与他同在。因此，智慧要我们明白：祂并不是因天主的工作而被造的，因为祂在预定的永恒工作中已经临在，并通过在形成万物时与天主同在而证明圣经并不虚假。\par

44. 异端者，你终于要从公教教导的启示中学习：基督为天主道路的开始并为祂的工作而被造的这句话含义是什么；并要由智慧本身的话语，受教于你那个不敬的愚钝的愚蠢。因为智慧如此开始：\Quote{若我要向你们宣告日常所行的事，我必记起那些从古时起的事。}因为智慧先前说过：\Quote{人类啊，我恳求你们，我向人子发出我的声音。愚昧人哪，你们要明白机敏；无知的人哪，你们要用心。}又说：\Quote{藉着我，君王统治，有权势者颁布公义。藉着我，王子被高举，藉着我，暴君拥有大地。}又说：\Quote{我行走在公平的道路上，漫游在正义的路径中间；为要分赐产业给那些爱我的人，用珍宝充满他们的宝库。}智慧并不沉默于祂的日常工作。首先劝诫众人，祂劝导愚昧人明白机敏，无知的人用心，好使热切勤勉的读者思考话语的不同和分别的意义。于是祂教导说，藉着祂的方法和法规，一切成功、知识、名声或财富的获得都得以成就：祂表明在祂之内包含了君王的统治、有权势者的谨慎、王子的著名事迹，以及拥有大地的暴君的公义；祂更不与恶行混杂，不参与不义的行为；所有这些都是智慧所行，为的是藉着参与每一件公平和正义的工作，祂能为那些爱祂的人提供永恒之物的财富和不朽坏的珍宝。因此，智慧在宣告祂要讲述日常之事后，应许祂也要记起那些从古时起的事。如今，何等盲目，竟以为在万世开始前已有事情成就，而这些事明确宣称只是从万世开始算起！因为凡从万世开始算起的工作，其本身都在那开始之后；相反，万世之前的事则先于万世的秩序，后者迟于它们。所以，智慧在宣告祂要记起那些从万世开始算起的事之后，说：\Quote{主为祂工作的道路的开始创造了我}，这些话指的是从万世开始算起所成就的事。因此，智慧所教导的，并非那宣称在万世之前的受生，而是那与万世本身同时开始的\OriginalTerm{分施}{dispensation}。\par

45. 我们还必须探究，那在万世之前受生的天主，又为天主道路的开始并为祂的工作而被造，这说法是什么意思。这的确是说：凡在万世开始之前有受生的，便有无穷无尽受生的永恒；但凡同一受生从万世开始起被视为一创造，为天主的工作和道路，它便被应用到工作和道路作为创造的起因。首先，既然基督是\OriginalTerm{智慧}{Wisdom}，我们必须看祂本身是否是天主工作的道路的开始。对此，我想没有疑问；因为祂说：\BibleQuote{我是道路}，又说：\BibleQuote{除非经过我，谁也不能到父那里去。}\BibleRef{若 14:6}道路是行路者的向导，奔跑者的标线，无知者的护卫，可以说是未知和渴慕之事的导师。因此祂为道路的开始、为天主的工作而被造，因为祂是道路，引领人们到父那里。但我们必须寻求这创造的目的是什么，它从万世开始起算。因为这也是末后分施的奥秘，其中基督再次在形体上被造，并被宣告是天主工作的道路。再者，祂从万世开始起为天主的道路而被造，那时祂使自己服从于有形的受造物形态，取了受造者的形状。\par

46. 那么，让我们看看为了天主的什么道路，为了天主的什么作为，智慧从世代的开始就被受造，尽管她是在万世之前由天主所生。亚当听到了在乐园中行走的一位的声音。你认为，如果祂没有取受造物的外貌，祂的临近能被听见吗？祂行走时被听见，这难道不是证明祂以受造的形象临在吗？我不问，祂以什么外貌向加音、亚伯尔和诺厄说话，以及祂以什么外貌接近哈诺客并祝福他。一位天使向哈加尔说话，而祂无疑也是天主。当祂像天使一样显现时，祂是否具有与祂作为天主的本性相同的形象？当然，天使的形象被显示出来，随后才提到天主的本性。但我为什么提到天使呢？祂作为一个人来到亚巴郎面前。在人的外貌下，在那受造物的形状中，难道基督不是以祂作为也是天主所拥有的本性临在吗？一个人说话，以身体临在，被食物滋养；然而天主被崇拜。确实，那位曾是天使的现在也是人，为了拯救我们，免得我们设想这些多样化的方面——同一种受造物状态的不同方面——中的任何一个就是祂作为天主的自然形象。再者，祂以人的形状来到雅各伯面前，甚至抓住他摔跤；祂用手抓住，用四肢搏斗，弯曲腰身，采取我们的每一个动作和姿势。但祂再次显现，这次是给梅瑟，如同火焰；以便你学会相信，这受造的本性要为祂提供外在的外貌，而不是体现祂本性的实在。那一刻，祂具有燃烧的能力，但祂并没有采取火的本性所固有的毁灭性质，因为火显然燃烧了，而荆棘却没有受损。\par

47. 纵观整个时间进程，认清祂以什么外貌出现在农的儿子若苏厄面前——这位先知以祂的名字命名——或以什么外貌出现在依撒意亚面前（他向祂叙述说看见了祂，正如福音所见证的\BibleRef{若 12:41}），或以什么外貌出现在厄则克耳面前（他甚至被接纳认识复活），或以什么外貌出现在达尼尔面前（他在世代永国中承认人子），或以什么外貌出现在所有其余的人面前（祂以各种受造物的形象向他们显现），\Emph{为天主的道路和天主的作为}，也就是说，为了教导我们认识天主，并有益于我们的永恒状态。为什么这种专为人类救恩设计的方法，在现今却导致对祂永恒诞生的如此不敬的攻击？你所谈论的受造物，是从世代的开始就开始的；但祂的诞生是无终的，且在万世之前。你当然可以坚持说我们在强解词语，如果一位先知，或主，或一位宗徒，或任何神谕曾用\Quote{受造}之名来描述祂永恒神性的诞生。在所有这些显现中，作为吞噬之火的天主，以受造的方式临在，以至于祂可以用祂用以取它的同一能力卸下这受造的形象，能够再次毁灭那仅仅为被观看而存在的事物。\par

48. 但是，宗徒将那在童贞女体内成胎的、有福而真实的肉身的诞生既称为创造也称为制造，因为那时我们受造之物的本性和形象都诞生了。因此，祂是天主自己的儿子，以人的形态和人的起源而被制造；祂不仅是制造，而且也是受造，正如经上所说：\BibleQuote{但是，当时期一满，天主就派遣了自己的儿子，生于女人，生于法律之下，为要把在法律之下的人赎出来，使我们获得义子的身份。}\BibleRef{迦 4:4} 因此，祂是天主自己的儿子，以人的形态和人的起源而被制造；祂不仅是制造，而且也是受造，正如经上所说：\BibleQuote{正如真理在耶稣内，你们应按从前的的生活方式，脱去那因私欲的迷惑而败坏的旧人。然而，你们应在心思的灵里更新，并穿上那按天主受造的新人。}\BibleRef{弗 4:21} 因为那曾是子的，也生为人子。这不是神性的诞生，而是肉身的创造；那新人取了族类的名称，且按那在万世之前受生的天主受造。至于那新人如何按天主受造，他在后面解释，加上：\Emph{在正义、圣德和真理中}。因为祂内没有诡诈；祂为我们成了正义、圣化，并且祂本身就是真理。那么，这就是基督，按天主受造的新人，我们所穿上的那位。\par

49. 那么，如果智慧在说它记念从世代的开始以来所成就的事时，说它是为天主的作为和天主的道路而受造的；然而，在说它受造的同时，又教导说它在万世之前已被建立，免得我们以为那受造形象——如此多样且频繁地取用——的奥秘涉及其本性的某种变化；——因为其建立的稳固性不允许任何能够推翻它的动乱，然而，唯恐建立似乎意味着比诞生少一些，智慧宣称自己在万有之前受生：——如果是这样，那么为什么\Quote{受造}一词现在被用于那既在万有之前受生，又在万世之前被建立之物的诞生呢？因为那在万世之前被建立的，从世代的开始起，为了天主的道路和祂的作为的开端而被重新创造。我们必须在这个意义上理解从世代的开始开始的受造与那先于世代和万有的诞生之间的区别。不敬至少没有这个借口，即它可以以错误为由为自己的亵渎辩护。\par

50. 因为虽然理解的软弱可能妨碍虔敬之人的感知，即使在如此解释之后，他仍可能无法领会\Quote{受造}的意义；尽管如此，宗徒话语的字面本身，当他将\Quote{造成}一词用于真正的诞生时，本应足以让一个真诚的（即使不是聪慧的）信仰相信：\Quote{受造}一词旨在导向对生成的信仰。因为当宗徒意欲宣称那从一位父所生的一位（即主从一位童贞女所生，没有因人类情欲而产生的受孕）时，他显然有一个明确的目标，称祂为\Quote{由女人所生}，而祂知道自己并曾多次宣称祂是所生的。祂希望\Quote{出生}指向生成的真实性，而\Quote{造成}则证明那从一位父所生的一位，因为\Quote{造成}一词排除了通过人类交合受孕的观念，它被明确声明祂是由童贞女所\Emph{造成}的，尽管同样肯定祂是所生而非造成的。但是，看哪，异端者，你是多么不虔敬。没有任何先知、福音作者或宗徒的话曾说耶稣基督是从天主受造，而非从祂所生；而你否认出生，却断言受造，但不是按照宗徒的意义，当祂说祂是受造的，免得有任何怀疑祂是作为从一位父所生的一位而出生。你以最不虔敬的意义作出断言，暗示天主不是通过传达本性的出生而获得祂的存在；尽管受造物本应从无中产生。这是你不幸心灵的首要感染，不是因为你称出生为受造，而是因为你使你的信仰适应受造而非出生的观念。然而，虽然这会标志一个贫乏的理解力，但仍不会标志一个完全不虔敬的人，如果你称基督为受造的，以便人们能认识到祂从天父那不可能的出生，即作为从一位所生的一位。\par

51. 但是，坚定的宗徒信仰不允许任何这些说法。因为它知道基督是在何种时间秩序中被受造的，以及祂是在何种永恒中被生的。此外，祂是天主从天主所生，祂真实出生和完美生成的神性是不容置疑的。因为关于天主，我们只承认两种存在方式：出生和永恒；而且，出生不是在任何事物之后，而是在万有之前，因此出生只见证了存在的根源，并不断言后裔与存在根源之间的任何不一致。尽管如此，按照共同的承认，这个出生，因为它来自天主，暗示了一个相对于存在根源的次级位置，但又不能与该根源分离，因为任何试图超越接受出生事实的思想，也必然要穿透生成的奥秘。因此，这是唯一虔敬地谈论天主的方式：认识祂是圣父，并同祂认识那由祂所生的圣子。确实，关于天主，我们没有被教导任何东西，除了祂是天主独生子的圣父和创造者。所以，不要让人的软弱超越自身；且让它只作这个宣认，其中唯有它的救恩——就是在降生成人的奥秘之前，它始终确信关于主耶稣基督的这一事实：祂是所生的。\par

52. 至于我，只要我借着你赐予我的圣神而拥有能力，圣父，全能的天主，我将不仅永远宣认你是天主，也永远宣认你是圣父。我也绝不会陷入如此的愚妄和不虔敬，竟使自己成为你全能和你奥秘的审判者，也不会让这软弱的理解力傲慢地寻求超越那虔敬的信仰，即教导我的对你无限的信仰和对你永恒的信仰。我不会宣称你曾没有你的智慧、你的能力和你的圣言，没有独生天主，我的主耶稣基督。那限制我们本性的软弱而不完美的语言，不会主导我对你的思想，以致我言语的贫乏窒息信仰为沉默。因为虽然我们拥有我们自己的话语、智慧和能力，是我们自由内在活动的产物，然而你的是完美天主的绝对生成，祂是你的圣言、智慧和能力；因此祂永远不会与你分离，祂在这些你永恒属性的名字中显示为从你所生。然而，祂的出生只被显示到足以显明你是祂存在的根源；但足以确证我们对祂无限的信仰，因为记载说祂是在永恒之前所生的。\par

53. 在人间的事务中，祢在我们面前设置了许多事物，我们虽不知其因由，但效果并非未知；而敬畏在我们本性固有的无知处灌输信德。因此，当我向祢的天空举起我这双微弱的眼睛时，我对它的确切认知仅限于它是祢的。因为我在其中看到这些安置星辰的轨道，它们的年周运转，以及昴星团、大熊星和晨星，各有其被指派的服务中不同职责，我认识到祢在其中的临在，哦天主，即使我无法对它们有任何清晰的理解。当我观看祢的海洋那奇妙的涨落时，我知道我不仅未能理解水的起源，甚至未能理解这变幻水域的运动；但我仍把握住对某个合理原因的信德，尽管它是我不能看见的，并且在这些我也不认识的事物中也不失于认识祢。此外，当我转念思虑大地时，它借着隐秘力量的力量，使之腐烂它所收到的所有种子，在腐烂后使之活起来，在活起来后使之繁衍，在繁衍后使之强壮；在所有这些变化中，我找不到任何我的理性所能理解的东西，但我的无知有助于认识祢，因为虽然我对那服侍我的本性一无所知，我却借着实际经验到我拥有的好处而认识祢。再者，虽然我不认识自己，但我察觉如此之多，以致我因不认识自己而更加惊异祢。因为在不理解的情况下，我察觉到我心智在发挥其能力时的某种运动、秩序或生命；而这察觉本身我归功于祢，因为祢虽拒绝我理解我本性的最初起源的能力，却赐予我察觉本性及其魅力的能力。而且，既然在自己身上我认识祢，尽管无知，如此认识祢，我就不愿在有关祢的事上，因不理解而对祢的全能抱持较弱的信德。我的思想将不试图把握并掌控祢独生子起源，我的官能将不为超越真理即祂是我的造物主和我的天主的真理而费力。\par

54. 祂的诞生是在永远之前。若有什么先于永远，那将是，当永远被理解时，仍避而不解的东西。而这东西是祢的，并且是祢的独生子；并非部分，也不是延长，也不是为适合祢行动方式的某种理论而捏造的空名。祂是子，是由祢、天主圣父所生的子，祂自己是真天主，由祢在祢本性的统一中受生，当在祢之后被承认，却与祢同在，因为祢是祂永恒起源的永恒作者。因为祂来自祢，祂次于祢；但既然祂是祢的，祢不可与祂分离。因为我们绝不可断言祢曾一度无祢的子而存在，免得我们指责祢或有缺陷，因当时不能生育，或在生育后有多余。因此，对于我们而言，永远受生的确切意义就是，我们知道祢是祢独生子的永远之父，祂在永远之前由祢所生。\par

55. 但是，就我而言，我不能满足于以我的信德和声音的服务，否认我的主和我的天主、祢的独生子耶稣基督是受造物；我还必须否认\Quote{受造物}这名号属于祢的圣神，因为祂由祢发出，并透过祂被派遣，我对一切属于祢的事物的敬畏是如此之大。而且，因为我知道唯独祢是未生的，独生子是由祢所生，我不会拒绝说圣神是被生的，或断言祂曾被创造。我惧怕那些因使用这名号\Quote{受造物}而暗示于祢的亵渎，这受造物是我与祢所造的其他存有共有的。祢的圣神，如宗徒所说，搜查并知道祢的深奥事，作为我的代言者，向祢说出我无法言说的话；而我难道要用\Quote{受造物}这名号来表达，更确切地说，羞辱祂本性的能力，这能力永远自存，由祢透过祢的独生子而来？除了属于祢的，没有什么能穿透到祢内；也不能有一个异于并陌生于祢的权能之能动性，来度量祢无限尊威的深度。凡进入祢内的是属于祢的；凡以搜查的能力居住在祢内的，没有什么对祢是陌生的。\par

56. 但我无法描述祂，因为我无法描述祂为我恳求的情形。正如在启示中，祢的独生子在永恒之前由祢所生，当我们停止与语言上的含糊和思想上的困难搏斗时，祂诞生的唯一确定性依然存在；同样，我在意识中紧握真理，即祢的圣神是出于祢并透过祂，尽管我无法凭理智理解它。因为在祢属神的事上，我反应迟钝，正如祢的独生子所说：\BibleQuote{\Emph{你不要惊奇，因我给你说了：你们必须重生。风随意向那里吹，你听到它的响声，却不知道它从哪里来，往哪里去：凡由水和圣神而生的，就是这样。}} \BibleRef{若 3:7} 我虽怀有重生的信德，却是在无知中相信；我拥有现实，却不理解它。因为我的意识并未参与导致这新生命的诞生，而它的效果却显而易见。此外，圣神没有界限；祂按祂的意愿说话，说祂所愿意的，并在祂愿意的地方说话。既然祂来去的原因无人知晓，尽管观察者意识到这一事实，我岂能把圣神的本质归于受造之物，并用确定祂起源的时间来限制祂？祢的仆人\Person{若望}确实说，万物是藉着圣子而造的，圣子作为天主圣言，在起初就与祢同在，哦天主。\Person{保禄}又列举万有都是在祂内受造的，在天上的和在地上的，看得见的和看不见的。\BibleRef{哥 1:16} 并且，当他宣布一切都是藉着基督并在基督内受造时，关于圣神，他认为只需称祂为祢的圣神就已足够。与这些特选的子民一道，我将如此思考这些问题；正如效法他们的榜样，关于祢的独生子，我将不说任何超出我理解的话，只简单地宣告祂是由祢所生，同样，我也将效法他们的榜样，不逾越人类理智所能认识的关于祢的圣神的知识，只简单地宣告祂是祢的圣神。愿我的份不是无益的争辩，而是坚定不移的信德之坚定的宣认！\par

57. 我恳求祢，保持我这虔诚的信德纯洁无瑕，甚至直到我的灵魂离去，赐予我能如此表达我的信念：使我永远紧握我在重生之信经中所宣认的，即当我因圣父、圣子和圣神之名受洗时所宣认的。简言之，让我朝拜祢、我们的圣父，并与祢一同朝拜祢的圣子；让我得蒙祢的圣神的恩宠，祂是由祢藉着祢的独生子而发出的。因为我有一位令人信服的信德证人，祂说：\BibleQuote{\Emph{父，凡是我所有的，都是祢的；祢所有的，都是我的}} \BibleRef{若 17:10}，就是我的主耶稣基督，祂常居于祢内，出于祢，并与祢同在，永远是天主：祂是永受赞美的，直到永远。阿们。\par

\SourceRecord{Translated by E.W. Watson and L. Pullan.；Nicene and Post-Nicene Fathers, Second Series；Vol. 9.；Edited by Philip Schaff and Henry Wace.；Buffalo, NY: Christian Literature Publishing Co.,；1899.；<http://www.newadvent.org/fathers/330212.htm>.}
